\documentclass[11pt]{article}
\usepackage[T1]{fontenc}
\usepackage{lmodern}
\usepackage[margin=1in]{geometry}
\usepackage{amsmath,amssymb,amsthm,mathtools}
\usepackage{aliascnt}
\usepackage{booktabs,longtable,array}
\usepackage{enumitem}
\usepackage{graphicx,tikz}
\usetikzlibrary{arrows.meta,positioning,calc,fit}
\usepackage{microtype}
\usepackage{placeins}
\usepackage{needspace}
\usepackage[numbers,sort&compress]{natbib}
\usepackage{xurl}
\usepackage[colorlinks=true,linkcolor=blue!45!black,citecolor=blue!45!black,urlcolor=blue!45!black]{hyperref}
\hypersetup{pdftitle={Finite congruence lattices: characterization and undecidability},pdfauthor={OpenAI}}
\usepackage[capitalise,noabbrev,nameinlink]{cleveref}
\usepackage{etoolbox}
\setcounter{tocdepth}{2}
\setcounter{secnumdepth}{3}
\setlength{\emergencystretch}{2em}
\setlist{itemsep=3pt,topsep=6pt}
\allowdisplaybreaks[1]
\newtheorem{theorem}{Theorem}[section]
\newaliascnt{lemma}{theorem}
\newtheorem{lemma}[lemma]{Lemma}
\aliascntresetthe{lemma}
\newaliascnt{proposition}{theorem}
\newtheorem{proposition}[proposition]{Proposition}
\aliascntresetthe{proposition}
\newaliascnt{corollary}{theorem}
\newtheorem{corollary}[corollary]{Corollary}
\aliascntresetthe{corollary}
\theoremstyle{definition}
\newaliascnt{definition}{theorem}
\newtheorem{definition}[definition]{Definition}
\aliascntresetthe{definition}
\newaliascnt{convention}{theorem}
\newtheorem{convention}[convention]{Convention}
\aliascntresetthe{convention}
\newaliascnt{example}{theorem}
\newtheorem{example}[example]{Example}
\aliascntresetthe{example}
\crefname{convention}{Convention}{Conventions}
\Crefname{convention}{Convention}{Conventions}
\theoremstyle{remark}
\newaliascnt{remark}{theorem}
\newtheorem{remark}[remark]{Remark}
\aliascntresetthe{remark}
\providecommand*{\lemmaautorefname}{Lemma}
\providecommand*{\propositionautorefname}{Proposition}
\providecommand*{\corollaryautorefname}{Corollary}
\providecommand*{\definitionautorefname}{Definition}
\providecommand*{\conventionautorefname}{Convention}
\providecommand*{\exampleautorefname}{Example}
\providecommand*{\remarkautorefname}{Remark}
\numberwithin{equation}{section}
\DeclareMathOperator{\Con}{Con}
\DeclareMathOperator{\Aut}{Aut}
\DeclareMathOperator{\End}{End}
\DeclareMathOperator{\Sym}{Sym}
\DeclareMathOperator{\Alt}{Alt}
\DeclareMathOperator{\Soc}{Soc}
\DeclareMathOperator{\core}{core}
\DeclareMathOperator{\GL}{GL}
\DeclareMathOperator{\SL}{SL}
\DeclareMathOperator{\PGL}{PGL}
\DeclareMathOperator{\PSL}{PSL}
\DeclareMathOperator{\AG}{AG}
\DeclareMathOperator{\Fr}{Fr}
\DeclareMathOperator{\rank}{rank}
\DeclareMathOperator{\im}{im}
\DeclareMathOperator{\id}{id}
\newcommand{\F}{\mathbb F}
\newcommand{\N}{\mathbb N}
\newcommand{\Z}{\mathbb Z}
\title{Finite congruence lattices:\\characterization and undecidability}
\author{OpenAI}
\date{September 24, 2026}
\begin{document}
\maketitle
\begin{abstract}
We give an explicit colored-graph characterization of the finite nonempty lattices
that occur as full congruence lattices of finite algebras, and prove that
deciding this representation property is undecidable. In particular, the
finite lattice representation problem has a negative answer. We also prove
that recognition of full subgroup intervals in finite groups is undecidable.
\end{abstract}
\tableofcontents
\section{Introduction}\label{sec:introduction}

Gr\"atzer and Schmidt proved that every algebraic lattice is the full
congruence lattice of an algebra \citep[Theorem~I]{GratzerSchmidt1963}.
The carrier in that theorem need not be finite.  The finite lattice
representation problem asks whether every finite lattice is
the congruence lattice of a finite algebra.  The operations of the algebra are
unrestricted: their number, arities, and tables may depend on the lattice.
Thus a negative answer must exclude all finite algebras, rather than only
groups, rings, or algebras in a fixed variety.

There is also a decision question: given the finite order table of a
lattice, can one determine whether such a finite representation exists?
McNulty records this precise question separately from the universal
representation problem \citep[p.~30]{McNulty2015}.  We prove that no such
decision procedure exists, and hence that not every finite lattice has a
finite representation.  We also construct a nonrepresentable family
directly.  Both conclusions must be distinguished from the effective
search for affirmative witnesses, which is possible.

The relation with finite groups is particularly close.
P\'alfy and Pudl\'ak proved that every finite lattice is a congruence lattice
of a finite algebra if and only if every finite lattice is an interval in a
subgroup lattice of a finite group \citep[Theorem~2]{PalfyPudlak1980}.
This is an equivalence between two universally quantified statements.
It does not assert that these two kinds of representation are equivalent
for each individual lattice.  Our reductions below are proved for the
particular lattice families to which they are applied.

Finite partition embeddings do not settle the problem: every finite
lattice embeds in the equivalence lattice of a finite set
\citep[Theorem~7.10]{PudlakTuma1980}, but the image need not be the full
congruence lattice of an algebra on that set.  Positive representation
results include all finite lattices of width at most two
\citep[Theorem~4.2]{Snow2000}.  Small examples likewise give little
indication of the obstruction.
DeMeo's analysis leaves only one possible exception among lattices with at
most seven elements \citep[Theorem~6.1.1 and Figure~6.2]{DeMeo2012}.
We represent that exception explicitly in \Cref{sec:small}.
The obstruction in this paper instead uses long systems of affine
subgeometries and the restrictions that their subgroup labels impose
simultaneously in several ambient simple groups.

\subsection{Statements}

A finite algebra means a nonempty finite set \(A\), together with a finite
list of total operations of finite nonnegative arity.  The list may be
empty.  An operation of arity zero specifies a constant and imposes no
additional condition on an equivalence relation.
Write \(\Con(A)\) for the lattice of all congruences of the algebra.
For a nonempty finite lattice \(L\), let
\[
 \mathcal R(L)\quad\Longleftrightarrow\quad
 L\cong\Con(A)\text{ for some finite algebra }A .
\]
Every representation in this paper is onto the entire congruence lattice.
An embedding into a proper sublattice is not a representation in this sense.

\Needspace{8\baselineskip}
\begin{theorem}\label{m:main}
The following statements hold.
\begin{enumerate}[label=\textup{(\roman*)}]
\item A finite nonempty lattice is representable if and only if it satisfies
the finite colored-graph criterion in \Cref{f:criterion}.
For each candidate graph size the criterion is decidable by a finite test.
\item There is no deterministic Turing machine that, given the order table
of a finite nonempty lattice, always halts and decides \(\mathcal R(L)\).
\item There is a nonrepresentable finite lattice \(L_0\) of minimum
cardinality, specified by the finite formulas in \Cref{sec:minimum}.
Every finite lattice of smaller cardinality is representable, and
\(\lvert L_0\rvert>7\).
\item Given the order table of a finite nonempty lattice \(L\), it is
undecidable whether \(L\cong[H,G]\) for some finite group \(G\) and
subgroup \(H\le G\), where \([H,G]\) is the full subgroup interval.
\end{enumerate}
\end{theorem}

Part~\textup{(ii)} resolves the finite lattice representation problem
negatively, since otherwise the constant affirmative algorithm would be a
decider.  We also construct a nonrepresentable family directly, before the
undecidability argument.  Part~\textup{(iii)} specifies an actual finite
order table by an unambiguous finite description.  Its constants include
one ordinary Busy Beaver integer.  The description does not evaluate
that integer, the minimum cardinality, or the entries of the order table.

The distinction between the three congruence-lattice assertions is useful.
The criterion
in Part~\textup{(i)} gives a search that terminates on representable
lattices; it does not provide a computable stopping rule for negative
instances.  The proof of Part~\textup{(ii)} establishes precisely that
there is no uniform such rule.  In Part~\textup{(iii)}, a single finite
search range is specified mathematically, using the Busy Beaver
constant; no uniform algorithm for evaluating these constants is asserted.

\subsection{Main ideas and proof organization}

We first make the affirmative search explicit.  Edge colors of a finite
complete graph encode principal congruences, and label-decreasing
selfmaps replace the choice of an algebraic signature.  A path condition
ensures that the recovered relations exhaust the congruence lattice.
For fixed graph size this is a finite test.  Enumerating graphs therefore
finds every affirmative witness, and a decider exists if and only if a
computable bound on witness sizes exists (\Cref{sec:graph}).  This is the
effective bridge used in the undecidability proof.

The graph method itself is established background.  Finite unary closure
already gives positive semidecidability
\citep[Section~3.2, Theorem~3.2.1]{DeMeo2012}.  Pudl\'ak's
valuation-colored construction uses graph paths and stable maps, which
preserve an edge color or collapse the edge; its general representation
can have an infinite carrier \citep[pp.~270--272]{Pudlak1976}.
Kenney's graphical-algebra framework already uses label-decreasing maps,
principal-congruence complete graphs, edge-image paths, and reconstruction
by all unary graph endomorphisms
\citep{Kenney2010}; see also \citet[pp.~51--58]{Kenney2009Slides}.
We present the finite criterion directly in terms of the lattice and one
finite graph, with the proof and bounded arithmetic needed for the final
minimum specification.

The first structural construction produces a negative family, without
using a hypothetical decider.  We assemble a simple self-dual lattice
from private truncated affine charts.  On a least carrier, every unary
operation is a permutation or a constant, but the permutation action
need not be transitive.  A full-diagonal construction in a semidirect
power represents the dual lattice as an ordinary subgroup interval;
self-duality gives the required orientation.  Normal-subgroup tests then
leave either an almost simple interval or extension-kernel data.  A
four-dimensional affine grid excludes the latter, yielding an ordinary
almost simple interval for the whole lattice or an initial point link
(\Cref{sec:affine}).

The ordinary-interval step develops methods used by P\'alfy and
Pudl\'ak, including restriction along idempotents on a least carrier and,
in the later subdirect descent, the correspondence between transitive
actions and stabilizer intervals
\citep[Lemmas~1 and~3]{PalfyPudlak1980}.  The subsequent extension
analysis is related to the twisted-wreath and signalizer approaches of
Baddeley and Lucchini, Baddeley, B\"orner, and Aschbacher
\citep{BaddeleyLucchini1997,Baddeley1998,Borner1999,Aschbacher2008};
P\'alfy compares their hypotheses and interval reductions
\citep[Section~4]{PalfyTwisted}.  These results concern different classes
of lattices and do not supply a pointwise reduction for the present
family.  Here the normal tests, compatibility of actual simple factors,
and extension-kernel grid provide that reduction directly.

Recognizing abstract simple-group types at consecutive labels would not
suffice for the obstruction.  The affine patterns force actual nested
simple subgroups, while exclusions are imposed relative to every earlier
ancestor at once (\Cref{sec:patterns,sec:classical,sec:alternating}).
Bounded Lie rank is excluded by algebraic subgroup tests and the
sufficiently-general subgroup method of Larsen and Pink
\citep[Theorem~0.5]{LarsenPink2011}; \Cref{sec:quantitative} supplies an
effective version for the later size ceiling.  In high rank, one must
further exclude field roots and tensor decompositions in every ancestor.
These protections make the nested actions compatible with rational
composition.  Highest-weight theory then produces universal cuts in the
full weight multiset, and multiplicities rule out a sufficiently long
protected path (\Cref{sec:bounded,sec:weights}).  Choosing more affine
generations than this bound gives the nonrepresentable family.

Undecidability requires a different uniformity statement.  A weighted
template is a subgroup interval built from finitely many permutation
sets, called units, whose cardinalities encode integer values.  A
positive solution of an arithmetic circuit equation supplies such a
template, already finitely representable by the coset construction.
The number of units depends only on the circuit's syntax, not its
solution.  We prove that any finite representation yields a new reading
of all the encoded values in one common partition into cells.  Equal
totals encode addition, and square cardinalities encode multiplication;
the new values need not equal the defining ones.

\Needspace{3\baselineskip}
Two independent bounds make this useful.  First, cell sizes are bounded
computably in the size of a representation.  Second, the number of
elements of a template is bounded computably in its number of units,
independently of their weights.  A hypothetical decider would thus bound
the sizes of affirmative witnesses for every template arising from a
fixed circuit.  Reading such a bounded witness would replace any
positive solution by a computably bounded one.  Searching the resulting
finite box contradicts Diophantine undecidability.  No algorithm is
required to construct the unknown-weight template from the circuit
(\Cref{sec:counting,sec:diophantine}).  The same replacement argument,
using bounds on top-group orders instead of algebra carriers, proves
undecidability of full subgroup-interval recognition.

The common reading is the structural burden of this second branch.
Reverse alternating tests reuse the pruning argument through
elementary-linear charts.  A forward marker theorem instead uses marker
omissions and their three-coatom constraints to reach the protected-path
bound (\Cref{sec:reverse,sec:forward,sec:forward-completion}).  Together
with the weighted normal tests, they force the
remaining units into an actual alternating power.  Normal and parity
tests then recover one cell alignment in which every total and square
comparison holds, even after the nontrivial lattice images needed by the
least-carrier reduction (\Cref{sec:weighted}).  The separate counting
argument concerns the literal templates, not those readings.  It uses
nonabelian coordinate compression, effective integer profiles,
bounded-tuple encodings of diagonal matchings, and Mar\'oti's
primitive-group order theorem \citep[Theorem~1.1]{Maroti2002}.

Finally, \Cref{sec:quantitative} gives an ordinary computable size ceiling
for one member of the nonrepresentable affine family.  This step is
independent of both the hypothetical decider and the later
Busy Beaver constant.  The latter constant bounds the witness sizes
returned by the prescribed affirmative graph searches up to the ceiling.
Explicit bounded arithmetic then
selects the first negative order table, ordered first by cardinality
(\Cref{sec:minimum}).

\begin{figure}[tbp]
\centering
\begin{tikzpicture}[
  box/.style={draw,rounded corners=2pt,align=center,text width=4.7cm,
    inner sep=6pt,font=\small},
  arr/.style={-{Stealth[length=2mm]},thick},
  node distance=10mm and 13mm]
\node[box,text width=8.2cm] (engine)
  {Simultaneous pruning in actual simple groups\\
   and the protected-path bound};
\node[box] (affine) at (-3.2,-1.7)
  {Private affine charts:\\a nonrepresentable family};
\node[box] (marker) at (3.2,-1.7)
  {Reverse tests and\\the forward marker theorem};
\node[box,below=of affine] (ceiling)
  {Effective parameter control:\\a computable size ceiling\\for one negative lattice};
\node[box,below=of marker] (weighted)
  {Weighted representations:\\one cell alignment for\\all arithmetic readings};
\node[box,below=of ceiling] (graph)
  {Finite graph witnesses:\\affirmative search and\\conditional size bounds};
\node[box,below=of weighted] (count)
  {Independent template count:\\lattice size bounded by units,\\not numerical weights};
\node[box,below=of graph] (minimum)
  {An ordinary Busy Beaver cutoff\\and finite selection of a\\minimum negative order table};
\node[box,below=of count] (undec)
  {Bounded replacement of\\Diophantine solutions:\\no representation decider};
\draw[arr] (engine.south west) -- (affine.north);
\draw[arr] (engine.south east) -- (marker.north);
\draw[arr] (affine) -- (ceiling);
\draw[arr] (marker) -- (weighted);
\draw[arr] (ceiling.west) -- ++(-0.35,0) |- (minimum.west);
\draw[arr] (graph) -- (minimum);
\draw[arr] (graph.south east) -- (undec.north west);
\draw[arr] (count) -- (undec);
\draw[arr] (weighted.east) -- ++(0.35,0) |- (undec.east);
\end{tikzpicture}
\caption{The two uses of the structural argument.  The affine branch
gives a negative lattice and a computable size ceiling independently of
any hypothetical decider.  The weighted branch combines simultaneous
arithmetic readings with a separate, weight-independent count; a
decider would then bound positive Diophantine solutions.  Finite graph
witnesses also supply the bounded searches used, with the size ceiling,
in the minimum-table specification.}
\label{fig:proof-overview}
\end{figure}


\subsection{Conventions and dependencies}

All finite lattices and algebra carriers are nonempty.  The one-element
lattice is included.  A subgroup interval \([H,G]\) always means the
whole interval.  Invariant-subgroup intervals in a normal direct power
will be identified explicitly and are not silently substituted for
ordinary intervals.  The notation \(\core_G(H)\) denotes the largest
normal subgroup of \(G\) contained in \(H\).
Quotienting by a normal subgroup already contained in the base preserves
the full interval above the base.

Labels such as \(K_X\) always denote actual subgroups in the realization
or quotient currently under discussion.  In a reverse affine chart,
a smaller flat has a larger subgroup label.
A normal subgroup is \emph{retained} when it lies in the current base;
it is \emph{missing} when the relevant normal test forces the base
to supplement it.  An element \(y\) of a lattice is left-modular if
\[
 x\vee(y\wedge z)=(x\vee y)\wedge z
 \qquad\text{whenever }x\leq z.
\]
The field of an affine chart and the field of a group's natural module
need not be the same.  Group bases, field letters and local ranks are
reintroduced when their ambient objects change.

We use the classification and standard structure of finite simple
groups, algebraic groups and their representations.  The exact
conventions and applicable ranges are recorded in
\Cref{sec:affine} and at their later uses.
Established results are cited; the structural tests, compatibility
arguments, weighted readings and effectivity deductions specific to
the constructions are proved below.
\FloatBarrier

\section{A finite colored-graph characterization}\label{sec:graph}

We first give a criterion whose individual witnesses can be checked by finite
combinatorial calculations. Throughout this section, $L$ is a finite nonempty
lattice, with least element $0$.

The graph language follows the valuation and path approach of
Pudl\'ak \citep[Definitions~1--3 and pp.~271--272]{Pudlak1976}; the
unary-closure criterion is recorded in
\citet[Theorem~3.2.1]{DeMeo2012}.  We give the full finite argument,
including its witness-size consequence.  Kenney's graphical-algebra
formulation contains the label-decreasing-map and edge-image-path
reconstruction used here
\citep{Kenney2010}\citep[pp.~51--58]{Kenney2009Slides}.
Our statement isolates a finite test without an auxiliary graphical
composition operation.

\begin{definition}[Admissible placements]\label{f:def-placement}
An $L$-colored complete graph consists of a finite nonempty vertex set $A$ and a
symmetric function $d:A\times A\to L$, with $d(p,p)=0$. Loops are included among
its edges. A placement into another such graph $(A',d')$ is a map $f:A\to A'$;
it is \emph{admissible} if
\[
 d'(f(p),f(q))\le d(p,q)\qquad(p,q\in A).
\]
For a nonempty set $S$ of edges of $(A,d)$, let $\Gamma_S$ be the undirected
graph on $A$ whose marked adjacencies are the pairs
$\{f(p),f(q)\}$, where $\{p,q\}\in S$ and $f:A\to A$ is admissible.
Collapsed pairs are allowed and do not affect connectivity. Write
$a_S=\bigvee_{\{p,q\}\in S}d(p,q)$.
\end{definition}

\begin{theorem}[Graph criterion]\label{f:criterion}
The lattice $L$ is the full congruence lattice of a finite algebra if and only
if there is an $L$-colored complete graph $(A,d)$ satisfying all three
conditions below.
\begin{enumerate}[label=\textup{(\roman*)}]
\item\label{f:graph-triangle} If $p\ne q$, then $d(p,q)\ne0$, and
\[
 d(p,q)\le d(p,u)\vee d(u,q)\qquad(p,q,u\in A).
\]
\item\label{f:graph-reflection} Whenever $a\not\le b$ in $L$, there are vertices
$p,q$ with $d(p,q)\le a$ and $d(p,q)\not\le b$.
\item\label{f:graph-connectivity} For every nonempty edge set $S$, vertices
$p,q$ satisfying $d(p,q)\le a_S$ lie in the same connected component of
$\Gamma_S$.
\end{enumerate}
For each fixed $|A|$, these conditions involve only lattice inequalities,
finite vertex labelings, and paths.
\end{theorem}

\begin{proof}
Suppose first that $\phi:L\to\operatorname{Con}(\mathcal A)$ is an isomorphism,
where $\mathcal A$ has finite carrier $A$. Define
\[
 d(p,q)=\phi^{-1}\bigl(\operatorname{Cg}_{\mathcal A}(p,q)\bigr),
\]
where $\operatorname{Cg}_{\mathcal A}(p,q)$ is the least congruence containing
$(p,q)$. Equality is the bottom congruence, so the diagonal colors are zero
and the other colors are nonzero. The congruence
$\operatorname{Cg}_{\mathcal A}(p,u)\vee
\operatorname{Cg}_{\mathcal A}(u,q)$ contains $(p,q)$ by transitivity, proving
\ref{f:graph-triangle}. If $a\not\le b$, order reflection of $\phi$ gives a
pair in $\phi(a)\setminus\phi(b)$, and that pair proves
\ref{f:graph-reflection}.

A \emph{one-variable translation} of an operation of $\mathcal A$ is obtained
by leaving one argument free and fixing all its other arguments. Every such
translation preserves every congruence, and is therefore an admissible
self-placement. Admissible self-placements are closed under composition.
Consequently a translation sends a marked adjacency of $\Gamma_S$ to another
marked adjacency, possibly a loop. It follows, by applying the translation
to paths, that the equivalence $\rho_S$ of belonging to the same connected
component is stable under all translations. An equivalence with this
property is a congruence: change the entries of an operation's input tuple
one at a time. Nullary operations impose no further condition.

The identity placement is admissible, so $\rho_S$ contains each pair in $S$.
Its corresponding lattice element is thus at least $a_S$. If
$d(p,q)\le a_S$, then $(p,q)\in\rho_S$, which proves
\ref{f:graph-connectivity}.

Conversely, suppose $(A,d)$ satisfies the three conditions. Give $A$ one unary
operation for each admissible self-placement. There are finitely many such
maps, so this defines an allowed finite algebra $\mathcal B$. For $a\in L$ set
\[
 \theta_a=\{(p,q)\in A^2:d(p,q)\le a\}.
\]
The diagonal convention, symmetry, and \ref{f:graph-triangle} show that
$\theta_a$ is an equivalence relation. Admissibility makes it a congruence of
$\mathcal B$. The map $a\mapsto\theta_a$ is order preserving, and
\ref{f:graph-reflection} makes it order reflecting.

Let $\theta$ be any congruence of $\mathcal B$, and take for $S$ all its edges,
including all loops. This is a nonempty set, even if $A$ has one vertex.
Put $a=a_S$. By construction $\theta\subseteq\theta_a$. Each marked adjacency
of $\Gamma_S$ belongs to $\theta$, since every admissible placement is an
operation of $\mathcal B$; hence so does every connected pair. Condition
\ref{f:graph-connectivity} gives $\theta_a\subseteq\theta$. Thus
$\theta=\theta_a$, proving surjectivity and the desired isomorphism. This
argument includes the one-element lattice, represented by the one-vertex
graph and the one-element algebra.
\end{proof}

\begin{corollary}[Enumeration and bounds]\label{f:cor-witness-bound}
Representability is semidecidable. Moreover, a total algorithm deciding
representability of finite lattice order tables exists if and only if there
is a total computable function $b:\mathbb N_{>0}\to\mathbb N_{>0}$ such that
every representable lattice $L$ has a graph witness in \Cref{f:criterion} with
at most $b(|L|)$ vertices.
\end{corollary}
\begin{proof}
For each positive integer $k$, enumerate all colorings on a $k$-element vertex
set and check the three finite conditions. Running through $k=1,2,\ldots$
halts affirmatively exactly on representable lattices. A computable bound
$b$ makes this search finite and decides representability. Conversely, given
a decider, enumerate the finitely many order tables on each fixed
$n$-element set, discard tables that are not lattices, and use the decider to
retain the representable ones. Run the witness search to completion for each
retained table, and take the maximum of the resulting witness sizes and $1$.
This computes the required $b(n)$.
\end{proof}

The later undecidability argument will therefore also exclude every such
computable witness bound.

\section{Representations through seven elements}\label{sec:small}

DeMeo's small-lattice reduction \cite[Theorem~6.1.1 and Chapter~6]{DeMeo2012}
provides finite-algebra representations for every lattice of at most seven
elements, with one possible exception. We now represent that remaining
lattice directly.  The closure constructions used for the other small
lattices include Snow's parallel-sum construction: take a disjoint union
and add a new bottom and top.  A one-element summand is allowed
\citep[Lemma~3.9]{Snow2000}; in particular, the parallel sum of a
four-element chain and a singleton is represented
\citep[Lemma~4.1]{Snow2000}.

\begin{definition}\label{f:def-l7}
Let $L_7$ have elements $0,1,a,b,c,d,x$, with $0$ least, $1$ greatest, and
with the comparabilities among the five interior elements exactly
\[
 a<c,\qquad b<c,\qquad b<d.
\]
In DeMeo's group-interval labeling these elements are
$J_1,J_2,M_2,M_1,K$, respectively.
\end{definition}

\begin{figure}[htbp]
  \centering
  \begin{tikzpicture}[
    x=1cm,y=1cm,
    lattice point/.style={circle,fill,inner sep=1.7pt},
    every label/.style={font=\small},
    line width=0.55pt
  ]
    \node[lattice point,label=below:{$0=H$}] (zero) at (0,0) {};
    \node[lattice point,label=left:{$a$}] (a) at (-2.4,1.5) {};
    \node[lattice point,label=below right:{$b$}] (b) at (-0.4,1.5) {};
    \node[lattice point,label=left:{$c=D$}] (c) at (-1.6,3) {};
    \node[lattice point,label=below right:{$d=\operatorname{PGL}_2(\mathbb F_8)$}]
      (d) at (1.3,3) {};
    \node[lattice point,label=right:{$x=\operatorname{PGL}_2(\mathbb F_4)$}]
      (x) at (5.5,2.15) {};
    \node[lattice point,label=above:{$1=G$}] (one) at (0.6,4.6) {};
    \draw (zero)--(a)--(c)--(one)
          (zero)--(b)--(c)
          (b)--(d)--(one)
          (zero)--(x)--(one);
  \end{tikzpicture}
  \caption{The lattice $L_7$ realized as the subgroup interval $[H,G]$,
    where $G=\operatorname{PGL}_2(\mathbb F_{64})$ and
    $H=\operatorname{PGL}_2(\mathbb F_2)$. The group $D$ is the
    dihedral eigenpair stabilizer of order $126$; its subgroups $a$ and $b$
    containing $H$ have orders $42$ and $18$, respectively. Among the five
    interior elements, the only strict comparabilities are
    $a<c$, $b<c$, and $b<d$.}
  \label{f:fig-l7}
\end{figure}


\begin{proposition}\label{f:prop-l7-interval}
With their usual subfield inclusion, the groups
\[
 H=\operatorname{PGL}_2(\mathbb F_2)
 \le G=\operatorname{PGL}_2(\mathbb F_{64})
\]
have subgroup interval $[H,G]\cong L_7$.
\end{proposition}
\begin{proof}
We first classify the overgroups $H\le K\le G$ by their Sylow
two-subgroups and odd-order eigenpair stabilizers, and then determine
their intersections.
We use the projective-line actions, extending scalars to the algebraic
closure when speaking of eigenlines. For any finite field $Q$, counting
invertible matrices and dividing by scalar matrices gives
\[
 |\operatorname{PGL}_2(Q)|=|Q|(|Q|-1)(|Q|+1).
\]
In particular $|H|=6$ and $H$ induces $S_3$ on its subfield line. It has no
common fixed point: $z\mapsto z+1$ fixes only infinity, and its conjugate by
inversion fixes only zero. Also
\begin{equation}\label{f:eq-centralizer-h}
 C_G(H)=1.
\end{equation}
Indeed, a transformation commuting with these two involutions fixes zero and
infinity individually, so is a scaling; commuting with $z\mapsto z+1$ then
makes that scaling the identity.

Fix an overgroup $H\le K\le G$. In characteristic two, a nonidentity element
of $G$ either has one Jordan block up to scalar, and has order two, or has
two distinct eigenlines, and has odd order. The latter assertion follows by
diagonalizing over a finite extension and taking the ratio of the two
eigenvalues. For each unordered eigenpair, its individual stabilizer in
$K$ is cyclic of odd order, since the diagonal ratio embeds it in the
multiplicative group of a finite field. The nonidentity parts of the
nontrivial such stabilizers partition the nonidentity odd-order elements:
each such element determines its unique eigenpair. Choose representatives
$I_i$ for their $K$-conjugacy classes and put $t_i=|I_i|\ge3$. Their
normalizers stabilize the eigenpairs, so
\[
 |N_K(I_i)|=n_i t_i,\qquad n_i\in\{1,2\}.
\]

The full translation group is a Sylow two-subgroup of $G$, by the order
formula; all Sylow two-subgroups are its conjugates. A Sylow two-subgroup of
$K$ has order $s\ge2$ and hence lies in a translation group, fixing a unique
point of the $\mathbb F_{64}$-line. Distinct Sylow two-subgroups have distinct
fixed points: at any point all involutions fixing it lie in its one
translation group, whose intersection with $K$ is itself a two-group. Let
$\Omega$ be the set of these fixed points. It is a $K$-orbit of size $h>1$,
since otherwise $H$ would fix its unique point.

At a point of $\Omega$, take coordinates putting that point at infinity.
The stabilizer $K_\infty$ is affine. Its translation kernel has order $s$;
its multiplier image has odd order $j$ and preserves the additive subgroup
of translation parameters. Nonidentity multipliers act freely on nonzero
parameters, so $j\mid s-1$. Therefore $|K|=hsj$, and the number of involutions
is $h(s-1)$. Counting the identity, the involutions, and the odd-order
elements yields
\begin{equation}\label{f:eq-overgroup-count}
 \sum_i\frac1{n_i}\left(1-\frac1{t_i}\right)
 =1-\frac{s-1}{sj}-\frac1{hsj}
 =1-\frac1j+\frac{h-1}{hsj}.
\end{equation}
The sum is nonempty because $H$ contains an element of order three.

\emph{The case $j=1$.}
The right side of \Cref{f:eq-overgroup-count} is less than $1/s$; each
summand is at least $1/3$. Thus $s=2$, and there is exactly one summand. It
must have $n_i=2$, because a summand with $n_i=1$ is at least $2/3$.
Substitution gives $t_i=h$, so
$|N_K(I_i)|=2h=|K|$. Hence $K$ lies in the setwise stabilizer of one eigenpair.
The order-three subgroup of $H$ must fix the pair individually, since it
cannot act nontrivially on two points. This eigenpair is therefore the one
of that order-three subgroup.

\emph{The case $j>1$.}
Now $s\ge4$. The number $j$ is one of the $t_i$: a generator of the multiplier
image lifts to an affine map with nonidentity multiplier, two fixed points,
and order $j$. The individual stabilizer of its two fixed points embeds in
the multiplier image and contains that cyclic subgroup, so has order $j$.
If the sum in \Cref{f:eq-overgroup-count} had only this term, it would be at
most $1-1/j$, contrary to the equation. The whole sum is less than $1$.
Since each summand is at least $1/3$, there are exactly two summands, and
both have $n_i=2$; otherwise their sum would be at least $2/3+1/3$.
Writing $t$ for the other order gives
\begin{equation}\label{f:eq-overgroup-parameters}
 1-\frac jt=\frac{2(h-1)}{hs}.
\end{equation}
In particular $t>j$, and $s(t-j)<2t$ implies
\[
 0<t-j<\frac{2j}{s-2}\le\frac{2(s-1)}{s-2}\le3.
\]
The two orders are odd, so $t=j+2$. Equation
\eqref{f:eq-overgroup-parameters} now gives $h(j+2-s)=j+2$.
Since $h>0$, we have $j>s-2$, and $j\le s-1$ therefore forces
$j=s-1$ and $h=s+1$.

The multiplier realization gives $s-1\mid63$. As $s$ is a power of two with
$4\le s\le64$, it follows that $s\in\{4,8,64\}$. The multiplier group is
then exactly $\mathbb F_s^\times$ inside $\mathbb F_{64}^\times$. The
translation-parameter subgroup, of size $s$, is stable under these scalars,
so is a one-dimensional $\mathbb F_s$-space, say $v\mathbb F_s$. Translations
act freely on the other $s$ points of $\Omega$, and hence those points form
one coset of $v\mathbb F_s$. An affine coordinate change thus makes
$\Omega=\mathbb P^1(\mathbb F_s)$.

Its full setwise stabilizer in $G$ is $\operatorname{PGL}_2(\mathbb F_s)$:
the images of infinity, zero, and one determine a projective transformation,
and every ordered triple of distinct subfield points is the image of that
triple under a subfield transformation. The equality
$|K|=(s+1)s(s-1)$ now gives
$K=g\operatorname{PGL}_2(\mathbb F_s)g^{-1}$ for some $g\in G$.
This conjugate is in fact the standard subfield group. Let $\sigma$ be the
entrywise $s$-power automorphism. It fixes both $H$ and $g^{-1}Hg$ pointwise.
Consequently $g\sigma(g^{-1})$ centralizes $H$, so
\Cref{f:eq-centralizer-h} gives $\sigma(g)=g$ in the projective group.
Normalize one nonzero matrix entry of a representative of $g$ to $1$;
all its entries are then fixed by $\sigma$. Thus
$g\in\operatorname{PGL}_2(\mathbb F_s)$, as required.

It remains to record the incidences between the groups in this list. The
eigenpair of the order-three subgroup of $H$ splits over $\mathbb F_{64}$.
Its full setwise stabilizer $D$ is dihedral of order $2\cdot63$ and contains
$H$, generated there by a rotation of order three and a flip. Subgroups
between $H$ and $D$ correspond to rotation subgroups containing the
order-three one, after adjoining that same flip. Their orders are
$6,18,42,126$, so there are precisely two strict intermediate groups, of
orders $18$ and $42$.

The other proper overgroups are the standard subfield groups over
$\mathbb F_8$ and $\mathbb F_4$, whose intersection is $H$: normalize one nonzero matrix entry to $1$
and use $\mathbb F_8\cap\mathbb F_4=\mathbb F_2$. Over
$\mathbb F_4$ the eigenpair is split, so its setwise stabilizer has order
$2(4-1)=6$; consequently
$D\cap\operatorname{PGL}_2(\mathbb F_4)=H$. Over $\mathbb F_8$ the pair is
nonsplit, because $3\nmid7$. Its individual stabilizer has order $8+1=9$.
To see the latter count directly, matrices fixing the two eigenlines
centralize a matrix with irreducible quadratic minimal polynomial. Its
invertible centralizer is the multiplicative group of the quadratic field
extension, and quotienting by subfield scalars gives order
$(8^2-1)/(8-1)=9$. The flip from $H$ doubles this order. Hence
$D\cap\operatorname{PGL}_2(\mathbb F_8)$ is exactly the order-$18$ group.
We obtain the asserted interval by taking
\[
 c=D,\quad d=\operatorname{PGL}_2(\mathbb F_8),\quad
 x=\operatorname{PGL}_2(\mathbb F_4),
\]
and taking $a,b$ to be the groups of orders $42,18$, respectively.
\end{proof}

\begin{lemma}[Coset-action correspondence]\label{f:lem-coset}
For finite groups $H\le G$, the full congruence lattice of the unary algebra
given by the left $G$-action on $G/H$ is isomorphic to $[H,G]$.
\end{lemma}
\begin{proof}
For a $G$-invariant equivalence $\theta$ on $G/H$, put
$K_\theta=\{g\in G:gH\mathrel\theta H\}$. This contains $H$. If
$g,g'\in K_\theta$, invariance applied to $gH\mathrel\theta g'H$ gives
$H\mathrel\theta g^{-1}g'H$, so $K_\theta$ is a subgroup. Conversely,
$H\le K\le G$ gives the invariant equivalence obtained by pulling equality
back along $G/H\to G/K$. These constructions are inverse and preserve order.
\end{proof}

\begin{theorem}\label{f:small-lattices}
Every finite nonempty lattice of at most seven elements is representable.
\end{theorem}
\begin{proof}
\Cref{f:prop-l7-interval,f:lem-coset} represent $L_7$ by a finite unary
algebra. DeMeo's small-lattice reduction supplies the other lattices
\cite[Theorem~6.1.1]{DeMeo2012}.
\end{proof}

\section{Affine test lattices and almost simple intervals}\label{sec:affine}

The negative examples will come from two opposite orientations of a finite
tree of truncated affine geometries. The common lattice is simple and
self-dual. Its reverse links have several horizontal branches, and their
bottom diamonds exclude particular normal-subgroup configurations. We first
construct these lattices and reduce a hypothetical representation to an
ordinary almost simple subgroup interval. Subsequent sections prune paths
through alternating and classical groups; bounded Lie rank gives a separate
contradiction, and additional path protections bound the number of remaining
stages.

\begin{theorem}\label{f:negative-family}
There are choices of the finite parameters in \Cref{f:def-test-lattice} for
which the resulting finite lattice is not the congruence lattice of any
finite algebra.
\end{theorem}

The construction and reductions in this section are uniform under the
inequalities stated below. The successive choices of larger parameters and
the completion of the proof of \Cref{f:negative-family} are given in
\Cref{sec:patterns,sec:classical,sec:alternating,sec:bounded,sec:weights}.
This theorem gives negative examples and, after bounding the parameters, a
cardinality ceiling. It does not assert that the chosen family member has
minimum cardinality; the minimum problem is treated separately.

\subsection{Structural inputs and conventions}

We use the classification of finite simple groups and its standard
structural consequences, including Schreier's Theorem, the automorphisms
and central covers of finite groups of Lie type, and the classification of
finite two-transitive groups. The algebraic-group inputs include
Lang--Steinberg, highest-weight theory, Steinberg's theorems, and the
algebraic Peter--Weyl filtration. The bounded-rank argument will use the
sufficiently-general subgroup Theorem of Larsen--Pink
\cite{LarsenPink2011}. To specify the scope of these inputs, we fix the
following conventions.

\begin{convention}[Lie-theoretic realizations]\label{f:conv-lie-inputs}
We use split root-system realizations over algebraic closures, with both
simply connected and adjoint root data, the Chevalley root and parabolic
structure, and the isogeny and Steinberg-map classification. Field maps,
pinned diagram maps, and, in the relevant special characteristics, the
special isogenies of bounded types are combined with inner automorphisms.
A Steinberg map of positive height supplies the finite fixed group.
For a simple adjoint group at sufficiently large parameter, or in the
high-rank cases used here, the associated simple finite group is the derived
fixed group. Outside the finitely many exceptional ranks and parameters for
simplicity and Schur multipliers, the simply connected fixed group is its
perfect universal central cover. Inner-diagonal, field, graph, and applicable
special actions induce the usual automorphisms of the simple projective
group. No uniform covering assertion is made at the excluded small cases.
For high-rank natural signed types, type $B$ is used only in odd
characteristic; in characteristic two we use the symplectic realization of
type $C$. These are the usual structure and automorphism conventions
\cite{GLS1994,Steinberg1968}.
\end{convention}

\begin{convention}[Representation-theoretic inputs]\label{f:conv-representation-inputs}
Over the algebraic closure in defining characteristic, we use rational
highest-weight theory: Weyl groups, Weyl and induced modules, weight convex
hulls, and the triangular and distribution-algebra properties of highest
weights. We also use the restricted-simple theorem for simply connected
groups and their Frobenius kernels, Steinberg's Tensor Product Theorem and
finite-group restriction and parametrization theorems on the simply
connected finite cover, and the two-sided algebraic Peter--Weyl good
filtration of the regular function algebra
\cite{Steinberg1963,Jantzen2003,Humphreys2006}.
The precise finite $P$-restricted convention at high rank is given in
\Cref{sec:weights}. These inputs concern rational modules or the specified
full finite covers; they do not assert that arbitrary almost simple subgroup
inclusions arise from algebraic-group inclusions.
\end{convention}

On the permutation side we use Jordan's prime-cycle Theorem: a primitive
group of degree $d$ containing a cycle of prime length at most $d-3$
contains $A_d$. We use Bochert's primitive-index bound, in particular an
exponential lower bound on the index in $S_d$ of a primitive group not
containing $A_d$ in sufficiently large degree, and Burnside's prime-degree
Theorem. The finite two- and multiply-transitive classification gives the
following consequences: a sufficiently large alternating simple group has
only its natural degree as a nontrivial two-transitive action; in high
degree a six-transitive group contains the alternating group; and a
non-affine two-transitive group of degree $37$ contains $A_{37}$
\cite{DixonMortimer1996}.
For the last assertion, the projective degrees $1+s+\cdots+s^j$, with $s$ a
prime power, never equal $37$: for $j=1$, one would need $s=36$; for
$j=2,3,4$ consecutive possible integer values bracket $37$ without equality;
and for $j\ge5$ even $s=2$ gives at least $63$. The unitary, Suzuki, Ree, and
symplectic-over-two degrees $s^3+1$, $s^2+1$, and
$2^{e-1}(2^e\pm1)$ also give no permitted parameter of degree $37$, and the
additional small and sporadic actions on the classification list do not
have that degree. The simple-group identifications in the classification
exclude spurious high-degree alternating overlaps.
We use ordinary finite subset Ramsey theory, finite multidimensional
Hales--Jewett for a fixed alphabet, number of colors, and number of free
slots, and later the finite vector-space Ramsey Theorem for linear
subspaces \cite{Ramsey1930,HalesJewett1963,GrahamLeebRothschild1972,GrahamLeebRothschild1973}.
More specialized uses of these inputs are stated where they occur.

\subsection{Charts, branches, and prepared links}

\begin{definition}[Affine test lattice]\label{f:def-test-lattice}
Let $q\ge5$ be a prime power such that $q-1$ is not a prime power, and let
$r,D,g,m,n,N$ be positive integers satisfying
\begin{equation}\label{f:eq-parameters}
 D\ge2,\qquad g\ge2,\qquad 3\le m<n-2,\qquad Dr<m,\qquad N\ge n+2.
\end{equation}
There are arbitrarily large allowable $q$, for example powers of $7$.
A \emph{chart} is a copy of $\operatorname{AG}(N,q)$. The rank of a nonempty
affine flat is its affine dimension plus one; the empty flat has rank zero.
To attach a fresh chart at a specified flat $W$, identify a flat of the new
chart affinely with $W$, and take all its other points fresh. Its intersection
with the union of the old point sets is exactly $W$. Attachment at the empty
flat makes no identifications.

Build the trunk charts recursively, beginning with the path ending at
$W=\varnothing$.
\begin{enumerate}[label=\textup{(\roman*)}]
\item At a path ending at $W$, attach $g$ fresh \emph{private charts} at $W$.
\item In each such chart, reserve a child path extending this path by $U$
for every flat $U\supset W$ with rank increment between $1$ and $r$.
Perform stages for $D$ generations: the empty path has a stage, as does
every child path reached in fewer than $D$ steps.
\item Finally, at every nonempty flat of rank at most $m$ in a trunk chart,
attach at least one additional fresh \emph{leaf chart}, with no stages of
its own.
\end{enumerate}
Repeated reservations may use distinct fresh charts. There are only finitely
many charts. The primal interior labels are the point sets that are
nonempty affine flats of dimension at most $n$ in some chart, ordered by
inclusion. Adjoin the empty bottom and one formal top above all these
labels, obtaining a bounded poset $L^+$. Let $L^-$ be an order-reversed copy
of $L^+$. The \emph{test lattice} $L$ is their horizontal sum: identify the
two bottoms and the two tops, with no other cross-comparabilities.
The interior labels on $L^-$ are called \emph{reverse flats}, named by their
affine point sets; their lattice order reverses inclusion. On this side the
top may accordingly be called the empty flat.
\end{definition}

A comparability component below means a connected component of the
undirected comparability graph on the strict interior of a bounded lattice.
Distinct such components are horizontal: two elements from different
components have meet zero and join one.

\begin{proposition}[Geometry of the test lattice]\label{f:prop-chart-properties}
The construction in \Cref{f:def-test-lattice} gives a finite self-dual
lattice $L$. It is graded of length $n+2$, atomistic, coatomistic, and simple.
Any two charts intersect in a flat of each, of rank at most $m$, with the
same affine structure on the intersection. Within one chart the primal
join of two flats is their affine span if that span has dimension at most
$n$. Attachment preserves all old meets and joins. In particular, flats in
two distinct private charts attached at the same $W$, neither contained in
$W$, have no common proper primal upper bound. There are at least three
horizontal comparability components in the interior of $L$, each containing
a chain of distinct comparable interior elements.
\end{proposition}
\begin{proof}
Induct on the attachments. When a chart is attached at $W$ in an old chart,
its intersection with any other old chart is exactly that chart's
intersection with $W$. By the inductive hypothesis this is a flat in the
old chart containing $W$ and in the other old chart, with agreeing affine
structures. It is therefore a flat of the new chart as well. This proves
the pairwise-intersection assertion. Trunk attachment flats have rank less
than $Dr<m$, and leaf attachment flats have rank at most $m$, so every
intersection of distinct charts has rank at most $m$. Affine subflats of a
shared flat are consequently identified consistently.

It follows that intersections of labels are again flats, allowing the
empty set. Moreover, whenever one nonempty label is contained in another,
it is an affine subflat in every chart containing the larger label. Thus
$L^+$ has finite meets and a top, which suffice to give all joins. A new
proper upper bound of two old labels, not itself already an old flat, can
contain those old labels only if both lie in the attachment flat $W$; it
then contains their old affine span in $W$. Hence no old join changes.
Meets are set intersections and do not change either. This also proves
the claim about private charts, first at their attachment and then after
all subsequent attachments. Within a chart the stated affine span is the
least proper upper bound. The horizontal sum with the opposite lattice
is a self-dual lattice.

Every proper chain in one orientation can be refined one affine dimension
at a time, up to dimension $n$, before reaching the formal top. Thus $L$
is graded of length $n+2$. A flat is generated by its points, so $L^+$ is
atomistic; intersecting its $n$-dimensional extensions inside one chart
recovers the flat, so it is coatomistic. The same properties hold in the
opposite orientation and hence in the horizontal sum.

The $g\ge2$ initial charts are disjoint. An attachment meets all old points
in a flat contained in one old chart; it cannot connect two old initial
branches by comparabilities. Each orientation therefore has at least two
horizontal interior components, and each component contains long affine
chains. In particular $L$ has at least three such components.

Finally, let a nonidentity lattice congruence identify two different
comparable elements $u<v$; such a pair exists whenever any two distinct
elements are identified, by meeting and joining them. By atomisticity there
is an atom $a\le v$ with $a\not\le u$. Meeting the identification with $a$
gives $a\equiv0$. Choose an interior element $b$ in another component.
Joining with $b$ gives $1\equiv b$. Meeting this with an interior element $c$
in a third component gives $c\equiv0$. Since $a\vee c=1$, we conclude
$1\equiv0$. The congruence is total, proving simplicity.
\end{proof}

\begin{proposition}[Prepared reverse links]\label{f:prop-prepared-links}
If $X$ is a nonempty flat of rank at most $m$ in a trunk chart, the reverse
ideal with top $X$, denoted $[0,X]$, is graded, atomistic, and coatomistic.
Its strict interior has at least two horizontal comparability components,
each containing comparable distinct elements. The following additional
properties hold for $L$ and for every such prepared link.
\begin{enumerate}[label=\textup{(\roman*)}]
\item\label{f:link-diamonds} Every bottom interval of length two is a diamond
$M_b$ with $b\ge3$ atoms, where $b-1\ge2$ is not a prime power.
\item\label{f:link-upper-diamonds} In every interior component there is an
element $y$ two ranks below the interval's top whose upper interval is a
diamond with at least two middle elements.
\end{enumerate}
In addition, $L$ has the following separation property:
\begin{enumerate}[label=\textup{(\roman*)},resume]
\item\label{f:link-separation} For any atom $u$ and any proper upper bound
$\omega\ge u$, there is a rank-two element $z\ge u$ with
$z\not\le\omega$.
\end{enumerate}
\end{proposition}
\begin{proof}
A leaf attached at $X$ supplies a branch of strict superflats of $X$
incomparable with the previous strict superflats. Later attachments cannot
connect two old branches in this reverse link: any connection through a
new chart passes through superflats of $X$ lying in its attachment flat;
if any exist, they are already connected within that flat. Every
superflat of $X$ lies in a chart containing $X$, and in that chart it has
extensions and intermediate flats. Each branch therefore contains
comparable distinct elements and chains out to $n$-flats. Gradedness follows
from these dimension chains. The flats extending $X$ by one dimension, and
the intersections of $n$-extensions, establish atomisticity and
coatomisticity in the reversed order.

On the primal side of $L$, a bottom interval of length two is a line, with
$b=q$ points; the choice of $q$ gives the claimed condition on $b-1$.
On the reverse side, including a prepared link, its top flat has dimension
$n-1$. Such a flat belongs to exactly one chart, since distinct charts meet
in rank at most $m<n-2$. Its $n$-dimensional extensions are counted by
one-dimensional subspaces of a vector space of dimension $N-n+1$, so
\[
 b=\frac{q^{N-n+1}-1}{q-1},\qquad
 b-1=q(1+q+\cdots+q^{N-n-1}).
\]
The second factor is greater than one because $N\ge n+2$, and is congruent
to $1$ modulo the characteristic prime. Thus $b-1$ has at least two prime
divisors. This proves \ref{f:link-diamonds}.

For \ref{f:link-upper-diamonds}, choose a primal $(n-1)$-flat, or, on the
reverse side, a flat whose rank exceeds that of the empty flat or the
prepared $X$ by two. The appropriate upper interval is a diamond, with at
least two middle flats. Such choices can be made in every branch.

For \ref{f:link-separation}, a primal atom is a point. In one chart through
it there is a line escaping the given proper flat $\omega$. A reverse atom
is an $n$-flat, belonging to a unique chart; its proper upper bound
corresponds to a nonempty subflat. Choose a hyperplane of the $n$-flat that
does not contain that subflat. It is the required rank-two reverse label.
\end{proof}

\subsection{From algebras to ordinary group intervals}

We write $[H,G]$ for the ordinary interval of all subgroups between finite
groups $H\le G$. The adjective \emph{ordinary} distinguishes such intervals
from invariant-subgroup intervals and from the extension posets introduced
below.

\begin{lemma}[Subdirect products of simple groups]\label{f:lem-subdirect}
A subdirect subgroup of a finite power of a nonabelian simple group is a
product of full, possibly automorphism-twisted, diagonal subgroups on
disjoint parts partitioning the coordinates. Containment between two such
subgroups means refinement of the tied parts, with agreement of the
diagonal identifications on each refined part. If a subgroup contains the
untwisted full diagonal, all its diagonal identifications are untwisted.
\end{lemma}
\begin{proof}
Induct on the number of coordinates. Project away the last coordinate.
The kernel in the last simple factor is normal in that factor, by
surjectivity, so is either the whole factor or trivial. In the first case
the last factor splits off. In the second, the subgroup is the graph of a
surjective homomorphism from the shorter subdirect product to the last
simple group. By induction the shorter group is a product of nonabelian
simple diagonal factors. The images of those factors are commuting normal
subgroups of a simple group; exactly one has nontrivial image, and on that
factor the map is an isomorphism. This ties the last coordinate to just
that part. The description of containment follows by projecting to pairs
of coordinates: a required tie must hold in the smaller subgroup with the
same identifying automorphism. Containing the untwisted full diagonal
forces each such automorphism to be the identity.
\end{proof}

\begin{proposition}\label{f:prop-ordinary-interval}
If the test lattice $L$ is representable by a finite algebra, it is
isomorphic to an ordinary subgroup interval of a finite group.
\end{proposition}
\begin{proof}
We begin with the least-carrier restriction argument of
\citet[Lemma~1]{PalfyPudlak1980}. The final diagonal construction below
will turn the resulting permutation action into an ordinary interval.
Choose a representing algebra with smallest possible carrier $A$. Let $F$
be the monoid of all maps $A\to A$ preserving its represented congruences.
The unary algebra $(A;F)$ has exactly those congruences: each represented
congruence is preserved by definition, while $F$ contains all one-variable
translations of the original operations, whose preservation forces an
equivalence to be an original congruence.

Let $e\in F$ be idempotent. On $eA$ use the unary operations
$ef|_{eA}$, $f\in F$. Restriction induces a surjective lattice homomorphism
\begin{equation}\label{f:eq-retraction}
 \operatorname{Con}(A;F)\longrightarrow
 \operatorname{Con}(eA;\{ef|_{eA}:f\in F\}).
\end{equation}
Meets are preserved by restriction. To check joins, retract an equivalence
path by $e$: its vertices then lie in $eA$ and each step remains in the
same congruence. For surjectivity, if $\delta$ is a congruence on the displayed
algebra on $eA$, define
\[
 x\mathrel{\widehat\delta}x'
 \quad\Longleftrightarrow\quad
 (efx,efx')\in\delta\ \text{for every }f\in F.
\]
This is an $F$-invariant equivalence. Its restriction to $eA$ is exactly
$\delta$: one direction uses its compatibility with all $ef|_{eA}$, and
the other uses $f=1$ and $e|_{eA}=1$.

If $|eA|>1$, the target of \eqref{f:eq-retraction} is nontrivial. Simplicity
of $L$ makes the homomorphism injective, and minimality of $|A|$ forces
equality of carrier sizes. A surjective idempotent is the identity.
Thus every nonunit idempotent of $F$ is constant.

Let $I$ be the ideal of nonunits of $F$. Every sufficiently long product
of elements of $I$ is constant. Indeed, otherwise take more than $|F|$
successive prefixes in order of application of a product whose final
image has size greater than one. Every prefix image then has size greater
than one, and two prefix maps coincide. The intervening nonempty product
fixes their common image pointwise. An idempotent power of that product
is a nonunit idempotent fixing at least two points, a contradiction.

For $\theta\in\operatorname{Con}(A;F)$ let $\Phi(\theta)$ be the equivalence
generated by
\[
 \{(f(x),f(y)):f\in I,\ (x,y)\in\theta\}.
\]
This is a congruence, by the ideal property of $I$, and
$\Phi(\theta)\le\theta$. It preserves joins: both equivalence generation and
images of paths show that applying it to a join gives the join of the
images. Iterating $\Phi$ generates precisely the pairs obtained from
products of the corresponding length in $I$. The inserted equivalence
closures contribute nothing extra, since an image of a path is a path;
diagonal pairs accommodate collapsed steps. The product observation makes
some iterate of $\Phi$ identically zero. For an atom $\alpha$, the deflation
$\Phi(\alpha)\le\alpha$ must therefore be zero, since the only other
possibility would fix $\alpha$ forever. Atomisticity now gives
$\Phi(1)=0$. Hence every map in $I$ sends every pair to a diagonal pair:
all nonunits are constant.

The units consequently form a permutation group $J$ with exactly $L$ as
its lattice of invariant equivalences; constants impose no conditions.
Choose any finite nonabelian simple group $C$, and form the coordinate
permutation semidirect product $C^A\rtimes J$. Then
\begin{equation}\label{f:eq-diagonal-interval}
 [\Delta(C)\rtimes J,\ C^A\rtimes J]
\end{equation}
is dual to that invariant-equivalence lattice. Indeed, intersection with
$C^A$ identifies this interval with the $J$-invariant subgroups above the
full diagonal. By \Cref{f:lem-subdirect} these are the diagonal products
determined by partitions of $A$, and subgroup containment reverses
refinement of the corresponding equivalences. Since $L$ is self-dual,
\eqref{f:eq-diagonal-interval} represents $L$ itself.
\end{proof}

\subsection{The monolith and extension description}

For a subgroup $H\le G$, its core is
$\operatorname{core}_G(H)=\bigcap_{g\in G}gHg^{-1}$. A group with a unique
minimal nontrivial normal subgroup is called \emph{monolithic}; that
minimal normal subgroup is its \emph{monolith}.
The normal-supplement and unique-monolith argument follows
\citet[Lemma~3.6]{PalfyTwisted}. The horizontal components of our prepared
intervals give the required lattice obstruction directly.

\begin{proposition}[Normal supplements]\label{f:prop-monolith}
Suppose $[H,G]$ is isomorphic either to $L$ or to a prepared reverse link.
After quotienting by $\operatorname{core}_G(H)$, every nontrivial normal
subgroup is supplemented by $H$. The quotient has a unique minimal normal
subgroup $M=S^s$, with $S$ nonabelian simple, and $H$ acts transitively on
its simple factors. Moreover $C_G(M)=1$ in the core quotient. In that
quotient the subgroup interval identifies with
\begin{equation}\label{f:eq-invariant-interval}
 [P,S^s]^H,\qquad P=H\cap S^s,
\end{equation}
the $H$-invariant subgroups in this range, by intersection with $S^s$;
the inverse sends $Y$ to $HY$.
\end{proposition}
\begin{proof}
If $R\trianglelefteq G$, then $HR$ permutes as a set with every overgroup of
$H$. Dedekind's law consequently makes $HR$ left-modular in $[H,G]$.
There is no proper interior left-modular label: for a putative such label
$u$, choose comparable distinct interior labels $v<w$ in another horizontal
component. Then
\[
 (v\vee u)\wedge w=w\ne v=v\vee(u\wedge w).
\]
Thus either $R\le H$ or $HR=G$. In the core quotient the first alternative
forces $R=1$.

An abelian minimal normal supplement would identify the interval with
an invariant-subgroup interval in that abelian group. Such an interval is
modular, contradicting the displayed failure of the modular identity.
If two distinct minimal normal subgroups existed, they would commute.
Using one to supplement $H$, the conjugation action of $H$ on the other
would induce all its inner actions. The interval would again identify
with a modular interval, now of invariant normal subgroups of that other
minimal normal subgroup. Both possibilities are impossible. The unique
minimal normal subgroup is therefore a power $S^s$ of a nonabelian simple
group. Its factor orbits under $H$ are normal in $HM=G$, so minimality makes
the action transitive. Its centralizer is normal; if nontrivial it would
contain $M$, which is impossible because $M$ is centerless. Finally $HM=G$,
and the elementary intersection--product correspondence gives
\eqref{f:eq-invariant-interval}.
\end{proof}

For a subgroup $Y\le S^s$, write $c(Y)$ for the product of its coordinate
projections. It is $H$-invariant whenever $Y$ is. We will need a precise
parametrization of the invariant subdirect groups. We use the
invariant-subdirect/extension correspondence in
\citet[Proposition~2.2 and Lemma~2.3]{PalfyTwisted}, expressed here for
the transitive action on the actual simple coordinates.

\begin{lemma}[Extension description]\label{f:lem-extension-description}
Suppose $H$ acts by automorphisms on $S^s$, transitively on the simple
coordinates. Fix a coordinate $o$, identify its group with $S$, and let
$A_0\le H$ be its stabilizer, with induced action
$\alpha:A_0\to\operatorname{Aut}(S)$. The $H$-invariant subdirect subgroups
of $S^s$ are in bijection with extensions
\begin{equation}\label{f:eq-extension-data}
 A_0\le T\le H,\qquad
 \beta:T\longrightarrow\operatorname{Aut}(S),\qquad
 \beta|_{A_0}=\alpha.
\end{equation}
This bijection reverses extension order: increasing the extension means
increasing its domain and agreeing on the smaller domain. If one restricts
to subdirect groups containing a specified $P$, then above any retained
group in actual subgroup order all restrictions of its extension to
intermediate domains still occur.
\end{lemma}
\begin{proof}
By \Cref{f:lem-subdirect}, a subdirect group has tied diagonal parts that
form an $H$-invariant partition of the coordinates. The part containing $o$
is a block for the transitive action; it equals $To$, where $T$ is its
setwise stabilizer and $A_0\le T$. The action of $T$ on that diagonal,
identified with $S$ by projection to $o$, gives $\beta$ extending $\alpha$.

Conversely, given \eqref{f:eq-extension-data}, define the diagonal on $To$
using left-action notation. At coordinate $ho$, for $h\in T$, send $u\in S$
to the transport by $h$ of $\beta(h)^{-1}(u)$ from coordinate $o$. Replacing
$h$ by $ha$, with $a\in A_0$, gives the same element because
$\beta(a)=\alpha(a)$. This defines a $T$-invariant full diagonal. Transport
it to the other blocks and multiply the resulting disjoint diagonal
factors. The result is $H$-invariant and subdirect, and recovers the original
data. Refinement of tied parts, with agreement of identifications, is
exactly restriction of the extension. The last assertion follows because
such restrictions give larger actual subgroups, still containing $P$.
\end{proof}

\begin{proposition}[Cases for a prepared interval]\label{f:prop-interval-cases}
Use the notation of \Cref{f:prop-monolith}, in the core quotient. Let
$\Lambda\le\operatorname{Aut}(S)$ be the image of a coordinate stabilizer
in $H$, and put $I=\Lambda\cap S$, identifying $S$ with its inner
automorphisms. The following are the only possibilities.
\begin{enumerate}[label=\textup{(\roman*)}]
\item\label{f:case-subdirect} $c(P)=S^s$. All labels are subdirect. If the
interval represents the whole $L$, its opposite has an ordinary subgroup
representation with top-group order at most $|H|<|G|$.
\item\label{f:case-product} $P=c(P)\ne1$. The interval is isomorphic to the
ordinary almost simple interval $[\Lambda,S\Lambda]$.
\item\label{f:case-extension} $P=1$ and $I=S$. Every nonbottom label is
subdirect, and the dual interval is the extension poset of
\Cref{f:lem-extension-description} with one artificial top adjoined.
\end{enumerate}
In particular the cases $P<c(P)<S^s$, $P=1$ with $I=1$, and $P=1$ with
$1<I<S$ are impossible.
\end{proposition}
\begin{proof}
We check all the cases, retaining the prepared-interval hypotheses.

\emph{Subdirect bottom.}
If $c(P)=S^s$, every label is subdirect and refines the fixed diagonal ties
of $P$. Restrict to one tied part of $P$. Its setwise stabilizer in $H$
acts transitively on its coordinates. A refinement of its coordinates
that is invariant under this stabilizer extends uniquely to the other
parts by $H$, keeping the twists fixed by $P$. Thus the interval is dual
to the invariant-equivalence lattice of that transitive action. By
\Cref{f:lem-coset}, the latter is an ordinary interval above a point
stabilizer in the part stabilizer. Its top has order at most $|H|<|G|$,
proving \ref{f:case-subdirect}.

\emph{A proper nonproduct bottom.}
Suppose $P<c(P)<S^s$. Choose a horizontal interior component different from
that of $c(P)$. Every label $Y$ in it has
$Y\vee c(P)=S^s$ and $Y\cap c(P)=P$. Its coordinate projections already
contain those of $P$, so adjoining $c(P)$ does not enlarge them. The join
condition therefore makes $Y$ subdirect. In its diagonal description,
intersection with $c(P)$ restricts each coordinate to the corresponding
nontrivial projection of $P$, with the same ties. These restrictions
project onto every such subgroup, because the intersection contains $P$.
There are comparable distinct labels in the chosen component. Their
strict containment changes at least one diagonal tie, and releasing a tie
strictly enlarges the intersection: its tied coordinate subgroup is
nontrivial and may now be chosen independently on the separated parts.
This contradicts the fact that both intersections equal $P$.

\emph{A nontrivial product bottom.}
Suppose $P=c(P)\ne1$, necessarily a proper product. A subdirect group
containing independent nontrivial subgroups in each coordinate cannot
have any nontrivial ties. Thus the only subdirect label above $P$ is
$S^s$. For a coatom $Y$, the product $c(Y)$ is a label containing $Y$;
maximality and the preceding observation imply $Y=c(Y)$. Since the
interval is coatomistic, every label is a product.
At one coordinate let $Q$ be the projection of $P$. Then
$1<Q\le I$: the coordinate elements of $P\le H$ induce their inner
automorphisms. Products identify the interval with $[Q,S]^\Lambda$.
For any coatom $T$ of this invariant interval, $I$ normalizes $T$.
If $TI=S$, then $T$ is normal in $TI$, hence in $S$, impossible because
$1<Q\le T<S$. Thus $TI<S$, and maximality gives $I\le T$.
Intersecting all coatoms yields $I\le Q$, so $I=Q$.
Intersection with $S$ and multiplication by $\Lambda$ now identify
$[Q,S]^\Lambda$ with $[\Lambda,S\Lambda]$ inside
$\operatorname{Aut}(S)$, proving \ref{f:case-product}.

\emph{Trivial bottom and no inner operators.}
Suppose $P=1$ and $I=1$. Schreier's Theorem makes $\Lambda$ soluble, since
it embeds in $\operatorname{Out}(S)$. We first show that a soluble operator
group on a finite nonabelian simple group has a nontrivial proper invariant
subgroup. We may replace it by its effective image. If that image is
trivial, choose any subgroup of prime order. Otherwise choose a minimal
normal elementary abelian $p$-subgroup $A$ of the operator group. Its fixed
group is invariant, and if nontrivial it is proper by effectiveness.
If its fixed group is trivial, orbit counting on the simple group gives
$p\nmid|S|$. For each prime divisor $\ell$ of $|S|$, the number of Sylow
$\ell$-subgroups divides $|S|$ and is prime to $p$, so the $p$-group $A$
fixes at least one such Sylow subgroup. There is at most one: the set of
conjugators between two $A$-invariant Sylow subgroups is $A$-invariant and
has size $|N_S(P_\ell)|$, prime to $p$, so has an $A$-fixed conjugator.
The trivial fixed group forces that conjugator to be $1$. Uniqueness
makes the Sylow subgroup invariant under the whole operator group.
This proves the assertion.

Choose a minimal nontrivial $\Lambda$-invariant subgroup $Q$ of $S$.
Its characteristic subgroups are invariant, so $Q$ is characteristically
simple. If it were a power of a nonabelian simple group, apply the
preceding assertion to the soluble operator image of a factor stabilizer,
and transport the resulting proper invariant subgroup through the
transitive factor action. This would give a nontrivial proper invariant
subgroup of $Q$, a contradiction. Therefore $Q$ is elementary abelian.
Its normalizer is proper because $S$ is simple. It strictly contains $Q$:
if $Q$ is not Sylow, use the normalizer property in a Sylow subgroup
containing it; if it is Sylow and were self-normalizing, Burnside transfer
would give a normal prime complement, again contradicting simplicity.

Transport $Q$ and $N_S(Q)$ to all coordinates and take their products.
They are labels, the latter strictly above the former. If the product of
the $Q$'s has lattice rank at least two, choose a rank-two label below it
and any atom in that interval. Otherwise the product itself is an atom;
choose a rank-two label above it and below the normalizer product. In both
cases the chosen atom $A$ is an elementary abelian normal subgroup of the
rank-two group $R$. Every other atom is an $H$-invariant complement to
$A$ in $R$, and conversely every such complement is an atom: intersection
with $A$ is trivial, and the join with $A$ is $R$ in the diamond.
Fix one invariant complement $C$. In the semidirect product $R=A\rtimes C$,
other complements are graphs of derivations $C\to A$. Since $H$ preserves
both factors, the invariant complements correspond precisely to the
$H$-equivariant derivations. These form a vector space over the prime field
of $A$. Their number is a prime power, allowing the zeroth power.
This is the operator-invariant version of the complement-counting
argument in \citet[Theorem~3]{PalfyPudlak1980}.
The number equals $b-1$ in the bottom diamond $M_b$, contradicting
\Cref{f:prop-prepared-links}\ref{f:link-diamonds}.

\emph{Trivial bottom and proper nontrivial inner operators.}
Suppose $P=1$ and $1<I<S$. Transporting $I$ gives a proper interior product
label. A label $Y$ in another horizontal component is subdirect: its
nontrivial coordinate projection $R$ satisfies $RI=S$, because its join
with the product of the $I$'s is full. The subgroup $I$ normalizes $R$,
so $R\trianglelefteq RI=S$ and hence $R=S$.
Write $(T,\beta)$ for the extension describing $Y$ in
\Cref{f:lem-extension-description}. Its image contains all inner
automorphisms. Otherwise $\beta(T)\cap S$ is a nontrivial proper invariant
subgroup of the diagonal, containing $I$. Transporting it over the tied
parts and multiplying gives a nonbottom label below $Y$ with proper
coordinate projections. It would lie in the same component, contradicting
the preceding conclusion about all its labels.

Choose $y$ in this component two below the top with a diamond above it.
The filter above $y$ identifies dually, by restriction of $\beta_y$, with
$[A_0,T_y]$. Every domain strictly larger than $A_0$ has image containing
$S$. Apply this observation to $A_0\ker\beta_y$: its image is still
$\Lambda$, which does not contain $S$. Hence $\ker\beta_y\le A_0$.
After quotienting by this kernel, every proper increase over $\Lambda$
contains $S\Lambda$, so there is a unique minimal proper increase. This
contradicts the diamond's having at least two middle elements.

\emph{Trivial bottom and all inner operators.}
Finally, if $P=1$ and $I=S$, every nonbottom label has nontrivial coordinate
projections invariant under all inners, so those projections equal $S$.
The nonbottom labels are exactly the invariant subdirect groups, and
\Cref{f:lem-extension-description} identifies their dual with the extension
poset. The bottom contributes the artificial top. This proves
\ref{f:case-extension} and completes the case analysis.
\end{proof}

Choose now an ordinary representation of the whole $L$ with smallest
possible top-group order, which exists under the hypothesis of
\Cref{f:prop-ordinary-interval}. Self-duality and
\Cref{f:prop-interval-cases}\ref{f:case-subdirect} exclude the subdirect-bottom
case. Consequently either $L$ already has an ordinary almost simple
representation, or it is isomorphic to an extension poset with an artificial
top. We eliminate the latter alternative unless a prepared reverse point
link has the desired almost simple representation.

\subsection{Pure extensions and their kernels}

The full-inner extension case and its kernel description appear in
\citet[Proposition~2.5 and Remark~2.6]{PalfyTwisted}; see also the
attribution to Aschbacher's signalizer formulation there in Section~4.
We now use the affine lattice to constrain these extensions further.

\begin{definition}\label{f:def-pure-extension}
An extension datum consists of a subgroup $H$ of a finite ambient group
$K$, a nonabelian simple group $B$, and a homomorphism
$\psi:H\to\operatorname{Aut}(B)$ whose image contains $B$, identified with
its inner automorphisms. Its extension poset consists of all homomorphisms
$\beta:T\to\operatorname{Aut}(B)$, $H\le T\le K$, restricting to $\psi$,
ordered by extension, together with an artificial top. Put
\[
 D_0=\ker\psi,\qquad P_0=\psi^{-1}(B).
\]
An extension is \emph{pure} if $\beta(T)=\psi(H)$, equivalently
$T=H\ker\beta$.
\end{definition}

\begin{lemma}[Purity and kernel joins]\label{f:lem-purity}
If the full extension poset with artificial top of
\Cref{f:def-pure-extension} is isomorphic to $L$, every proper label is pure.
These labels are then parametrized by all $H$-normalized subgroups $E\le K$
satisfying
\begin{equation}\label{f:eq-pure-kernel}
 D_0\le E,\qquad E\cap H=D_0.
\end{equation}
The domain is $T=HE$, and its entire lower ideal is the ordinary interval
$[H,T]$. Proper joins of kernel labels are generated subgroups. If two
kernel labels join to the artificial top, their generated subgroup
contains $P_0$.
\end{lemma}
\begin{proof}
Subextensions of pure extensions are pure. If finitely many pure labels
have a common extension, their join is its restriction to the subgroup
generated by their domains, and is pure as well: its image is generated
by copies of $\psi(H)$.

Suppose an atom $(T,\beta)$ is not pure. The intermediate restriction to
$H\ker\beta$ is pure, so atomhood forces $H\ker\beta=H$ and
$\ker\beta=D_0$. Since $H$ already maps onto the inner subgroup $B$, it
also follows that $\beta^{-1}(B)=P_0$. Thus $T$ normalizes both $P_0$ and
$D_0$. The simultaneous normalizer
\[
 N_K(P_0)\cap N_K(D_0)
\]
acts on $P_0/D_0\cong B$. Its conjugation homomorphism restricts to $\psi$
on $H$, since action on the identified inner automorphisms determines an
automorphism of the centerless group $B$. On each nonpure atom it similarly
agrees with $\beta$. This gives one extension label $\omega$, strictly below
the artificial top, bounding every nonpure atom.

Choose a nonpure atom $u$. By
\Cref{f:prop-prepared-links}\ref{f:link-separation}, there is a rank-two
label $z\ge u$ not below $\omega$. At most one atom of its diamond lies
below $\omega$, since two distinct atoms join to $z$. The diamond has at
least three atoms, so at least two are pure. Their proper join is $z$,
which is consequently pure; its subextension $u$ is pure, a contradiction.
All atoms, and hence all proper labels by atomisticity and purity of proper
joins, are pure.

A pure extension has a kernel satisfying \eqref{f:eq-pure-kernel} and
domain $HE$. Conversely, such a subgroup $E$ is normal in $HE$, and
$HE/E\cong H/D_0$ defines the unique extension killing $E$ and agreeing with
$\psi$. This proves the parametrization. If $H\le U\le HE$, then
$U=H(U\cap E)$, so restricting the extension identifies its lower ideal
with all of $[H,HE]$.
For a proper join, restriction of a common extension to the generated
domain shows that the kernel is the subgroup generated by the kernels:
their product with $H$ is the generated domain, and their intersection
with $H$ is $D_0$. If two labels instead join to the artificial top and
their generated kernel $E$ still met $H$ in $D_0$, the same construction
would give a common proper extension, impossible. Therefore $E\cap H$
properly contains $D_0$. It is normal in $H$, so its image is a nontrivial
normal subgroup of the almost simple group $\psi(H)$; it contains $B$.
As $D_0\le E$, this implies $P_0\le E$.
\end{proof}

\subsection{The lower affine-grid contradiction}

We compare pure extensions indexed by the points, lines, and planes of
an affine four-space. Four dimensions allow two planes to meet in exactly
one point; after using this to choose a suitable plane, the contradiction
is obtained inside that plane. The next lemma requires purity only for
the labels of dimension at most two.

\begin{lemma}[Lower-grid obstruction]\label{f:lem-lower-grid}
Fix an extension datum as in \Cref{f:def-pure-extension}, and an affine
four-space over a finite field with at least three elements. Suppose each nonempty
flat $X$ of dimension at most two supplies a pure extension with domain
$T_X>H$ and kernel $E_X$, ordered in reverse containment, with the following
properties.
\begin{enumerate}[label=\textup{(\roman*)}]
\item If $X\cap Y=\{i\}$, their labels have proper join the point label $i$;
thus $T_i=\langle T_X,T_Y\rangle$ and
$E_i=\langle E_X,E_Y\rangle$. Distinct point labels have incompatible
join, and their generated kernels contain $P_0$.
\item For each $X$, modulo $L_X=\operatorname{core}_{T_X}(H)$ the interval
has the normal-supplement and unique nonabelian monolith properties of
\Cref{f:prop-monolith}, including transitivity of the base on the simple
factors and the trivial monolith centralizer. Denote
its monolith by $W_X$, with preimage $V_X$ in $T_X$.
\item For every point $i\in X$, the subgroup
$(T_X\cap V_i)/L_i$ is subdirect in $W_i$ and is strictly above the
intersection of the base with $W_i$.
\end{enumerate}
These data cannot exist.
\end{lemma}
\begin{proof}
We organize the argument into the compatible monolith embeddings, the
selection of a plane with small cores, and the final factor action.

\emph{Monolith embeddings.}
For $i\in X$ put
\begin{equation}\label{f:eq-zxi}
 Z_{X,i}=(T_X\cap V_i)/L_i\le W_i.
\end{equation}
The base $H$ is transitive on the simple factors of $W_i$. It is therefore
transitive on the tied diagonal factors of its invariant subdirect group
$Z_{X,i}$. This makes $Z_{X,i}$ a nontrivial minimal normal subgroup of
\[
 T_X/L_i=(H/L_i)Z_{X,i}.
\]
It is not contained in $H/L_i$, by the strictness in the hypothesis.
Since $T_X\le T_i$, we have $L_i\le L_X$. Minimal normality and
$L_X\le H$ now give $L_X\cap V_i=L_i$. Thus the image of $Z_{X,i}$ in
$T_X/L_X$ is still injective and minimal normal, so equals $W_X$.
We obtain a canonical $H$-equivariant subdirect injection
$f_{X,i}:W_X\to W_i$, and
\begin{equation}\label{f:eq-monolith-preimages}
 V_X=(T_X\cap V_i)L_X.
\end{equation}

For $i\in X\subseteq Y$ we have
$H\le T_Y\le T_X\le T_i$ and $L_i\le L_X\le L_Y$.
Intersecting \eqref{f:eq-monolith-preimages} with $T_Y$ gives
\[
 V_X\cap T_Y=(V_i\cap T_Y)L_X,\qquad
 V_Y=(V_X\cap T_Y)L_Y,\qquad
 L_Y\cap V_X=L_X.
\]
The last equality follows also by using $L_Y\cap V_i=L_i$.
These formulas give injections $f_{Y,X}:W_Y\to W_X$ composing with the
$f_{X,i}$ as expected: lift representatives in $V_i\cap T_Y$.
They are subdirect. Indeed, every simple factor of $W_X$ is visible by
projection to some simple coordinate of $W_i$ through $f_{X,i}$, and the
composite image of $W_Y$ projects onto every such coordinate. The kernel
of the action of $H$ on $W_X$ is exactly $L_X$, by the trivial centralizer
of the monolith in $T_X/L_X$.

If $X\cap Y=\{i\}$, the proper-join condition gives
$T_i=\langle T_X,T_Y\rangle$. Writing each domain as $H$ times its
intersection with $V_i$ shows that
\begin{equation}\label{f:eq-core-intersection}
 W_i=\langle Z_{X,i},Z_{Y,i}\rangle,
 \qquad L_i=L_X\cap L_Y.
\end{equation}
For completeness, the subgroup generated by the two $Z$'s contains
$(H\cap V_i)/L_i$, is normalized by $H/L_i$, and multiplying it by $H/L_i$
gives $T_i/L_i$. Intersecting that product with $W_i$ therefore recovers
all of $W_i$. An element of $H$ acts trivially on this generated group
exactly when it acts trivially on both embedded monoliths; this proves
the asserted equality of cores.

\emph{A plane whose cores are killed by $\psi$.}
For any affine plane $U$, write
\[
 L_* =\bigcap_{i\in U}L_i.
\]
Every line domain $T_Y$, for a line $Y\subset U$, normalizes $L_*$. To see
this, fix a target point $j\notin Y$ and choose two distinct lines through
$j$ meeting $Y$, at points $k_1,k_2$. An element of $T_Y\le T_{k_t}$
normalizes $L_{k_t}$, so sends $L_*\le L_{k_t}$ into the core of the line
$jk_t$, because $L_{k_t}\le L_{jk_t}$. Equation
\eqref{f:eq-core-intersection} for those two lines puts the resulting
image inside $L_j$. For $j\in Y$, use $T_Y\le T_j$ directly. The same
argument for inverses proves normalization. Two distinct lines through a
point generate its domain, so $L_*$ is normal in the group generated by
all point domains of $U$.

Since $L_*\le H$, we have $L_*\cap E_i=L_*\cap D_0$ for every point $i$.
Both $E_i$ and $L_*$ are normal in $T_i$, so
$[E_i,L_*]\le L_*\cap D_0$. Thus every $E_i$ centralizes
$L_* /(L_*\cap D_0)$. The group generated by kernels of distinct points
contains $P_0$, by hypothesis. Therefore $\psi(L_*)$ centralizes $B$,
whose centralizer in $\operatorname{Aut}(B)$ is trivial. Hence
\begin{equation}\label{f:eq-plane-intersection-killed}
 L_*\le D_0.
\end{equation}

There is a point $i$ with $L_i\le D_0$. Otherwise every $\psi(L_i)$ is a
nontrivial normal subgroup of $\psi(H)$ and contains $B$. Take an iterated
commutator of the original normal subgroups $L_i$ inside $H$. It lies in
$\bigcap_{i\in U}L_i$ and its $\psi$-image still contains $B$, because
$B$ is perfect.
That contradicts \eqref{f:eq-plane-intersection-killed}.
Choose two affine planes in the four-space meeting exactly at this point
$i$. Their cores intersect in $L_i$ by
\eqref{f:eq-core-intersection}, so their two core images cannot both
contain $B$, by the same commutator argument. Hence one of these planes,
henceforth denoted $U$, has $L_U\le D_0$. All point cores of $U$ then lie
in $D_0$, since $L_i\le L_U$ for $i\in U$.

Quotient the group generated by the point domains of $U$ by
$\bigcap_{i\in U}L_i$. This subgroup is normal there as shown above,
lies in the base, and is contained in all the point cores and all extension
kernels in the plane. The core in each quotient is the quotient of the
old core. Thus the inclusions, join statements, and monolith identifications
survive. We retain the notation and may assume
\begin{equation}\label{f:eq-zero-common-core}
 \bigcap_{i\in U}L_i=1,\qquad L_U\le D_0.
\end{equation}

\emph{Compatible perfect subgroups.}
For a nonempty affine subflat $X\subseteq U$ set
\begin{equation}\label{f:eq-mx}
 M_X=E_X\cap\bigcap_{i\in X}V_i.
\end{equation}
Every factor $V_i\cap T_X$ is normal in $T_X$, as is $E_X$, so $M_X$ is
normal in $T_X$. It lies in $V_X$ by
\eqref{f:eq-monolith-preimages}. It maps onto $W_X$: each
$V_i\cap T_X$ maps onto $W_X$, while the normal image of $E_X$ modulo $L_X$
is nontrivial, since $HE_X=T_X>H$ and $L_X\le H$, and therefore contains
the monolith. Taking nested commutators of these finitely many normal
subgroups gives a subgroup in their intersection whose image still
contains $W_X$, by perfection of $W_X$. As $M_X\le V_X$, its image is exactly
$W_X$. The supplement condition then yields
\begin{equation}\label{f:eq-domain-mx}
 HM_X=T_X,\qquad D_0M_X=E_X.
\end{equation}
The second equality follows by intersecting the first with $E_X$.

For $X=U$, the kernel of $M_U\to W_U$ is contained in every $L_i$: it lies
in $L_U\cap V_i=L_i$. Equation \eqref{f:eq-zero-common-core} makes it trivial.
Consequently $M_U\cong W_U$ is an embedded product of nonabelian simple
groups; write its factors as $C_a$. The normal subgroups $L_U$ and $M_U$
of $T_U$ have trivial intersection, so $L_U$ centralizes every $C_a$.

The subdirect embedding $f_{U,X}$ assigns to each factor $C_a$ a disjoint
set of factors of $W_X$ on which it has its tied diagonal, and these sets
partition all factors of $W_X$. Let $M_{X,a}$ be the preimage in $M_X$ of
the entire product of $W_X$-factors assigned to $a$. These groups generate
$M_X$, are permuted by $H$, and each is normalized by $L_U$: the latter
centralizes $M_U$, so fixes its factors and their assigned sets.
For $i\in X$, the images of $M_{X,a}$ in $W_i$ have support only over the
same index $a$. For $a\ne b$ it follows that
\begin{equation}\label{f:eq-cross-commutators}
 [M_{X,a},C_b]\le L_i\le L_U.
\end{equation}
In fact $M_{X,a}$ centralizes $C_b$. Since $L_U$ centralizes $C_b$,
\eqref{f:eq-cross-commutators} and the three-subgroups lemma imply
$[M_{X,a},[C_b,C_b]]=1$; perfection gives the claim. Likewise, for meeting
lines $X,Y$ (allowing $X=Y$) and $a\ne b$, a common point $i$ gives
\begin{equation}\label{f:eq-line-commutators}
 [M_{X,a},M_{Y,b}]\le L_i\le L_U.
\end{equation}

\emph{The final action on the simple factors.}
Choose four pairwise nonparallel lines in $U$, two through a point $i$ and
two through a distinct point $j$. There are at least four directions,
since the field has at least three elements; the lines can be chosen
with the required distinct directions. Set
\[
 R'=\langle M_Y:Y\text{ is one of these four lines}\rangle.
\]
It is normalized by $H$. The two proper point joins and then the
incompatible join of $i,j$, together with \eqref{f:eq-domain-mx}, imply
$P_0\le D_0R'$. Hence
\begin{equation}\label{f:eq-surjects-b}
 \psi(P_0\cap R')=B.
\end{equation}

It remains to show that this image is soluble. We first show that
$P_0\cap R'$ fixes every factor $C_a$. Schreier's Theorem then controls
its outer action, and we show that $\psi$ kills the subgroup inducing
inner automorphisms.

Fix an index $b$.
A word in the generating groups $M_{Y,a}$ can be reordered as
\[
 \ell u v,
\]
where $\ell\in L_U$, all letters of $u$ have indices $a\ne b$, and all
letters of $v$ have index $b$. For adjacent letters in the wrong order,
put $e=xyx^{-1}y^{-1}\in L_U$, by
\eqref{f:eq-line-commutators}. If $p$ is the preceding word, then
\[
 pxy=peyx=e(e^{-1}pe)yx.
\]
Conjugating the preceding letters by $e$ preserves their indices because
$L_U$ normalizes every indexed generating subgroup. Thus the swap reduces
the inversion count. Induction gives the required reordering.

In right-conjugation notation, $C_b^{\ell u v}=C_b^v$, since $\ell$ and $u$
centralize $C_b$. The element $v$ centralizes every other $C_a$, and $C_b$
already commutes with them, so $C_b^v$ commutes with every $C_a$ for
$a\ne b$. If the word lies in $P_0\cap R'\le H$, it permutes the factors
of $M_U$; its image of $C_b$ cannot be another $C_a$, because that would
make a nonabelian simple group commute with itself. It must therefore be
$C_b$, proving the assertion.

The action of $P_0\cap R'$ on $W_U\cong M_U$ consequently has soluble
quotient modulo its inner action, by Schreier's Theorem on each fixed
simple factor. An element inducing an inner action on $W_U$ lies in
$V_U$: after multiplying by the corresponding monolith element its image
in $T_U/L_U$ centralizes $W_U$, and that centralizer is trivial. On the
other hand, purity and $L_U\le D_0$ give
\[
 H\cap V_U\le H\cap(E_UL_U)\le D_0L_U=D_0,
\]
because the nontrivial normal image of $E_U$ contains $W_U$ and hence
$V_U\le E_UL_U$. Thus the inner-action kernel within $P_0\cap R'$ is
killed by $\psi$. Consequently $\psi(P_0\cap R')$ is a quotient of a soluble group,
contradicting \eqref{f:eq-surjects-b} and the nonabelian simplicity of $B$.
\end{proof}

\begin{theorem}[Ordinary almost simple starting interval]\label{f:thm-almost-simple-reduction}
If the test lattice $L$ is representable, then either the whole $L$ or a
prepared reverse link with top an initial trunk point has an ordinary
almost simple subgroup-interval representation. On the reverse side,
there are at least $D-1$ further generations of prepared stages starting
at that top. This conclusion uses only the field condition and
\eqref{f:eq-parameters}, independently of later choices of larger
parameters.
\end{theorem}
\begin{proof}
Choose the least-top-order ordinary representation furnished by
\Cref{f:prop-ordinary-interval}. The reduction after
\Cref{f:prop-interval-cases} either already gives the whole $L$ as an
ordinary almost simple interval, or gives an extension poset with
artificial top isomorphic to $L$. In the latter case all proper labels
are pure by \Cref{f:lem-purity}. For each nonempty reverse flat $X$, write
$E_X,T_X$ for its kernel and domain. Its lower ideal is the ordinary
interval $[H,T_X]$.

Take an affine four-subspace of one initial trunk chart and consider all
its points, lines, and planes. Their ranks are at most three, hence at
most $m$, so all their reverse links are prepared. The proper and
incompatible joins have exactly the domain and kernel properties of
\Cref{f:lem-lower-grid}, by \Cref{f:lem-purity}. Each ordinary domain
interval has the normal-supplement and monolith properties of
\Cref{f:prop-monolith}. If at some point $i$ the nontrivial-product case
\ref{f:case-product} of \Cref{f:prop-interval-cases} occurs, it gives the
required almost simple representation of that point's link. Otherwise,
the remaining possibilities at every point are the subdirect-bottom and
the trivial-bottom full-inner cases. In both, every domain label strictly
above the base intersects the point monolith subdirectly modulo the core,
as required in \Cref{f:lem-lower-grid}. The strictness in that lemma follows
from $HV_i=T_i$: every incident domain satisfies
$T_X=H(T_X\cap V_i)$, so $T_X>H$ gives strict containment above the
point-base intersection after quotienting by $L_i$.
That lemma gives a contradiction.

For the whole lattice take the empty reverse flat as the starting top.
For a point in an initial trunk chart, its rank increment from the empty
flat is one, which lies between $1$ and $r$; its child path was therefore
reserved at the first stage. In either case at least $D-1$ generations
remain by the construction.
\end{proof}

\section{Simultaneous path pruning and affine witness patterns}\label{sec:patterns}

We now organize the local group arguments into a finite pruning construction on the test lattices of \Cref{sec:affine}. The construction is simultaneous over all ancestors of a path. Its output will be used with the classical and alternating exclusions in \Cref{sec:classical,sec:alternating}, \Cref{w:bounded-exclusion}, and \Cref{w:path-bound}.

\subsection{Path labels and compatible simple groups}

We work in the ordinary interval supplied by
\Cref{f:thm-almost-simple-reduction}, with fixed base group $H$.
Write $K_X$ for the subgroup labeled by a reverse flat $X$. Thus
$X\subseteq Y$ means $K_Y\le K_X$. A downward path chooses, at its current
flat $W$, a private trunk chart and then a flat $U\supsetneq W$ of rank
increment between $1$ and $r$ in that chart. The local group arguments
will be applied both in the current group's almost simple quotient and
in the quotients belonging to all earlier path nodes.

At a path node $W$ we maintain an actual nonabelian simple group
$B_W\trianglelefteq K_W$ such that
\[
 HB_W=K_W,\qquad C_W:=C_{K_W}(B_W)\le H,
\]
where $C_W$ is soluble and $K_W/C_W$ is almost simple with socle the
faithful image of $B_W$. The starting interval has this property. The
next lemma explains why the local conclusions in several such quotients
give one subgroup in the original interval.

\begin{lemma}[Compatibility of ancestor socles]\label{c:ancestor-socles}
Suppose $H\le K\le K_W$, and a local conclusion in $K_W/C_W$ supplies
a nonabelian simple normal subgroup $\overline B\trianglelefteq K/C_W$
contained in the image of $B_W$. Suppose also that
$(H/C_W)\overline B=K/C_W$, and that its centralizer in $K/C_W$ is
soluble and contained in $H/C_W$. Then $\overline B$ has a unique lift
$B\le B_W$, with
\[
 B\trianglelefteq K,\qquad HB=K,\qquad C_K(B)\le H,
\]
and $C_K(B)$ is soluble. The quotient by this centralizer is almost simple.
If the earlier actual socles contain $B_W$, the image of $B$ in every
earlier ancestor quotient agrees with any nonabelian simple normal
subgroup supplied there with soluble centralizer. All the simultaneous
structural conclusions therefore concern this same $B$.

Whenever two comparable candidate labels have these simple and
centralizer properties, their actual simple groups are nested, strictly
if the labels are distinct.
\end{lemma}

\begin{proof}
The quotient map is faithful on $B_W$, since $B_W\cap C_W=1$.
Its inverse on the image of $B_W$ defines $B$, which lies in $K$ and is
normal there. The supplement identity lifts because $C_W\le H$.
A lift of an element centralizing $\overline B$ centralizes $B$: its
commutators with $B$ lie both in $B$ and in $C_W$, so are trivial.
Hence the new centralizer is the inverse image of the quotient
centralizer and is soluble and contained in $H$. Quotienting by it gives
a faithful automorphism group of $B$ containing its inner automorphisms.

For an earlier node $W'$, the inclusion $B\le B_W\le B_{W'}$ makes its
image faithful in $K/C_{W'}$. The same commutator argument identifies
its centralizer there with $C_K(B)/C_{W'}$, which is soluble. Distinct
nonabelian simple normal subgroups would commute, contradicting that
solubility. Thus the simple supplied relative to $W'$ must be this image.
All these quotient maps preserve the exact interval, since their kernels
lie in $H$.

For comparable candidates with $K_U\le K_L$, the simple $B_U$ acts on
$B_L$. Schreier's Theorem makes its outer action trivial, so
$B_U\le B_L C_{K_L}(B_L)$. This is a direct product because $B_L$ is
centerless. Projection of $B_U$ onto the soluble centralizer is trivial,
giving $B_U\le B_L$. If the labels differ, their simples differ because
$HB_U=K_U$ and $HB_L=K_L$.
\end{proof}

\subsection{One pruning round and its failure witnesses}

A structural exclusion forbids containment in specified proper full
stabilizers in an ancestor's almost simple quotient. Such a full
stabilizer is pulled back to an actual subgroup of the original interval.
Avoiding it is preserved when a tested subgroup is enlarged.

At the current path flat $W$, fix one category of these stabilizers,
allowing finitely many ancestor and case choices. A round tests flats
$U\supsetneq W$ of a prescribed rank increment $s$ in each available
private chart. A successful flat has a label $K_U$ contained in none of
the stabilizers under test. After the first round, tests take place
inside a previously selected successful flat $V$, with
$s<\operatorname{rk}V-\operatorname{rk}W$. The subgroup $E=K_V$ then
satisfies $H<E<K_U$ and retains all preceding exclusions. We call it the
\emph{anchor} for that chart.

\begin{lemma}[Prepared full-stabilizer witnesses]\label{c:prepared-witness}
Suppose $K_W$ already avoids the final exclusion list relative to every
strict ancestor. A failed test at $U$ supplies an actual label $K_X$ with
\[
 W'\subseteq X\subseteq U,\qquad X\not\subseteq W,
\]
where $W'$ is the ancestor at which the failure occurs. If all trial
ranks lie below $m$, this witness has a prepared reverse lower interval.
In a round with anchor $E$, it also satisfies $H<E<K_X$.
\end{lemma}

\begin{proof}
The proper full stabilizer contains $K_U$, so its label lies between
$K_U$ and $K_{W'}$. It is therefore a subflat $X$ of $U$ containing $W'$.
If $X\subseteq W$, its subgroup contains $K_W$, contrary to the earlier
exclusions; for $W'=W$, properness gives the same conclusion.
Such an $X$ is nonempty and lies in the trunk chart containing $U$, so
\Cref{f:prop-prepared-links} applies. Finally $E<K_U\le K_X$.
\end{proof}

\subsection{Patterns forced by a failed round}

\begin{lemma}[Affine patterns in failed charts]\label{c:failed-pattern}
Suppose every increment-$s$ test in a sufficiently large available flat
fails, all trial ranks lie below $m$, and a witness
$X(U)\subseteq U$, with $X(U)\not\subseteq W$, is chosen for each test.
Color it by its ancestor, category and case. Given any finite number of
parameter directions, sufficient available dimension produces a family
$X_t$, indexed by an affine space over $\mathbb F_q$, with these properties:
\begin{enumerate}[label=\textup{(\roman*)}]
\item All witnesses have the same color and rank. For distinct $t,t'$,
their span has rank one larger, contains every $X_u$ on the parameter
line through $t,t'$, and gives the subgroup intersection
$K_{X_t}\cap K_{X_{t'}}$. Each witness label covers this intersection.
\item An affinely independent parameter list increases the span rank by
one at each new point, giving strictly descending subgroup intersections
for as long as the spans remain within the truncation.
\item In an anchored round every label contains the same $E>H$.
Families in different private charts have pairwise subgroup
intersections exactly $H$.
\end{enumerate}
The first round can supply more labels than the full interval height
plus any fixed constant, without requiring their total span to be proper.
\end{lemma}

\begin{proof}
Choose a fixed model for the horizontal increment and many unused
vertical coordinates, with coordinate space $V_{\rm vert}$.
For nonempty $W$, take an origin in $W$ and graph
linear maps from its $s$ new horizontal directions to the vertical
directions; these maps vanish on $W$. For empty $W$, use $s-1$ horizontal
affine dimensions and graph affine maps instead. Every graph is an
increment-$s$ test flat.

Color each map by its witness's failure color and the pullback of $X(U)$
into the fixed horizontal model. There are finitely many colors for the
fixed horizontal data. View the maps as words whose letters are
individual-coordinate linear or affine functionals. Multidimensional
Hales--Jewett \citep{HalesJewett1963} gives a monochromatic family with
fixed coordinates and any prescribed number $k$ of disjoint nonempty
variable batches. Thus the pulled-back witness is one fixed flat
$X_{\rm hor}\not\subseteq W$.

Choose a horizontal functional $\lambda$ nonzero at some point of $X_{\rm hor}$;
in the nonempty case it vanishes on $W$. Restrict the functional in each
variable batch to a scalar multiple of $\lambda$. Let
$\iota:\mathbb F_q^k\to V_{\rm vert}$ send each parameter coordinate
to the sum of the vertical basis vectors in its batch. The disjoint
nonempty batches make $\iota$ injective. The witness family is
\begin{equation}\label{c:graph-witnesses}
 X_t=\{(x,f_0(x)+\lambda(x)\iota(t)):x\in X_{\rm hor}\},
 \qquad t\in\mathbb F_q^k,
\end{equation}
where $f_0$ vanishes on the variable batches and records the fixed graph
coordinates elsewhere. At a point where $\lambda\ne0$, distinct parameters
have different vertical coordinates, so the witness flats are distinct.
The span of two graph flats adds exactly the vertical direction
$\iota(t'-t)$, and therefore contains the witnesses on the parameter
line. More generally,
\begin{equation}\label{c:graph-span-rank}
 \operatorname{rk}\langle X_{t_0},\ldots,X_{t_h}\rangle
 =\operatorname{rk}X_{\rm hor}+
   \dim\operatorname{aff}\{t_0,\ldots,t_h\}.
\end{equation}
The chart join rule converts these spans, when proper, into the asserted
subgroup intersections and covers. Pair spans are proper because the
trial ranks are below $m<n-2$. The anchor inclusion follows from the
round setup. Across private charts the witness flats lie outside $W$
and have no common proper primal upper bound, so their reverse meet is
the base $H$.

Taking sufficiently many parameter directions supplies the first-round
family of more labels than the interval height. Only pair spans and the
bounded independent lists used later must fit the truncation; the span
of the whole family need not be proper. In subsequent rounds the number
of parameter directions is bounded independently of $n$.
\end{proof}

\begin{figure}[tbp]
\centering
\begin{tikzpicture}[x=1cm,y=1cm,
  witness/.style={font=\small,inner sep=2pt},
  inclusion/.style={densely dashed},
  panel/.style={draw=black!30,rounded corners=3pt,inner sep=7pt}]
\node[witness] (h) at (0,0) {\(H\)};
\node[witness] (t1) at (-2.4,1.05) {\(T_1\)};
\node[witness] (t2) at (2.4,1.05) {\(T_2\)};
\node[witness] (k11) at (-3.7,2.2) {\(K_{X_1}\)};
\node[witness] (k12) at (-2.4,2.2) {\(K_{X_2}\)};
\node[witness] (k1q) at (-1.1,2.2) {\(K_{X_q}\)};
\node[witness] (k21) at (1.1,2.2) {\(K_{Y_1}\)};
\node[witness] (k22) at (2.4,2.2) {\(K_{Y_2}\)};
\node[witness] (k2q) at (3.7,2.2) {\(K_{Y_q}\)};
\node[font=\scriptsize] at (-1.75,2.2) {\(\cdots\)};
\node[font=\scriptsize] at (3.05,2.2) {\(\cdots\)};
\foreach \target in {k11,k12,k1q}{\draw (t1)--(\target);}
\foreach \target in {k21,k22,k2q}{\draw (t2)--(\target);}
\draw[inclusion] (h)--(t1);
\draw[inclusion] (h)--(t2);
\node[panel,fit=(t1)(k11)(k12)(k1q)] (p1) {};
\node[panel,fit=(t2)(k21)(k22)(k2q)] (p2) {};
\node[font=\small,above=3pt of p1] {First private chart};
\node[font=\small,above=3pt of p2] {Second private chart};
\node[font=\small,below=7pt of h] {\(K_{X_i}\cap K_{Y_j}=H\quad\text{for every }i,j\)};
\end{tikzpicture}
\caption{The subgroup relations supplied by \Cref{c:failed-pattern}.
In the first chart, the witnesses on one parameter line satisfy
\(\Sigma_1=\langle X_i,X_j\rangle\) for every \(i\ne j\),
\(\operatorname{rk}\Sigma_1=\operatorname{rk}X_i+1\), and
\(T_1=K_{\Sigma_1}=K_{X_i}\cap K_{X_j}\); each \(K_{X_i}\) covers \(T_1\).
The second chart has the analogous relations with \(Y_j,\Sigma_2,T_2\).
Solid segments denote these covers, and dashed segments denote inclusions,
with intervening labels omitted. The diagram is schematic in rank and
shows only the stated relations.}
\label{c:pattern-figure}
\end{figure}


Assign each failed chart the ancestor/category/case of one homogeneous
pattern extracted by \Cref{c:failed-pattern}.
Suppose the exclusion argument for each such choice
prohibits either one such pattern or a pair in distinct charts. In the
first case no chart can fail in that case; in the second at most one can.
Thus the number of charts lost in a round is at most the number of cases
requiring two charts. These are the group-theoretic failure colors; the
pullback-flat colors used inside Hales--Jewett do not enlarge the
chart-loss count. Retaining one successful flat in every surviving chart
provides the anchors for the next round.

If $F$ is successful after the last round, every
$W\subsetneq X\subseteq F$ has $K_X\ge K_F$ and therefore avoids every
completed-category stabilizer relative to every ancestor tested.
The local conclusions consequently apply to all comparable candidates
in $F$, before a next path node is chosen. \Cref{c:ancestor-socles} makes
their simple groups actual and nested. This is the property used for
simultaneous root and tensor protection in \Cref{w:selection}.

\begin{proposition}[Finite schedule of simultaneous exclusions]\label{c:pruning-schedule}
Suppose that the local exclusions below prohibit their prescribed failed patterns with absolute case and chart-loss bounds. For a bounded path length, finitely many rounds then produce a flat of final increment \(r\) such that the subgroup label of every proper extension in it is contained in none of the prescribed proper full stabilizers relative to any high-type ancestor. The next local simple has the corresponding transitivity or full-spanning properties relative to all those ancestors.
\end{proposition}

\begin{proof}
Use a simultaneous initial round for all initial classical subspace exclusions relative to any classical ancestors (the ones allowing a more-than-height requirement); their pattern proofs need no earlier anchor. The remaining categories of the structure arguments can then be ordered as there, requiring only bounded line/orbit or independent-list bounds: e.g. direct line systems and then other direct systems, prime fields then ordered factors for linear type, partial ordered factors then prime fields then paired idempotents and remaining ordered factors for isometry types, elementary subgroups then tensor systems; in alternating ancestors subsets and block partitions then affine, Cartesian and diagonal structures. Subcases rendered impossible by an already forbidden structure need no round of their own. There are only boundedly many categories/ancestor instances altogether for bounded path length. We thus obtain a pruned flat of final increment \(r\) from \(W\); all labels at proper extensions within it avoid the exclusions, giving the needed next almost simple socles by the structural arguments, and giving transitivity or natural full spanning as applicable relative also to earlier ancestors. Selection of which extension to take (additional closure-field tensor/root protections within a same-characteristic run) uses \Cref{w:selection}. If the current parent instead is of bounded Lie rank, use \Cref{w:bounded-exclusion} (same graph pattern with a bounded independent list and no prior anchor, relative just to this parent).
\end{proof}

\begin{proposition}[Order of the finite parameter choices]\label{c:parameter-order}
The test-lattice parameters can be chosen in an order that meets all the pattern requirements and all the later path requirements. In particular the initial height-dependent pattern requirement does not create a circular dependence between \(m,n,N\). Once the local group results are applied, these choices force either the bounded-type contradiction or a path longer than the protected-path bound.
\end{proposition}

\begin{proof}
Here is one parameter order, always finite. Fix the high thresholds, the resulting bounded-type constants, and the absolute bound on protected high-type path lengths from the group arguments. Choose \(D\) exceeding this bound by several stages (allowing the displaced point start and a stage for the bounded-type contradiction). Choose \(q\), a power of 7 for instance, above all absolute line-count requirements. Choose \(r\) as in the final flag selection, large for bounded-type requirements and the exponential-gap union bound over at most \(D\) steps, so the path ranks are at most \(Dr\). Choose larger intermediate trial increments backwards from \(r\), the increment after the first round thus a constant independent of \(m,n\): each preceding increment merely must allow the Hales--Jewett bound for the next bounded-slot round, at any path rank up to \(Dr\). Choose \(m\) above the ranks of these trials, then \(n\) sufficiently larger, including for residual interval length below all of them and bounded required spans. Now \(N\) can allow the first-round Hales--Jewett bound including the larger initial slot requirement depending on \(n\), and any no-anchor bounded-type trial. Finally take \(g\) sufficiently large for the failed-chart bounds at every stage. Thus the separate arguments imply contradiction on this finite lattice: either a bounded-type parent gives the direct contradiction, or too long a protected path is forced.
\end{proof}

\section{Pruning inside a high-rank classical ancestor}\label{sec:classical}

All assertions in this section concern the prepared configuration of \Cref{sec:patterns}. A structure is excluded when its full stabilizers cannot supply the failed-chart patterns required by \Cref{c:failed-pattern}; this is precisely the input needed by \Cref{c:pruning-schedule}. The statement concerns these prepared configurations, and includes the preceding exclusions specified at each stage.

\subsection{Standing interval and matrix conventions}

\begin{convention}[Prepared classical configuration]\label{c:classical-configuration}
Here we give the high-rank classical structure exclusions for the reverse-flat path argument. We work in one of the current almost simple quotients, with interval base again denoted by \(H\), top group \(M\), and simple socle \(S\le M\le{\rm Aut}(S)\). We have \(HS=M\): \(S\) is not in the base (the outer quotient is soluble and the base-core quotient must have a nonabelian monolith), and use the normal supplement property. (For a strict ancestor the top here still means that ancestor group in its almost simple quotient.) A bad structure has for label its \textbf{full stabilizer in \(M\)}. We consider labels containing the tested group and proper in \(M\); by irreducibility/fullness of the ambient simple groups below (of sufficiently high rank), the stabilizers we need are proper. Equivalently one can discard full-group stabilizers, checking properness when using the exclusions. Here is the same preparedness observation as for the witness flats in the path construction: any such proper full stabilizer containing a chart test label, or the chosen previously good anchor in the chart, for a structure on the full list of intended exclusions, also labels a trunk subflat within the relevant trial-rank range. Indeed for a strict ancestor the current path node has already avoided \emph{all} these final exclusions relative to it (even if the pruning order now starts over); for the current ancestor use properness. This applies also when testing or using a case condition on the structure before its exclusion is complete (such as which tensor side is retained). Each witness label \(K_i\) thus has full interval \([H,K_i]\) a prepared reverse link. Thus a normal subgroup of \(K_i\) is either retained (contained in \(H\)) or missing (supplemented by \(H\)), soluble normals are retained, and two commuting normals cannot both be missing (also valid modulo normal subgroups contained in \(H\)). These consequences of \Cref{f:prop-monolith} will be denoted by (\(*\)).

Recall the patterns supplied when pruning fails: many distinct labels in a private chart, including a line of \(q\) distinct labels with \(K_i\cap K_j\le K_l\) for all \(l\) on the line and \(i\ne j\) there. Moreover \(K_i\) covers \(K_i\cap K_j\), by the hyperplane-in-span property. We may require a bounded number of independent pattern points with strictly descending successive intersections, and in the initial round more distinct labels than the total height plus any fixed constant. We may use another such pattern in a separate chart for the same category, with \(K_i\cap K_j=H\) across the two. Any needed bounded subdivision into colors will be indicated; all patterns used together for a category can be taken homogeneous. Except in the initial round we use a common previously good anchor \(H<E<K_i\) for labels within a chart, with \([H,E]\) itself a prepared reverse link. Earlier exclusions imply \(E\) lies in none of the previously forbidden stabilizers. The prepared links have height exceeding any required absolute constants, interior horizontal branching, coatomisticity, and the bottom diamond parameters used in the interval reductions. Constants and bounded choices here are independent of the field and dimension of the ambient finite group; \(q\) denotes the \textbf{lattice} field size and will exceed all required absolute orbit/count bounds.
\end{convention}

\begin{convention}[Natural matrix groups and their small parameters]\label{c:classical-matrices}
Let \(V\) have dimension \(v\) over the finite ground field \(F\). The natural types are linear (\({\rm PSL}\)), unitary, symplectic, and orthogonal (nonsingular, using odd dimension only in odd characteristic); in isometry types write \(b\) for the defining nondegenerate bilinear or hermitian form, and include the defining quadratic form \(Q\) for orthogonal groups. Adjoints for \(b\) are denoted by \(*\). We use the standard descriptions of finite simple classical groups, their covering matrix groups (\({\rm SL},{\rm SU},{\rm Sp},\Omega\)), orders and automorphisms. Here \(v\) is above a sufficiently large absolute threshold. In the isometry cases \(M\) acts by semilinear similarities modulo scalars. In the linear case it acts on \(\mathcal A={\rm End}_F(V)\) by semilinear automorphisms and possibly anti-automorphisms (transpose with conjugation/field action); on the projective group the latter action uses also inverse. There is also the corresponding action on subspaces, order-reversing for an anti-action. Write \(+\) for the collineation/automorphism part (the entire group in the isometry cases). Isometry actions on \(\mathcal A\) respect the adjoint involution. All of this uses sufficiently high rank (so no triality or small exceptional automorphism issue).

We give a few more specifics on the standard classical-group facts being used (reference points are \citet{Taylor1992} and \citet{KleidmanLiebeck1990} for classical groups and their natural geometry, together with \Cref{f:conv-lie-inputs,f:conv-representation-inputs}). In recalling the group lists use parameter \(s=|F|\) except \(s=\sqrt{|F|}\) in unitary type. Forms and subspaces considered over the ground field use the finite-field Witt theory; ``nondegenerate'' for a subspace used in an orthogonal decomposition requires nondegeneracy of the restricted polar pairing also in characteristic two.

\begin{itemize}
\item In orthogonal type use \(\Omega\) as the kernel in the isometry group of determinant and spinor norm in odd characteristic, and of the Dickson invariant (rank of \(u-1\) modulo two) in even characteristic with nondegenerate quadratic space of even dimension. In particular ordinary isometry commutators are available inside it. These kernels are preserved by semilinear similarities (for the spinor norm on determinant one, scaling the form does not change the norm modulo squares, by the reflection formula with an even number of reflections). In the other types use determinant one in linear and unitary, and the full symplectic group. The quotients of linear similarity groups over the respective special groups are soluble; adjoining field automorphisms still gives soluble remaining quotients. For unitary and high orthogonal types the automorphisms needed in \(M\), including graph actions if present, are covered by the ground semilinear similarity description (no \(D_4\) here).
\item Recall for natural degree at least two the usual simplicity/perfection list: for linear and unitary the special group is perfect, simple projectively, outside degree two at \(s=2,3\) and the unitary degree three at \(s=2\); the analogous exceptions for symplectic are degree two at \(2,3\) and \({\rm Sp}_4(2)\). \(\Omega\) in the above convention/range of polar spaces has this high structure (perfect, simple projectively) in degree at least five. In degree three (odd characteristic) use the projective rank-one structure via the symmetric square of \({\rm SL}_2(s)\); in four the split determinant space has the two independent rank-one tensor actions (projectively), and nonsplit uses the natural module and its field-\(s\) twist tensoring for \({\rm SL}_2(s^2)\). Thus the split \(\Omega_4^+\) is perfect with two projective simple factors for \(s\ge4\), and \(\Omega_4^-\) is perfect, simple projectively. The plane and soluble cases needed will also be noted below.
\item The natural actions of the special groups in these assertions, apart from orthogonal planes, are absolutely irreducible, so their matrix lifts generate the full natural matrix algebra. Equivalently one can use the standard irreducible defining-characteristic vector modules and Steinberg restriction (restricted highest weight at the natural parameter; using \(B\) only for odd characteristic); in degree four orthogonal type use the rank-one tensor descriptions, and \({\rm SU}_2(s)\) gives the split \({\rm SL}_2(s)\) module over the extended field. For this last description one can write the defining hermitian form with matrix a suitable scalar times the standard nondegenerate alternating matrix in degree two (choose the scalar to make it hermitian and use equivalence of nondegenerate hermitian forms); the special group then includes \({\rm SL}_2(s)\), of the full special order. When using a derived replacement at \({\rm Sp}_4(2)\) we check it separately. The size comparisons below use only ordinary classical-group order formulas.
\end{itemize}
\end{convention}

The exclusion order is as follows, within the prepared configuration and
under the full hypotheses of the cited statements.  Initial subspace
exclusions precede direct-system pruning
(\Cref{c:linear-subspaces,c:isometry-subspaces,c:direct-pruning}).
In linear type, fields precede ordered factors
(\Cref{c:linear-fields,c:linear-factors}).  In isometry types, partial
ordered-factor pruning leaves the two stated retained-plane cases;
field transfer and prime-degree fields precede paired idempotents and
completion of the ordered-factor exclusions
(\Cref{c:partial-isometry-factors,c:field-transfer,c:isometry-fields,c:remaining-factors}).
Full spanning and elementary normalizers then precede permuted tensor
systems and the next actual simple socle
(\Cref{c:full-spanning,c:elementary-normalizers,c:tensor-pruning,c:classical-conclusion}).

\begin{lemma}[Properness of the stabilizers used below]\label{c:proper-stabilizers}
For sufficiently large ambient natural degree, each nontrivial structure used in the exclusions below has proper full stabilizer. The same properness check applies when an earlier exclusion is invoked for a structure found from the anchor.
\end{lemma}

\begin{proof}
For properness throughout, the high-rank covering matrix group of \(S\) is perfect and absolutely irreducible. It cannot preserve the initial subspace, signed-flag or individual-projector structures below (its own action is linear), nor even an adjoint-exchanged complementary pair of projectors up to swap in isometry type (perfection). It cannot normalize a proper direct-sum system below: irreducibility would force a uniform split by transitivity, then apply the direct-split order comparison below on \(V\) itself. Similarly use the tensor order comparison for ordered tensor factors or tensor systems; normalizing a field would force centralizing it by perfection, and a nontrivial elementary abelian subgroup inside \(S\) cannot be normal there. Thus when an earlier exclusion is invoked for a structure found using an anchor, its full stabilizer is indeed proper.
\end{proof}

\subsection{Initial subspace exclusions}

\begin{proposition}[Linear subspaces, signed flags and idempotents]\label{c:linear-subspaces}
In linear type, the initial round excludes stabilizers of proper nonzero subspaces when there are no anti-actions. In the presence of anti-actions it excludes signed flags \(0<I\le J<V\) and individual proper nontrivial idempotents. These exclusions require no prior anchor and make the collineation part of every subsequent good group irreducible.
\end{proposition}

\begin{proof}
First suppose linear type, without anti-actions. Exclude subspace stabilizers of nonzero proper \(U\). In each such label the shears from \(V/U\) to \(U\) are in the normal unipotent radical, hence in \(H\). Any two subspaces in the pattern must therefore be comparable. Stabilizers of a descending list along that chain have strictly descending intersections (special linear maps inside the smallest previous space can move a new proper subspace). The total interval height bounds their number, contradiction.

If \(M\) in linear type contains anti-actions, so does \(H\); fix one, \(a\in H\). Exclude the signed flag stabilizers for \(0<I\le J<V\) (automorphisms fix the two terms, anti-actions exchange them). For labels containing \(H\), \(J=aI\). Again the normal unipotent radical of the corresponding parabolic in \(S\) (preserved also by signed anti-actions) contains the shears from \(V/I\) to \(I\), and is in \(H\). Thus the \(I\)'s form a chain, the \(J\)'s the reversed chain above it. Intersections are strict taking descending \(I\)'s, just as before.

Also exclude stabilizers of individual nontrivial proper idempotents \(e\) (in the linear case with anti-actions). The two supported special linear groups on the complementary spaces are individually normal, commute projectively, and the quotient over their product is soluble. Thus exactly one is retained by (\(*\)); designate the retained space \(U\), with complement determined by \(a\) (anti-action on an idempotent sends its image to the subspace-dual transform of its kernel). In a pattern with \(\dim U\ge2\), the supported group acts irreducibly nontrivially there and fixes the complement pointwise. Any other retained space, being \(H^+\)-invariant, projects to zero along the complement or contains \(U\) (use differences under the supported group). By mutual irreducibility and distinctness, all retained spaces are pairwise in one another's complements. Their projectors are thus orthogonal idempotents in the ring sense. Intersecting the stabilizers strictly decreases until possibly the last one, by taking special linear maps on the unsplit remainder. Again this is height-bounded.

\textit{The retained line case.} For the color with retained dimension 1, take a cross-chart pair. There is a subspace \(B_0\) of codimension at most 4 in both designated complementary hyperplanes (that is, contained in both and of codimension at most 4 in \(V\)), disjoint from the span of both image lines; take a complement \(C_0\) containing both lines. The supported \({\rm SL}(B_0)\), fixing \(C_0\), belongs to \(K_i\cap K_j=H\). Thus every designated line in these patterns lies in \(C_0\) and every corresponding hyperplane contains \(B_0\), by large-dimension irreducibility. Both the total line span and the total span of dual functional lines have bounded dimension. Now a boundedly long list of independent pattern points already contradicts strict descent of intersections. Indeed descent must be strict also on intersection with \(S\), since \(HS=M\) and all groups contain \(H\). Track the growing two spans along the list; they each increase only boundedly often. Between increases, invariance is already imposed on the current spans, and further constraints are linear conditions on their endomorphism spaces (restrict matrices to the first and their transposes to the second). Each strict descent shrinks the solution subspace, by testing a witness matrix above the descent in \(S\). This bounds the length absolutely.

These exclusions imply irreducibility of the collineation part: if it preserves \(U\), then with anti-actions also \(aU\); the full group in question preserves these two with signed exchange. Nonzero intersection gives a signed flag (intersection and sum), and zero intersection a stabilized idempotent.
\end{proof}

\begin{proposition}[Singular and nondegenerate subspaces]\label{c:isometry-subspaces}
In isometry type, the initial round excludes stabilizers of nonzero totally singular subspaces, proper nonzero nondegenerate subspaces, and nonsingular radical lines in orthogonal characteristic two. It requires no prior anchor. The resulting good groups act irreducibly by semilinear transformations.
\end{proposition}

\begin{proof}
For isometry types first exclude stabilizers of nonzero totally singular subspaces \(U\) (totally isotropic for unitary or symplectic). The standard unipotent radical is again in \(H\). Here are details of the resulting comparability. Use Witt coordinates \(U\oplus Z\oplus U^-\) with last coordinate \(\lambda\) paired dually with \(U\). The radical acts by
\[
 (u,z,\lambda)\mapsto (u+f(z)+h\lambda,\ z+g\lambda,\ \lambda)
\]
with arbitrary \(f\) and accordingly adjoint \(g\), and a form-preserving correction \(h\). Such corrections exist by representing the quadratic correction as the diagonal of a bilinear pairing (orthogonal), filling an alternating correction (symplectic), or by surjectivity of quadratic-field trace on the diagonal in the hermitian case. With \(f=g=0\) the available \(h\) include the symmetric shears in symplectic type, skew-hermitian in unitary type, alternating in orthogonal type. These unipotent groups are inside the defining special groups (the additional determinant/spinor restrictions have order prime to the characteristic in the relevant unitary or odd-characteristic orthogonal cases). In orthogonal characteristic two the Dickson rank parity is even: the rank of the displacement is \(2\operatorname{rank}g\) plus the rank induced by \(h\) from \(\ker g\) to its dual \(U/{\rm im}\,f\), and the latter is alternating by the quadratic isometry condition.

Another invariant totally singular space within \(U^\perp\), if not within \(U\), must contain \(U\) by the arbitrary \(f\). A vector in such an invariant space with \(\lambda\ne0\) outside \(U^\perp\) is impossible in unitary/symplectic type by the pure \(h\)-shears (one can have nonzero pairing with \(h\lambda\)). In orthogonal type it forces the space to contain the annihilator hyperplane of \(\lambda\) in \(U\); also \(Z=0\), else choose a shear with nonsingular \(g\lambda\) and take the difference vector. Thus the only incomparable possibility consists of two split maximal singular spaces adjacent in a hyperplane. At most two distinct maxima can so occur (through their common hyperplane there are just two; a third adjacent to both and not through it contains two vectors with nonzero pairing modulo that hyperplane). The rest are in a chain. Intersections down a descending chain of singular spaces are strict by elementary special linear transvections inside the smallest previous singular space preserving the larger flag, extended dually to isometries (also in \(S\)). This again contradicts the pattern size.

Next exclude stabilizers of proper nonzero nondegenerate subspaces in isometry types. On the two orthogonal parts there are commuting normal supported special groups with soluble quotient over their product, so designate the unique retained side. For sufficiently large retained dimension its supported derived classical group acts irreducibly there, hence distinct retained spaces are pairwise orthogonal by the same mutual argument as above. Intersections strictly decrease until possibly the last (use the large irreducible supported group on the remaining orthogonal space). For bounded retained dimension use a cross-chart pair. Their two small spaces embed in a nondegenerate subspace of bounded dimension (complete the radical of their span by dual partners); the supported large irreducible classical group on its complement lies in \(H\). All the small spaces in the patterns thus have bounded total span, and strict descent for independent pattern points is bounded by the growing-span linear-constraint argument above. Use this same small-space argument for the stabilizers of nonsingular radical lines in orthogonal characteristic two. These exclusions ensure ordinary semilinear irreducibility: an invariant degenerate space has an invariant radical, totally singular except possibly in that quadratic case, where the kernel of the quadratic form on the radical is totally singular and otherwise the radical is a nonsingular line.

All initial categories needing more than height-boundedly many labels required no previously good anchor.
\end{proof}

\subsection{Direct systems of imprimitivity}

\begin{convention}[Direct systems after the subspace exclusions]\label{c:direct-configuration}
Next exclude stabilizers of sets of complementary nontrivial projectors giving a direct decomposition (at least two spaces), orthogonal nondegenerate decompositions in the isometry cases. Use the irreducible previously good \(E^+\). It is transitive on the blocks, all of common dimension \(k\), and a block stabilizer in \(E^+\) is semilinearly irreducible on its block (transport any supposed subspace). In isometry type the blocks consist of at most two isometry classes, with independent alternating groups on each class induced by the inner stabilizer: permute isometric copies, then take commutators. Similarly there is at least the alternating permutation group in linear type. If two isometry classes occur, they have equal sizes and the class partition is preserved by the full stabilizer (this uses only finite-field form types and semilinear similarities).
\end{convention}

\begin{lemma}[Two rich permutation actions]\label{c:orbit-bound}
Suppose a finite group has two permutation actions, each preserving a partition into at most two classes and inducing at least the independent alternating groups on its classes. The number of orbits on ordered pairs of pairs, with one index from each action in each pair, is bounded absolutely. Replacing indices by equivariant fibers of absolutely bounded size preserves this conclusion.
\end{lemma}

\begin{proof}
Here and below we use an orbit bound. Suppose a group has two permutation actions, each preserving a partition into at most two classes and inducing on the points at least the independent alternating groups on its classes. Then the number of orbits on ordered pairs of pairs (one index from each action per pair) is bounded absolutely. Indeed choose a size gap among the at most four class sizes, with all small sizes bounded and every large size at least 8 with alternating order exceeding the product of the quotient-index bounds from the small factorials, class permutations and parities (choose five successive absolute thresholds rapidly increasing, or more to leave a gap). Intersect the joint action with the product of all large alternating factors. This is normal of bounded index, projecting onto each large alternating factor by the index comparison and simplicity (one can restrict for this check to the preimage of the single alternating factor on its action side, whose image there is onto by hypothesis). By the subdirect description these actions split as independent diagonals, with identifications natural up to relabeling (automorphisms of alternating groups in large degree). The orbits on the specified positions are then bounded by equality patterns and small-class choices. Bounded equivariant fibers keep the bound absolute.
\end{proof}

The orbit bound will control the number of invariant partitions, rather
than the number of points. Suppose a group has at most $C$ orbits on
$\Pi\times\Pi$, and $\Omega\to\Pi$ is an equivariant map with fibers of
size at most $d$. Over a representative of each orbit on $\Pi\times\Pi$
there are at most $d^2$ ordered pairs, so there are at most $Cd^2$ orbits
on $\Omega\times\Omega$. An invariant equivalence relation is a union
of these orbits. Consequently
\begin{equation}\label{c:partition-bound}
 \#\{\text{invariant partitions of }\Omega\}\le 2^{Cd^2}.
\end{equation}
In the applications, $\Omega$ consists of common refinement cells or
tensor factors. We must recover each original system by grouping those
actual objects before this bound can count its full stabilizer labels.

\begin{lemma}[Direct line systems]\label{c:direct-lines}
Under \Cref{c:direct-configuration}, the full stabilizers of direct line systems are excluded. The only possible distinct linear systems reduce to matchings over \(\mathbb F_3\); the corresponding odd-characteristic orthogonal case also reduces to matchings. In both cases an absolute orbit bound contradicts the pattern-line size.
\end{lemma}

\begin{proof}
First exclude line systems, \(k=1\). In linear type the determinant-one diagonals in \(S\) are normal soluble in a label, thus in \(H\). A line of another system with support \(s\) on this one has diagonal orbit size \((|F|-1)^{s-1}\) (divide by \(\gcd(v,|F|-1)\) for full support), at most \(s\) by independence. For \(|F|=2\) the monomial collineations would preserve the all-ones vector line contrary to irreducibility. Thus distinct systems force \(|F|=3\) and support two throughout (\(v\) large and use \(E^+\)-transitivity). They then pair the lines of the first system by a perfect matching, putting both the sum and difference lines on each pair, uniquely. Indeed each support pair has orbit spanning its full plane, so distinct pairs are disjoint. In unitary type use determinant-one norm-one diagonals similarly; the base \(\sqrt{|F|}+1\ge3\) allows only support one. In odd-characteristic orthogonal type use sign flips even in each norm-square class, lying in \(\Omega\), again a normal soluble group. Orbit size at least \(2^{s-3}\) forces \(s\le5\); each occurring type is large, allowing arbitrary support signs, so actually \(s\le2\). Distinct systems again give matchings, uniquely per matching since the two lines on a pair must be orthogonal: they are generated by \(x+y,x-y\) from the two first-system lines with equal nonzero norms, forcing matching within one square class and the relative identification up to sign. There are no nondegenerate lines for alternating polar form. In the matching cases fix two pattern points \(i,j\) on a pattern line. Their intersection group contains alternating permutations of matched configurations independently in each of the (at most two) isometry types, so has a subgroup preserving both with absolutely bounded number of orbits on ordered pairs of first-system indices. Every other system on the pattern line gives a matching invariant by that subgroup. This bounds the number absolutely, contradiction for large lattice \(q\).
\end{proof}

\begin{lemma}[Supported block groups and exceptional planes]\label{c:direct-block-groups}
For direct blocks of dimension \(k\ge2\), the supported special groups described below give a normal projective product \(Y\), with retained block centers. Modulo those centers its nonsoluble parts are simple powers, with at most two simple sites per block. The orthogonal-plane cases in which the supported rotation group has order at most two are either already impossible or are excluded by the graph-isometry matching argument below.
\end{lemma}

\begin{proof}
We explain the supported inner group choices for \(k\ge2\). Use products across blocks of their natural special groups, actually supported (for orthogonal spaces we can take \(\Omega\) in the sense of the determinant-one spinor kernel in odd characteristic, and the kernel of the Dickson parity in even characteristic; thus in split four-space over two we use the group with the two independent natural \({\rm SL}_2\) tensor actions); replace \({\rm Sp}_4(2)\) by its simple derived group. Denote the resulting projective product by \(Y\trianglelefteq K\), and include its normal block-center subgroup among the soluble normals retained by (\(*\)). Except for soluble cases, modulo these centers \(Y\) is a direct product of isomorphic nonabelian simple groups, one per block, or two per block in orthogonal split dimension four. We recall details to control small cases. The standard natural special modules are absolutely irreducible in these dimensions except the orthogonal planes. This follows also from the usual natural modules in defining characteristic (orthogonal \(4^-\) uses the two Frobenius-conjugate natural \({\rm SL}_2\) factors tensoring after scalar extension, \(4^+\) the external tensor; orthogonal odd dimension here is in odd characteristic). For \({\rm Sp}_4(2)'\cong A_6\), restriction stays irreducible by normality in \({\rm Sp}_4(2)\) and the orders too large for \({\rm GL}_3(2)\); the end field cannot enlarge by the order in \({\rm GL}_2(4)\) either. Soluble low exceptions include natural \({\rm SL}_2\) (also \({\rm SU}_2\) or \({\rm Sp}_2\)) at parameters 2,3, \({\rm SU}_3(2)\) (its cycle and determinant-one cube-root diagonals act absolutely irreducibly), \(\Omega_3(3)\), and \(\Omega_4^+\) at 2,3; the indicated irreducibility still holds. For an orthogonal plane \(\Omega\) is the rotation group of order \((|F|\mp1)/\gcd(2,|F|-1)\), acting semisimply without fixed space or invariant nondegenerate line when that order exceeds two.

\textit{Small orthogonal planes.} For order at most two in those planes: split over two gives a canonical nonsingular vector per block, impossible by irreducibility; split over three gives a canonical orthogonal nonsingular line refinement (the two nonsingular lines in the split plane), now excluded. The remaining cases (split over five and nonsplit over three) have individual supported signs \(z_a=-1\) on block \(a\) in \(H\). Treat patterns of these cases separately. If two systems are compatible (commuting projectors), distinctness and transitivity give a common orthogonal line refinement, impossible. Otherwise some sign of one moves blocks of the other (if they all fix them the projections split the blocks). By rank two it transposes precisely two and is identity on the rest. Those two planes must be graphs \(\pm f\) of an isometry from the sign block \(U_a\) to a plane in its orthogonal complement: their sum contains \(U_a\), they project isomorphically in both directions there, and they are orthogonal. Every other sign preserves the graph pair by commutation with \(z_a\), so the partner plane is precisely some \(U_b\) (each sign outside \(a\) acts on it as a scalar). Transitivity gives a perfect matching and a graph-isometry choice on each edge. The intersection group for two fixed systems contains alternating permutations of identically aligned matched configurations (all are isometric). Every further system on the pattern line must give a matching invariant by this group, boundedly many possibilities, and boundedly many isometry choices because the field here is bounded and only one choice per edge orbit needs specifying. The edge orbit bound again follows from bounded rank on the block indices. This contradicts large lattice \(q\).
\end{proof}

\begin{lemma}[Supported groups in a prepared anchor]\label{c:supported-anchor-groups}
For every remaining direct system, each block \(a\) has an actually supported matrix subgroup \(R_a\), projectively contained in \(E^+\), of order greater than two. It is normalized semilinearly by the block stabilizer of \(E^+\). Its block module is semisimple and has no trivial constituent. This includes the case in which the supported simple power is missing from the base.
\end{lemma}

\begin{proof}
In the remaining nonexceptional direct systems, there exist for every block \(a\) actually supported matrix subgroups \(R_a\) projectively in \(E^+\), normalized semilinearly by the block stabilizer in \(E^+\), of order greater than two acting nontrivially there. For soluble \(Y\) or retained \(Y\), this follows from the full supported choices above. If \(HY=K\) is missing, \(H\) induces the full top permutation image (on blocks and simple sites), since \(Y\) acts trivially there. If two simple-site orbits occur (two per block), at most one orbit product is missing modulo centers, and the other supplies supported groups as required. Otherwise there is a transitive orbit, and we have the exact interval over \(P\) in the \(H\)-invariant subgroups of the simple power \(Y/(\text{block centers})\), where \(P\) is the base intersection image. In fact after the centers have been quotiented this transitive normal simple power is minimal normal, hence meets the base core trivially (being missing) and maps to the monolith in the core quotient. By \Cref{f:prop-interval-cases}, either \(P\) is a nontrivial product and all overgroups here are products, or all nonbottom labels are subdirect. The latter contradicts the height: strict comparable subdirect groups give strict refinement of invariant ties, whose chain length is absolutely bounded by the top orbit bound (simple sites have fibers of size at most two over blocks). Thus \(E>H\) gives an intersection product with all projections strictly larger than the corresponding nontrivial projections of \(P\). Lifting with the retained centers supplies the required singly supported \(R_a\) (one can take the full available preimage on that block in the specified supported group, including both sites if they are permuted there). Indeed all centers lost in passing from the supported choices in \(Y\) to the simple-site images have projective images in \(H\); so independent tuples in this intersection product lift with their supported representatives projectively in \(E\). Conjugating by a semilinear collineation of \(E^+\) preserving the block preserves the supported matrix choice as well, acting on representatives which are the identity on the other blocks. In all cases the block module for \(R_a\) is semisimple with no trivial constituent, since its module socle and its fixed space are invariant under the semilinearly irreducible block stabilizer.
\end{proof}

\begin{lemma}[Compatibility of direct systems]\label{c:direct-compatibility}
Any two remaining direct systems in one homogeneous pattern have commuting projectors and hence a common direct refinement.
\end{lemma}

\begin{proof}
Any two systems in such a pattern must be compatible. If all \(R_a\)'s of one fix the other's blocks individually, each projection onto \(a\) of such a block lies within the block (use differences and absence of a trivial quotient), giving compatibility. Thus incompatibility would force movement each way by supported matrices. A supported matrix moving blocks has displacement rank at most its support-block dimension and at least the moved block dimension. Hence both dimensions are equal and equality holds. On the moved system it then transposes just one pair of blocks with square the identity, and acts identically elsewhere (from the fixed dimension). Its displacement has full rank on the support block, so an involution there is impossible in characteristic two and can only be \(-1\) in odd characteristic. But a group of order greater than two cannot have only one element inducing a nonidentity block permutation, contradiction.
\end{proof}

\begin{proposition}[Completion of direct-system pruning]\label{c:direct-pruning}
The remaining direct systems are excluded, both when their supported normal product is retained and when it is missing. The bounds are absolute, independently of the ground field and natural degree.
\end{proposition}

\begin{proof}
Separate retained and missing colors. For retained systems, any two compatible systems must refine one another in some direction: incidence degrees are constant on each block side by transitivity, and if both exceed one, each touching block extends outside, hence is fixed setwise by the other's supported group there (fixed outside points). This gives an invariant proper nonzero block intersection, nondegenerate when relevant, impossible by the retained actions above. A coarse block of large dimension cannot admit the resulting uniform finer split under its full supported special group by orders. Uniformly, with base size \(s_0=|F|\) (\(\sqrt{|F|}\) in unitary type), and \(c=1\) in linear/unitary, \(c=1/2\) in orthogonal/symplectic, the large \(m\)-dimensional special group order is at least \(s_0^{c m^2-O(m)}\), whereas the isometry (or linear) stabilizer of a corresponding uniform \(t\)-block direct split is at most \(t!s_0^{c m^2/t+O(m)}\) by standard orders, smaller uniformly for large \(m\) and \(2\le t\le m\). Bounded coarse dimensions permit only boundedly many comparable refinements. This excludes the retained patterns.

For missing systems take all points of a pattern line. Compatibility
gives their common refinement into nonzero fine cells. Fix two systems
$i,j$ and group the fine cells by their pair of block indices. Each
aggregate cell has a supported special group inside $K_i\cap K_j$.
If its dimension is large, no further system can split it: a splitting
block extending outside the cell would give a proper invariant
intersection under the supported group, whereas a block contained in
the cell would give a uniform finer split contrary to the same order
bound. A small aggregate has at most its dimension many nonzero fine
cells. Thus the number of fine cells over each occupied pair is bounded
by an absolute constant $d$.

The map from a fine cell to its block-index pair is $H$-equivariant.
Indeed $H$ stabilizes both systems and permutes their common refinement;
anti-actions also permute products of commuting projectors. The two rich
top actions and \Cref{c:orbit-bound} give an absolute bound $C$ on the
orbits of ordered pairs of occupied pairs. Each original system is
recovered by grouping the fine cells and taking their direct sums.
Distinct full stabilizer labels therefore give distinct $H$-invariant
partitions of the fine cells. By \eqref{c:partition-bound} there are at
most $2^{Cd^2}$ such partitions, contradicting the chosen pattern-line
size. This completes direct-system pruning.
\end{proof}

\subsection{Tensor geometry and order comparisons}

\begin{lemma}[Ordered factors and their forms]\label{c:tensor-forms}
Let \(A\) be a unital proper nonscalar central simple \(F\)-subfactor of \(\mathcal A=\operatorname{End}_F(V)\), normalized by the current previously good anchor. Then \(V=X\otimes_FY\), with \(A=\operatorname{End}(X)\otimes1\) and commutant \(A'=1\otimes\operatorname{End}(Y)\). Both factors have dimension at least two. If \(A=A^*\), the ambient form is the tensor of factor forms. In bilinear characteristic two the factor forms are alternating; an orthogonal quadratic form is the canonical form with that polar pairing and zero on pure tensors.
\end{lemma}

\begin{proof}
We next treat subfields and ordered tensor structures. All algebras below inside \(\mathcal A\) are unital \(F\)-algebras (``prime'' extension or field algebra in these exclusions refers to prime relative degree, not to the characteristic field). Recall that a central simple matrix subfactor \(A\) gives
\[
 V=X\otimes_F Y,\qquad A={\rm End}(X)\otimes1,\qquad A'=1\otimes{\rm End}(Y)
\]
with \(A'\) its commutant (central simple algebras over finite fields split). For \(A\) proper and nonscalar both factor dimensions are at least two. Linear normalization acts by tensor matrices up to scalars, by the matrix algebra automorphism theorem; similarly with field automorphisms. If \(A=A^*\), both factors have defining forms, with \(b\) their tensor form up to a scalar: use the restrictions of the matrix involution to the two factors. These are both hermitian for unitary type, and each symmetric/alternating bilinear otherwise. Semilinear similarities preserving the ordered tensor act by semilinear similarities on the factors. In bilinear characteristic two, an ordered factor normalized by such an irreducible \(E\) cannot have nonalternating form (the diagonal contraction has an invariant hyperplane, then tensor with the other space). With both alternating, any orthogonal quadratic form in question must be the canonical one for the polar tensor, vanishing on pure tensors. Such a quadratic form exists (impose zero on basis tensors and expand, using alternation on both factors: for a pure tensor the potentially nonzero pair terms have both indices different and cancel by interchanging which second indices are matched to the two first indices); it transforms under \(E\) like the given one, and a nonzero difference, the square of a linear functional, contradicts irreducibility. We will use these facts whenever there is a common previously good anchor.
\end{proof}

\begin{lemma}[Uniform tensor order obstruction]\label{c:tensor-order}
For sufficiently large dimension \(m\), a full supported perfect natural classical group cannot normalize a proper nonscalar tensor factor or a system of at least two nontrivial tensor factors, even allowing permutations of the factors.
\end{lemma}

\begin{proof}
Useful also is the uniform order obstruction: a full supported perfect natural classical group on a space of sufficiently large dimension \(m\) cannot normalize a proper nonscalar central simple tensor factor, or a tensor system splitting the space into at least two nontrivial slots even up to permutations. Its order (also modulo scalars) is at least \(|F|^{m^2/2-O(m)}\); the linear tensor-system order is bounded by
\[
 t!\, |F|^{\sum_j x_j^2},\qquad \prod_{j=1}^t x_j=m,\quad t\ge2,\quad x_j\ge2 ,
\]
which is uniformly smaller (\(t\le\log_2 m\), \(\sum x_j^2\le m^2/4+4\), e.g. by merging factors). Use the linear action of the covering matrix group for this comparison. This proves also properness of such stabilizers in the large ambient group. Properness for field algebras follows directly from perfection (fix them pointwise if normalized by the socle) and absolute irreducibility on the natural module.
\end{proof}

\subsection{Field and ordered-factor exclusions in linear type}

\begin{proposition}[Prime-degree field algebras in linear type]\label{c:linear-fields}
After the direct-system exclusions, the full normalizers of field algebras \(D/F\) of prime relative degree are excluded in linear type, including when the ambient group contains anti-actions.
\end{proposition}

\begin{proof}
Exclude stabilizers of field algebras \(D_i/F\) of prime degrees \(r_i\). Their soluble normal projective torus \({\bf P}(D_i^*)\cap S\) is in \(H\), of order
\[
 \frac{|F|^{r_i}-1}{|F|-1}\ \big/\ \frac{\gcd(v,|F|-1)}{\gcd(v/r_i,|F|-1)}
\]
by the norm determinant. This exceeds \(r_i\) except possibly \(|F|=3,r_i=2\), when it is at least two. If a pair of fields do not commute, each torus injects by normalization into the other's relative field Galois group (a nontrivial torus element has lifts generating its prime field extension). This forces the exceptional case on both sides. If two fields commute, their generated algebra is commutative semisimple, necessarily a field since the primitive idempotent system would otherwise contradict the direct-system exclusion at \(E\). Distinct commuting fields therefore have distinct prime degrees. In a nonexceptional line take their compositum; the intersection group contains the natural determinant-one group over it, perfect if module dimension is at least two, which forces any third prime field on the line into the compositum by centralization (use elementary matrices). If dimension one, pairwise commutation alone does the same. This contradicts distinctness. In a line of the exceptional type, distinct fields must be noncommuting. They have generators \(J_i^2=-1\) with projective class in the indicated torus (norm one), and pairwise normalization gives anticommutation. Two generate a split quaternion algebra \(M_2(F)\); their intersection group includes the large determinant-one group on its tensor complement. Every line field lies in that bounded pair algebra by centralization, impossibly many. We separate the types by color here.
\end{proof}

\begin{proposition}[Ordered factors in linear type]\label{c:linear-factors}
After \Cref{c:linear-fields}, the full stabilizers of proper nonscalar ordered central simple factors are excluded in linear type.
\end{proposition}

\begin{proof}
Now exclude ordered proper nonscalar central simple factors \(A\). The projective supported special linear groups on both tensor spaces are commuting normals in the label, even allowing anti-actions, with soluble quotient. Exactly one is retained; designate its factor algebra. Its natural projective group is simple except dimension two over two or three. Over two, the characteristic order-three subgroup would give an invariant quadratic field, impossible. Over three, the \(A_4\) has a unique Klein four subgroup, normal and spanning the factor algebra by its quaternion lifts. Treat this last color separately from the simple case. For a pattern of designated algebras the groups are normal in \(H\); intersection nontrivial gives equality (or equality of the Kleins in the exceptional case), hence equality of algebras, since their lifts span. Otherwise they commute projectively and also linearly on special-group lifts (perfection, or no scalar characters for \({\rm SL}_2(3)\)). Thus distinct designated algebras commute elementwise. For a pair on a line the new algebras lie in the centralizer matrix factor, on whose space the intersection contains the full supported special linear group. In large dimension this prohibits any proper new tensor subfactor there by orders (the whole centralizer factor then permits at most one distinct new designated algebra); in bounded dimension only boundedly many distinct commuting nonscalar central simple factors can fit there anyway (their tensor-product matrix algebra embeds there by simplicity, so the product of degrees is bounded). Contradiction.
\end{proof}

\subsection{Partial ordered-factor pruning in isometry type}

\begin{proposition}[Ordered factors apart from two retained plane types]\label{c:partial-isometry-factors}
In isometry type, after the direct-system exclusions, the full stabilizers of ordered proper nonscalar selfadjoint central simple factors are excluded except possibly when the retained support is an odd-characteristic bilinear orthogonal plane or a symplectic plane over \(\mathbb F_2\). These exceptions concern the retained side of the ordered factorization.
\end{proposition}

\begin{proof}
First try to exclude ordered proper nonscalar selfadjoint central simple factors \(A\), leaving two possible remaining types temporarily. In the stabilizer the natural projective special groups supported on \(X,Y\) are commuting normals with soluble quotient. Again use \({\rm Sp}_4(2)'\) in its small case. The supports indicated lie in \(S\). Besides perfection in the usual cases, note unitary determinant one, and in orthogonal type: in characteristic two both factors are alternating as above and the tensor action from a factor has even displacement rank; in odd characteristic \(\Omega\) on a symmetric factor lies in \(\Omega\) on the tensor (diagonalize the other symmetric factor and use orthogonal sums; scaling does not change the spinor norm on determinant one). If both factors alternating in odd characteristic, the symplectic supports also lie projectively there by perfection or absence of a 2-quotient for \({\rm SL}_2(3)\). Quotient solubility follows by the two tensor projections and the standard similarity groups. Exactly one support is retained, and we designate its factor algebra.

We exclude the labels for which this retained type is neither a bilinear orthogonal plane in odd characteristic nor a symplectic plane over two. Here are the details of the supported projective groups within the designated algebra:

\begin{itemize}
\item Usually it is a single simple group with perfect special lift and absolutely irreducible spanning natural action. In split orthogonal dimension four (odd characteristic here) there are two projective \({\rm PSL}_2\) components on symplectic slots, with their matrix algebras tensoring to the factor. They are permuted by the label, even at parameter three: in the determinant-form model the singular rank-one quadric and its two rulings give these two tensor slots up to exchange (the maximal linear rank-one spaces are precisely of the two fixed-line kinds). Indeed the endomorphisms preserving each member of a ruling individually form the matrix algebra on the other coordinate (use two independent first-coordinate lines and their sum-vector line). Nonsplit four has the standard simple projective group with absolutely irreducible tensor-twist action instead.
\item The degree-two symplectic or special unitary groups at parameter three (also the symplectic slots just described at three) give projective \(A_4\) with Klein lifts spanning the slot and special lift \({\rm SL}_2(3)\) having no ground-field scalar characters. Indeed the Klein preimage in the rank-one special group at three is quaternion of order eight with commutator \(-1\) acting as negative identity in degree two, thus absolutely irreducible there (semisimplicity for this 2-group, no linear constituent possible), and the abelianization of \({\rm SL}_2(3)\) has order three whereas the ground scalar group here has order two or eight. A degree-two unitary factor at parameter two is impossible: its characteristic order-three subgroup in the projective special has two orthogonal nondegenerate eigenlines over four, giving a forbidden nondegenerate direct decomposition on tensoring. An orthogonal 3-space over three is likewise impossible: for the form up to similarity a sum of three squares, the three norm-one lines are the axes preserved as a set even by similarities, again a direct decomposition.
\item For the special unitary 3-space over four (parameter two) the projective special \(C\) of order 72 has a unique normal Sylow 3-group \(D_9\) elementary of order nine, contained in every nontrivial normal subgroup, and its lifts span the slot. Indeed in an orthonormal basis use the cycle and \({\rm diag}(1,\omega,\omega^2)\) for a cube root \(\omega\ne1\). Their commutator generates the scalar special center; projectively they give \(D_9\). In the projective full unitary group of order 216, the Fourier matrix and \({\rm diag}(1,1,\omega)\) normalize it inducing generators of \({\rm SL}_2(3)\) (rotation and shear). Thus the whole 216 normalizes it with centralizer the nine. The special image contains the nine and induces order eight acting irreducibly there (a line stabilizer in \({\rm SL}_2(3)\) has order six). Thus the nine is minimal normal and a normal subgroup disjoint from it must centralize it, proving the assertion. Its linear special lift has no character to the scalars of order three by this action and the center commutator.
\end{itemize}

These lists follow from the same natural classical groups as for direct blocks, now with odd characteristic on symmetric factors. The individual component matrix algebras here are selfadjoint.

Color according to simple components, \(A_4\) components, or the single \(C\). The projective components lie in \(H\), which permutes at most two of them in each system. In a simple or \(A_4\) pattern they cannot swap the other's pair (no quotient of order two), so components across systems mutually normalize. A nontrivial intersection, normal in each, gives equality or common Klein, hence the same component algebra; otherwise they commute projectively and linearly on special lifts as above. Likewise for the \(C\) pattern, using a common \(D_9\) if intersection is nontrivial (it contains the normal Sylow on each side, which must then be the same). Thus all the distinct component algebras are commuting tensor slots. For two labels \(i,j\) on a line their components (at most four) have a tensor centralizer factor, itself normalized by \(E\); any new component lies there. In the pair intersection is the large supported perfect classical group on that centralizer space if of large dimension, by the tensor form observations. It must stabilize individually the at most two components for each further point of the line. The order bound prohibits a proper nontrivial central simple subfactor there being new (the whole factor itself would add only one); if of bounded dimension only boundedly many commuting new ones fit. Thus the list of distinct component algebras on the whole line is absolutely bounded (at most four from the original two labels plus those new ones). Each designated algebra consists of one or two such components tensoring and determines the stabilizer (it and its tensor commutant determine the ordered tensor). Only boundedly many labels then fit on the line, contradiction. Consequently any remaining previously good-normalized proper ordered selfadjoint factorization must have the \textbf{retained} side one of the two remaining plane types.
\end{proof}

\subsection{Transfer to normalized field algebras}

The remaining field obstruction is commutation. In linear type the
retained scalar tori nearly settled this directly. In isometry type
we instead use classical groups over the extension fields. The next
setup identifies the algebra whose center determines which transfer
argument is needed.

\begin{convention}[A pair of field witnesses]\label{c:field-pair-setup}
Let $D_i,D_j$ be selfadjoint field algebras of prime relative degree,
with distinct full normalizer labels $K_i,K_j$ on a pattern line, and
put $T=K_i\cap K_j$. The algebra $R=F[D_i,D_j]$ is selfadjoint
semisimple: its radical is normalized by the common irreducible $E$
and must annihilate the space. Its primitive central idempotents are
either just one or exactly two exchanged by $*$, since sums on
$*$-orbits are orthogonal projectors contrary to the earlier
direct-system exclusion if there is more than one orbit.
\end{convention}

When the center is a field larger than $F$, or has two exchanged
components, we will find a perfect group inside $T$ whose centralizer
forces $D_i$ and $D_j$ to commute. In the scalar-center case we will
rule out $R=\mathcal A$ and use the resulting ordered tensor factor
to return to the two retained plane types.
We first establish the form and centralizer calculations needed for
these alternatives. The two-dimensional unitary exception in the
centralizer calculation will be handled by a norm-one scalar count;
exchanged components require the paired special linear group below.

\begin{lemma}[Transferred forms]\label{c:field-transfer}
Let \(D/F\subseteq\mathcal A\) be a normalized selfadjoint field algebra, of arbitrary finite degree, and put \(\delta=*|_D\). In the prepared irreducible-anchor configuration there is a nondegenerate transferred form \(B\) satisfying
\[
 b(dx,y)=\operatorname{Tr}_{D/F}\bigl(dB(x,y)\bigr).
\]
The associated bilinear, hermitian and quadratic structures, and their transformations under the normalizer, are as described below. The perfect natural special group of the transferred form, when available, lies in \(S\).
\end{lemma}

\begin{proof}
For a normalized field (degree possibly composite), let \(\delta=*|_D\). There is a nondegenerate transferred form \(B\) given by
\[
 b(dx,y)={\rm Tr}_{D/F}(d B(x,y)).
\]
It is linear over \(D\) on the first slot and \(\delta\)-semilinear on the second. It is symmetric/skew or hermitian/skew-hermitian as appropriate (skew-hermitian may be rescaled). In particular if \(b\) is alternating and \(\delta=1\) then so is \(B\), also in characteristic two by testing all square scalars. For orthogonal characteristic two and \(\delta=1\) transfer \(Q(dx)={\rm Tr}_{D/F}(d^2 Q_D(x))\), valid since the left side is additive in \(d\) (the cross term vanishes by alternation of \(B\)) and scales by squares over \(F\), so gives a linear function of \(d^2\) over \(F\); testing the trace against these same squares gives \(D\)-quadratic homogeneity of \(Q_D\) and polar \(B\). For orthogonal characteristic two and \(\delta\ne1\) use instead
\[
 Q(x)={\rm Tr}_{D_0/F}(B(x,x)),\qquad D_0=D^\delta .
\]
Indeed the right side has polar \(b\) and transforms the same way under any previously good anchor normalizing \(D\). The difference from \(Q\) has zero polar form, hence is the square of a linear functional over the finite, therefore perfect, ground field. If nonzero, its kernel would be a proper nonzero invariant subspace for that anchor. Irreducibility therefore makes the difference zero. All form transformation assertions follow from trace uniqueness under semilinear maps normalizing \(D\). For instance for a similarity \(u\) with multiplier \(\mu\) on \(b\) and induced field power \(\tau\) on \(D\), testing at \(d^\tau\) gives \(B(ux,uy)=\mu B(x,y)^\tau\). Thus the stabilizer acts projectively by corresponding \(D\)-semilinear similarities (extra scalar kernel soluble), and the standard \(D\)-linear special group, when perfect, lies inside \(S\).
\end{proof}

\begin{lemma}[Centralizers of perfect transfer groups]\label{c:transfer-centralizer}
Suppose the field normalizer in \Cref{c:field-transfer} is nonsoluble and the total ground-field dimension is sufficiently large. There is a perfect natural classical group \(P_D\), absolutely irreducible over \(D\). Its centralizer in ground-field endomorphisms is exactly \(D\), except possibly for a unitary \(D\)-module of dimension two with \(\delta\ne1\) fixing \(F\). This exception does not occur for prime extension degree in sufficiently large total dimension.
\end{lemma}

\begin{proof}
If this normalizer is nonsoluble and total ground-field dimension large, we can use a perfect such classical group \(P_D\) on the natural \(D\)-module, absolutely irreducible over \(D\). Indeed either module dimension or extension degree/field size is large, avoiding the small parameter exceptions; in the permanently soluble form dimensions nonsolubility of the normalizer would be impossible. Its centralizer in ground-field endomorphisms is exactly \(D\), except possibly the unitary module of dimension two over \(D\) with \(\delta\ne1\) fixing \(F\). To see this, extending from \(F\) to algebraic closure splits the module into copies twisted by the different embeddings of \(D\) over \(F\). They are irreducible and pairwise inequivalent since matrix traces generate \(D\) over \(F\): symplectic type uses an \({\rm SL}_2\) block; orthogonal type can use split 4-space with tensor trace there an arbitrary product, or norm traces from the \(4^-\) tensor twist, and in dimension three the symmetric-square traces (odd characteristic). Unitary type has traces generating at least \(D_0\) via \({\rm SU}_2\), and in dimension at least three a norm-one diagonal with entries \(a,b,(ab)^{-1}\), where \(a,b,ab\ne1\), has trace differing from inverse trace, giving all of \(D\). In dimension two, \(D_0F=D\) if \(\delta\) is nonidentity already on \(F\). Here on a split nondegenerate orthogonal 4-space use the two \({\rm SL}_2(D)\) tensor actions, supported there (available also within higher dimensions of Witt index at least two); in the nonsplit dimension-four case the trace on the tensor twist ranges over norms of arbitrary \({\rm SL}_2\)-traces, and in dimension three the values \(u^2-1\) from the symmetric square generate the field in odd characteristic. For \({\rm SU}_2\) the natural action over \(D\) realizes the split rank-one traces over \(D_0\). The indicated three-entry unitary diagonal exists even when the norm-one group has order three (take \(a=b\)); the difference of its trace from inverse trace, multiplied by \(ab\), is \(-(a-1)(b-1)(ab-1)\). Thus the embedding distinction (and absolute irreducibility in each copy, as for natural classical groups above) gives scalar centralizer on each over closure whenever claimed, hence the centralizer assertion over \(F\). In particular the exception does not arise for prime degree in sufficiently large ground-field dimension (it would require degree two and module dimension two).
\end{proof}

\begin{lemma}[Paired transfer spaces]\label{c:paired-transfer}
In the high-dimensional prime-degree setting of \Cref{c:field-pair-setup}, suppose an unordered complementary pair \(e,e^*=1-e\), normalized by \(E\), commutes with \(D_i\). Its two spaces are dually totally isotropic and, in orthogonal type, totally singular. Their simultaneous field-and-pair stabilizer contains the paired group \(\operatorname{SL}_m(D_i)\). In the nonsoluble case this group is perfect in \(S\), and every ground-field endomorphism centralizing it commutes with \(D_i\).
\end{lemma}

\begin{proof}
We will also use the following paired situation (still with the prime-degree \(D_i/F\) and high total dimension). Suppose two central or commuting projections to be used below are exchanged by \(*\) (a complementary pair \(e,e^*=1-e\) normalized as a pair by \(E\)), and commute with \(D_i\). Their spaces are dually totally isotropic, and totally singular also for \(Q\) when relevant (in characteristic two each quadratic restriction with zero polar is a linear square, and the sum of the two kernels would be a proper nonzero invariant space if a restriction were nonzero; the dimensions are large). The full simultaneous stabilizer of \(D_i\) and the pair thus includes paired \({\rm SL}_m(D_i)\) on these spaces, acting by inverse adjoint on the second, where \(m\) is their \(D_i\)-dimension; if \(m=1\) the stabilizer would be soluble. For \(m\ge2\) this group is perfect inside \(S\). Every endomorphism centralizing its matrices commutes with \(D_i\). Indeed over closure the constituents are natural at embeddings \(\tau\) over \(F\) and dual at \(\tau'\delta_i\) with \(\tau'\) over \(F\) (by the transferred pairing); the field scalars act there by \(\tau,\tau'\), respectively. Natural twists are distinguished by traces, as are dual twists; a natural and dual identification could only have \(\tau=\tau'\delta_i\) by traces on an \({\rm SL}_2\) block. It cannot occur for \(m\ge3\) (diagonals \(a,b,(ab)^{-1}\) again; \(|D_i|\ge4\)). If \(m=2\), either the two identified constituents have the same field-scalar action, or \(\delta_i\ne1\) fixes \(F\), forcing prime degree two and bounded total dimension. This proves the claim.
\end{proof}

\subsection{Prime-degree field pruning in isometry type}

\begin{proposition}[Prime-degree selfadjoint fields]\label{c:isometry-fields}
After direct-system pruning and the partial ordered-factor exclusion of \Cref{c:partial-isometry-factors}, the full stabilizers of selfadjoint field algebras of prime relative degree are excluded. The pattern is colored by whether the adjoint restriction on the field is the identity.
\end{proposition}

\begin{proof}
\textit{A larger center field.} If the center of \(R\) is a field larger than \(F\), take a central prime extension \(D'\) invariant under \(*\) and \(T\). If \(K_i\le N_M(D')\), then \(P_{D_i}\) forces \(D'\subseteq D_i\) by centralization, so \(D_i=D'\) commutes with \(D_j\). Otherwise the covering property gives
\[
 T=K_i\cap N_M(D')
\]
exactly. Consider the commutative semisimple \(L_f=F[D_i,D']\). If a field, the degrees must be distinct primes with compositum degree their product, and the perfect transfer group over it is available in \(T\) (the subfields are unique, and \(T\) is nonsoluble since it contains \(E\)). This perfect group centralizes \(D_j\), forcing commutation with \(D_i\) by the centralizer assertion. In the exceptional unitary dimension two case here, the compositum degree is \(2s\) with \(s\) a large odd prime (total dimension \(4s\)). Write \(a=|F|\) and identify the compositum with \(C=\mathbb F_{a^{2s}}\). Work in the actual scalar group
\[
 U=\{z\in C^\times:z^{a^s+1}=1\},\qquad |U|=a^s+1.
\]
The subgroup \(U_S\) whose projective classes lie in \(S\) has index at most two in \(U\): these are actual determinant-one isometries, with at most the spinor or Dickson restriction still to impose. Their classes belong to \(T\). Let \(U_0\le U_S\) be the scalar preimage of the kernel of their action on \(D_j\). Since the matrices are \(F\)-linear, this action has order at most \([D_j:F]\). That prime degree divides \(4s\), so it is either \(2\) or \(s\), and
\[
 |U_0|\ge\frac{a^s+1}{2s}.
\]
The two maximal proper relative subfields of \(C/F\) are \(\mathbb F_{a^s}\) and \(\mathbb F_{a^2}\). Their intersections with \(U\) have sizes at most \(2\) and \(a+1\), respectively: the first satisfies \(z^2=1\), while on the second the oddness of \(s\) turns the norm-one condition into \(z^{a+1}=1\). For sufficiently large \(s\), the displayed lower bound exceeds \(2+(a+1)\). Hence \(U_0\) contains a generator of \(C/F\). Ground-field scalar changes do not affect conjugation on \(D_j\), so this is a count in the scalar preimage throughout. Such a generator centralizes \(D_j\), forcing all of \(C\), and in particular \(D_i\), to commute with \(D_j\).

\textit{A split compositum.} If \(L_f\) here splits, there are exactly two exchanged components, and \(D_i,D'\) must have the same prime degree \(r\) (tensor over \(F\)). The possible primitive components are indexed by relative identifications of the two degree-\(r\) fields; \(E\) acts on these by translations of a cyclic \(r\)-set (the two field automorphisms extend the same base automorphism), transitively on the occupied subset. Thus \(r=2\). Now apply the paired situation to the two components: the indicated paired special linear group also centralizes \(D'\) (on each component that field acts through \(D_i\)). Hence it lies in \(T\), centralizes \(D_j\), and gives commutation with \(D_i\).

\textit{Two exchanged central components.} If the center of \(R\) instead has two exchanged components, use their projection-pair stabilizer in place of \(N_M(D')\). It cannot contain \(K_i\) by perfection and the centralizer of \(P_{D_i}\), so its full intersection with \(K_i\) is \(T\), by covering. The paired situation applies directly, yielding the same commutation.

\textit{Scalar center.} It remains to consider center \(F\). Here \(R\ne\mathcal A\), else \(T\)'s actions on the two fields have trivial projective kernel and soluble image. Thus \(R\) gives a proper ordered selfadjoint factorization, necessarily with one of the retained remaining plane factors by the partial exclusion. If \(R\) itself is that plane matrix algebra, then for bounded field size the intersection contains the high-dimensional perfect classical group on its tensor complement, forcing every line field into bounded \(R\), contradiction. For an odd orthogonal plane factor at large field size the projective supported rotations (from \(\Omega_2\)) have order greater than four, and lie in \(H\) as a normal soluble group for the stabilizer label of \(R\) (which contains \(E\)). Normalizing both \(D_i,D_j\) (quadratic within \(R\)), some nonscalar rotation centralizes both. In the plane matrix algebra this gives their equality (the centralizer of a nonscalar matrix there has dimension two), impossible. If instead \(R'\) (the matrix commutant of \(R\)) is the remaining plane algebra, \(K_i\) cannot preserve \(R\), since then \(P_{D_i}\) would act trivially on \(R'\) by its soluble similarity projection, contrary to its field centralizer. Thus again \(T\) is the full intersection \(K_i\cap N_M(R)\). On the large \(R\)-factor space apply transfer to \(D_i\) there, selfadjoint for its factor form. Its field normalizer there must be nonsoluble (otherwise both factor projections of \(T\) soluble). We get its supported perfect transfer group in \(T\), with centralizer in \(R\) exactly \(D_i\); in characteristic two here use the alternating factor form, the total tensor preserving its given orthogonal quadratic if necessary as explained earlier. This group centralizes \(D_j\subset R\), impossible.

\textit{Completion of the line argument.} Consequently the fields on a pattern line all commute unless we already obtained a contradiction. Two distinct ones of the same prime degree would give a split commutative algebra with exactly two exchanged components as above, degree two, and the two adjoint restrictions must differ (to exchange the relative-identification idempotents). Use homogeneous color of the adjoint type (identity or not on the prime-degree field algebras) to prevent this. Thus two have distinct primes and generate a field. Its transfer group in the nonsoluble intersection forces any third into the compositum, impossible by prime degrees. There is no unitary dimension-two exception here: if the involution fixes \(F\) but not the compositum, the two prime fields have different adjoint types. This completes prime-field pruning.
\end{proof}

\subsection{Paired idempotents and the remaining ordered factors}

\begin{proposition}[Completion of the algebra exclusions]\label{c:remaining-factors}
After \Cref{c:isometry-fields}, the full stabilizers of unordered proper idempotent pairs \(e,e^*=1-e\) are excluded. The two remaining retained-plane types in \Cref{c:partial-isometry-factors} are then impossible, completing the ordered-factor exclusions.
\end{proposition}

\begin{proof}
Exclude stabilizers of unordered pairs \(e,e^*=1-e\) of proper idempotents. For two such labels, choosing idempotents \(e,f\), the generated selfadjoint algebra is again semisimple. A center field larger than \(F\) is now forbidden. If there are two central components exchanged, their projector pair commutes with each of the given pairs. Distinct commuting such pairs are impossible: they would give all four projector cells (or already be the same pair, using adjoint exchange), with forbidden orthogonal direct sums on adjoint orbits. For center \(F\), note \((e-f)^2\) is central, and the algebra is spanned over this scalar algebra by \(1,e,f,ef\), since the relation rewrites \(fe\) in terms of \(e,f,ef\) and scalars and the generators are idempotent. It is therefore \(M_2(F)\), which must have remaining plane type. An orthogonal plane has at most one such pair. For the bounded symplectic plane case, the high perfect tensor-complement group in the intersection fixes individually each further pair on the line, forcing them into the bounded plane algebra. Thus all paired stabilizers are excluded.

Finally the remaining ordered factors themselves cannot occur: an odd orthogonal plane factor has the unique pair of isotropic lines over the closure, giving a forbidden projector pair if split, or a forbidden quadratic field (the diagonal algebra descended to \(F\)) if nonsplit, selfadjoint and normalized by all the factor similarities. A symplectic plane over two gives a normalized selfadjoint quadratic field via its characteristic order-three subgroup. This completes the stated field and ordered-factor exclusions.
\end{proof}

\subsection{Full spanning and elementary normalizers}

\begin{lemma}[Full spanning]\label{c:full-spanning}
After the subspace, direct-system, field, ordered-factor and paired-idempotent exclusions, every nontrivial projective subgroup \(P\le S\) normalized by a previously good anchor satisfies \(F[\text{lifts of }P]=\mathcal A\).
\end{lemma}

\begin{proof}
After the preceding exclusions we have the following full spanning property: any nontrivial projective subgroup \(P\le S\) normalized by a previously good anchor has
\[
 F[\text{lifts of }P]=\mathcal A .
\]
Indeed the radical vanishes by anchor irreducibility, and in isometry type the algebra is selfadjoint. Multiple center components are now impossible by the direct-system exclusions (and the exchanged pair exclusion for adjoints). Center-field enlargement and a proper nonscalar central simple factor are also excluded. Semilinear anti-actions in linear type respect the algebra as well (their group action on lifts uses inverse, harmless here).
\end{proof}

\begin{proposition}[Elementary abelian normalizers]\label{c:elementary-normalizers}
Under the full-spanning conclusion of \Cref{c:full-spanning}, full normalizers of nontrivial elementary abelian subgroups of \(S\) are excluded.
\end{proposition}

\begin{proof}
Exclude full normalizers of nontrivial elementary abelian subgroups \(P\le S\). In a pattern each such subgroup is contained in \(H\le E\) by (\(*\)). Its scalar commutator pairing from linear lifts is nondegenerate by full spanning. Thus it identifies \(P\) with its ground-field character dual (nondegeneracy already forces the exponent prime to divide \(|F^\times|\)). Lifts are linearly independent by their distinct simultaneous conjugation characters, and their linear span is the generated algebra. Hence \(|P|=v^2\) and its projective linear centralizer is itself (match the scalar conjugation character using the pairing, then use full spanning). There cannot be two distinct ones here normal in \(E\): a nontrivial intersection would itself satisfy the same properties, and a trivial intersection would give projective commutation. This excludes the pattern.
\end{proof}

\subsection{Tensor systems and missing normal powers}

\begin{lemma}[Reduction to one simple group per tensor slot]\label{c:tensor-simple-slots}
After the preceding exclusions, a tensor system has transitive equal-dimensional slots and the appropriate tensor forms. Its stabilizer induces at least the alternating group on the slots. To exclude all tensor systems, it suffices to exclude those having a normal supported projective product \(Y=B^t\le K\cap S\), with one nonabelian simple group per slot whose lifts span the slot algebra.
\end{lemma}

\begin{proof}
Now exclude stabilizers of tensor systems of at least two nontrivial matrix factors tensoring to \(\mathcal A\), allowing permutations of slots, with factors individually selfadjoint in isometry type. For systems under consideration the anchor acts transitively, else an orbit product contradicts ordered-factor exclusion. Thus we have \(t\) slots of common dimension \(k\). The tensor forms in isometry type can be aligned to identical forms on the factors up to scalars by transitivity/similarity (finite-field form classes). In bilinear characteristic two all are alternating (by transitivity and alternation on \(V\)); again any orthogonal quadratic here is the canonical one zero on simple tensors by irreducibility. Consequently the stabilizer label \(K\) induces at least \(A_t\) on the slots by taking commutators of actual tensor-slot permutations inside the full linear or isometry group.

We can reduce to the case of a normal supported projective product \(Y=B^t\) inside \(K\cap S\), one nonabelian simple group per slot (embedded projectively), whose linear lifts span each slot matrix algebra. Use the natural special groups as above, including \({\rm Sp}_4(2)'\). Here is the treatment of other possibilities. Small degree-two parameter cases at two and three in the linear, symplectic and unitary families give canonical nontrivial elementary abelian projective groups per slot (order three or the Klein, characteristic in the projective special). Their product is elementary abelian projectively inside \(S\), now forbidden; the same applies for \({\rm PSU}_3(2)\) using \(D_9\). Containment uses the same special supports as for ordered factors. \(\Omega_3(3)\) again gives direct orthogonal line cells from the axes. Odd-characteristic orthogonal planes of split type give pairs of singular lines per slot, whose tensor cells paired with their duals give a forbidden nondegenerate direct decomposition. In nonsplit type use the canonical quadratic fields per slot; their tensor algebra has \(2^{t-1}\) components, each central projector selfadjoint (conjugation on all the fields simultaneously), again forbidden. Split orthogonal four-space factors canonically refine into pairs of symplectic plane tensor slots as above, giving a normalized tensor system of the latter instead. Thus it suffices to exclude the single-simple-per-slot systems (first treating just these labels; the refinement observation then handles the four-space case).
\end{proof}

\begin{lemma}[Retained and missing tensor products]\label{c:tensor-products}
Retained products from \Cref{c:tensor-simple-slots} give at most one system in a pattern. For a homogeneous missing pattern, every \(N_i=E\cap Y_i\) is a nontrivial slotwise product \(\prod_aR_{i,a}\), strictly above the base intersection.
\end{lemma}

\begin{proof}
For a pattern with \(Y\) retained, these products are minimal normal subgroups of \(E\) by slot transitivity. Two must be equal (commutation impossible by perfection and spanning), which recovers the same system by the lifts of the simple factors. For missing patterns \(HY=K\), \(H\) induces the full top action, transitive and containing the natural alternating group. Exactly as for missing direct simple-site products, the interval over \(H\cap Y\) in \(Y\) of invariant subgroups must be the nontrivial product case by the reductions and subdirect tie-chain bound. In particular on a pattern line, for each system \(i\),
\[
 N_i=E\cap Y_i=\prod_a R_{i,a}\trianglelefteq E,
\]
a slotwise projective product strictly above the base intersection.
\end{proof}

\begin{lemma}[Common refinement of missing tensor systems]\label{c:tensor-refinement}
A finite pattern line of the missing systems in \Cref{c:tensor-products} has a common tensor refinement formed from the lift algebras of tuple intersections \(\bigcap_iR_{i,a_i}\). After scalar factors are removed, those algebras tensor to \(\mathcal A\) and group to recover every original slot. They are permuted transitively by \(E\). In bilinear characteristic two their forms are alternating; if an orthogonal quadratic form is present, it is canonical also for every grouping into two nonempty batches.
\end{lemma}

\begin{proof}
Write bars for the quotient by $\operatorname{core}_E H$, and let $W$
be the nonabelian simple-power monolith of $\bar E$. Each $\bar N_i$
contains $W$. The simple factors of $W$ partition among the
$\bar R_{i,a}$: their commutators with $W$ generate $W$ by perfection,
lie in the respective intersections with $W$, and are normal there and
mutually commuting.

We need these assignments to give nontrivial intersections in the actual
group, not only compatible images in the core quotient. Put
$N=\bigcap_iN_i$. Since the $N_i$ are normal in $E$, an iterated
commutator of them lies in $N$; its image contains $W$ by perfection.
Thus $\bar N\supseteq W$. Fix a simple factor $C$ of $W$, and let $a_i$
be its assigned slot in system $i$. Because $R_{i,a_i}\trianglelefteq N_i$
and $N\le N_i$, we have
\[
 [N,R_{i,a_i}]\le N\cap R_{i,a_i}.
\]
The image of this commutator contains $C$. Each $N\cap R_{i,a_i}$ is
normal in $N$, so a second iterated commutator shows that
\[
 M_\gamma:=\bigcap_iR_{i,a_i},\qquad \gamma=(a_i),
\]
has image containing $C$. In particular it is nontrivial.

Now take the intersections $M_\gamma$ for all slot tuples, allowing
trivial intersections. They are permuted by $E$ and generate a
nontrivial $E$-normalized subgroup of $S$, whose lifts span
$\mathcal A$ by \Cref{c:full-spanning}. For distinct tuples the lifts
commute linearly, because they lie in distinct tensor slots of at least
one original system. Their lift algebras therefore commute and jointly
generate $\mathcal A$. The center of each algebra centralizes all of
$\mathcal A$, hence is scalar. Its radical, multiplied by the commuting
remaining algebras, would give a nilpotent ideal of $\mathcal A$, hence
vanishes. The non-scalar lift algebras are thus central simple; their
tensor product maps isomorphically onto $\mathcal A$.

This also recovers the original embedded factors. Group the tuple
algebras by their assigned slot in any one system. Each grouped algebra
is contained in that slot, and the product of all grouped algebras is
still $\mathcal A$. Comparing dimensions in the original tensor
decomposition forces equality in every slot. We have obtained a common
tensor refinement, not merely abstract isomorphic factors.

The fine factors are selfadjoint when appropriate and are permuted
transitively by $E$, by the ordered-factor exclusion. Thus the same form
and quadratic observations apply. In particular in bilinear
characteristic two the fine-slot forms are all alternating, by
transitivity and alternation of the total form. Grouping them into two
nonempty batches still gives alternating forms on both. The orthogonal
quadratic, if present, is zero also on pure tensors for every such
two-batch grouping: its canonical quadratic with the given polar form
exists by alternation and is zero on fine pure tensors too. The two
quadratics therefore agree on a basis and, having the same polar form,
agree everywhere.
\end{proof}

\begin{proposition}[Completion of tensor-system pruning]\label{c:tensor-pruning}
The missing tensor systems of \Cref{c:tensor-products} are excluded by an absolute bound on the possible invariant partitions of their common fine slots. Hence all tensor systems are excluded.
\end{proposition}

\begin{proof}
Fix two systems $i,j$ and send each nontrivial fine factor to its pair
of original slot indices $(a_i,a_j)$. Grouping factors with the same
pair gives an aggregate matrix factor. On a large aggregate space,
$K_i\cap K_j$ contains the full supported perfect classical group.
No further system on the line can split this aggregate: its fine
factors would give a tensor system normalized by that group, contrary
to \Cref{c:tensor-order}. Thus a large aggregate contains only one
nontrivial fine factor. On a small aggregate of degree $h$, there are
at most $\log_2h$ such factors, because their degrees multiply to $h$
and each is at least two. The fibers of the displayed map are therefore
bounded absolutely.

The map is $H$-equivariant, since $H\le E$ stabilizes both original
systems and permutes all tuple intersections; discarding scalar factors
preserves this action. By \Cref{c:orbit-bound}, the two rich top actions
give an absolute bound on ordered-pair orbits of occupied pairs, hence
on ordered-pair orbits of fine factors. Each original system is recovered
by grouping those embedded factors, as proved in
\Cref{c:tensor-refinement}. Distinct stabilizer labels therefore give
distinct invariant partitions. The bound \eqref{c:partition-bound}
contradicts the pattern-line size and completes tensor-system exclusion.
\end{proof}

\begin{theorem}[The next classical-ancestor socle]\label{c:classical-conclusion}
Let \(K\) be a tested descendant with a prepared interval, good for all the exclusions in this section, or containing the final good anchor. There is an actual nonabelian simple subgroup \(B\trianglelefteq K\) inside the ancestor socle \(S\), with full natural spanning, such that
\[
 HB=K,\qquad C_K(B)\le\operatorname{core}_K H,
\]
and \(C_K(B)\) is soluble. The quotient by this centralizer is almost simple with socle \(B\), and the exact subgroup interval is preserved.
\end{theorem}

\begin{proof}
Finally consider a tested descendant \(K\) with prepared interval, good for these exclusions (or containing the final good anchor). It has \(K\cap S\ne1\) since \(K\) is nonsoluble. Take a minimal \(K\)-normal subgroup there. It cannot be elementary abelian, and if a power of nonabelian simples, their lift algebras commute linearly by perfection, and jointly span fully. As in the refinement argument they are central simple tensor factors, individually selfadjoint when relevant; multiple ones would be a forbidden tensor system. Thus there is an actual simple \(B\trianglelefteq K\) in \(S\), absolutely irreducible projectively. Its projective linear centralizer is trivial by perfection and spanning, so \(C_K(B)\) injects in the soluble outer quotient and lies in \({\rm core}_K H\) by (\(*\)). Also \(B\) is not retained (else base-core quotient soluble), hence \(HB=K\). Quotienting its centralizer gives the required next almost simple group with socle \(B\), keeping the exact interval.
\end{proof}

\section{Pruning inside an alternating ancestor}\label{sec:alternating}

We use the same failed-chart construction to force the next almost simple group inside a high-degree alternating ancestor. The successive anchors become transitive, primitive, non-affine and non-Cartesian as the exclusions proceed.

\begin{convention}[Prepared alternating configuration]\label{c:alternating-configuration}
Use the pattern, prior-anchor and full-stabilizer-label conventions of the structure exclusions. Thus the present almost simple ancestor quotient is \(A_v\le M\le S_v\) of high degree on \(\Lambda\); labels under test and witnesses are in the interval above \(H\). For a witness \(K_i>H\), \([H,K_i]\) has the prepared link properties, including (\(*\)): a normal subgroup is retained in \(H\) or supplemented by \(H\) (missing), soluble normals are retained, and two commuting normals cannot both be missing (in the core quotient they would both contain the nonabelian monolith). Where earlier exclusions are used we have the previously good \(H<E<K_i\) common within a pattern, with link properties at \(E\) as well. Patterns can have a whole line of \(q\) distinct labels with \(K_i\cap K_j\le K_l\) for distinct \(i,j\) and all \(l\) on the line; cross-chart pairs when invoked have \(K_i\cap K_j=H\).
\end{convention}

\subsection{Permutation orbit bounds and primitivity}

\begin{lemma}[Two alternating permutation images]\label{c:alternating-orbits}
If a group acts on two finite sets, inducing at least the alternating group on each, then its number of orbits on ordered pairs of pairs, one coordinate from each set in each pair, is bounded absolutely.
\end{lemma}

\begin{proof}
We use the following orbit bound for two finite set actions on sizes \(s,t\), each inducing at least the alternating group on that set. The number of orbits on ordered pairs of pairs (one coordinate from each set per pair) is absolutely bounded. Indeed the joint action is subdirect on the two induced groups. In large degrees either the product of the alternating groups is contained, or the two alternating simples are linked diagonally, hence the same degree and natural identification up to relabeling. This follows by the normal kernels of subdirect projections and simplicity (normal subgroups not containing the alternating simple must centralize it); use also the usual automorphism theorem for alternating groups in high degrees. If only one degree is large the large alternating factor acts independently by the same argument. Equality patterns and the bounded small sizes give the assertion.
\end{proof}

\begin{proposition}[Subsets and equal-block partitions]\label{c:alternating-primitivity}
In the prepared alternating configuration, full stabilizers of proper nonempty subsets and nontrivial equal-block-size partitions are excluded. Subsequent good anchors are primitive.
\end{proposition}

\begin{proof}
First exclude stabilizers of proper nonempty subsets. For a cross-chart
pair, the intersection contains supported alternating groups on at most
four cells. Thus $H$ has boundedly many point orbits, and every subset
whose stabilizer is a witness is a union of these orbits. Only boundedly
many subsets, and hence witness labels, are possible.

Next exclude nontrivial equal-block-size partitions, which suffice for
primitivity once transitivity is known. Write $s$ for the number of
blocks and $k$ for their common size. In a witness label, the
intersection $Y$ of $M$ with the base block group is normal. If $Y$ is
retained and $k\ge3$, every individually supported $A_k$ lies in $H$.
Its primitivity makes the restriction of any other tested
$H$-invariant equivalence to that block either universal or discrete.
In the discrete case a point has an equivalent point outside the block,
since the other partition also has block size at least two. The
supported transitive group fixes that outside point and would make it
equivalent to every point of the original block, a contradiction.
Consequently each original block lies in a block of the other partition.
Applying this in both directions to two retained systems with block
sizes at least three forces their equality. For retained $k=2$, the
even simultaneous flips on any two blocks lie in $H$. A distinct
matching fails invariance by flipping a block while leaving one of its
points' outside partner fixed, using the high degree.

If $HY=K_i$, then $H$ induces the full symmetric top action on the
blocks. For a cross-chart pair of missing partitions,
\Cref{c:alternating-orbits} bounds the orbits on ordered pairs of block
cells. Moreover $H$ contains the supported alternating groups on the
points of each cell, so each such orbit has boundedly many point-pair
orbits above it, including when both points lie in the same cell.
Thus $H$ has boundedly many orbits on $\Lambda^2$. Every invariant
equivalence is a union of those orbits, giving the absolute partition
bound used in \eqref{c:partition-bound}. This is too small for all the
distinct pattern labels. Here and below the chart field size exceeds
these absolute bounds, after coloring by the indicated finitely many
cases.
\end{proof}

\begin{proposition}[Elementary abelian regular groups]\label{c:alternating-affine}
For a primitive previously good anchor, normalizers of elementary abelian regular groups are excluded.
\end{proposition}

\begin{proof}
Exclude next normalizers of elementary abelian regular groups. Such groups lie in \(A_v\) (parity of translations, high degree) and would thus be retained soluble normals. They would be minimal normal in the primitive previously good \(E\), by transitivity of every nontrivial normal subgroup and regularity. Distinct ones must commute, impossible for regular self-centralizing groups.
\end{proof}

\subsection{Cartesian decompositions}

\begin{lemma}[Retained Cartesian products]\label{c:cartesian-retained}
For a primitive anchor avoiding the affine exclusion, a normalized nontrivial Cartesian decomposition has transitive slots of common size \(k\ge5\). Its full stabilizer has a normal product \(Y=A_k^s\le A_v\) and induces at least \(A_s\) on slots. Retained products determine at most one system in a pattern.
\end{lemma}

\begin{proof}
Next exclude stabilizers of nontrivial Cartesian decompositions of the point set (product coordinates \(\Lambda\simeq\prod_{a=1}^s \Lambda_a\), \(s\ge2,\ |\Lambda_a|\ge2\); stabilize their system, i.e. equality-of-coordinate partitions, allowing slot permutations). Primitivity of \(E\) gives slot transitivity (an orbit of slots otherwise gives an invariant nontrivial equivalence), and thus common slot set size \(k\), say on \(s\ge2\) slots. Sizes \(2,3,4\) would give canonical elementary abelian regular groups (\(S_2,A_3,V_4\), respectively, multiplied over slots). For \(k\ge5\) let \(Y_i=A_k^s\) acting slotwise, inside \(A_v\) and normal in the label \(K_i\). The induced slot action of \(K_i\) contains \(A_s\) by commutators of true slot permutations. These stabilizers are proper (Cartesian coordinate agreement prevents full 2-transitivity). If \(Y_i\le H\), it is minimal normal in \(E\), and has trivial centralizer on \(\Lambda\) (point stabilizer in \(Y_i\) self-normalizing). Thus all retained ones in the pattern must coincide, determining the slots themselves by the orbits of factor subproducts.
\end{proof}

\begin{lemma}[Missing normal powers with primitive factor action]\label{c:primitive-product-interval}
Let \([H,K_i]\) be a prepared interval and \(Y_i\trianglelefteq K_i\) a direct power of nonabelian simple groups with \(HY_i=K_i\). If the action of \(H\) on its simple factors is primitive, then \(H\cap Y_i\) is a nontrivial coordinate product, and every group between \(H\) and \(K_i\) intersects \(Y_i\) in a coordinate product.
\end{lemma}

\begin{proof}
Here and in the next exclusion we use a simple consequence of \Cref{f:prop-interval-cases}. If \(HY_i=K_i\) for a direct normal power of nonabelian simples and \(H\)'s action on the factors is primitive, then \(P=H\cap Y_i\) is a nontrivial coordinate product and all groups between \(H,K_i\) intersect \(Y_i\) in coordinate products. Indeed \(Y_i\) is minimal normal, missed by the base so embeds as the minimal normal modulo core, and the exact interval is that of invariant subgroups above \(P\). In the earlier reductions, the subdirect case and the all-nonbottom-subdirect case are now both impossible: ties must be factor-index invariant, so primitivity allows no chain of comparable distinct proper nontrivial subdirect groups. This leaves precisely the asserted product case.
\end{proof}

\begin{lemma}[Common Cartesian refinement]\label{c:cartesian-refinement}
A finite pattern line of missing Cartesian labels has a common refinement into finite coordinate sets. It is constructed from the tuple intersections of the slotwise products \(E\cap Y_i\). The original Cartesian systems are recovered by grouping the fine coordinates, and \(E\) permutes the fine coordinate directions.
\end{lemma}

\begin{proof}
For missing Cartesian labels \(HY_i=K_i\), \(H\) induces the full top image, transitive and containing \(A_s\), hence primitive on slot indices. On a line of these labels write
\[
 N_i=E\cap Y_i=\prod_a R_{i,a}\trianglelefteq E
\]
a slotwise product strictly above $H\cap Y_i$. In
$\bar E=E/\operatorname{core}_E H$, each $\bar N_i$ contains the
nonabelian simple-power monolith $W$. A simple factor $C$ of $W$ lies
in one of the $\bar R_{i,a_i}$ for each $i$: the commutators of those
slot groups with $W$ generate $W$, are normal in $W$, and commute for
distinct slots.

Apply the two commutator steps from the proof of
\Cref{c:tensor-refinement}. First, since the $N_i$ are normal in $E$
and their images contain the perfect group $W$, their intersection
$N=\bigcap_iN_i$ still has image containing $W$. Next,
$[N,R_{i,a_i}]\le N\cap R_{i,a_i}$ has image containing $C$.
The subgroups $N\cap R_{i,a_i}$ are normal in $N$, so iterating
commutators among them shows that the actual tuple intersection
\[
 T_\gamma=\bigcap_iR_{i,a_i},\qquad \gamma=(a_i),
\]
has image containing $C$.

Take these intersections for all slot tuples. They are permuted by
$E$ and generate a nontrivial normal subgroup, which is transitive
because $E$ is primitive. Distinct tuple groups commute, since they
act in distinct slots of at least one original system. Consider their
action through the abstract product $\prod_\gamma T_\gamma$ and fix
a point. If a product of tuple elements fixes that point, retain only
the components assigned to one prescribed slot $a$ of an original
system $i$. This partial product still fixes the point: it acts only
on coordinate $a$, where its action agrees with that of the original
fixing product. Repeating this restriction through every original
system isolates one tuple component at a time. Thus the stabilizer
in the abstract product is exactly the product of the individual
point stabilizers.

Transitivity now identifies the point set with the product of the
individual coset sets. Each original coordinate depends only on the
tuple components assigned to its slot. The maps from these separate
batches to the original coordinate sets have a bijective product map,
so each is itself a bijection. Grouping the fine coordinates therefore
recovers the original equality-of-coordinate partitions. Discard the
trivial coset sets; the remaining coordinate directions are still
permuted by $E$.
\end{proof}

\begin{proposition}[Completion of Cartesian pruning]\label{c:cartesian-pruning}
Missing Cartesian labels are excluded. Together with \Cref{c:cartesian-retained}, this excludes all nontrivial Cartesian decompositions.
\end{proposition}

\begin{proof}
Fix two distinct systems $i,j$ and group the fine coordinates by the
occupied pairs of their original slot indices. On an aggregate set of
size $u$, the supported $A_u$ lies in $K_i\cap K_j$, hence in every
label on the pattern line. If another system split that aggregate,
then on a slice with all outside coordinates fixed its nonconstant
coordinate partitions would give a nontrivial Cartesian decomposition
invariant under the natural $A_u$. Such a split has $u\ge4$ and
distinguishes pairs agreeing in some coordinate from pairs agreeing
in none, contrary to 2-transitivity. Thus there is only one nontrivial
fine coordinate over each occupied pair.

The map from fine coordinates to occupied pairs is $H$-equivariant:
$H\le E$ permutes the tuple groups and their coordinate directions,
and stabilizes both original systems. By
\Cref{c:alternating-orbits}, the two slot actions bound the number of
ordered-pair orbits on this set. Each original system is recovered
by an $H$-invariant grouping of the fine coordinates, as proved in
\Cref{c:cartesian-refinement}. Distinct full stabilizer labels therefore
give distinct invariant partitions. The bound
\eqref{c:partition-bound}, with fibers of size one, contradicts the
number of labels on the pattern line.
\end{proof}

\subsection{Diagonal actions and the almost simple conclusion}

\begin{lemma}[Primitive groups after Cartesian pruning]\label{c:primitive-diagonal}
Let \(E_0\) be a primitive group satisfying the preceding affine and Cartesian exclusions. If it is not almost simple, it normalizes a transitive group \(Y=B^s\), with \(B\) nonabelian simple and \(s\ge2\), whose point stabilizer is a full possibly twisted diagonal. Its degree is \(|B|^{s-1}\). This includes the pair of opposed regular simple groups.
\end{lemma}

\begin{proof}
Recall a short primitive-group observation after these exclusions. A minimal normal \(P=C^u\) of a primitive group \(E_0\) here is transitive with \(C\) nonabelian simple (an abelian one would be regular), \(E_0=P(E_0)_x\) at a point \(x\). Its point stabilizer \(Q=P_x\) has no proper intermediate increase towards \(P\) normalized by \((E_0)_x\). Taking its coordinate product closure, \(Q\) is either a product or subdirect. In the product case \(u>1\) would give a Cartesian coset product; in the subdirect case the same holds for multiple tie parts (uniform sizes by factor transitivity of \((E_0)_x\)). Indeed the point stabilizer of \(E_0\) is transitive on factors (minimal normality, with \(P\) itself not permuting them). In the product situation, including when taking batches given by multiple tie parts, the cosets \(P/Q\) accordingly split as a product of coset sets. Their sizes are uniform and nontrivial, \(P\) respects this product coordinatewise, and \((E_0)_x\) respects it by its conjugation/transport of the factors or batches with their respective stabilizers in \(Q\). Thus for \(u>1\) we have just a full diagonal, possibly with twists, and \(v=|C|^{u-1}\). For \(u=1\) with nontrivial \(C_{E_0}(P)\), that centralizer must be transitive so both are regular opposite simple groups, again giving a transitive power \(C^2\) with full diagonal point stabilizer. Consequently any non-almost-simple case here normalizes a transitive \(Y=B^s,\ s\ge2\), \(B\) nonabelian simple, with full diagonal point stabilizer.
\end{proof}

\begin{proposition}[Diagonal normalizers]\label{c:diagonal-pruning}
Under the successively primitive, non-affine and non-Cartesian anchor hypotheses, the full normalizers of the groups \(Y\) in \Cref{c:primitive-diagonal} are excluded. Both retained and missing normal powers are covered.
\end{proposition}

\begin{proof}
Finally exclude the normalizers \(K_i\) of all such \(Y\). They are proper since \(Y\le A_v\) by perfection, and induce at least \(A_s\) on factors by diagonal-coordinate factor permutations and commutators. For retained \(Y\le H\), transitivity of \(E\) on factors, when present, makes \(Y\) minimal \(E\)-normal with trivial centralizer (the diagonal is self-normalizing), leaving no distinct such system. If instead \(E\) has multiple factor orbits, their products are each transitive so use at least \(s-1\) factors each, forcing just two opposed regular simple groups, the only minimal normals in \(E\), again determining \(Y\). If \(K_i\)'s top action is intransitive, likewise \(s=2\), and at least one individual factor is retained by (\(*\)); such factors are regular simple minimal normals of \(E\), determining \(Y\) by taking also their opposite regular group (distinct such minimal normals commute so determine the same pair).

\textit{The missing transitive case.} It remains that \(HY=K_i\) and \(H\) induces a transitive top action containing \(A_s\). By the interval product consequence, \(N=E\cap Y=\prod R_a\) is a nontrivial slot product with each projection proper (since \(E<K_i\)). Choose a minimal \(E\)-normal \(P'=C^u\) inside it. Each factor of \(P'\) is contained actually in one \(R_a\), since \([P',R_a]\le P'\cap R_a\), normal in \(P'\), and those commutators generate \(P'\) by perfection. This assignment is \(E\)-equivariant with some common positive number \(h\) per slot, so \(u=sh\). The primitive observation applied to \(E\) gives \(Q=P'_x\) a full diagonal. If \(h>1\), take inside \(P'\) the product of the projections of \(Q\) along the factor blocks of the assignment. It is a proper intermediate increase, normalized by \(E_x\), impossible. If \(h=1\), point degrees give \(|C|^{s-1}=|B|^{s-1}\), contrary to containment in the proper \(R_a\).
\end{proof}

\begin{theorem}[The next alternating-ancestor socle]\label{c:alternating-conclusion}
After the exclusions, every tested descendant with its prepared interval is primitive almost simple. Its actual simple socle is transitive, lies in \(A_v\), and is supplemented by \(H\).
\end{theorem}

\begin{proof}
These arguments use only successively primitive, non-affine, and non-Cartesian prior anchors as needed and finite case subdivisions. Thus after the exclusions the tested descendant with its prepared interval is primitive almost simple, with an actual simple transitive socle in \(A_v\), supplemented by \(H\) via (\(*\)) (it cannot be retained, since then factoring by the base core leaves a soluble quotient by Schreier, inconsistent with the prepared interval), as desired.
\end{proof}

\section{Excluding bounded algebraic types}\label{sec:bounded}

Other than the high alternating and classical cases, simple socles of bounded order can be discarded immediately by the available interval length: \(M=HS\) at an almost simple current parent with socle \(S\), so \(|S|\ge|M:H|\) tends uniformly above any prescribed constant. By CFSG the remaining cases use a finite list of root systems, with parameter as large as needed, \(S=(G^{F_A})'\) for a simple adjoint algebraic group over \(k=\overline{\mathbf F}_p\).

Throughout this section we work at a current almost simple parent in the prepared reverse-flat interval of \Cref{sec:patterns}. Thus $H\le M\le\operatorname{Aut}(S)$, $S$ is the simple socle, and $M=HS$. The root systems under consideration belong to a fixed finite list. Every complexity bound below depends only on this list and on the stated degree bound; it does not depend on the characteristic or the finite-field parameter.

\subsection{Uniform algebraic input}

\begin{definition}[Bounded equation complexity]\label{w:complexity}
Fix closed affine split models for the root systems in the finite list, with a bounded number of coordinates and bounded-degree group operations. A closed subset has equation complexity at most $d$ if it is the common zero set, inside the chosen group, of polynomials of degree at most $d$. There need not be a previously specified bound on the number of equations: their linear span lies in the finite-dimensional space of polynomials of degree at most $d$.
\end{definition}

\begin{theorem}[Sufficiently general finite subgroups]\label{w:lp}
For the fixed finite list of simple adjoint types, there is a uniform degree bound $d_0$ with the following property. If a finite subgroup $\Gamma\le G(k)$, where $k=\overline{\mathbf F}_p$, is contained in no proper closed subset of equation complexity at most $d_0$, then there is a Steinberg map $\Psi$ such that
\[
 (G^\Psi)'\le\Gamma\le G^\Psi.
\]
Here the degree bound follows from the uniform constructible-family
formulation of the sufficiently-general subgroup theorem of
Larsen--Pink \citep[Theorem~0.5 and Metadefinition~2.2]{LarsenPink2011}:
the exceptional family has uniformly bounded defining equations in
the fixed affine models.
\end{theorem}

This formulation applies to the fixed split models in all characteristics. Its hypothesis is sufficient generality, expressed by avoidance of bounded-complexity proper closed subsets. The alternative argument in \Cref{sec:quantitative} will also establish the implication needed over the indicated algebraic closures.

\begin{convention}[Point maps and height]\label{w:point-maps}
We also use the standard classification of the Steinberg maps and of automorphisms and orders of the finite Lie type groups. On points take the group of maps generated by inner and diagram automorphisms, split Frobenius powers and inverses, and special isogeny powers where applicable (square split \(p\)-Frobenius in special small characteristics). It has a height homomorphism to \(\mathbf Z\) with unit of parameter growth \(p\), or \(\sqrt p\) in the special case. Height zero maps are ordinary algebraic automorphisms, and positive height maps are Steinberg maps with fixed groups, of parameter \(P=(\text{unit})^{\rm height}\). The normalizer of \(S\) here induces the automorphisms required for our large-parameter \(S\). Modulo inners one uses the standard pinned maps here: split Frobenius commutes with diagram maps, and in the cases with the special self-isogeny one uses its powers (square as above, no other nontrivial ordinary diagram needed in those types). Thus arbitrary words simplify to an inner twist of these combinations; the positive-height ones have standard fixed-group types (inner twists removed by Lang--Steinberg). The exceptional generators only concern fixed small characteristics. All coordinates, variety complexities etc. may use fixed closed affine split models of bounded size/degree, with the group operations and pinned ordinary diagram maps of bounded complexity uniformly in the characteristic. Images of bounded-degree-equation varieties by the whole point-map group have bounded equation degree (field powers just transform coefficients, inner and diagram have bounded complexity, special map of fixed degree has inverse using itself and inverse field power). Positive maps are given algebraically of degree \(O(P)\). The classification and fixed-group facts used here are those of \citet{Steinberg1968,GLS1994}. We write $\mathcal A(G)$ for this group of maps on $G(k)$, $\operatorname{ht}$ for its height, and $\delta$ for the parameter unit, so $\delta=p$ or $\sqrt p$.
\end{convention}

\begin{lemma}[Order, degree, and pointwise rigidity]\label{w:height-rigidity}
For the finite type list, let $b=\dim G$ and $B(\Psi)=(G^\Psi)'$. For sufficiently large parameter $P=\delta^{\operatorname{ht}(\Psi)}$, the following hold.
\begin{enumerate}
\item $B(\Psi)$ is simple, and its order exceeds $|G^\Theta|$ for every $\Theta\in\mathcal A(G)$ with $0<\operatorname{ht}(\Theta)<\operatorname{ht}(\Psi)$.
\item For each fixed equation complexity $d$, $B(\Psi)$ lies in no proper closed subset of complexity at most $d$, once $P$ is sufficiently large for $d$.
\item If a positive-height point map fixes $B(\Psi)$ pointwise, it is a positive integral power of $\Psi$.
\end{enumerate}
\end{lemma}
\begin{proof}
Two useful uniform comparisons follow. Write \(b=\dim G\). For \(P\) large, \(B(\Psi)=(G^\Psi)'\) is simple of order at least \(P^b(1-o(1))/z\) inside \(G^\Psi\) of order \(P^b(1+o(1))\), with index at most a center bound \(z\) from the root system. In particular \(|B(\Psi)|\) exceeds the entire fixed-group order at any smaller positive height, since \(z<(\text{unit})^b\) by the usual center sizes. This applies uniformly to comparison with smaller heights: below parameter \(\sqrt P\) even the rough order bounds suffice, and otherwise the two product formulas both have small relative error. Also \(B(\Psi)\) lies in no proper closed subvariety of bounded degree, for \(P\) sufficiently large for that bound. On any irreducible component of dimension \(e<b\) and bounded degree, the fixed-point equations of degree \(O(P)\) cut a finite set of size \(O(P^e)\), by affine B\'ezout (combine the equations successively generically to cut to dimension zero). Hence a point map of positive height fixing \(B(\Psi)\) pointwise must be a power of \(\Psi\): divide its height by the given height and compose off that power on points. Positive remainder is impossible by orders, zero-height nonidentity by the bounded equations of its fixed locus.

For the last assertion, if $\alpha$ has height $q\operatorname{ht}(\Psi)+r$ with $0\le r<\operatorname{ht}(\Psi)$, the point map $\Psi^{-q}\alpha$ still fixes $B(\Psi)$ pointwise. If $r>0$, its fixed group is too small. If $r=0$, a nonidentity algebraic automorphism has a proper fixed locus with uniformly bounded equations, contradicting the second assertion. Hence $r=0$ and $\alpha=\Psi^q$.
\end{proof}

\subsection{Witness stabilizers and the divisor contradiction}

\begin{lemma}[A bounded-complexity failure witness]\label{w:witness}
Let $H\le K<M$ be a tested overgroup with $\Gamma=K\cap S\ne1$, and suppose $\Gamma$ fails the sufficient-generality hypothesis of \Cref{w:lp}. For sufficiently large parameter of $S$, there is a proper full stabilizer label $K'$ with $K\le K'<M$ such that
\[
 K'\cap S=N_G(Z)\cap S
\]
for a proper closed subset $Z\subset G$ whose equation complexity, and that of its conjugation normalizer $N_G(Z)$, are uniformly bounded.
\end{lemma}
\begin{proof}
Consider a proper \(K\) in the interval, above a tested large-index label, and \(\Gamma=K\cap S\ne1\). If insufficiently general for the sandwich, it lies in a proper bounded-degree-equation closed \(X\subset G\) by Larsen--Pink as used above. Take the subgroup of our point-map group preserving \(S\) and inducing \(M\)-elements on \(S\); it maps onto \(M\) by the automorphism description. Let \(\tilde K\) be the full preimage of \(K\) there (so each such map normalizes \(\Gamma\)). Then
\[
 Z=\bigcap_{\alpha\in\tilde K,\ \gamma\in\Gamma} \alpha(\gamma X)
\]
contains \(\Gamma\), is proper closed with equations of uniformly bounded degree, left-\(\Gamma\)-invariant and \(\tilde K\)-invariant. Even for the whole intersection one can just span all the equations occurring in the fixed bounded-degree space. Its point-map stabilizer (within that preimage of \(M\)) thus defines a label \(K'\ge K\), since the stabilizer before projection contains the kernel there (already in \(\tilde K\)). Identifying \(S\) with its inner actions on itself, \(K'\cap S=N_G(Z)\cap S\) by using the actual inner lifts to \(G\). The conjugation normalizer here has bounded-degree equations (test the equations of \(Z\) on all conjugates of its points by \(g,g^{-1}\)). It is proper since otherwise all conjugates of \(\Gamma\) left-stabilize \(Z\), while their normal closure is Zariski-dense by simplicity and trivial center. The degree comparison for \(S\) thus makes \(K'\) proper.

Here $N_G(Z)$ denotes the conjugation stabilizer, not the left stabilizer. To justify its properness explicitly, if every $g\in G(k)$ normalizes $Z$, left invariance under $\Gamma$ implies left invariance under every $g\Gamma g^{-1}$. The Zariski closure of the group generated by these conjugates is a nontrivial closed normal subgroup of the simple adjoint group $G$, hence is $G$. Since $Z$ is closed and nonempty, left invariance by this dense subgroup forces $Z=G$, a contradiction. The nonidentity of $\Gamma$ suffices because the adjoint group has trivial center.
\end{proof}

\begin{proposition}[No bounded-type parent with a stage remaining]\label{w:bounded-exclusion}
Fix the finite list of bounded root systems. There is a fixed trial increment $t$ such that a prepared current parent with a remaining private-chart stage cannot have socle of one of these types, provided the residual interval length is sufficiently large. All choices are uniform in the characteristic and parameter.
\end{proposition}
\begin{proof}
At the current parent flat $W$, test a sufficiently large fixed
increment $t$ in a private chart. We first find a flat $U$ of that
increment for which every $K\in[K_U,M]$ admits the sandwich for
$K\cap S$. If every trial failed, \Cref{w:witness} would supply full
stabilizer witnesses above the trial labels. The affine-pattern
construction gives an independent parameter list whose successive
label intersections decrease strictly for more steps than the dimension
of the relevant polynomial space. Descent remains strict after
intersection with $S$, since $M=HS$ and every overgroup of $H$ is
determined by its intersection with $S$. Each strict descent must add
a defining polynomial outside the span of the preceding equations.
But all conjugation-normalizer equations lie in one bounded-dimensional
space, a contradiction. This step needs only a bounded independent
list, no prior anchor, and tests relative to the current parent.

Put $E=K_U$ and $\Gamma_K=K\cap S$ for $E\le K\le M$. We now have
\[
 B(\Psi_K)\le\Gamma_K\le G^{\Psi_K}.
\]
The reserved interval length below $E$ makes every parameter here
large enough for \Cref{w:height-rigidity}. Indeed, if that residual
height is $h$, then
$|\Gamma_K|\ge|\Gamma_E|\ge[E:H]\ge2^h$, whereas all fixed groups
below a prescribed parameter have bounded order for the fixed type
list. Write $B_E=B(\Psi_E)$. Perfection and the abelian quotient of a
fixed group by its derived subgroup give $\Gamma_E'=B_E$.

Since $B_E\le\Gamma_K$, the map $\Psi_K$ fixes $B_E$ pointwise.
Rigidity therefore gives $\Psi_K=\Psi_E^{a_K}$ for a positive integer
$a_K$. The ambient map $F_A$ likewise fixes $B_E$, so $F_A=\Psi_E^d$.
It also fixes $B(\Psi_K)$ pointwise, and rigidity applied to that
simple group makes $F_A$ a positive power of $\Psi_K$. Consequently
\[
 F_A=\Psi_E^d,\qquad \Psi_K=\Psi_E^{a_K},\qquad a_K\mid d.
\]
These are equalities of point maps, so each exponent determines an
actual embedded fixed group.

Every lift of $E$ normalizes $\Gamma_E$ and hence $B_E=\Gamma_E'$.
Conjugating $\Psi_E$ by such a lift gives a map of the same height
fixing $B_E$ pointwise; rigidity makes it equal to $\Psi_E$. Thus the
lifts commute with $\Psi_E$. For each divisor $a$ of $d$, put
$B_a=B(\Psi_E^a)$. This group lies in $S$: its elements are fixed by
$F_A$, and perfection puts it in $(G^{F_A})'$. It is normalized by
$E$, so $E B_a$ is an overgroup label. Moreover
\[
 (E B_a)\cap S=\Gamma_E B_a\subseteq G^{\Psi_E^a},
 \qquad E B_1=E,\qquad E B_d=M.
\]
If $a\mid b$ with $a<b$, then $B_a\le B_b$, and
\Cref{w:height-rigidity} gives $|B_b|>|G^{\Psi_E^a}|$.
Hence $E B_a<E B_b$. A saturated divisor chain therefore gives
$\Omega(d)$ strict steps in the exact height-$t$ interval $[E,M]$,
where $\Omega(d)$ counts prime factors with multiplicity. Thus
\[
 \tau(d)=\prod_{\ell^a\parallel d}(a+1)
 \le2^{\Omega(d)}\le2^t.
\]

For a fixed exponent $a_K$, the sandwich leaves only boundedly many
choices for $\Gamma_K$: its quotient by $B_{a_K}$ is a subgroup of
the bounded-order group $G^{\Psi_E^{a_K}}/B_{a_K}$. If $z$ bounds
that order, $C=2^z$ bounds the number of choices. Since
$K=H\Gamma_K$, no additional choices of outer automorphisms occur.
We conclude that $|[E,M]|\le C2^t$.

The reverse-flat interval, however, contains at least
$q^{\lfloor(t-1)^2/4\rfloor}$ intermediate subflats, by the ordinary
linear-subspace count; this bound also covers an empty starting flat.
For sufficiently large fixed $t$ it exceeds $C2^t$, already using
$q\ge5$. This is the required contradiction, uniform in characteristic
and parameter, whenever a bounded-type parent has a stage remaining.
\end{proof}

\section{Protections and the length of high-type paths}\label{sec:weights}

The structural exclusions give nested simple groups whose actions remain
transitive or absolutely irreducible in every earlier ancestor. We now
need two further conclusions. First, within each consecutive
same-characteristic classical run, the chart construction can select
descendants whose actions on each ancestor's natural module admit no
nontrivial factorization into individually invariant tensor factors
over the algebraic closure, and that are fixed by no prime compositional
root of its defining field map. Second, any path with these protections
has absolutely bounded length. The selection uses the affine geometry;
the length bound uses only the group properties stated below, so it can
also be applied to the later reverse and forward constructions.

All high-degree and high-rank thresholds in this section are sufficiently
large absolute constants. Bounded types have already been treated in
\Cref{sec:bounded}.

\begin{definition}[Prepared paths and ancestor properties]\label{w:prepared-path}
A \emph{prepared path} is a strict descending chain of finite groups
$K_W$ containing a common subgroup $H$, together with actual strictly
nested nonabelian simple groups $B_W\trianglelefteq K_W$ such that
\[
 HB_W=K_W,\qquad C_{K_W}(B_W)\le H,
\]
where each centralizer is soluble and the quotient by it is almost
simple with socle $B_W$. Fix the usual natural realizations of the
alternating and classical socles. For every ancestor--descendant pair
require the following properties.
\begin{enumerate}
\item In an alternating ancestor, the descendant socle is transitive on the natural set.
\item In a high classical ancestor, the descendant socle is projectively absolutely irreducible on the natural module. Its label normalizes no proper ground-field tensor system of the endomorphism algebra of the kind excluded in \Cref{sec:classical}; in isometry types the factors are individually selfadjoint.
\end{enumerate}
\end{definition}

The pruning construction of \Cref{sec:patterns} supplies these properties
by \Cref{c:classical-conclusion,c:alternating-conclusion}.
\Cref{c:ancestor-socles} makes the simple groups actual and compatible
in all ancestor quotients, whose centralizer kernels lie in $H$.
Consequently a full stabilizer containing a tested label pulls back to
a label of the exact original interval. In the selection argument the
same properties hold for all comparable candidate subflats of a
successful pruned flat before the next node is chosen: if
$W\subset L\subset U$, then $B_U\le B_L\le B_W$.

\subsection{Classical conventions and changes of type}

\begin{convention}[High classical groups]\label{w:classical-conventions}
We use the standard structure theory of finite simple classical groups, simply connected coverings, Schur multipliers and automorphisms. In sufficiently high rank there are no small covering exceptions: the universal central extension is the usual fixed group \(\widehat S=\widehat G^\Phi\) of simply connected type. Write \(k=\overline{\mathbf F}_p\) in characteristic \(p\) defining the group. We also use the adjoint algebraic group \(G\) over \(k\), with \(S=(G^\Phi)'\), and regard it on points as a projective group on a natural \(v\)-space. Here \(\Phi\) is standard of parameter \(P=p^f\) (split Frobenius composed possibly with diagram/inner automorphisms); ground field \(F\) on the finite natural module is of size \(P\), or \(P^2\) for unitary type. We use types A (linear and unitary), C, odd orthogonal B only in odd characteristic, and even orthogonal D; by the bilateral (bilinear) form types we mean B,C,D, not the unitary case. The usual automorphisms of \(S\) lift to point maps in the group on adjoint \(G(k)\) generated by inner and diagram automorphisms and integral powers of a split Frobenius \({\rm Fr}_p\) (including inverse on points). Modulo inner this group is abelian, with at most two diagram choices and field exponent called height. A positive-height element is an ordinary Steinberg map (inner conjugate to standard form by Lang--Steinberg). We shall use that in this adjoint group the order of the derived fixed simple group at positive height \(f_2\) exceeds the order of the whole fixed group at \(0<f_1<f_2\), uniformly in high rank. Indeed the classical order formulae give order \((p^f)^{\dim G}\) times factors bounded above and away from zero (products of \(1\pm(p^f)^{-j}\)); the derived index costs at most a polynomial in the rank. When finite types \(B_e,C_e\) coincide in characteristic two we use the symplectic realization (type C); the odd-dimensional \(B_e\) convention in this section is only for odd characteristic. In particular all Steinberg maps needed \emph{in this high classical section} are ordinary graph/field ones, not the special square-root-field maps. We use the standard high-rank multiplier and automorphism results in \citet{GLS1994,KleidmanLiebeck1990}.
\end{convention}

We use the standard highest weight theory of simply connected semisimple algebraic groups, including Steinberg's Tensor Product Theorem and his parametrization of the irreducibles of \(\widehat S\) in defining characteristic as restrictions of \(L_{\widehat G}(\lambda)\) for \(P\)-restricted dominant weights (including standard diagram twists). Projective representations of \(S\) lift to ordinary ones on its universal cover, uniquely once the projective matrices are fixed (perfection). 

More explicitly, we use the universal-central-extension and automorphism description in the high-rank range of the standard multiplier/automorphism theorems for finite groups of Lie type. Thus lifting a projective representation means pulling back the scalar central extension and then using universality on \(\widehat S\). The Steinberg theorems we use here are the ordinary \(p\)-adic tensor product for rational simples and the defining-characteristic Restriction Theorem on the full simply connected fixed group: highest weights with fundamental coefficients \(0,\ldots,P-1\) give irreducibles, pairwise inequivalent, and give all of them, also for these standard twisted forms. With pinned map a diagram twist times split \({\rm Fr}_p^f\), a further field twist rotates digit positions on restriction modulo \(f\), with wrapping applying a diagram (since the defining map is the identity on the finite points). One obtains this by twisting the tensor product digit factors separately. Equivalence of linear lifts up to a projective intertwiner on the perfect cover already implies ordinary equivalence (the scalar discrepancy would be a character). These are the ordinary tensor-product and finite-restriction theorems, including ordinary diagram twists, of \citet[Theorems 1.1, 7.4 and 9.3]{Steinberg1963}; see also \citet{Jantzen2003}.

\begin{lemma}[A lower bound for projective degree]\label{w:min-degree}
There is an absolute constant $c_0>0$ such that every nontrivial projective representation in defining characteristic of a sufficiently high-rank simple classical group with natural dimension $v$ has dimension at least $c_0v$.
\end{lemma}
\begin{proof}
Any nontrivial irreducible projective representation of \(S\) in characteristic \(p\) has dimension at least a positive absolute constant times \(v\) in high rank: its highest-weight Weyl orbit already has at least that size (permutations of a nonconstant weight word in type A; signed permutations, possibly with even sign changes, in B,C,D). This bound suffices here and also bounds any nontrivial projective representation, using perfection if the lift has only trivial composition factors. In the last case a lift with all composition factors trivial has unipotent, hence soluble, image. A perfect group has no nontrivial soluble image, so a nontrivial representation must have a nontrivial composition factor.
\end{proof}

\begin{lemma}[Alternating subsequences]\label{w:alternating-switch}
A prepared path contains at most two alternating socles of sufficiently large degree.
\end{lemma}
\begin{proof}
Any alternating subsequence of the path has bounded length. Suppose there are three socles \(A_n>A_m>A_e\) as subgroups (\(n>m>e\)), with the transitivity for each ancestor pair. For large \(e\), \(m\ge\binom e2\). In fact an intransitive point stabilizer in natural \(A_e\)-coordinates giving index \(>e\) has index at least \(\binom e2\) (if within a point group \(A_{e-1}\) properly, use simplicity for minimum proper index at least \(e-1\)). Proper transitive subgroups of \(A_e\) have index at least exponential in \(e\) for large \(e\), by Bochert's Theorem in the primitive case, and the block partition count in the imprimitive case. For instance with \(b\) blocks of size \(a\), both at least two, pair the blocks disjointly and interchange corresponding positions arbitrarily in each pair, giving at least \(2^{(a-1)\lfloor b/2\rfloor}\) partitions. Thus \(|A_e|=\exp(o(m))\). If \(D\) is the stabilizer in \(A_m\) on the \(n\)-set, \(A_m=A_eD\) by transitivity, so \(D\) cannot be transitive on the natural \(m\)-set by the same index bound. Nor can it have just orbits of sizes 1 and \(m-1\): the bound on proper transitive subgroups on \(m-1\) would force \(D=A_{m-1}\), contrary to \(n>m\). It thus stabilizes an \(h\)-set, \(2\le h\le m-2\), making \(A_e\) \(h\)-homogeneous on \(m\). This implies 2-homogeneity: the count inside an \(h\)-set of edges from any one pair orbit is constant; comparing sets consisting of \(i\) or \(j\) and \(h-1\) other points shows the adjacency rows at \(i,j\) agree away from both (all such sums of differences vanish). Thus there is just one edge orbit. Even order then gives 2-transitivity by an involution flipping an ordered pair. By the classification of finite 2-transitive groups, alternating simple groups of sufficiently large degree have no such nonnatural action. This proves the bound. The permutation-group inputs used here are the primitive subgroup bound and the classification of finite doubly transitive groups; see \citet{DixonMortimer1996}.
\end{proof}

\begin{lemma}[Changes of defining characteristic]\label{w:characteristic-switch}
For a classical ancestor $i$ and classical descendant $j$ in a prepared path with distinct defining characteristics, there is an absolute $c>0$ such that
\[
 v_i\ge P_j^{c v_j}.
\]
There cannot be three high classical nodes $i<j<l$ with $p_i\ne p_j$ and $p_j\ne p_l$. Consequently, after splitting at alternating nodes and changes of defining characteristic, the path is the union of an absolutely bounded number of consecutive classical runs of one characteristic.
\end{lemma}
\begin{proof}
For high classical pairs in distinct characteristics, absolute constants give \(v_i\ge P_j^{c v_j}\) for ancestor \(i\), descendant \(j\), with \(c>0\). Here is a direct bound. The high classical \(S_j\) contains images of paired special linear groups \({\rm SL}_h(P_j)\), \(h\ge c' v_j\), using a subspace in linear type or paired maximally isotropic coordinates in the form types (one may restrict scalars of matrix entries in the unitary coordinates to \(\mathbf F_{P_j}\)). These act nontrivially including on column shear elements. The restricted projective representation lifts linearly on that \({\rm SL}_h(P_j)\), centrally closed for such high \(h\). Its abelian column shears of size \(P_j^{h-1}\) are diagonalizable here and have some nontrivial character occurring; the \({\rm SL}_{h-1}\) action is transitive on all their nontrivial characters. This proves the assertion. Thus we cannot have classical nodes \(i<j<l\) (indices increasing downwards) with \(p_i\ne p_j,\ p_j\ne p_l\), since then
\[
 \log|S_l|=O(v_l^2\log P_l)=O(\log^2 v_j)
\]
whereas \(v_i\) is at least exponential in \(v_j\). But \(|S_l|\ge v_i^2\) by its absolutely irreducible projective action in \(i\): lift that action to its perfect central cover. Absolute irreducibility implies that its matrices span the full $v_i\times v_i$ matrix algebra. Central lifts of one projective element are scalar multiples, so at most $|S_l|$ matrices are needed in this span. This contradicts the two exponential lower bounds for sufficiently high rank.

Hence there are boundedly many interruptions of consecutive high-classical runs in the \emph{same} characteristic, on paths avoiding bounded types. It suffices to bound those runs. We now work in such a run, notation \(S_i=B_{W_i}\).
\end{proof}

\subsection{Canonical tensor factors over the algebraic closure}

\begin{lemma}[Canonical tensor atoms]\label{w:atoms}
Let $S_i$ be a high classical ancestor, and let $B$ be a candidate simple subgroup of high classical type in the same defining characteristic. Assume that $B$ acts projectively absolutely irreducibly on the ancestor natural module $V$ over $k=\overline{\mathbf F}_p$. Among decompositions of $\operatorname{End}_k(V)$ into nontrivial full matrix factors fixed individually by $B$, there is a unique finest one as a collection of embedded matrix subalgebras. Call its factors the \emph{atoms}. There is at least one atom for each nonzero ordinary $p$-digit in the restricted highest weight of the lifted representation of $\widehat B$. The only possible further separation within a digit is across the multiple bond in types B and C. Every semilinear algebra action, or algebra anti-action with inversion on the group, normalizing the projective $B$-action permutes these atoms.
\end{lemma}
\begin{proof}
To verify this, in any tensor factorization into nontrivial slots with \(B\) fixing the factor algebras individually, lift each projective action to \(\widehat B\). They tensor to the original lift by perfection, and are all irreducible. Regard them as restrictions of rational simples with restricted highest weights for the finite-group parameter of \(B\). Their tensor is then irreducible rationally as well. By Steinberg decomposition, even the subproducts of the contributing digit modules at any one exponent (untwisted) must be irreducible.

Recall the top Levi face of a rational simple \(L(\lambda)\) for nodes \(I\), using weights differing downwards from highest only by \(I\)-roots, is a simple highest module for the Levi derived group with the corresponding coefficients. Indeed it is generated by the highest vector there by distribution-algebra triangular decomposition (PBW); a nonzero Levi submodule on that face has a highest unipotent-fixed vector, fixed also by outside raising groups since they go off the weights. In the simple module this forces the highest line. The face of a tensor product with the sum as highest weight tensors the corresponding faces. Now at any one digit, no two modules from different slots can have active coefficients in the same simple-bond connected component of the Dynkin diagram. Otherwise take a shortest simple-bond path between active nodes of two different slots, possibly just one node (minimize over such pairs in a single simple-bond component). For the two modules the faces on that path have endpoint highest weights \(a\omega_1,b\omega_\ell\) of type \(A_\ell\), \(1\le a,b<p\), with no other supports there by minimality. These are symmetric powers of natural and dual modules (the same natural type for \(A_1\)), simple for these exponents. For example simplicity follows using torus monomials of distinct weights and elementary root transfers of powers with nonzero coefficients at degree less than \(p\). Their tensor is reducible by contraction (one derivative on each side paired), a nonzero equivariant map to a strictly smaller space. This contradicts the face property.

Thus coefficients are disjoint with no carry, and the digitwise sum highest weight is restricted and must be exactly the given tuple for \(V\). Take one atom per nonzero digit except for the following possible refinement in B,C: if the regular chain and the node beyond the multiple bond both contribute and can be separated in a factorization, split them. Indeed in such a separation their two rational pieces tensor irreducibly to the digit module, so all such splits at all digits can be made together. This gives a factorization, and any other one as above has its slot representations given by grouping the atoms (by the coefficient argument and uniqueness of the simples). This is true of the actual matrix subalgebras: equivalent factorizations with fixed slot representation classes have their tensor identifications differing, after intertwiners, by a scalar since \(V\) is absolutely irreducible for \(\widehat B\). In more detail, separated active components at a digit give actual rational factor modules by taking the digit pieces of the slot modules of that separation; their subproduct, untwisted, is irreducible as above, with the individual component weights and their sum all forced. Thus it realizes the indicated split even as rational modules of that digit, independently of the other digits. This proves the claim. Any semilinear or anti-action on the algebra normalizing the projective \(B\)-action (using inverse with an anti-action) permutes the atoms, by the same intrinsic description.
\end{proof}

\begin{lemma}[Descent and the single Frobenius cycle]\label{w:atom-cycle}
Retain the hypotheses of \Cref{w:atoms}, and assume the candidate label has the ancestor tensor-system exclusion in \Cref{w:prepared-path}. Let $\sigma$ be arithmetic Frobenius for the ancestor ground field, except in unitary type where it is the extended semilinear adjoint anti-action with coefficient exponent $P_i$. Then $\sigma$ permutes the atoms in one cycle. If its length is $t=t(B)$, all atoms have a common degree $u\ge2$ and
\[
 u^t=v_i.
\]
The lifts of the candidate normalizer used to act on the endomorphism algebra commute with $\sigma$. In bilateral form types the factors are individually selfadjoint, with alternating factor forms in characteristic two.
\end{lemma}
\begin{proof}
The map $\sigma$ permutes the atoms because it fixes the projective
$B$-action, using inversion with the anti-action in unitary type. We
first specify lifts of the ancestor automorphisms. Use finite-field
projective semilinear representatives, similarities in the form types,
and also graph anti-actions in linear type. Lift their field exponents
by integral powers of split Frobenius on $k$. These lifts commute with
$\sigma$: in unitary type this is the hermitian-similarity identity
extended to $k$, and otherwise it follows from the finite-field matrices.

The finite projective representatives are faithful on the ancestor
simple group. Indeed, if a collineation with coefficient exponent
$0\le a<\log_p|F|$ centralizes it, perfection identifies its natural
finite lift with the coefficient twist of highest weight $p^a\omega_1$.
Restricted-weight uniqueness forces $a\equiv0\pmod{f_i}$. In unitary
type it also rules out $a=f_i$, since wrapping produces the distinct
diagram-dual natural weight. With the field part removed, absolute
irreducibility leaves only a scalar. The same weight comparison rules
out a linear-type duality in the kernel. Thus closure lifts differ only
by arithmetic ground-field Frobenius powers. A system stable under those
powers as a collection has a well-defined full stabilizer in the finite
ancestor. In unitary type the representatives here are collineations;
$\sigma$ itself is the semilinear adjoint anti-action and $\sigma^2$ is
arithmetic Frobenius for $F$.

We next check the forms needed for descent. In bilateral types the
perfect lift of $B$ preserves the ambient form. Restricted-weight
uniqueness makes its digits self-dual, and any pieces separated across
the B or C multiple bond are self-dual as well. Their invariant forms
tensor to the ambient form by uniqueness, so the atom algebras are
individually selfadjoint. In characteristic two each factor form is
alternating: a nonalternating symmetric invariant form would give the
nonzero invariant functional $x\mapsto\sqrt{b(x,x)}$ on a nontrivial
irreducible module.

If a quadratic form is present in characteristic two, every grouping
into at least two aggregate factors has the canonical quadratic with
the given polar form and value zero on simple tensors. To construct it,
prescribe zero on tensor-basis vectors. On a simple tensor, polar terms
with a repeated index vanish and the remaining terms cancel in pairs
by exchanging one index. This quadratic is invariant under the lifting
group. Its difference from the given invariant quadratic is the square
of an invariant linear functional, hence zero.

Suppose $\sigma$ has more than one orbit on atoms, and tensor together
the atoms in each orbit. The aggregate algebras descend under arithmetic
Frobenius to central simple $F$-factors, which split over the finite
field. In unitary type first descend under $\sigma^2$; stability under
$\sigma$ then gives selfadjointness for the hermitian adjoint. Bilateral
selfadjointness was just proved. Candidate lifts commute with $\sigma$
and therefore permute these aggregate factors. This is exactly a
ground-field tensor system forbidden by the preceding exclusions.

Its full stabilizer is proper. Otherwise $S_i$ acts trivially on at most
$\log_2v_i$ slots, by simplicity and comparison of its order with the
corresponding factorial. Absolute irreducibility makes each slot action
nontrivial. A smallest slot has degree at most $\sqrt{v_i}$, contrary to
\Cref{w:min-degree} at the fixed high threshold. There is therefore one
$\sigma$-cycle. Its atoms have equal degree $u$, and $u^t=v_i$.
\end{proof}

\subsection{Counting proper containing stabilizers}

The next two counts supply the finite lists needed for selection. Each
full stabilizer contains the candidate label, so it is an actual label
in the ancestor interval. To count distinct primes, we construct a
subgroup lying in all previously imposed stabilizers but not in the
next one. Their intersections then descend strictly, and interval
height bounds the number of such steps. Tensor coarsenings leave a
bounded small-slot remainder; fixing roots do not.

\begin{proposition}[Counting prime tensor coarsenings]\label{w:coarsening-count}
Under the hypotheses of \Cref{w:atom-cycle}, each prime $s\mid t(B)$ defines a canonical congruence coarsening $\mathcal P_s$ into $s$ matrix factors, normalized by the candidate label. Its full stabilizer in the ancestor is proper. There is an absolute constant $c_1$ such that
\[
 \#\{s:s\text{ is prime and }s\mid t(B)\}
 \le \operatorname{ht}[K_{\rm cand},K_{W_i}]+c_1.
\]
The height is taken in the actual subgroup interval, or equivalently in its prepared ancestor quotient.
\end{proposition}
\begin{proof}
For each prime $s\mid t$, group the atoms by their residue classes modulo
$s$ around the $\sigma$-cycle. Candidate lifts commute with the cycle,
so they normalize this canonical coarsening $\mathcal P_s$. Its full
stabilizer is well defined by \Cref{w:atom-cycle}. If the entire socle
$S_i$ preserved it, its permutation action on the $s$ slots would
commute with an $s$-cycle. That image is cyclic, hence trivial by
perfection. A slot of degree $u^{t/s}\le\sqrt{v_i}$ would then carry a
nontrivial projective action, contradicting \Cref{w:min-degree}.
Thus each coarsening gives a proper containing label.

Let $d$ be the product of distinct primes already imposed. Their common
refinement is $\mathcal P_d$, with $d$ slots of degree $h=u^{t/d}$.
In $G_i(k)$ take the product $X_d$ of the supported connected projective
groups on these slots: $\operatorname{PGL}_h$ in ambient type A, and
the connected projective classical groups of the factor forms in
bilateral types. Put $X_1=G_i(k)$ itself. In particular, at $d=1$ we
keep the orthogonal ambient group in characteristic two. For $d>1$
the tensor construction embeds these products in the connected ambient
group; the canonical-quadratic observation in \Cref{w:atom-cycle}
justifies the characteristic-two orthogonal case.

The map $\Phi_i$ cycles the $d$ factors through $\sigma$, using inversion
on units in the anti-case. On one factor $\Phi_i^d$ is an ordinary
positive Steinberg map of parameter $P_i^d$: after this coefficient
twist, its matrix-algebra action is an isomorphism or anti-isomorphism
preserving the factor group. The fixed group $X_d^{\Phi_i}$ therefore
projects isomorphically to that one-factor fixed group. While $h$ is
sufficiently large, its derived group gives a full high-rank natural
simple $Y$ on the chosen slot. It lies in $S_i$ by taking commutators
and preserves every previously imposed coarsening.

For a new prime $s$, suppose $Y$ also preserves $\mathcal P_s$.
Perfection and commutation with the induced $s$-cycle make it fix those
slots individually. Inside the chosen $d$-slot, the two coarsenings
intersect in $s$ tensor factors of degree $h^{1/s}$, by the Chinese
remainder description of their atom indices. At least one factor
action is projectively nontrivial, since their tensor is the natural
action of $Y$ on that slot. Its degree contradicts \Cref{w:min-degree}.
Thus the next stabilizer makes the intersection strictly smaller.

Interval height bounds the primes added while $h$ is high. At stopping,
$h\ge2^{t/d}$ and $t/d$ is divisible by every unused distinct prime.
Bounded $h$ leaves only absolutely many unused primes. This gives the
claimed additive constant $c_1$ in the actual interval, whose height is
unchanged by the centralizer quotient.
\end{proof}

\begin{proposition}[Counting fixing prime roots]\label{w:root-count}
Let $B$ be as in \Cref{w:atom-cycle}, in the high classical ancestor with adjoint group $G_i$ and ordinary Steinberg map $\Phi_i$ of height $f_i$. Consider pairs $(s,\Theta)$ where $s$ is prime,
\[
 \Theta^s=\Phi_i,
 \qquad \Theta|_B=1,
\]
and $\Theta$ belongs to the inner--diagram--Frobenius point-map group. The full normalizer labels of $S(\Theta)=(G_i^\Theta)'$ are proper, contain the candidate label, and number at most
\[
 2\operatorname{ht}[K_{\rm cand},K_{W_i}].
\]
There are at most two maps at each possible height, one for each diagram choice, and all such roots commute.
\end{proposition}
\begin{proof}
The inner centralizer of $B$ is trivial by projective absolute
irreducibility and perfection. Modulo inner automorphisms the point-map
group is abelian, with at most two diagram choices at any height.
Consequently two maps of the same height and diagram choice fixing $B$
agree, and all fixing roots commute. Conjugating a root by a candidate
lift changes no outer class and still fixes $B$ pointwise. The same
argument shows that the lift commutes with the root and normalizes
$S(\Theta)$.

An inner element normalizing $S(\Theta)$ likewise commutes with $\Theta$,
by applying this centralizer argument to the natural absolutely
irreducible fixed simple. Since $G_i$ is adjoint, that element itself
lies in $G_i^\Theta$. The root has height $f_i/s<f_i$, so the order
comparison makes its normalizer proper in the ancestor. We have thus
obtained actual proper labels containing the candidate.

Choose one root $\Theta_s$ for each occurring prime. For a product $d$
of primes already chosen, select integers $a_s$ with
$\sum_{s\mid d}a_s d/s=1$ and put
\[
 \Theta_d=\prod_{s\mid d}\Theta_s^{a_s}.
\]
The roots commute, and $\Theta_s^s=\Phi_i$ gives
\[
 \Theta_d^{d/s}=\Theta_s,\qquad
 \operatorname{ht}(\Theta_d)=f_i/d.
\]
For example, raising the product to $d/s_0$ gives
$\Theta_{s_0}^{a_{s_0}d/s_0}$ times
$\Phi_i^{\sum_{s\ne s_0}a_s d/(ss_0)}$, which is exactly
$\Theta_{s_0}$. This verifies the identity also with diagram components.
For the empty product set $\Theta_1=\Phi_i$.

The group $S(\Theta_d)$ is fixed by each previous root and lies in all
their normalizer labels. If it normalized the fixed simple for a new
prime root $\Theta_l$, the preceding inner-normalizer argument would
make it fixed by $\Theta_l$ as well. It would then lie in the fixed group
of the combined root $\Theta_{dl}$, contradicting
\[
 |S(\Theta_d)|>|G_i^{\Theta_{dl}}|.
\]
Each new prime therefore makes the label intersection strictly smaller.
There are at most as many primes as the interval height, and at most
two roots at each associated height. The number of normalizer labels
is consequently at most twice that height.
\end{proof}

\begin{corollary}[Equal natural modules force a fixing root]\label{w:natural-root}
Suppose $B<S_i$ is a same-series high classical candidate of the same natural dimension as $S_i$, and its projective action on $V_i$ is its own natural module, allowing a Frobenius twist and, in type A, a diagram natural/dual twist. Then some prime root of $\Phi_i$ fixes $B$ pointwise. Thus the absence of such fixing roots excludes this case.
\end{corollary}
\begin{proof}
We record why root-free protection precludes equal-natural embeddings within one root series. If the action of \(B\) on \(V_i\) is projectively its own natural module (allowing Frobenius twists, also the diagram natural/dual twist in A), both groups in the same classical series and dimension, then \(B=S(\Psi)\) inside \(G_i\) for a standard Steinberg map up to conjugacy. Indeed identify the natural spaces over \(k\) by the representation equivalence (split Frobenius twists preserve the standard finite images); forms match up to scalar by uniqueness, also the quadratic when needed by square-difference uniqueness on the finite group. This identifies the corresponding algebraic adjoint groups. Divide the height of \(\Phi_i\) by that of \(\Psi\); after composing off that power of \(\Psi\) we still fix \(B\) pointwise, thus cannot have positive remainder by the order comparison. At height zero, an inner centralization is trivial; so is a possible D-diagram centralization (acting in \({\rm PGL}(V_i)\)); and a non-inner A-graph centralization would give projective self-duality of the finite natural representation, impossible by the distinct restricted weights and perfection. So \(\Phi_i=\Psi^z\), \(z>1\) by strict inclusion, and some prime root fixes \(B\).
\end{proof}

\subsection{Simultaneous protection in a pruned chart}

\begin{definition}[A protected classical run]\label{w:protected}
A consecutive high classical run of one defining characteristic is \emph{protected} if for every ancestor--descendant pair within the run the descendant has one tensor atom on the ancestor natural module and no prime root of the ancestor Steinberg map fixes it pointwise. The path also retains all properties of \Cref{w:prepared-path}.
\end{definition}

\begin{proposition}[Selecting all ancestor protections]\label{w:selection}
Assume that the current same-characteristic classical run is already protected through its current node $W$. Fix an upper bound $D$ on the number of path steps and a fixed number $c$ of continuation colors. There is a finite extension rank $r$, depending only on $D,c$ and earlier fixed requirements, such that from every suitably pruned private chart at a current high classical node one can select a next proper extension. If that extension continues the current same-characteristic run, it has both protections relative to every ancestor in that run, including the current node. Earlier structural exclusions are preserved.
\end{proposition}
\begin{proof}
Work inside a successful pruned flat of extension rank $r$ at the current
node $W$. All its comparable candidate labels have the prepared
properties and actual nested simples, by \Cref{sec:patterns}.
Color its proper extensions with at most $c$ colors, recording in
particular whether their socle continues the current same-characteristic
high-classical run. Consider a uniform random flag at extension ranks
\[
 \ell,2\ell,\ldots,(c+1)\ell,\qquad r=(c+1)\ell.
\]
Every flag has an equal-colored pair. Hence some two ranks $a<b$ have
matching probability at least $\binom{c+1}{2}^{-1}$. Conditioning on the
$b$-extension gives a particular $U$ for which at least that fraction
of the $a$-extensions $L\subset U$ have its color. If this color does
not continue the run, any such $L$ needs no additional protection.

Otherwise apply the two counts to $B_U$ in each ancestor $i$ of the
current run, including $W$ itself. These lists detect every failure of
protection at a matching $L$. Indeed, if $t(B_L)>1$, a prime congruence
coarsening of its atoms has a stabilizer containing $K_L$. Since
$B_U\le B_L$ fixes its factors individually, each factor groups atoms
of $B_U$. The partition is invariant under the single Frobenius cycle;
a cycle has a unique invariant partition into $s$ congruence classes.
It is therefore one of the coarsenings counted for $B_U$. A prime root
fixing $B_L$ likewise fixes $B_U$ and occurs on its root list.

None of the counted labels contains $K_W$. For an earlier ancestor,
such containment would make $B_W$ fixed by the root, or individually
fix the coarsening slots: its slot permutation commutes with the cycle
and is trivial by perfection. Either conclusion contradicts the
protection already imposed on $W$. For the current ancestor, properness
of the counted labels gives the same noncontainment. Thus each bad label
is $K_Z$ for a subflat $Z\subseteq U$ with $Z\not\subseteq W$.

For a uniform $a$-extension $L$ in $U$, the probability of $Z\subseteq L$
is at most $2q^{a-b}$: it suffices to require one fixed point of
$Z\setminus W$ in the ordinary affine-flat count. Here $q$ is the chart
field size. The two counts give $O(1+\operatorname{ht}[K_U,K_{W_i}])$
bad labels per ancestor. Reverse rank differences bound this height by
$Dr$ over at most $D$ path steps. The total number is therefore at most
$CD(1+Dr)$ for an absolute $C$ after the colors and thresholds are fixed.
Since $b-a\ge\ell$, their union occupies a fraction at most
\[
 2CD(1+D(c+1)\ell)q^{-\ell}
\]
of those extensions. This tends to zero with $\ell$, whereas the
matching-color fraction is at least $\binom{c+1}{2}^{-1}$. Choose $\ell$
large enough also for the earlier fixed requirements. A matching $L$
outside the union has one atom and no fixing prime root relative to
every ancestor in the run. It retains the structural exclusions because
it lies in the successful pruned flat.

There is no circular choice of parameters: the following path bound
concerns paths already satisfying these group protections. Its absolute
constant is obtained first, then $D$ is fixed, and only then are $\ell$
and $r$ chosen.
\end{proof}

\subsection{From finite inclusions to rational compositions}

The tensor protection makes each edge a single ordinary $p$-digit.
To compose two edges rationally, the lower algebraic group must preserve
the middle group's defining form. Bilinear forms follow from
restricted-weight uniqueness. Quadratic forms in characteristic two
require the next two lemmas. Root protection will have a separate role:
it excludes natural coefficient-one edges by \Cref{w:natural-root}.

\begin{lemma}[Restriction is injective in degree-one cohomology]\label{w:h1-injective}
Let $J$ be a simply connected simple classical algebraic group in the
high-rank range of \Cref{w:classical-conventions}, over
$k=\overline{\mathbf F}_p$. Let
$\Phi=\tau\operatorname{Fr}_p^f$ be an ordinary pinned diagram--field
Steinberg map with $f\ge1$, parameter $P=p^f$, and diagram action $\tau$
on weights. If $D'=L_J(\xi)$ is a nontrivial $P$-restricted simple
rational $J$-module, then
\[
 H^1(J,D')\longrightarrow H^1(J^\Phi,D')
\]
is injective. The finite group here is the full simply connected fixed group.
\end{lemma}
\begin{proof}
Put $\mathcal G=k[J]^{J^\Phi}$, with right finite invariance and left
translation. It suffices to show
\[
 (D'\otimes(\mathcal G/k))^J=0.
\]
Indeed, a finite-fixed vector splitting an extension of $k$ by $D'$
defines the equivariant function sending $g$ to its translate of that
vector. It is right finite invariant, and its image in the quotient
module is constant. Modulo constant functions it therefore lies in the
displayed invariant space. Vanishing makes the function itself constant,
giving a rationally fixed lift.

The Lang map $g\mapsto g\Phi(g)^{-1}$ identifies the right quotient
$J/J^\Phi$ with $J$ \citep{Lang1956,Steinberg1968}
\citep[Definition~17.95 and Theorem~17.96]{Milne2022}.
It is surjective and étale since $d\Phi=0$, hence faithfully flat.
Its fibres are exactly the right cosets of the finite reduced
fixed-point scheme: $h=gx$ with $x=\Phi(x)$. Descent of regular functions
identifies $\mathcal G$ with the regular function algebra carrying the
twisted left-right action.

The Donkin--Koppinen filtration begins with constants and has tensor
pairs of induced highest-weight modules as its sections
\citep[Theorem~4.3 and its preceding construction]{Mathieu2000}
\citep[Section~2.4]{BendelNakanoPillen2011}. Pulling this two-sided
filtration back to the twisted diagonal gives an exhaustive filtration
of $\mathcal G/k$ with sections
\[
 \nabla(\nu)\otimes\nabla(\nu^*)^{[\Phi]},
 \qquad \nu\ne0\ \text{dominant}.
\]
Here $*$ denotes highest-weight dual, $[\Phi]$ pullback, and $\Delta$
below is a Weyl module. A nonzero invariant after tensoring a section
by $D'$ gives a nonzero map
\[
 \Delta(\nu)^{[\Phi]}\longrightarrow L(\xi)\otimes\nabla(\nu).
\]
The source is still generated by its highest vector, since $\Phi$ is
surjective on algebraic-group points. The map is nonzero there, so
\[
 P\tau(\nu)\le\xi+\nu
\]
in dominance order. The source is also trivial on the $f$-th Frobenius
kernel $J_f$. Its nonzero image therefore gives
$\operatorname{Hom}_{J_f}(L(\xi)^*,\nabla(\nu))\ne0$, with kernel
and invariance understood as group schemes.

Since $\xi$ is $P$-restricted, the restricted-simple theorem and
Steinberg decomposition force a rational composition factor of
$\nabla(\nu)$ whose highest weight has the form
\[
 \xi^*+P\gamma\le\nu,\qquad \gamma\text{ dominant}.
\]
Indeed, every rational simple restricts to copies of the $J_f$-simple
for its $P$-restricted part, and those restricted simples are pairwise
distinct. Pairing the two inequalities with $\rho^\vee$ gives the
diagram- and dual-invariant height comparisons
\[
 (P-1)\operatorname{ht}(\nu)
 \le\operatorname{ht}(\xi)\le\operatorname{ht}(\nu).
\]

Only $P=2$ and equality remain. Then $p=2$, $f=1$, and $\xi$ has an
active coefficient one. Equality gives $P\tau(\nu)=\xi+\nu$, so the
image of the source highest vector lies on the unique tensor highest
line. Let $v_\xi\otimes v_\nu$ span it and choose $\alpha$ with
$\langle\xi,\alpha^\vee\rangle=1$. The lowering operator has value
\[
 (f_\alpha v_\xi)\otimes v_\nu
   +v_\xi\otimes(f_\alpha v_\nu).
\]
The first summand is nonzero by the simple-root bracket, and the
first-factor weights $\xi-\alpha$ and $\xi$ differ. The summands cannot
cancel, including in characteristic two. The line is not $J_1$-fixed,
contradicting the kernel-trivial source action. Thus no section has
an invariant. Every invariant would occur in the exhaustive filtration
and have a nonzero successive image, so the required vanishing follows.

The restricted-kernel and regular-bimodule results are the simply
connected theorems in \citet[Part II, Chapter 3 and II.4.20]{Jantzen2003};
see also \citet[Section~1.4]{LinNakano2006} and
\citet[Theorem~4.3]{Mathieu2000}. The pullback to the ordinary
diagram-twisted diagonal uses no complete reducibility.
\end{proof}

\begin{lemma}[The quadratic form in characteristic two]\label{w:quadratic}
Let $J$ be simply connected of sufficiently high type $D_e$ in characteristic two, with ordinary pinned map $\Phi=\tau\operatorname{Fr}_2^f$, $f\ge1$. Let $\eta\ne0$ be $2$-restricted, and consider $L_J(2^y\eta)$. Suppose a quadratic form $Q$ on this module is invariant under the full finite cover $J^\Phi$, and its polar form is rationally $J$-invariant. Then $Q$ itself is rationally $J$-invariant.
\end{lemma}
\begin{proof}
Untwist $2^y$ to $L_J(\eta)$; split Frobenius permutes $J^\Phi$
because $\Phi$ is a standard pinned graph/field map. Exact invariance
under the chosen finite linear lift holds by perfection: its similitude
scalar is a character. Rational invariance of the polar form makes
\[
 c(g)=gQ-Q
\]
a regular cocycle with values in the subspace of squares of linear
forms. We will take its coefficientwise square root to obtain a cocycle
in $L_J(\eta)^*$.

The target Lie action on squares is zero, so the differential of $c$
at the identity kills commutators. The Lie algebra of simply connected
type D is perfect also in characteristic two: simple coroots span the
torus part and are opposite-root brackets, and every D-root pairs oddly
with some coroot. Thus this differential, and by translation every
differential of $c$, vanishes.

Each coordinate of $c$ is consequently a square in the function field,
using the kernel of differentials over a perfect field. Its square root
is regular by normality. Taking roots of the coefficients in a fixed
dual basis therefore gives a regular linear-form cocycle, zero on
$J^\Phi$. Additivity, injectivity and equivariance of squaring on
$k$-points give its cocycle identity.

The square-root cocycle takes values in the nontrivial restricted dual module after untwisting. By \Cref{w:h1-injective}, its class is zero, so it has the form $g\ell-\ell$ for a linear functional $\ell$. Its vanishing on $J^\Phi$ makes $\ell$ finite-fixed. The restricted finite module is nontrivial and irreducible, so $\ell=0$. Thus the square-root cocycle, and therefore $gQ-Q$, is identically zero. The perfect-Lie-algebra argument above is used here only for type D in characteristic two.
\end{proof}

\begin{proposition}[Composing the rational highest-weight actions]\label{w:rational-composition}
Let $i<j<l$ lie in one protected run, and suppose these three nodes have the same series A, B, C, or D. Use consistent natural realizations and simply connected covers of the actual nested simple subgroups. Write $G_s=\widehat G_j$ and $J=\widehat G_l$. The lifted edge representations $ij$ and $jl$ have highest weights
\[
 p^x\lambda,\qquad p^y\eta,
\]
where $\lambda,\eta\ne0$ are $p$-restricted. The latter representation induces an algebraic homomorphism $J\to (G_s)_{\rm ad}$ which lifts to $J\to G_s$. The resulting rational composition of $L_{G_s}(p^x\lambda)$ is irreducible. Its highest weight $\mu$ is one nonzero $p$-digit times a power of $p$, and that digit has the shape, up to diagram, of the direct edge $il$.
\end{proposition}
\begin{proof}
Consider path nodes \(i<j<l\) in one protected run, of the same root series A, B, C or D. Use consistent finite-group identifications at each node (actual simple \(B_{W_j}\), its projective natural realization, and simply connected cover, for both its incoming and outgoing edges). Thus the projective restriction along \(i\to j\to l\) is just via inclusion of the same nested simples as \(i\to l\); the weight colors below also respect diagram changes. Write \(G_s=\widehat G_j\) and \(J=\widehat G_l\) temporarily for the simply connected algebraic groups. Choose pinned coordinates with standard graph/field Frobenius on each (conjugating the finite realizations accordingly if necessary). The edges \(ij,jl\), evaluating natural modules on descendants projectively, have highest weights \(p^x\lambda,p^y\eta\) respectively for the lifts, \(\lambda,\eta\ne0\) \(p\)-restricted, since each has one atom. We claim \(L_J(p^y\eta)\) on the natural space for \(j\) gives a morphism \(J\to (G_s)_{\rm ad}\) inducing the given projective action on finite points. This is immediate in A. In the remaining types the finite representation preserves the bilateral form (perfection); by restricted-weight uniqueness it is rationally self-dual, and the invariant form rationally therefore agrees up to scalars. This suffices except for D in characteristic two, where \Cref{w:quadratic} supplies rational invariance of the quadratic form as well.

The morphism to the adjoint group (in form types one can first map by the given linear representation to the connected classical isometry group, using connectedness of \(J\), then take its adjoint quotient) lifts to \(G_s\) by simply-connectedness. For example pull back the central isogeny; the reduced identity group in the pullback maps surjectively to \(J\) (finite surjection, finitely many components), by a central isogeny with kernel of multiplicative type; its radical has trivial image and hence dimension zero, so it is semisimple and maps isomorphically by the simply connected root datum of \(J\). Its natural action agrees linearly with the above one (no scalar characters). On finite points the lift differs at most by central elements from lifts to the usual finite covering group at node \(j\). Indeed both give the prescribed projective matrices on the natural space (the image of the finite cover surjects projectively onto the given socle), and the kernel of the projective natural map on simply connected algebraic-group points is central. Thus composing \(L_{G_s}(p^x\lambda)\) yields on restriction to \(J^\Phi\) projectively the direct-edge action \(il\). The rational composition is irreducible by finite irreducibility. Its highest weight \(\mu\) for \(J\) is again a single \(p\)-digit times a power of \(p\), by Steinberg decomposition and the protected individual tensor indecomposability on finite points. Indeed multiple rational digits here would give multiple nontrivial-dimensional factors fixed individually by the finite cover, hence also as matrix factor algebras by the finite projective \(S_l\)-action (and this projective action already agrees as just explained). Moreover the digit shape agrees up to diagram with the direct edge's: reduce its Frobenius twist exponent on the finite group using the defining graph/field map and apply restricted-weight uniqueness (linear lifts agree by perfection).
\end{proof}

\subsection{Universal cuts of the full weight multiset}

\begin{lemma}[Chamber maxima and universal cuts]\label{w:universal-cuts}
Use the triple, groups, and weights of \Cref{w:rational-composition}. Let $\Omega$ be the full weight multiset of $L_J(\eta)$, with every multiplicity retained. After mapping a maximal torus of $J$ into one of $G_s$, the natural weights restrict to $p^y\Omega$. One Weyl translate of $\lambda$ maximizes throughout the positive cocharacter chamber of $J$, and on each active fundamental summand of $\lambda$.

In type A an active node $a$ of the upper weight $\lambda$ therefore
gives a cut selecting $a$ entries of $\Omega$, each dominating every
unselected entry, with no equal-vector tie across the cut. In B and C
the analogous signed fundamental cut selects $a$ entries, including
the last node up to the B spin scale; its selected weights are nonzero
and in the positive cone. This also holds at the nonspin D nodes when
$\lambda$ has no spin support. At an active D spin node the maximum
instead chooses one orientation $u_a$ of each natural signed pair, and
\[
 u_a+u_b\text{ is nonnegative on the chamber whenever }a\ne b.
\]
Distinct ambient orbit weights cannot restrict to the same highest weight. More generally, every weight of a simple highest module differing from highest by a multiple of one simple root has multiplicity at most one.
\end{lemma}
\begin{proof}
Map a maximal torus of $J$ into one of $G_s$. The restricted natural
multiset is $p^y\Omega$. By irreducibility of the composition and the
highest-weight convex hull property, the maximum pairing on the orbit
of $p^x\lambda$ is the linear function given by $\mu$ throughout the
positive cocharacter chamber of $J$. Choose a maximizing orbit weight
at a strictly interior cocharacter. Its restriction must equal $\mu$
as a vector, so that same Weyl translate $w\lambda$ maximizes everywhere
in the chamber.

It also maximizes each active fundamental summand. For any one
cocharacter a common Weyl translate maximizes all fundamental summands;
equality in their positive sum therefore forces equality for each.
Two distinct ambient orbit weights cannot both restrict to the
composed highest weight, whose multiplicity is one.

Suppress the positive scales $p^x,p^y$ for the cut description. In type
A an active upper node $a$ selects $a$ natural coordinates at $w\lambda$.
Exchanging one selected and one unselected coordinate shows, throughout
the lower chamber, that every selected restriction dominates every
unselected restriction. An equal-vector tie would give another
maximizing ambient orbit weight across a strict coefficient boundary,
contradicting highest multiplicity one.

In B and C the analogous choice takes $a$ entries of the full signed
multiset, also at the last node up to the B spin scale. In D the same
argument applies before the spin nodes when $\lambda$ has no spin
support. Sign changes show that selected restrictions lie in the
positive cone; in D compensate by an unused coordinate. A zero-vector
restriction would permit a sign change in the full maximizing orbit
weight, using a zero-coefficient coordinate for D parity, and again
contradict multiplicity one. Thus the selected values are nonzero.
Signed coordinate exchanges compare them with every unused value,
including negative choices and the extra zero in B. An equal nonzero
tie across the cut also changes the ambient orbit weight; parity can
be preserved by changing both directed signs when needed. These cuts
are therefore strict as asserted.

For D spin support, take an active spin fundamental summand instead.
Its maximizing signed choices $u_i$ satisfy the stated pair-sum
condition, because flipping any two choices cannot increase the
pairing on the chamber.

These comparisons concern all entries of $\Omega$, including nonextremal weights and repeated entries. A cut of its Weyl orbit alone is only a necessary consequence. The single-root multiplicity assertion follows because the corresponding weight space of the highest-weight module is generated by the one divided-power monomial in that root; it does not assert multiplicity one for arbitrary weights.
\end{proof}

\subsection{The word argument in type A}

\begin{lemma}[Type A orbit cuts]\label{w:a-word-cuts}
Let $J$ have sufficiently high type $A_{e-1}$, and let $0\ne\eta$ be $p$-restricted. Suppose a nonempty proper universal cut of the full weight multiset of $L_J(\eta)$ is given as in \Cref{w:universal-cuts}. Replace the cut by the longest-element image of its complement if necessary so its selected part of the extremal orbit has size at most half that orbit. Unless $\eta$ is supported at one endpoint, the selected orbit words consist of the highest word alone, possibly its single boundary swap for a binary word, or the positive half of the orbit of the exact highest root. In the final case the highest weight is exactly $\epsilon_1-\epsilon_e$, not a larger restricted multiple.
\end{lemma}
\begin{proof}
Take \(J\) of type \(A_{e-1}\). Its orbit on \(\eta\) consists of permutations of a descending nonconstant word, ordered for dominance by prefix sums (coordinate sum zero). In a universal cut of \(\Omega\), the upper orbit words form a nonempty proper subset \(U\), since the orbit contains the full highest and lowest weights. We may study the smaller orbit side, replacing the full cut by the longest-element reversal of its complement when needed (this is again an upper cut of the weight multiset, now of complementary size) so \(|U|\) on the orbit is at most half. Then \(U\) is an initial segment both for lexicographic decreasing order and for lexicographic increasing order starting from the right. This follows by the first differing coordinate comparisons from strict cone domination and equality of total sums.

Here is the description needed of such segments. Let \(M\) be the largest letter, multiplicity \(b\), and \(m\) the smallest, multiplicity \(c\). If \(U\) contains a word not starting with \(M\), it contains all starting with \(M\). If \(b\ge2\) this includes one also ending with \(M\), forcing all ending with lower letters; the two fractions (and extra words) contradict the half bound. Thus \(b=1\). With second-largest letter \(T\), a starting-\(M\) word ending in \(T\) forces in all ending below \(T\). If there are at least two such lower letters counting multiplicity, these include one also starting below \(T\), forcing all starting at or above \(T\). Again the fractions contradict the bound. Namely if \(l\) letters are below \(T\), the proportions ending below \(T\) and starting at or above \(T\) are \(l/e\) and \(1-l/e\); in the indicated situation \(U\) contains each such class entirely and strictly. We obtain only a natural two-value shape with singleton extreme against constant, or \(M,T^{e-2},m\). Symmetrically these are the exceptions if some selected word does not end with \(m\).

Otherwise all words in $U$ start with $M$ and end with $m$. Until the
first word not ending with $m$, left lex order fixes the first $e-c-1$
coordinates of the descending word: the greatest word not so ending
already has those coordinates. Right lex order confines changes to
the first $b+1$. Only positions in $[e-c,b+1]$ can therefore change.
This interval has at most one position unless $b+c=e$, the binary
case; then its two positions permit only the single boundary swap.
Thus, outside the exceptional shapes, $U$ consists of $\eta$ alone or
of $\eta$ and that swap.

For \(M,T^{e-2},m\), use positions \(a,d\) of the largest and smallest letters. Both orders begin with the half \(a<d\), one prioritizing \(a\) increasing there and the other \(d\) decreasing. A proper nonempty common prefix of the half has only \((1,e)\) (taking \(a>1\) forces \((1,2)\), then the full half; similarly for \(d<e\)). Dominance across the cut consisting of the half requires \(M-T=T-m\), by comparing words with both defects early and both defects late, in both choices of early side. At a prefix separating the two defects of one word (both before the boundary) from those of the other (both after) this compares \((M-T)+(m-T)\) to zero, in both directions. Thus \(\eta=z(\epsilon_1-\epsilon_e)\) in that case (\(\epsilon_i\) standard coordinates). In the full weight set each selected weight then has all prefix sums nonnegative (dominance over negative scaled simple roots on the other side), and each unselected one has them nonpositive. If \(z>1\), the weight \(\eta-\alpha_1\) occurs by simple-root lowering (\(z<p,\ \alpha_1=\epsilon_1-\epsilon_2\)), and can be permuted to mixed prefix signs. So this half exception requires \(z=1\).
\end{proof}

For edges whose descendant has type A, color the nonzero restricted digit by the following five possibilities:
\begin{enumerate}
\item one active endpoint with coefficient one;
\item one active endpoint with coefficient greater than one;
\item one active interior node;
\item multiple active nodes, excluding the exact highest root;
\item the exact highest root.
\end{enumerate}
This palette is invariant under the A diagram reversal.

\begin{proposition}[No homogeneous type A configuration]\label{w:a-obstruction}
Four nodes of a protected run of type A cannot have all pair edges of one color in the preceding palette.
\end{proposition}
\begin{proof}
Take the first three of four homogeneous nodes, with the notation of
\Cref{w:rational-composition}. For each color we first constrain the
extremal orbit cut, then use actual full weights or their multiplicities.

\paragraph{Multiple support and the highest root.}
In the multiple nonexact color, the orbit lemma leaves only $\eta$ on
the smaller side of every cut. A full selected weight $\nu$ must
dominate every active simple reflection $s_h\eta$. Thus the nonnegative
difference $\eta-\nu$ is supported only at $h$ for each active $h$.
Multiple support forces $\nu=\eta$, whose multiplicity is one. Every
active node of the upper digit $\lambda$ must consequently be an
endpoint; since it has multiple support, both endpoints are active.
Apply the same argument to the last three nodes to obtain that
conclusion for $\eta$ as well. In the exact highest-root color both
digits already have these two active endpoints.

Choose natural coordinates at a maximum of $\lambda$ whose first,
second, penultimate and last restrictions, before the lower twist, are
\[
 \eta,\quad \eta-\alpha_1,\quad
 w_0\eta+\alpha_1,\quad w_0\eta,
\]
respectively. Endpoint activity supplies these distinct weights in
the lower module. The middle coefficients of $\lambda$ are constant,
so their placement preserves the maximum. Write $\theta=w\lambda$
for this maximizing weight, and $\beta_1,\beta_{v_j-1}$ for the two
endpoint simple roots in that coordinate order. Restricted activity
makes $\theta-\beta_1$ and $\theta-\beta_{v_j-1}$ distinct weights of
the upper module. Their restrictions satisfy
\[
 \operatorname{res}\bigl(p^x(\theta-\beta_1)\bigr)
 =\operatorname{res}\bigl(p^x(\theta-\beta_{v_j-1})\bigr)
 =\mu-p^{x+y}\alpha_1.
\]
The two ambient weight spaces stay independent on restriction. They
give multiplicity at least two along one simple-root ray from the
composed highest weight, contrary to \Cref{w:universal-cuts}.

\paragraph{One interior node.}
Write $\lambda=c\omega_a$ and $\eta=z\omega_h$, with $a$ and $h$
interior in their respective diagrams. Replace the full cut by the
longest-element image of its complement if required by
\Cref{w:a-word-cuts}. Its selected orbit words are $\eta$, possibly
together with $s_h\eta$.

Every full selected weight $\nu$ differs from $\eta$ only in the
$\alpha_h$ direction. If $s_h\eta$ is outside, domination of that
word proves this. Otherwise use the two outside words
$s_{h-1}s_h\eta$ and $s_{h+1}s_h\eta$: their downward supports intersect
only at $h$. The actual module also contains
\[
 \zeta=\eta-\alpha_h-\alpha_{h-1}.
\]
To obtain it, lower at $h$ and then across the simple bond to its left,
where the new coefficient is one. This weight is outside the cut,
because its downward support is not confined to $h$. Domination of
$\zeta$ leaves only $\eta$ and $\eta-\alpha_h$ on the selected side,
each of multiplicity at most one.

The cut has $a$ entries, or $v_j-a$ in the complemented case. Both
numbers are at least two, so these are exactly its two entries.
Since the full multiset $\Omega$ sums to zero, complementation and
the longest element replace a sum by its highest-weight dual. Hence
\[
 \mu=c\,p^{x+y}(2\eta-\alpha_h)
 \quad\hbox{or}\quad
 \mu=c\,p^{x+y}(2\eta-\alpha_h)^*.
\]
Both neighboring fundamental coefficients of $2\eta-\alpha_h$ are
positive. No common Frobenius rescaling removes either support node,
contradicting the one-node digit shape of the direct edge.

\paragraph{One endpoint.}
In the coefficient-one color, the edge is a natural or dual-natural
module of equal dimension, excluded by \Cref{w:natural-root}.
Otherwise the upper digit module, up to duality, is a symmetric power
of degree $c$ with $2\le c<p$. The lower natural or dual-natural
restriction has highest weight $\eta_0$, equal to $\eta$ or $\eta^*$,
and a highest simple-root string containing vectors $a_0,a_1,a_2$
of weights $\eta_0,\eta_0-\alpha,\eta_0-2\alpha$, since its restricted endpoint
coefficient is also greater than one. The distinct monomials
\[
 a_0^{c-1}a_2,\qquad a_0^{c-2}a_1^2
\]
both have weight $c\eta_0-2\alpha$ before the common twists. They
contradict the same single-root multiplicity bound. Thus every color
is impossible.
\end{proof}

\subsection{The signed series}

\begin{lemma}[Signed orbit cuts]\label{w:signed-cuts}
Let $J$ have sufficiently high type B, C, or D, and suppose $L_J(\eta)$ is self-dual, with $\eta\ne0$ restricted. In type D suppose for this lemma that no spin node of $\eta$ is active. Let a signed fundamental cut as in \Cref{w:universal-cuts} be given, including the last B and C nodes and excluding only the D-spin case, and let $U$ be its selected extremal orbit words. Use nonnegative dominant magnitudes, write $h$ for their nonzero count, and let $O^+$ be the unsigned permutation orbit.
\begin{enumerate}
\item If some unsigned word is unselected, then $U\subset O^+$. For $2\le h<e$, $|U|\le2$, with the binary possibility $\eta,s_h\eta$; for $h=e$, needed only in B and C, $U=\{\eta\}$.
\item If every unsigned word is selected and $h<e$, then either $h=1$ or $h=2$ with the two nonzero magnitudes equal.
\item If every unsigned word is selected and $h=e$ in B or C, all magnitudes are equal and $|U|\le3$.
\end{enumerate}
\end{lemma}
\begin{proof}
Now \(J\) is B,C or D of rank \(e\), and dominance in particular requires the prefix comparisons in the usual ordered signed coordinates. The orbit of \(\eta\) is antipodal by self-duality; make its dominant magnitudes nonnegative using the D diagram if necessary. First suppose no D spin nodes active (case of spin activity handled separately). For a universal fundamental cut above in \(\Omega\), upper orbit words \(U\) include \(\eta\), and each has all prefix sums nonnegative and is nonzero. Write \(h\) for the number of nonzero magnitudes in a word, and \(O^+\) for all unsigned permutations.

If some unsigned word is outside \(U\), comparison of total sums gives \(U\subset O^+\). If \(h<e\), we have all choices of signs (even in D, by zeros). Before the earliest negative word in decreasing lex order the first \(h-1\) positive coordinates are fixed, hence \(|U|\le e-h+1\). Indeed in this case all negative words are outside \(U\), and one can put the largest \(h-1\) magnitudes first positive, with the final nonzero one negative later; left prefix domination gives the lex bound. For \(2\le h<e\) in high rank this is at most half the unsigned words (there are at least \(\binom eh\)). In this range apply the A-word prefix description among those equal-total words. We obtain at most two upper words (the adj-shaped half for a word \(M,T^{e-2},m\) of three distinct magnitudes would be too large by \(e-h+1\), and the remaining binary singleton-extreme exception in this range has a singleton zero, already giving bound two by \(e-h+1\)). For a binary shape in that range with equal nonzero coordinates they are \(\eta\) and possibly \(s_h\eta\), also when \(h=e-1\) by left lex order. If instead \(h=e\), here only needed for B,C, the first \(e-1\) positives are fixed before the first negative word, giving just \(\eta\).

Suppose all unsigned words are selected. If $h<e$, every signed word
with positive first coordinate is forced in: otherwise it fails the
first-prefix comparison with a selected unsigned word starting with
zero. For $h\ge3$, put a smallest nonzero magnitude first positively
and two others negatively; the third prefix is negative. For $h=2$
with unequal magnitudes, the smaller positive followed by the larger
negative gives the same contradiction. Thus only $h=1$ or two equal
nonzero magnitudes can occur.

For $h=e$ in B or C, suppose the smallest and largest magnitudes are
$P<M$, and let $S$ be the sum of all $e$ magnitudes. The signed word
with $M$ first and all other entries negative must be selected, by
comparison with a selected unsigned word starting with $P$. The word
starting with $-P$ and having all other entries positive is unselected,
since its first prefix is negative. Domination of their total sums
would require
\[
 2M-S\ge S-2P,\qquad\text{hence}\qquad S\le M+P,
\]
impossible with at least three positive coordinates. All magnitudes
are therefore equal. A selected word cannot have two or more negatives,
since total-sum comparison would then force in every single-negative
word, including one with negative first prefix. Left lex order before
the first two-negative word now leaves only the all-positive word and
the words with one negative in either of the last two positions.
Thus $|U|\le3$.
\end{proof}

For a fixed series B, C, or D, use the following edge colors:
\begin{enumerate}
\item a single node $1$, with coefficient one or greater than one as separate colors;
\item a single interior nonspin node, separating index $2$ from larger indices;
\item the single last node in B or C;
\item any spin activity in D;
\item all remaining multiple support.
\end{enumerate}
Here ``interior'' in B and C excludes the last node. B and C are separate vertex types, and B's last-node color includes its spin weight. The D colors are invariant under interchange of its spin nodes.

\begin{proposition}[No homogeneous signed configuration]\label{w:signed-obstruction}
Four nodes of a protected run having one fixed series B, C, or D cannot have all pair edges of one color in the preceding palette.
\end{proposition}
\begin{proof}
Use the triple notation of \Cref{w:rational-composition}, and consider the possible common edge color.

\paragraph{D spin activity.} For D spin activity on the two edges, \(\eta\) has at least \(e-1\) nonzero coordinates. In its orbit take words starting with a smallest nonzero entry positive and two others negative, then many independent signs on the remaining nonzero ones respecting parity. If all $e$ coordinates are nonzero, parity leaves at least $2^{e-4}$ choices; if one is zero, it absorbs the parity adjustment and gives the same lower bound. Their first prefix is positive and third prefix negative. No two constructed words are antipodes, since the first coordinate is positive in each. Thus for $e\ge7$ this gives more than $e$ distinct antipodal pairs with mixed prefix signs. Self-duality supplies antipodality even in odd D rank, where it forces the last coordinate to be zero when opposition interchanges the spin nodes. Each pair appears (multiplicity one) as a natural ambient signed pair in \(\Omega\). In the spin maximum for \(\lambda\) as above, its chosen orientation therefore contributes a negative prefix somewhere. But at any one prefix at most one coordinate choice can have negative sum, by the pair-sum condition. Contradiction.

\paragraph{Remaining multiple support.} In the remaining multiple-support color, the orbit description forces \(U\) to have just \(\eta\). Indeed \(h=1\), two equal nonzero coordinates only, and all equal full magnitudes are all single-node shapes here. And a bound of two words in the description reduces to just \(\eta\) under multiple support: any additional one would have to dominate the active reflections outside, impossible (its difference down from highest must be supported only at each such node; if itself an active reflection use a different active node). Then every full selected weight also dominates all those reflections, giving just \(\eta\) with multiplicity one. This cannot accommodate the active indices of \(\lambda\).

\paragraph{One interior nonspin node.}
Write $\lambda=c\omega_a$ and $\eta=z\omega_h$, where $a$ and $h$
are interior nonspin nodes in their respective diagrams. Here
$2\le h<e$, with $h\le e-2$ in D. If all unsigned words were selected,
the two-equal-magnitude exception would force $h=2$. The palette then
also gives $a=2$, whereas the cut would contain at least
$\binom e2>2$ unsigned words. Thus only $\eta$, possibly together with
$s_h\eta$, occurs in the selected orbit.

As in type A, every full selected weight differs downward from $\eta$
only at $h$. If $s_h\eta$ is selected, use reflected words from two
distinct Dynkin neighbors; at the D fork either right neighbor suffices.
Intersecting their downward supports does not require the right bond
to be simple. The outside weight
$\eta-\alpha_h-\alpha_{h-1}$ is present by the left simple bond.
After the first lowering its vector is highest for that left rank-one
subgroup and has coefficient one there, so the second lowering is
nonzero in every characteristic. Domination of this outside weight
leaves at most the two entries $\eta,\eta-\alpha_h$, each of
multiplicity one. Since the cut has $a\ge2$ entries, it has exactly
two. The composed highest weight is a positive scalar multiple of
$2\eta-\alpha_h$, with both neighboring nodes active, contrary to
the direct edge's one-node digit shape.

\paragraph{The last B or C node.} For the last single node of B,C the fundamental maximum in the upper group orients every natural signed pair to the upper side. It therefore chooses half the orbit of \(\eta\), which has all equal positive magnitudes before sign changes. This contradicts the bound above.

\paragraph{Node one.} Coefficient one gives the prohibited natural embedding. For a single node 1 with both coefficients larger than one, write \(\lambda=z\omega_1\). Assign \(\epsilon_1,\epsilon_2,\epsilon_3\) at a maximum respectively to \(\eta,\eta-\alpha_1,\eta-2\alpha_1\), up to twist (distinct, nonzero and nonopposite, adjusting unused signs if needed). The top \(A_2\)-face in the upper module before twist is the natural symmetric power of degree \(z\). Thus the distinct weights \((z-1)\epsilon_1+\epsilon_3\) and \((z-2)\epsilon_1+2\epsilon_2\) both occur in these coordinates and restrict to $z\eta-2\alpha_1$ before the common twists. Their two weight spaces give multiplicity at least two along one simple-root ray from highest after restriction, contradicting \Cref{w:universal-cuts}.

All constructions require only fixed lower bounds on rank: four
coordinates suffice for the type A equal-gap test, three for the signed
total-sum comparison, and seven for the displayed D parity count.
They therefore apply after fixing one sufficiently large absolute
high-rank threshold.
\end{proof}

\subsection{The absolute path bound}

\begin{theorem}[Bound on protected high-type paths]\label{w:path-bound}
There is an absolute constant $D_0$ such that a prepared path all of whose socles are of high alternating or high classical type, and whose consecutive same-characteristic classical runs are protected, has fewer than $D_0$ nodes. This constant is independent of all chart dimensions, finite fields, and path-construction parameters.
\end{theorem}
\begin{proof}
Within a protected classical run, first color vertices by the four series A, B, C, and D. The edge palettes above have an absolute number of colors and are compatible with diagram changes. Finite Ramsey theory \citep{Ramsey1930} therefore bounds the number of nodes in the run: a sufficiently long run contains four same-series nodes whose pair edges are homogeneous, contrary to \Cref{w:a-obstruction,w:signed-obstruction}. All pairs of the selected nodes remain protected because the protection was imposed for every ancestor--descendant pair, not just consecutive nodes.

By \Cref{w:alternating-switch}, at most two high alternating nodes occur. By \Cref{w:characteristic-switch}, there cannot be two successive changes of characteristic among the classical nodes after ignoring the alternating interruptions. Thus there are only absolutely boundedly many same-characteristic classical runs. Applying the run bound to each gives $D_0$.
\end{proof}

\begin{proof}[Completion of the proof of \Cref{f:negative-family}]
Use the almost simple reduction of \Cref{f:thm-almost-simple-reduction} and the prepared path construction of \Cref{sec:patterns}. Choose the high thresholds and $D_0$ first, then take the number $D$ of reserved generations greater than $D_0$, with room for the initial displacement and one remaining bounded-type stage. Choose the lattice field size $q$ above the finitely many absolute line requirements. Next choose the terminal chart increment $r$ using \Cref{w:selection} and \Cref{w:bounded-exclusion}. Choose the earlier trial increments backwards from $r$ as in \Cref{sec:patterns}, then $m$ above their ranks and $n$ large enough for all residual interval lengths and bounded spans. Finally choose $N$ for the initial pattern demand and enough private charts $g$ for all finite losses.

At each stage, \Cref{sec:classical,sec:alternating} supplies a pruned chart with the ancestor properties of \Cref{w:prepared-path}. A bounded-order socle is excluded by the residual interval length. A bounded-type Lie socle is impossible by \Cref{w:bounded-exclusion}. In every remaining case the next high-type node can be selected, with \Cref{w:selection} adding the protections whenever the current classical run continues. This produces more than $D_0$ high-type nodes, contradicting \Cref{w:path-bound}. It proves the required finite large-family nonrepresentation statement. Together with \Cref{sec:small}, this gives a finite minimum strictly greater than seven. Its explicit size ceiling and the description of $L_0$ are supplied in \Cref{sec:quantitative,sec:minimum} after the uniform-decision argument.
\end{proof}

\section{Reverse alternating tests}\label{sec:reverse}

The undecidability argument will use exact subgroup intervals in both
orientations, and intervals of invariant subgroups in powers of a
nonabelian simple group.  We first prove the reverse exclusion needed for
that reduction.  Its hypothesis is an arbitrary exact realization of a
reversed alternating subgroup lattice.  From such a realization we recover
an ordinary site interval, build the linear charts of the path argument
from actual elementary subgroups, and reproduce the structural exclusions.
The forward marker theorem in the next section has a complementary
conclusion: after a bounded prefix, it supplies an actual alternating
simple subgroup with many point-stabilizer labels.

In \Cref{sec:weighted}, markers will be placed beside units carrying grids
and full permutation groups.  Normal and projection tests on nontrivial
lattice images of that template will allow these two theorems to force a
forward product configuration.  Intersections of cuts then identify one
common partition into actual orbit cells, in which equal totals and
squares supply arithmetic.  Finally, \Cref{sec:counting} bounds the number
of lattice elements computably in the number of units, independently of
the formal alphabet sizes.  A computable bound for carriers of realizable
lattices would consequently bound positive Diophantine solutions whenever
they exist.

Throughout, reversal means reversal of the entire displayed parameter
lattice.  Selecting its coordinate-product labels produces a separate
ordinary site interval; it need not preserve all labels of the original
invariant interval.  We continue to use property~\((*)\): every normal
subgroup of the upper endpoint is contained in, or supplemented by, the
base, and the quotient by the base core has a unique minimal nontrivial
normal subgroup, which is nonabelian.  Fixed absolute thresholds need not
be effectively evaluated: they may be fixed constants in an implication
between the existence of algorithms.  All groups in this section are
finite.

\begin{theorem}[Reverse alternating exclusion]\label{r:reverse-exclusion}
For all sufficiently large integers $v$, the lattice
$\operatorname{Sub}(A_v)^{\mathrm{op}}$ has no exact realization as
\[
 [P,S^b]^\Delta,
\]
where $S$ is nonabelian simple, $\Delta$ acts by automorphisms transitively
on the $b$ simple factors, $P$ is a $\Delta$-invariant coordinate product,
and either $P\ne1$ or $b=1$.
\end{theorem}

\subsection{Normal tests and the ordinary site interval}

We first establish the normal test needed at every prepared link.  The
test applies to the full invariant interval through a semidirect product;
we will then transfer its conclusion to the ordinary interval formed by
the coordinate-product labels.

\begin{lemma}[Semiregular alternating normal test]\label{r:alternating-normal}
Let $X\le A_n$ be semiregular and elementary abelian, allowing $X=1$.
For $n$ sufficiently large in terms of $|X|$, every exact finite group
interval isomorphic to $[X,A_n]$ or $[X,A_n]^{\mathrm{op}}$ has
property~\((*)\).
\end{lemma}

\begin{proof}
We first show that a normal subgroup can give only an endpoint label.
The unique nonabelian monolith will then follow from nonmodularity.
Suppose $[H_0,G_0]$ is such an interval, and let $D\in[X,A_n]$ label
$H_0N$, where $N\trianglelefteq G_0$.  This label satisfies the
left-modular identity, in either orientation.  Moreover, projecting
$H_0N$ to $[L,K]\subseteq[H_0,G_0]$, by meeting and joining in either
order, gives
\[
 L(N\cap K),\qquad N\cap K\trianglelefteq K.
\]
The identity
\begin{equation}\label{r:left-modular}
 z\vee(d\wedge y)=(z\vee d)\wedge y\qquad(z\le y)
\end{equation}
is self-dual.  Consequently neither $D$ nor any of these projections
can admit a test
\begin{equation}\label{r:forbidden-test}
 D\cap y\le z<y\le\langle D,z\rangle.
\end{equation}
We show that $D=X$ or $D=A_n$.

Choose $37$ points in distinct $X$-orbits.  Their $X$-translated
$37$-tuples are pairwise disjoint.  Let $T$ be the product of the
alternating groups supported on these tuples.  Thus $X\cap T=1$, and
$[X,XT]$ identifies with the $X$-invariant subgroups of $T$.  For
$L\le A_{37}$ let $T_L$ denote its product in the matched coordinates.
Put $D_0=D\cap T$, and write $V\le A_{37}$ for a coordinate projection
of $D_0$.  The following pair of conditions is impossible whenever
$B_0\le a<c\le A_{37}$:
\begin{equation}\label{r:local-test}
 \langle V,B_0\rangle\cap c\le a,
 \qquad
 c\le\langle a^{\langle V,a\rangle}\rangle.
\end{equation}
Indeed project once more to
$D'=\langle D_0,T_{B_0}\rangle$.  The first condition gives
$D'\cap T_c\le T_a$.  Conjugating independently supported copies of
$a$ by elements whose site projection belongs to $V$ gives
$T_c\le\langle D',T_a\rangle$, contrary to
\eqref{r:forbidden-test}.

Suppose first that $V$ is proper and transitive.  Burnside's prime-degree
Theorem \citep{Muller2005} and the classification of finite
two-transitive groups imply
that $V$ lies in the affine normalizer of a regular cyclic group of
order $37$; in particular it contains a $37$-cycle.  Here the
non-affine possibilities at degree $37$ all contain $A_{37}$.
For completeness, $37$ is not a projective point number
$1+s+\cdots+s^j$ for a prime power $s$: for $j=1$ one would have
$s=36$, for $j=2$ the adjacent possibilities $s=5,6$ give $31,43$,
and for $j\ge3$ the small remaining possibilities likewise miss $37$.
The other classical and exceptional two-transitive degrees give none
at $37$; we use the permutation-group classification in
\citet[Table~1]{MeagherSpigaTiep2016}.

Choose a cyclic $37$-group $a$ so that
$c=N_{A_{37}}(a)$ intersects $V$ trivially.  To justify the choice, both
$V$ and $c$ have order at most $37\cdot18=666$.  A nonidentity affine
permutation fixes at most one point.  From its cycles one can choose
nine disjoint directed moved pairs.  A commuting permutation is
determined on those eighteen points by the ordered images of the nine
sources, and has at most $19!$ choices on the rest.  Its centralizer
in $S_{37}$ therefore has order at most
$(37\cdot36\cdots29)19!$, so its $A_{37}$-class has size at least
$28!/(2\cdot19!)>666^2$.  Averaging the intersections of $V$ with
conjugates of $c$ now gives a conjugate with no nonidentity intersection.
The groups $V$ and $a$ contain distinct regular cyclic $37$-groups;
a proper transitive subgroup would be affine, where the Sylow
$37$-group is unique.  Hence $\langle V,a\rangle=A_{37}$.
Its simplicity gives \eqref{r:local-test} with $B_0=1$, a contradiction.

Suppose next that $V$ is nontrivial and intransitive.  Choose
complementary invariant sets $A,F$ as follows: if there are at least two
singleton orbits, use two of them as $F$; if there is exactly one, use
that one; if there are none, use a smallest orbit.  In all cases
$|A|\ge19$.  Remove two points from $A$ when $|F|\ge2$, and three
when $|F|=1$, leaving a set $Z$ meeting every orbit in $A$.  In the
first case with two fixed points removed to form $F$, this is possible
because a nontrivial even action has
$\sum_O(|O|-1)\ge2$ on its other orbits.  In the remaining cases the
orbits of $A$ are nonsingletons, so the required removals are possible.
Set $B_0=A_Z$.  The $V$-conjugates of $B_0$ generate $A_A$:
their supports cover $A$ and have large overlaps.

Take a $4$-cycle on the removed points and enough points of $F$ to
make four.  When two points come from each set, arrange the two from
$A$ consecutively and the two from $F$ consecutively.  Multiply this
cycle by a disjoint transposition in $Z$, obtaining an even permutation
$g$.  Set
\[
 a=\langle B_0,g^2\rangle,\qquad c=\langle B_0,g\rangle.
\]
No element of $c\setminus a$ preserves $A$, whereas
$\langle V,B_0\rangle$ does.  Thus the first condition in
\eqref{r:local-test} holds.  The normal closure of $a$ in
$\langle V,a\rangle$ contains $A_A$ and its conjugate by $g^2$.
Their overlapping supports generate the alternating group on $A$
with the used points of $F$ added.  This group contains $c$, proving
the second condition and yielding another contradiction.

The two exclusions show that every site projection $V$ is $1$ or
$A_{37}$.  We next rule out diagonal ties in the latter case.
Then $D_0$ is subdirect and in fact equals $T$.
Otherwise the description of subdirect simple powers as products of
twisted diagonals gives a tied pair of sites.  Choose
$a=C_2<c=C_6\le A_{37}$.  Independent nontrivial supports from
$T_a$ break all ties, so $\langle D_0,T_a\rangle=T$.
Modularity would give
\[
 T_c=\langle D_0\cap T_c,T_a\rangle.
\]
On a tied pair, quotient each $C_6$ coordinate by its $C_2$ subgroup.
The image of $D_0\cap T_c$ has size at most $6$, whereas that of the
left side has size $3^2=9$.  The displayed equality is impossible.
Thus every transversal product $T$ is either contained in $D$ or has
trivial intersection with $D$.

If one such $T$ is contained in $D$, then all are.  Change one point
of a transversal at a time, using spare $X$-orbits when needed.
Consecutive transversals share supported three-cycles, so the
alternative of trivial intersection is impossible at the next one.
The supported $A_{37}$'s on the transversals themselves have connected,
overlapping supports covering the letter set.  They generate $A_n$,
and hence $D=A_n$.

Otherwise $D\cap T=1$ for every transversal product, also for shorter
transversals after extending them to length $37$.  We will show that every
element of $D$ then belongs to $X$.  First suppose $d\in D$
sends a point outside its $X$-orbit.  Consider the directed non-loop
moves on the set of $X$-orbits.  If three have six distinct endpoint
orbits, use a three-cycle $u$ on their sources; its $d$-conjugate
$u'$ is supported on their destinations.  If not, endpoints of a
maximal disjoint matching comprise at most four designated orbits,
and every non-loop move meets one of them.  Since each designated
orbit has size $|X|$, only boundedly many orbits, in terms of $|X|$,
are touched.  On sufficiently many untouched orbits, two points in
different orbits have the same $X$-offset under $d$.  Multiplying $d$
by an element of $X\le D$, make it fix those two points.  Use a
three-cycle $u$ on them and one point moved outside its orbit, and
again put $u'=u^d$.  In either case one transversal product $T$
contains $u,u'$, and
\[
 z=XT_{\langle u\rangle}<y=\langle z,u'\rangle\le XT,
 \qquad D\cap y=X,\qquad y\le\langle D,z\rangle.
\]
This contradicts \eqref{r:forbidden-test}.

It remains to exclude $d\in D\setminus X$ preserving each $X$-orbit.
Some transversal five-tuple has its matching of $X$-translated tuples
disrupted by $d$; preservation of every such matching would force all
offsets to be one common element of $X$.  On this matching use
\[
 z=XT_{C_5}<y=XT_{A_5}.
\]
Inside $\langle D,z\rangle$, take the supported five-cycles on the
original matching, their $d$-conjugates, and all $X$-translates.
The connected components of their support hypergraph on the five
orbits have size at least ten: the original matching partitions
these orbits into five-tuples, a disrupted tuple joins at least two,
and the $X$-action is transitive on the original matching tuples.
On each component the generated group is primitive.  Indeed a
supported prime cycle cannot permute nontrivial blocks: if it moved
blocks of size at least two, it would move at least ten points.  Its
support must therefore be within a single block, and connectedness
of the support hypergraph leaves no nontrivial block system.
Jordan's prime-cycle Theorem \citep[Theorem~1.1]{Jones2014} now supplies the full alternating group
on each component, independently, because the generating cycles
are independently supported.  It follows that
$y\le\langle D,z\rangle$, again contradicting
\eqref{r:forbidden-test}.  Thus $D=X$ or $A_n$, proving the normal
endpoint assertion.

Pass to the quotient by $\operatorname{core}_{G_0}(H_0)$.
Every nontrivial minimal normal subgroup supplements the base.
An abelian one would identify the interval, by intersection, with
an invariant-subgroup lattice of an elementary abelian group, hence
a modular lattice.  But $[X,A_n]$ is nonmodular, as already witnessed
inside the transversal product labels by the intransitive test above.
If two distinct minimal normals occurred, they would commute.  Using
one as the supplement identifies the interval with its invariant
subgroups; supplementation by the other means that the base induces
all inner actions on the first.  All these intersections would then
be normal subproducts of its simple factors, again giving a modular
lattice.  The same contradictions hold for the reversed lattice.
The core quotient therefore has a unique nonabelian minimal normal
subgroup.
\end{proof}

\begin{lemma}[Product labels and the site normal test]\label{r:site-normal}
Suppose the realization in \Cref{r:reverse-exclusion} exists.
After restricting its parameter lattice to a sufficiently large
supported $A_n\le A_v$, coordinate-product labels form an ordinary
interval
\[
 [H,M],\qquad H=\Lambda Q,\quad M=\Lambda S,
\]
where $Q$ is the base projection at a fixed site and
$\Lambda\le\operatorname{Aut}(S)$ is the site-stabilizer action.
They are closed under meets and joins in the original interval.
Write $K_X$ for the ordinary label indexed by a prime-generated
$X\le A_n$.  For every bounded-size semiregular elementary $q$-group
$X$, including $1$, with $q>5$ prime, the interval $[H,K_X]$ has
property~\((*)\), has base equal to a meet of strictly intermediate
coatoms, and has arbitrarily large residual height when $n/|X|$
is sufficiently large.  These conclusions persist under the core
and centralizer quotients whose kernels lie in the base.
\end{lemma}

\begin{proof}
A proper subdirect label cannot contain a nontrivial invariant product
$P$: factor transitivity makes every base projection nontrivial, and
its independent supports normally generate each full simple site.
For $b=1$ the assertion is automatic.  A coatom therefore has proper
product closure and, by maximality, equals that closure.  A subgroup
of $A_v$ generated by prime-order elements is a join of atoms in
parameter order; its reversed label is a meet of coatoms, hence a
product.  In particular a supported $A_n$ is a product label.  It
can be used as the new base; we keep the name $P$.

Identify $S$ with its inner automorphisms and set $I=\Lambda\cap S$.
A proper coatom product projects to a nontrivial proper group $T$
normalized by $I$, maximal in the local proper invariant interval.
Since $TI=S$ would make $T\trianglelefteq S$, we have $TI<S$, whence
$TI=T$.  Intersecting these coatoms gives $I\le Q$, since the
corresponding prime-order groups generate $A_n$.  Consequently an
overgroup of $H$ in $M$ is determined by its intersection with $S$;
that intersection transports equivariantly to every site.  Conversely
each invariant site projection containing $Q$ gives such an overgroup.
Coordinate products are closed under intersection and generation,
proving the interval assertion.

Fix $X$ as in the statement and put $Z=K_X\cap S$.
The exact invariant interval from $P$ to its product label is an
ordinary interval with endpoints
\begin{equation}\label{r:semidirect-test}
 P\rtimes\Delta,\qquad (\textstyle\prod Z)\rtimes\Delta,
\end{equation}
in matched-site notation; replace $\Delta$ by its finite acting image
if necessary.  It has reversed parameter lattice $[X,A_n]$.
A $K_X$-normal subgroup $N\le Z$ transports well-definedly, by
site-stabilizer invariance, to a normal product at the upper endpoint
of \eqref{r:semidirect-test}.  By \Cref{r:alternating-normal}, unless
$N\le Q$, it supplements the lower endpoint.  Thus $QN=Z$, and
$HN=K_X$.

We also need actual product coatoms with meet $H$.  On three points
in distinct free $X$-orbits, the product of translated supported
$C_3$'s is elementary abelian and disjoint from $X$.  Coprimality
$q\ne3$ splits it into irreducible $X$-modules.  Adjoining one
nonzero irreducible summand to $X$ covers $X$, remains strictly
below $A_n$, and is prime-generated.  The joins of all these groups
contain every supported transversal three-cycle, hence generate
$A_n$.  On reversing, their product labels are strictly intermediate
coatoms whose meet is $H$.  Using disjoint triples of orbits gives
arbitrarily long chains of such prime-generated product labels.

Now let $R\trianglelefteq K_X$ satisfy $R\cap Z\le H$.
Then $[R,Z]\le H\cap Z$.  Therefore $R$ normalizes the intersection
with $Z$ of each indicated coatom.  If $R$ is not contained in one
of them, it and that coatom generate $K_X$, so the coatom's inner
intersection is normal in $K_X$.  That intersection properly contains
$H\cap Z$, but its product with $H$ is precisely the proper coatom.
This contradicts the preceding normal test for normals within $Z$.
Hence $R\le H$.

Every missing normal thus has a missing intersection with $Z$.
Two missing normals commuting modulo the base core, including a
single one abelian modulo the core, would transport their inner
intersections to missing normal products in
\eqref{r:semidirect-test}.  They still commute modulo the transported
product of the core's intersection with $Z$.  That product is normal
at the upper endpoint and lies in its base.  This contradicts the
unique nonabelian monolith from \Cref{r:alternating-normal}.
All assertions of property~\((*)\) follow.  Factoring a core or
centralizer kernel already in the base preserves the exact intervals,
the coatom meets, and these conclusions.
\end{proof}

\subsection{Elementary charts and supported normal products}

The site interval need not retain every label of the original invariant
interval.  We therefore build charts for which all comparable labels are
products, and identify the precise replacements for the cross-chart and
normal-product arguments used in the earlier pruning construction.

\begin{lemma}[Elementary charts with projective separation]
\label{r:projective-charts}
Fix a prime $q>5$, a bounded path depth $D$, a maximum increment $r$,
and finite chart dimensions and numbers.  There are arbitrarily large
integers $n$ for which one can assign semiregular elementary
$q$-subgroups of $A_n$ to all
nodes and private charts in the finite recursive construction,
compatibly at their attachments.  At a node $W$, each private chart is
$W\oplus\mathbf F_q^N$ and extends the prescribed action of $W$.
Two nonidentity root-coordinate translations from distinct private
charts attached at that same $W$ generate a common
$\operatorname{PSL}_2(q^u)$ factor.
Here $u$ may be chosen a sufficiently large odd prime at each
attachment.
\end{lemma}

\begin{proof}
At a stage over $W$ take a replicated regular action of
\begin{equation}\label{r:attachment-group}
 W\times\operatorname{PSL}_2(\mathbf F_{q^u}).
\end{equation}
Assign different projective points $e$ to the private charts, and
embed their new coordinates as $N$-dimensional $\mathbf F_q$-subspaces
in the corresponding unipotent root groups.  For distinct $e,f$,
unipotent determinant-one lifts with coordinates $x,y$ have product
trace
\[
 2+c_{ef}xy,\qquad c_{ef}\ne0.
\]
Choose the embeddings so that, for every pair of nonzero chart
vectors from different charts, this trace has square outside
$\mathbf F_q$.  Such a choice exists: independent random linear
maps from the fixed-dimensional chart spaces to $\mathbf F_{q^u}$
are injective with probability tending to one.  Nonzero input vectors
in different charts have independent uniform images.  For each
fixed nonzero $c_{ef}$ the probability that
$2+c_{ef}xy\in\mathbf F_q$ tends to zero, uniformly in the chosen
points.  There are only finitely many tests.  Since $u$ is odd, an
element of $\mathbf F_{q^u}$ whose square belongs to $\mathbf F_q$
already belongs to $\mathbf F_q$; thus the stated squared-trace
avoidance follows.

The Dickson subgroup Theorem for $\operatorname{PSL}_2$
\citep[Theorem~2.2]{Giudici2007} now implies
that the two translations generate the full projective group.
Their distinct fixed points exclude the point-stabilizer alternative;
their order $q$ excludes the dihedral groups, whose orders are
$q^u-1$ and $q^u+1$, and $q>5$ excludes the small groups
$A_4,S_4,A_5$.  Since $u$ is odd, the square-field
$\operatorname{PGL}_2$ alternative cannot occur.  Since both $q$ and
$u$ are prime, the only remaining proper subfield group is
$\operatorname{PSL}_2(q)$.  Squared traces of determinant-one lifts
in a conjugate of that group belong to $\mathbf F_q$, contrary to our choice.
The square is unchanged by changing the lift projectively.

Choose $n$ divisible by all orders needed for this finite recursive
construction and sufficiently large for every residual-height test.
At a new attachment, the given semiregular $W$-action and the
restriction of the replicated regular action
\eqref{r:attachment-group} are both disjoint unions of regular
$W$-sets of the same total size.  An equivariant bijection identifies
them, extending the existing action.  The actions lie in $A_n$:
$W$ has odd order, and the projective factor is perfect, so neither
has a nontrivial sign character.  No disjointness between later
subtrees is needed.  The recursive construction is finite, and the
orders and sizes can all be accommodated by the final choice of $n$.
\end{proof}

\begin{lemma}[Linear failure patterns]\label{r:linear-patterns}
Consider a pruning round at an elementary parameter $W$ in the charts
of \Cref{r:projective-charts}, with actual ancestor socles as in
\Cref{c:ancestor-socles}.  Assume $K_W$ avoids the final exclusion list
relative to every strict ancestor.  In a failed chart every trial
extension $U$ of the prescribed increment has a chosen proper
full-stabilizer witness at some tested ancestor $W'$, with parameter
\[
 W'\le X(U)\le U,\qquad X(U)\not\le W.
\]
In an anchored round, the available elementary extension $B$ has
$K_B$ previous-good and $\dim(B/W)$ greater than the trial increment.
Take $q$ above the absolute line bounds for the initial exclusions,
the available dimensions sufficiently large for the prescribed finite
patterns, and $n/|X|$ sufficiently large for \Cref{r:site-normal} at
every elementary parameter used.

Then the graph construction of \Cref{c:failed-pattern} supplies lines
and bounded independent patterns in every failed chart.  For two matched patterns
in distinct private charts attached at the same current $W$, there is
one ordinary subgroup $J$ such that every cross-pair intersection is
$J$, with $H\le J$, and $J$ lies in every witness label of both patterns.
Every comparable chart interval is an exact reversed linear-subspace
interval, and every tested witness and previous-good anchor has
property~\((*)\).  In the respective high-type classical and alternating
configurations, this replacement for a cross-chart intersection equal to
$H$ makes the initial exclusions apply using only a bounded line or
independent list.
\end{lemma}

\begin{proof}
All intermediate subgroups between comparable elementary chart
parameters are elementary and prime-generated.  By
\Cref{r:site-normal}, their labels are products, so these intervals
are exact even after passage to the ordinary site interval.  The
same is true between an ancestor and a comparable trial extension,
and after the indicated ancestor core or centralizer quotients.
Thus full stabilizers witnessing failure have labels $K_{X_i}$
with all prepared-link normal tests available.

Within an available extension of $W$, choose horizontal and vertical
complements from the new root coordinates.  Color graphs of linear
maps by the failure's ancestor, category, bounded case data, and
horizontal witness subspace.  The multidimensional Hales--Jewett Theorem supplies arbitrarily many affine parameter slots with all
this data fixed \citep{HalesJewett1963}.  The fixed horizontal
witness is not contained in $W$.  Choose a horizontal functional
nonzero on its image outside $W$, and vary each slot by multiples
of that functional.  Distinct parameters give distinct witness
subspaces $X_i$.  The span of two includes all witnesses on their
parameter line and adds exactly one dimension; an affinely
independent list adds a dimension at every step.  Upon reversal,
pair intersections lie below every label on that line, are covered
by the two labels, and bounded independent lists give strictly
decreasing successive intersections.  In particular all pairs of
distinct labels on one parameter line have the same intersection $E$,
and $E$ lies in every label on that line.

For two private charts attached at $W$, write $L_0,L_1\le W$ for
the projections of their witnesses to $W$.  These projections are
constant within their respective patterns, because only the vertical
root coordinates vary.  Put $G=\operatorname{PSL}_2(q^u)$ for the
attachment factor.  Any cross pair generates a subgroup of
$(L_0+L_1)\times G$ projecting onto $G$, by
\Cref{r:projective-charts}.  Its commutator subgroup lies in $G$ and
projects onto $[G,G]=G$, since $W$ is abelian.  Hence the generated
subgroup is exactly $(L_0+L_1)\times G$.  It is independent of the
chosen pair, contains every witness parameter in both patterns, and
is prime-generated.  Its ordinary reversed label $J$ is therefore
the common cross-pair intersection, and
\[
 H\le J\le K_{X_i}
\]
for every witness in either pattern.

A witness need only contain the subspace of its tested ancestor
and extend outside the current $W$; it need not contain $W$.
The fixed $W$-projections suffice for the preceding argument.
We now verify the uses of $J$ in the initial exclusions.  For subset
stabilizers, $J$ contains the supported alternating groups on the at
most four intersection cells, so it has boundedly many point orbits;
every subset in either pattern is $J$-invariant.  For missing
equal-block partitions, supplementation still gives $H$ its rich
actions on both block lists.  Thus $H$ has an absolute bound on
orbits of ordered pairs of intersection cells.  The supported
alternating groups within those cells lie in $J$; a cell-pair
stabilizer therefore has an absolute bound on point-pair orbits,
including coincident cells and cells of size at most two.  Hence $J$
has boundedly many point-pair orbits and boundedly many invariant
equivalences.  This gives the contradiction in
\Cref{c:alternating-primitivity}; the retained-block cases use the
original normal tests.  In the classical small-space cases, the
supported large group on the common complement lies in $J$ and
preserves every structure in both patterns.  The bounded-span and
linear-constraint arguments of
\Cref{c:linear-subspaces,c:isometry-subspaces} therefore apply to a
bounded independent list.

The remaining initial classical cases need only a parameter line.
For a descending list of distinct subspaces or signed flags, the
intersection of the first two stabilizers fails to preserve the
third, by the elementary maps in \Cref{c:linear-subspaces}.  This
contradicts $E\le K_{X_3}$.  For singular spaces, first discard the
at most two exceptional incomparable maximal spaces, then choose
three members of the remaining chain.  For mutually perpendicular
retained large spaces or orthogonal idempotents, choose four distinct
members.  After the first two structures are imposed, the third
cannot exhaust the unsplit remainder because the fourth nonzero
space remains.  The supported classical or special linear group on
that remainder gives the same strict third decrease.  Thus a fixed
number of labels on a line replaces every height-sized initial
comparison.
\end{proof}

We next recover the independently supported groups needed when a
normal simple power is missing.  We cannot assume that every
intermediate intersection with that power is a coordinate product:
the ordinary site interval may omit labels of the full invariant
interval.  The structural exclusions need only a previous-good
common intersection containing enough independently supported normal
subgroups.  We find it by bounding the codimension of product closures
and then selecting compatible affine parameters.  The selection uses
the following finite-vector-space observation.

\begin{lemma}[Simultaneous affine kernels]\label{r:affine-kernels}
Fix a finite field $\mathbf F_q$ and integers $h,m\ge0$.
For sufficiently large $d$, any family of linear maps
$\phi_i:\mathbf F_q^d\longrightarrow\mathbf F_q^h$, indexed by
$i\in\mathbf F_q^d$, has an affine $m$-space $A$ such that
\[
 \phi_i(j-i)=0\qquad(i,j\in A).
\]
\end{lemma}

\begin{proof}
Color every two-dimensional subspace, with its canonical reduced
row-echelon basis $(b_1,b_2)$, by the finite table
\[
 \bigl(\phi_{a b_1+b b_2}(c b_1+d b_2)\bigr)_{a,b,c,d\in\mathbf F_q}.
\]
The number of colors depends only on $q,h$.  The finite-vector-space Ramsey Theorem gives a homogeneous subspace of dimension $m+2$
\citep{GrahamLeebRothschild1972,GrahamLeebRothschild1973}.
Let $e_1,e_2,\ldots,e_{m+2}$ be its canonical basis, and put
\[
 A=e_1+\langle e_3,\ldots,e_{m+2}\rangle.
\]
Reduced row-echelon coordinates are compatible with this passage to
a subspace: multiplying a reduced row-echelon matrix in these
coordinates by the reduced row-echelon matrix of the containing
space again gives reduced row-echelon rows in ambient coordinates.
For $i\in A$ and $w\in\langle e_3,\ldots,e_{m+2}\rangle$, both
$(i,e_2)$ and $(i,e_2+w)$ are such canonical plane bases.
Their colors agree, so $\phi_i(e_2)=\phi_i(e_2+w)$.
Linearity gives $\phi_i(w)=0$, which proves the claim.
\end{proof}

\begin{lemma}[Supported products in a selected pattern]
\label{r:supported-products}
Let $X_i\le B$ be the elementary parameters of a failure pattern in
\Cref{r:linear-patterns}, with $K_B$ a previous-good common anchor.
Suppose the witness
$K_{X_i}$ has a missing normal simple power $Y_i$, possibly modulo
a retained normal subgroup, and $H$ is transitive on its simple sites
with an absolute bound on its ordered-pair orbit count.  After
selecting a line, or an affine subpattern of any prescribed positive
dimension, its common intersection label $E$ has property~\((*)\),
contains $K_B$ and hence remains previous-good, and, for each $i$,
contains an $E$-normal product
\[
 N_i=\prod_a R_{i,a}\ \le E\cap Y_i
\]
of independently supported proper subgroups, not contained in $H$.
Its image in $E/\operatorname{core}_E(H)$ contains the monolith.
The supported factors are nonsoluble, and the same statement with
preimages applies when retained centers have been factored out.
\end{lemma}

\begin{proof}
Write $Y=Y_i$ temporarily, so $HY=K_{X_i}$.
Every intermediate overgroup is $H$ times its intersection with $Y$.
Take the coordinate-product closure of this intersection and multiply
by $H$.  For $X_i\le Z\le B$ this is an overgroup
$K_{g_i(Z)}$, where
\[
 X_i\le g_i(Z)\le Z,
\]
and $g_i$ is monotone.  The resulting group's intersection with $Y$
is the product closure itself.  Therefore $g_i(Z)=X_i$ forces the
intersection at $K_Z$ to be subdirect in $Y$.  All intermediate
intersections from $K_Z$ to $K_{X_i}$ are then subdirect.  Their
tie relations are unions of $H$-orbits on ordered pairs of sites;
a strict inclusion forces strict refinement of ties, because
comparable subdirects with the same ties coincide.  The pair-orbit
bound gives an absolute $h$ with
\[
 g_i(Z)=X_i\quad\Longrightarrow\quad\dim(Z/X_i)\le h.
\]
Choose a complement $C$ of $g_i(B)/X_i$ in $B/X_i$.
Monotonicity gives $g_i(C)\le g_i(B)\cap C=X_i$ (using the preimage
of $C$ in $B$ in this notation), so
\begin{equation}\label{r:closure-codimension}
 \operatorname{codim}_B g_i(B)\le h.
\end{equation}
Exactness between comparable elementary labels ensures that every
closure label used here really is indexed by the indicated subspace.
The same calculation can be made modulo a retained normal subgroup;
taking its preimages does not alter the intervals.

In the graph pattern, $X_j$ differs from $X_i$ in the single vertical
direction determined linearly by the parameter difference $j-i$.
The preimage of $g_i(B)$ in these vertical directions has codimension
at most $h$, by \eqref{r:closure-codimension}.  Choose a map $\phi_i$
to $\mathbf F_q^h$ with that preimage as kernel.
Apply \Cref{r:affine-kernels}, asking beforehand for enough
Hales--Jewett slots.  In the resulting affine subpattern,
$X_j\le g_i(B)$ for every $i,j$.  Let $U$ be the span of its witnesses
and put $E=K_U$.  Then $U\le g_i(B)$ for all $i$, and consequently
\[
 H<K_{g_i(B)}\le E<K_{X_i}.
\]
Here $K_{g_i(B)}>H$ because $g_i(B)\le B<A_n$.
The label $E$ is above $K_B$, so it inherits all previous-good
exclusions, and \Cref{r:site-normal} gives property~\((*)\) for
$[H,E]$.

The intersection $K_{g_i(B)}\cap Y_i$ is a coordinate product not
contained in the base; otherwise supplementation would give
$K_{g_i(B)}=H$.  Each of its single-site components already lies
independently in $E\cap Y_i$.  Let $R_{i,a}$ be the full subgroup
of $E\cap Y_i$ supported on site $a$ and put
$N_i=\prod_aR_{i,a}$.  Conjugation by $E$ permutes these support
subgroups, so $N_i\trianglelefteq E$, and $N_i\not\le H$.
Site transitivity and $E<K_{X_i}$ imply that each factor is proper:
one full simple factor would force all of them and then $E=K_{X_i}$.
The factors are mutually conjugate under $H$, and if they were
soluble, $N_i$ would be a soluble missing normal, contrary to
property~\((*)\).  Thus they are nonsoluble.  The image of $N_i$
contains the common nonabelian monolith of
$E/\operatorname{core}_E(H)$, again by property~\((*)\).

We spell out the scope needed in the structure exclusions.
For direct systems calculated modulo block centers, each block
carries at most two simple sites, and the at-most-two block classes
have independent alternating top actions induced already by $H$
through supplementation; in linear type there is one list.
These give the required absolute pair-orbit bound.  In the case
of several site orbits one instead uses the supports of a retained
orbit as in \Cref{c:supported-anchor-groups}.  Lifting independent factors on a
block to its natural supported matrix group gives projective image
in $E$, since the lost centers were retained.  The lift is preserved
by the semilinear block stabilizer in $E^+$ and supplies the groups
of order greater than two required in the block-module argument.
If needed take products of the supported factors over one block.

For tensor and Cartesian systems the $Y_i$ already have independent
simple supports.  The common-monolith and iterated-commutator
arguments of \Cref{c:tensor-refinement,c:cartesian-refinement} apply to these
$E$-normal products $N_i$, without requiring
$N_i=E\cap Y_i$.  Their tuple intersections still generate a
nontrivial $E$-normalized subgroup.  It spans in the tensor case,
and is transitive in the Cartesian case, by the previous-good
exclusions inherited from $K_B$.  In the diagonal primitive case
choose a minimal $E$-normal subgroup inside the nontrivial product
of proper supported factors.  This is exactly the containment
needed by that argument.  Retained centers, when present, are
normal in $E$ and lie in $H$, so these conclusions survive their
removal and reinstatement.
\end{proof}

\begin{proof}[Proof of \Cref{r:reverse-exclusion}]
Prepare the charts in \Cref{r:projective-charts}, and use
\Cref{r:site-normal} for every elementary parameter encountered.
At every node the current almost simple quotient is the one from
\Cref{c:ancestor-socles}; initially it is $M$.  The normal, residual-height,
and exact comparable-link hypotheses have all been verified.
\Cref{r:linear-patterns} supplies the line, independent-list, and
cross-chart configurations for the pruning rounds relative to every
ancestor.  \Cref{r:supported-products} replaces the interval-product
consequence exactly at its uses for missing direct, tensor,
Cartesian, and diagonal systems.  Thus each finite round either
loses only the bounded number of charts allowed by the matched
failure arguments or produces the next actual almost simple socle,
with all ancestor-relative exclusions.  These are the structural
conclusions of \Cref{c:classical-conclusion,c:alternating-conclusion}.

If a parent has bounded type, its bounded-order case is excluded by
residual height.  In bounded Lie rank, the bounded-equation pattern
argument of \Cref{sec:bounded} applies unchanged with a bounded
independent list: strict subgroup descents are strict after
intersection with the socle, whereas the spaces of equations have
bounded dimension.  Thus one obtains a sandwich throughout the
interval over a selected increment $t$.  Its height is $t$ and the
same divisor count gives at most $C2^t$ labels.  The exact
$t$-dimensional vector-subspace interval has at least
$q^{\lfloor t^2/4\rfloor}$ labels, a contradiction for sufficiently
large fixed $t$.  The sufficiently-general input is used here in
precisely the scope established in \Cref{sec:bounded}
\citep{LarsenPink2011}.

For a high-type path, use linear flags in the final protection
selection of \Cref{w:selection}.  All bad containing labels belong
to the exact subspace intervals just proved.  The same
exponential-gap estimate for flag proportions applies; hence the
closure-field tensor and root protections hold simultaneously for
the required ancestor comparisons.  The absolute group-theoretic
bound \Cref{w:path-bound} on the length of such a protected path
then contradicts a
choice of depth $D$ larger than that bound.

There is no dependence of the earlier finite choices on $n$.
First fix the high thresholds and the protected path bound, and
choose $D$ beyond it.  Next choose a prime $q$ above every absolute
line bound, then $r$ for the bounded-type count and protection
estimates.  Choose trial dimensions backwards through the finite
pruning rounds, allowing at each round all Hales--Jewett and
\Cref{r:affine-kernels} slot requirements for the trial and bounded
path ranks.  The latter requirements depend only on $q$, desired
fixed dimensions, and absolute orbit bounds.  The initial
comparisons require only a line or a bounded independent list, as
proved in \Cref{r:linear-patterns}; none requires comparison with
$|A_n|$.  Choose $N$ and the number of private charts large enough
for the finite color and chart-loss bounds.  Only then choose the
fields of \eqref{r:attachment-group} and finally $n$.
All elementary groups $X$ used are now bounded in size, so $n/|X|$
can exceed all normal-test and residual-height requirements while
$n$ accommodates every replicated regular action.  This gives one
fixed forbidden degree $n$.  For every $v\ge n$, the supported
$A_n$ restriction of \Cref{r:site-normal} gives the same contradiction.
\end{proof}

\section{The forward marker theorem}\label{sec:forward}

This section concerns pure markers.  Its target is a product-interval
statement forcing a bounded prefix and many actual point stabilizers
in an ordinary site group.  The individual group-type arguments that
follow are exclusions under the assumption that this point configuration
is absent; their final contradiction again uses the protected path
bound.  The normal and mixed-projection tests below will also apply
internally when the simple sites change.

\begin{definition}[Pure marker lattice]\label{r:pure-markers}
For a finite set $T$ of $t$ markers, let $F_T$ have labels
\[
 (I,V),\qquad I\subseteq T,\quad V\le S_{T\setminus I}.
\]
Choose a set $Z$ disjoint from $T$, with $|Z|=t+1$ or larger, and
interpret this label as the subgroup
\[
 S_{Z\cup I}\times V\ \le S_{Z\cup T}.
\]
Order, meet, and join are subgroup containment, intersection, and
generation in this model.  The markers in $I$ are called absorbed.
We write
\[
 X_I=(I,1),\qquad x_i=X_{\{i\}},\qquad
 M_D=X_{T\setminus D},\qquad
 C_D=(T\setminus D,S_D)\quad(D\ne\varnothing).
\]
The bottom and top are $X_\varnothing$ and $X_T$.
We also write $F_t$ when only the cardinality matters.
\end{definition}

\begin{lemma}[Elementary marker properties]\label{r:marker-elementary}
The labels in \Cref{r:pure-markers} are closed under meet and join.
The coatoms are exactly the $C_D$, $\varnothing\ne D\subseteq T$.
The axis filter $[X_I,1]$ is $F_{T\setminus I}$, and the axis ideal
$[0,X_I]$ is $F_I$.  Moreover
\[
 \operatorname{ht}(F_t)=O\bigl(t\log\log(t+3)\bigr).
\]
\end{lemma}

\begin{proof}
The independent symmetric group on the absorbed part makes the
intersection again a group with independent full symmetric action
on its new absorbed part.  In a join, that part expands exactly
along the generated orbits meeting it: symmetric transpositions
there and their conjugates give the full symmetric group on the
expanded part, independently of the induced group on the remaining
markers.  Thus both operations stay within the indicated labels.
These rules depend only on the marker labels $(I,V)$, so the resulting
lattice is independent of the choice of $Z$ with $|Z|>t$.

If a proper label has unabsorbed set $D$, it lies below $C_D$.
The latter is the full stabilizer of the partition into $Z\cup(T\setminus D)$
and $D$.  These parts have unequal sizes, the first larger than
$t\ge|D|$; adjoining any element failing to preserve them generates
the full symmetric group by the same absorption argument.  Hence
$C_D$ is a coatom, proving the classification.  The descriptions
of the axis filter and ideal follow directly by treating $Z\cup I$
as the fixed absorbed set, or by restricting to $Z\cup I$,
respectively.

For the height use $|Z|=t+1$, embedding $F_t$ in the subgroup lattice
of $S_{2t+1}$.  Every strict step in a subgroup chain consumes at
least one prime factor, with multiplicity, of $(2t+1)!$.
Writing $n=2t+1$, the total number of such factors is
\[
 \sum_{p\le n}\sum_{a\ge1}\left\lfloor\frac{n}{p^a}\right\rfloor
 \le n\sum_{p\le n}\frac1{p-1}
 =O\bigl(n\log\log(n+2)\bigr),
\]
by the reciprocal-prime bound (the difference between $1/(p-1)$
and $1/p$ is summable).  This proves the estimate.
\end{proof}

\begin{lemma}[Marker normal test]\label{r:marker-normal}
For sufficiently large $t$, every exact finite group interval
isomorphic to $F_T$ or $F_T^{\mathrm{op}}$ has property~\((*)\).
\end{lemma}

\begin{proof}
Let $\omega=(I,V)$ label the base times a normal subgroup.
As in \eqref{r:left-modular}, it and all its projections to
subintervals satisfy the left-modular law in either orientation.
Indeed in actual subgroup order, with base $H_0$, normal
$N\trianglelefteq G_0$, and $H_0\le D\le E\le G_0$, the identity is
\[
 (DN)\cap E=D(N\cap E)=D\vee((H_0N)\cap E).
\]
Projection to $[D,E]$ gives $D(N\cap E)$, with
$N\cap E\trianglelefteq E$.  Self-duality of the left-modular
identity gives the reversed statement as well.  In particular the configuration
\begin{equation}\label{r:marker-forbidden}
 \omega\wedge b\le a<b\le\omega\vee a
\end{equation}
is impossible, also after any projection.

Suppose $I$ is nonempty and proper, with $|I|\ge2$.
Choose four positions $P_4$ in cyclic order such that membership
in $I$ is nonconstant, occurs in both parity classes, and is not
invariant under rotation by one.  Such a choice is possible for
large $t$: two chosen positions of $I$ can be adjacent, and at
least one of the other positions lies outside $I$; if the
complement has only one point, use three from $I$.
For the rotation $s$ put
\[
 a=(T\setminus P_4,\langle s^2\rangle),\qquad
 b=(T\setminus P_4,\langle s\rangle).
\]
The $I$-membership condition forces $\omega\wedge b\le a$.
Each parity orbit of $s^2$ meets the absorbed part of $\omega$,
so $\omega\vee a=1$.  Thus \eqref{r:marker-forbidden} holds.

If $I=\{j\}$ and $V$ moves some marker $k$, choose $P_4$ so that
$k$ has opposite parity from $j$ and its image under some element
of $V$ lies outside $P_4$.  That image is already absorbed in $a$;
the two parity classes are therefore absorbed in $\omega\vee a$.
Again the same rotation gives \eqref{r:marker-forbidden}.
If $V=1$, use three positions containing $j$, with $C_3<S_3$
in place of the two rotation groups and absorb the outside.
The join with the $C_3$ absorbs all three positions, while the
meet with the $S_3$ lies in $C_3$.  Thus nonempty proper $I$ is
impossible in all cases.

Now let $I=\varnothing$ and suppose $V$ moves $j$.
Project above $x_j$.  In the marker description of that axis
filter, $\omega\vee x_j$ has nonempty absorbed set because the
$V$-orbit of $j$ contains another point.  The preceding test makes
it the top.  Consequently $V$ is transitive on $T$.
Every marker atom $x_j$ complements $\omega$.

First take the forward orientation.  Complementation of a
base-normal multiple gives the interval isomorphisms
\[
 [0,x_j]\longrightarrow[\omega,1],\qquad u\longmapsto u\vee\omega.
\]
As $x_j$ is an atom, $\omega$ is a coatom; since its absorbed set
is empty, \Cref{r:marker-elementary} gives $\omega=C_T$.
Factor by the base core.  Choose a minimal missing normal subgroup
inside the putative normal subgroup; it still has a nonbottom
proper normal-multiple label, and the same argument identifies
that label with $C_T$.  Write this minimal normal as $N$ and the
base as $H$.  At a transposition label
\[
 0<d=(\varnothing,\langle(i\ k)\rangle)<C_T,
\]
the intersection $N\cap K_d$ is a nontrivial proper subgroup of
$N$, and $H(N\cap K_d)=K_d$.
For every $j$ the marker operations give
\[
 C_T\wedge(d\vee x_j)=d.
\]
Because $N$ is normal in the whole upper endpoint and lies in
$K_{C_T}$, this equality shows that $K_{x_j}$ normalizes
$N\cap K_d$.  The groups $K_{x_j}$ generate the upper endpoint.
Thus $N\cap K_d$ is a nontrivial proper normal subgroup of $N$
normal in that endpoint, contradicting minimality.

In the reverse orientation the normal projection instead yields
isomorphisms, in the displayed $F_T$ order,
\[
 [x_j,1]\longrightarrow[0,\omega],\qquad u\longmapsto u\wedge\omega.
\]
For $k\ne j$ the label $X_{\{j,k\}}$ is strictly above $x_j$;
therefore $X_{\{j,k\}}\wedge\omega>0$.
This meet is nontrivial only if $V$ contains the transposition
$(j\ k)$.  Hence $V=S_T$ and $\omega=C_T$.
But a supported three-cycle group containing $j$ cannot be in
the image of the indicated isomorphism.  Any label above $x_j$
whose meet with $C_T$ contains that cycle must absorb all three
positions; its meet then contains the full $S_3$, rather than
just the cycle group.  This contradicts surjectivity.

We have proved the normal endpoint assertion.  The lattice $F_T$
is nonmodular for large $t$: it contains the full subgroup lattice
of $S_T$, or one may use the rotation obstruction just exhibited.
In the faithful base-coset action after factoring the core, every
minimal normal subgroup is transitive because it supplements the
base.  An abelian minimal normal identifies the interval with an
invariant-subspace lattice and makes it modular.  Two distinct
nonabelian minimal normals commute; supplementation by one makes
the base induce all inner actions on the other, so intersections
there are normal subproducts of simple factors, again forming a
modular lattice.  These contradictions give the unique nonabelian
monolith required by property~\((*)\).
\end{proof}

\begin{lemma}[Exclusion of a mixed marker base]\label{r:marker-mixed}
For sufficiently large $t$, suppose that
\[
 F_T\cong[P,S^b]^\Delta
\]
in forward order, with $S$ nonabelian simple and $\Delta$
transitive on its factors.  If $c$ denotes coordinate-product
closure, then
\[
 P<c(P)<S^b
\]
is impossible.
\end{lemma}

\begin{proof}
Assume the mixed case and let $\omega=(I,V)$ label $c(P)$.
The projections of $P$ are nontrivial and proper at every site,
by factor transitivity.  A label joining with $\omega$ to the top
must be subdirect: adjoining $c(P)$ does not change its site
projections.  For two distinct comparable subdirect labels, their
intersections with $\omega$ are distinct.  Indeed their tied
coordinates and restricted twisting maps are realized using the
nontrivial full projections of $P$ to the factors of $c(P)$;
a strict refinement of ties gives a strict change of the
intersection.  Conversely, comparable subdirect groups with
unchanged ties and twists coincide.

Apply the rotations in the proof of \Cref{r:marker-normal}.
They give $\omega\wedge b\le a<b$ with both $a$ and $b$ joining
$\omega$ to the top.  Hence both are subdirect, whereas
\[
 a\wedge\omega=b\wedge\omega,
\]
contradicting the preceding strictness.  This eliminates nonempty
proper $I$, including the singleton cases of that proof.
The top absorbed set is impossible because $\omega$ is proper.

If $I=\varnothing$ and $V$ moves $j$, consider $c(x_j)$.
It contains $\omega\vee x_j$.  In the axis filter above $x_j$,
apply the nonempty-absorption test just used if the product
closure is still a proper mixed label.  If $x_j$ itself is already
a product, its closure equals it; this cannot happen because it
contains $\omega\vee x_j>x_j$.  Thus $c(x_j)$ must be the top.
So $x_j$ is subdirect.  Its join with $\omega$ is then the top:
$\omega$ has independent nontrivial supports, whose conjugates
under a subdirect group generate each simple site by simplicity.
In parameter terms this makes $V$ transitive.  All $x_j$ are
therefore subdirect and have $x_j\wedge\omega=0$.
Strictness for the comparable subdirect labels
$x_j<X_{\{j,k\}}$ gives
$X_{\{j,k\}}\wedge\omega>0$; as before, $V$ contains every
transposition, and so $V=S_T$.

Each meet $x_j\wedge\omega=0$ is the same subgroup $P$.
Write $c(P)=\prod_o T_o$ over the simple sites, where every $T_o$ is
nontrivial.  On a tied pair of sites, the two-coordinate projection
of $P$ is the graph of the identifying isomorphism restricted to
$T_o$; on an untied pair it is the whole product.  Thus $P$ determines
the ties and the restrictions of their twisting maps to the $T_o$,
although it need not determine the full maps on $S$.  Every invariant
overgroup of a subdirect label is obtained by an operator-invariant
refinement of its ties, retaining the twists on the remaining parts.
Its intersection with $c(P)$ depends only on that refinement and the
restricted twists.  Consequently meeting with $\omega$ has the same
image on every filter above $x_j$.

This is false for $\omega=C_T$.  Fix a triple of markers.  If $j$
lies outside it, joining its supported three-cycle label with $x_j$
and then meeting $C_T$ recovers exactly that three-cycle label.
If $j$ lies inside the triple, any label above $x_j$ whose meet with
$C_T$ contains the cycle must absorb all three positions; its meet
then contains the full $S_3$.  It cannot yield just the three-cycle
subgroup.  A triple and an outside marker give the contradiction.
\end{proof}

\begin{lemma}[Ordinary site interval for product markers]
\label{r:marker-site}
Suppose $F_T=[P,S^b]^\Delta$ exactly in forward order, with
$\Delta$ transitive on the simple factors and $P$ a coordinate
product satisfying $P\ne1$ or $b=1$.  For sufficiently large $t$,
all coatoms and all axis labels are products.  At one site, let
$Q$ be the projection of $P$ and let
$\Lambda\le\operatorname{Aut}(S)$ be the image of the site
stabilizer.  Then $\Lambda\cap S\le Q$, and the product labels
form the ordinary interval
\[
 [H,S\Lambda],\qquad H=\Lambda Q.
\]
This correspondence preserves meets and joins in the full invariant
interval.  We denote its subgroup at a product label $u$ by $K_u$.
\end{lemma}

\begin{proof}
A proper subgroup projecting fully at every simple site cannot
contain the independent nontrivial base projections: their
conjugates normally generate the full simple sites.  At $b=1$
this is automatic.  Hence each coatom has proper product closure
and equals that closure by maximality.
For a proper coatom let $T_1$ be its projection at the fixed site.
It is nontrivial and proper.  The group
$T_1(\Lambda\cap S)$ is still invariant and transports to every
site.  It is proper, since equality to $S$ would imply
$T_1\trianglelefteq S$.  Maximality consequently gives
$T_1(\Lambda\cap S)=T_1$.
The singleton-omission coatoms $C_{\{i\}}$ intersect to the base,
so their site projections give $\Lambda\cap S\le Q$.
The same intersections over $i\notin I$ give $X_I$; thus every
axis label is a product.

An invariant site projection containing $Q$ transports uniquely
to a $\Delta$-invariant coordinate product.  As
$\Lambda\cap S\le Q$, these projections are exactly the
intersections with $S$ of the overgroups of $H$ in $S\Lambda$.
An overgroup is recovered by multiplying its intersection with
$S$ by $H$.  Finally, coordinate products are closed under
intersection and generation in the original simple power,
proving the asserted compatibility.
\end{proof}

\begin{theorem}[Forward marker theorem]\label{r:forward-marker}
For every nonnegative integer $m$, there are integers $a(m)$ and
$t_0(m)$ such that the following holds whenever $|T|\ge t_0(m)$.
Every exact forward realization
\[
 F_T=[P,S^b]^\Delta
\]
with $S$ nonabelian simple, $\Delta$ transitive on the factors,
and invariant coordinate product $P\ne1$ (or $b=1$), has the
following configuration in the ordinary site interval of
\Cref{r:marker-site}.  There is $A\subset T$ with $|A|\le a(m)$
such that $K_{M_A}$ has an actual high-degree alternating simple
normal subgroup $B_A\le S$ satisfying
\[
 HB_A=K_{M_A},\qquad
 C_{K_{M_A}}(B_A)\le H,
 \qquad C_{K_{M_A}}(B_A)\text{ is soluble}.
\]
In the almost simple quotient
$K_{M_A}/C_{K_{M_A}}(B_A)$, at least $m$ distinct labels
$M_{A\cup\{i\}}$, with fresh indices $i\in T\setminus A$,
are full stabilizers of individual points in the natural
alternating action.  The constants $a(m)$ and $t_0(m)$ are independent
of $S,b,\Delta,P$ and the realization; the high-degree threshold is
the fixed one used in the structural exclusions.
\end{theorem}

The proof of \Cref{r:forward-marker} occupies the rest of this
section and \Cref{sec:forward-completion}.  Its initial site
reduction is \Cref{r:marker-site}.  The normal test and mixed-base
exclusion apply whenever new simple sites occur internally.
The bounded-type step uses exactly the sufficiently-general
bounded-complexity input established in \Cref{sec:bounded}; the
high-type steps retain the protected root and tensor conditions
of \Cref{sec:weights}.

\subsection{Pure links and simultaneous pruning}
\label{fa:links}
We prove \Cref{r:forward-marker} by successive exclusions, always stopping if the required point
stabilizers have already appeared.  The present subsection records the transfer and selection facts that
allow every exclusion to be made simultaneously below each bounded path prefix.  All labels refer to the
original marker set of \Cref{r:pure-markers}; ordinary labels are those represented by products at the
original simple sites.

\begin{lemma}[Exclusion of bounded types]\label{fa:bounded-types}
In the product configuration of \Cref{r:forward-marker}, the ambient simple group has unbounded order and
unbounded classical rank or alternating degree as the required marker reservoir grows.  More precisely,
any prescribed bounds on order and on Lie rank can be excluded by choosing the reservoir sufficiently
large.  This assertion also holds in every exact pure-marker filter obtained by the changes of simple
sites below.
\end{lemma}

\begin{proof}

Bounded order is immediately impossible by the length of an axis chain. In bounded Lie rank use the
argument of \Cref{w:lp,w:height-rigidity,w:bounded-exclusion}. For a prescribed sufficiently large constant \(k\) test \(M_U\),
\(|U|=k\). There must be some for which all ordinary labels above it admit the sandwich on intersection
with \(S\): otherwise pick many disjoint such sets and take proper bounded-equation stabilizer witnesses
as there. Each has omitted (nonabsorbed) set nonempty inside its tested set, hence intersections strictly
descend (also ordinary products, strict on intersection with \(S\)). The linear spaces of equations
cannot give such long descent. Now residual lengths under \(M_U\) are arbitrarily large and the same
height and order comparison for the sandwich maps bounds the number of ordinary labels above it by
\(C2^{h}\), \(h=O(k\log\log(k+3))\). But all partitions of \(U\) occur via intersections of cuts \(C_D\)
(individual symmetric group on every part), giving at least \(\exp((1/2+o(1))k\log k)\) for example by
pairings. This is impossible. In particular we can require arbitrarily high thresholds on the remaining
types here or at any use with a large enough marker set.

\end{proof}

\begin{convention}[Relative labels at an ancestor]\label{fa:relative-labels}
At a path prefix \(A\), suppose that we have the prepared-path data of
\Cref{w:prepared-path}: an actual simple subgroup
\[
 B_A\trianglelefteq K_{M_A},\qquad B_A\le S,\qquad HB_A=K_{M_A},
\]
whose centralizer in \(K_{M_A}\) is soluble and contained in \(H\).  Initially \(A=\varnothing\).  We
work in its almost simple quotient, and write \(T'=T\setminus A\) for the relative marker set.  Thus the
relative top is \(M_A\), even when a displayed relative label uses only \(T'\).
A proper full stabilizer containing a tested axis with a bounded number of further omissions has a product label
\begin{equation}\label{fa:witness-labels}
 r_D=(T'\setminus D,V_D),\qquad
 g_D=(\varnothing,V_D)=r_D\wedge X_D,
 \qquad \varnothing\ne D\subseteq T',\quad V_D\le S_D.
\end{equation}
Here \(D\) is a subset of those further omissions, and \(V_D\) fixes every position outside \(D\) in the
expression for \(g_D\).  A normal subgroup is \emph{retained} at this link if it is contained in
\(K_{g_D}\); otherwise it is \emph{missing}.
\end{convention}

We will use exact pure-marker links before passing to the ordinary site
interval.  These are different intervals: the ordinary one corresponds
exactly to the coordinate-product labels in the full formal link, and
need not contain every formal label.  Normal subgroups can nevertheless
be tested in the full link by transporting them across the original
simple sites.  We first prove this normal-subgroup transfer, then use
it to find independently supported groups inside intersections.

\begin{lemma}[Normal transfer on pure links]\label{fa:pure-links}
For the labels in \Cref{fa:witness-labels}, the full formal interval \([g_D,r_D]\) is \(F_{T'\setminus
D}\).  If this residual set is sufficiently large, its ordinary group interval \([K_{g_D},K_{r_D}]\) has
the normal property of \Cref{r:marker-normal}, including its nonabelian monolith modulo the base core.
For pairwise disjoint supports \(D_i\), put
\[
 J=\bigwedge_i r_{D_i},\qquad g_* =\bigvee_i g_{D_i},\qquad D_* =\bigcup_i D_i.
\]
The same conclusion holds for \([K_{g_*},K_J]\), and for distinct \(i,j\) one has \(g_{D_i}\le r_{D_j}\)
and \(r_{D_i}\vee r_{D_j}=M_A\).
\end{lemma}

\begin{proof}

The full formal link \([g_D,r_D]\) is \(F_{T'\setminus D}\):
every intermediate increase keeps \(V_D\) independent.  For disjoint
supports the label calculation gives \(g_{D_i}\le r_{D_j}\) and
\(r_{D_i}\vee r_{D_j}=M_A\).  The link \([g_*,J]\) is likewise
pure on \(T'\setminus D_*\), with independent side group
\(\prod_i V_{D_i}\).  We prove the normal assertion for either
pair of endpoints, denoted by \(g,r\).

Work before any ancestor-centralizer quotient, and let \(S\) denote
the original simple site.  Put
\[
 L=K_g,\qquad R=K_r,\qquad Q=L\cap S,\qquad Z=R\cap S.
\]
Let \(P_g,P_r\) be the corresponding products at the original sites,
with fixed-site projections \(Q,Z\).  They have these projections
because \(\Lambda\cap S\) lies in the original base.  With
\(\Delta\) replaced by its finite acting image if necessary, the
full formal link has the exact ordinary realization
\[
 [P_g\rtimes\Delta,\ P_r\rtimes\Delta].
\]
Indeed its intermediates all contain \(\Delta\), and hence correspond
to precisely the \(\Delta\)-invariant subgroups between the two
products.  The normal test applies to this full interval.

First let \(N\trianglelefteq R\) lie in \(Z\).  Since
\(\Lambda\le L\le R\), its transport to the other sites is
well-defined.  The resulting product is normal in
\(P_r\rtimes\Delta\).  By \Cref{r:marker-normal}, either
\(N\le Q\) or \(QN=Z\); in the second case \(LN=R\).
Thus every missing normal inside the original site supplements the
ordinary lower endpoint.

Now let \(N\trianglelefteq R\) satisfy \(N\cap Z\le L\).
Then \([N,Z]\le Q\), so \(N\) normalizes the intersection with
\(S\) of each singleton-omission axis coatom of the link.  These
coatoms are products, are strictly intermediate, and have meet \(L\).
If \(N\) escaped one of them, adjoining it to that coatom would
generate \(R\); the coatom's inner intersection would therefore be
normal in \(R\).  It strictly contains \(Q\), but its product
with \(L\) is only the proper coatom, contradicting the preceding
paragraph.  Hence \(N\le L\).  Together the two paragraphs prove
the normal endpoint assertion for every normal subgroup of \(R\).

Finally put \(C=\core_R(L)\).  Two missing normals commuting modulo
\(C\), or one missing normal abelian modulo \(C\), would transport
their intersections with \(Z\) to missing normal products in the
exact semidirect interval.  Their commutators lie in the transport of
\(C\cap Z\), a normal product contained in its base.  This contradicts
the unique nonabelian monolith supplied there by \Cref{r:marker-normal}.
The ordinary interval therefore has property~\((*)\).  Quotienting
an ancestor centralizer preserves it, since that kernel already lies
in \(L\).

\end{proof}

The lower endpoint for retention is thus \(K_{g_D}\) on the link
indexed by \(D\), and \(K_{g_*}\) on a simultaneous intersection link.
We keep these endpoints when applying the subsequent normal-product
arguments.

\begin{lemma}[The disjoint-support join restriction]\label{fa:disjoint-join}
For disjoint nonempty \(D_1,D_2\) and \(J=r_{D_1}\wedge r_{D_2}\), every formal label \(u\) satisfies
\begin{equation}\label{fa:join-restriction}
 J\le u<r_{D_1}\quad\Longrightarrow\quad u\vee r_{D_2}<M_A.
\end{equation}
\end{lemma}

\begin{proof}

Indeed, all positions outside \(D_1\cup D_2\) are already absorbed at \(u\). The part of \(D_2\) still
not absorbed there stays separate also from \(D_1\) since \(u\le r_{D_1}\), and is preserved by
\(V_{D_2}\) since \(J\le u\); if this part were empty one would already have \(r_{D_1}\).

\end{proof}

\begin{lemma}[Independent support in intersections]\label{fa:intersection-support}
In the setting of \Cref{fa:pure-links}, suppose that \(Y\trianglelefteq K_{r_D}\), with \(Y\le B_A\), and
that a retained normal subgroup \(Z\le Y\cap K_{g_D}\) makes \(Y/Z\) a missing power of nonabelian simple
groups.  Assume that \(K_{g_D}\) acts transitively on its simple sites with a bounded number of orbits on
ordered pairs.  At a sufficiently long residual link, singleton axis intersections with this power are
products and their simple sites have arbitrarily high type.  Moreover, in \(E=K_J\), the product of the
individual simple-site intersections of \((E\cap Y)/Z\) is \(E\)-normal and missing relative to
\(K_{g_*}\).  It has nonsoluble factors and its image contains the nonabelian monolith modulo the base
core of \(E\).  If \(J<r_D\), these singly supported factors are proper.
\end{lemma}

\begin{proof}

There are two sets of coordinates in this argument: the original \(b\)
simple sites, and the simple factors of \(Y/Z\) inside each of them.
We first obtain an exact interval at the new sites, then show that its
axis groups contain independent supports.

Lift \(Y,Z\) inside the actual \(B_A\le S\), before any ancestor
quotient, and transport them across the original sites.  Since
\(K_{g_D}Y=K_{r_D}\), we also have
\[
 (K_{g_D}\cap S)Y=K_{r_D}\cap S.
\]
Thus the transported product of copies of \(Y\) supplements the full
lower endpoint, while the transported product of copies of \(Z\)
is retained.  Every intermediate subgroup is its lower endpoint
times its intersection with that normal product.  After factoring
the retained product, \([g_D,r_D]\) is therefore an exact
invariant-subgroup interval in \((Y/Z)^b\).  The operators are the
original action together with conjugation by the label at \(g_D\);
they act transitively on the new simple sites.

Consider a singleton axis \(g_D\vee X_{\{i\}}\), with \(i\) outside
the used supports.  Its intersection with the new power remains a
product across the original \(b\) sites.  It has nontrivial projections
at the new sites, by strict increase from the lower endpoint and
transitivity.  These projections cannot be full.  Otherwise every
subgroup on a long axis chain above it would be subdirect.  Strict
axis increases give strict intersection increases by supplementation;
inside an original site these would require strict refinements of
\(K_{g_D}\)-invariant tie partitions.  The bounded number of orbits
on ordered pairs bounds the length of such a chain, as in
\Cref{c:orbit-bound}, whereas the residual marker set is arbitrarily
large.

The filter above this singleton axis is again an exact pure-marker
interval.  Its bottom has nontrivial proper projections to the new
sites, so \Cref{r:marker-mixed} makes it a coordinate product there.
\Cref{r:marker-site} then makes all subsequent axes products at those
sites as well.  Applying \Cref{fa:bounded-types} to this exact filter
forces its new site simple to have arbitrarily high type.  This uses
the bounded-type exclusion, not the conclusion of the forward theorem
being proved.

Now put \(E=K_J\).  The axis
\[
 g_D\vee X_{T'\setminus D_*}\le J
\]
supplies independent new-site supports in \(E\), modulo \(Z\).
Inside one original site take the product of all individual
simple-factor intersections of \((E\cap Y)/Z\).  This product is
\(E\)-normal.  It is missing relative to \(K_{g_*}/Z\): the axis
intersection just obtained is a product contained in it, and
supplementation recovers that axis label by multiplying its lift by
\(K_{g_D}\).  That axis is not below \(g_*\).

By \Cref{fa:pure-links} on \([K_{g_*},K_J]\), the supported normal
product has image containing the nonabelian monolith modulo the base
core.  It is nonsoluble, and transitivity makes every supported factor
nonsoluble.  If \(J<r_D\), each factor is proper: a full one would,
by the transitive \(K_{g_D}\)-action and supplementation, give
\(E=K_{r_D}\).  This proves precisely the supported-product conclusion
needed below; it does not identify that product with all of
\((E\cap Y)/Z\).  In direct-system applications the block centers
are retained, so lifting these supports with their centers gives
the required actual block-supported matrices.

\end{proof}

\begin{remark}\label{fa:quotient-labels}

When reasoning in the almost simple quotient at \(A\) we generally keep ordinary subgroup labels
(including \(H\)) for their images as well. Overgroups there containing the base correspond without
change of order to ordinary labels below \(K_{M_A}\), since the kernel was already in the base. The image
of \(B_A\) still identifies with the \emph{actual} \(B_A\) by injectivity. Thus even if normal-product
transport is used while reasoning in this quotient, the inner groups transported can be taken back inside
that actual simple group first. Ordinary labels throughout are products at the original sites; their
being a sublattice, rather than necessarily the full displayed label lattice, matters for cover/size
arguments. Within the ideal every relative cut \(C_D\) is product (global cut intersected with \(M_A\));
in particular a maximal ordinary proper label has to be a cut, by containment.

\end{remark}

\begin{lemma}[When a product cover is a full cover]\label{fa:product-cover}
Let \(\Delta\) act on a finite direct product by coordinate-permuting
automorphisms.  Let \(X=\prod_i X_i<Y=\prod_i Y_i\) be invariant
coordinate products, with no invariant coordinate product strictly
between them.  If an invariant subgroup \(U\) satisfies \(X<U<Y\),
then \(X_i\trianglelefteq Y_i\) at every site.  Consequently,
nonnormality at one site makes \([X,Y]^\Delta\) a cover in the full
invariant-subgroup interval as well.
\end{lemma}

\begin{proof}
The coordinate closure of \(U\) is \(Y\).  Conjugating the independent
copies of \(X_i\) inside \(U\), whose projections are all of \(Y_i\),
therefore gives
\[
 N=\prod_i\langle X_i^{Y_i}\rangle\le U.
\]
This is an invariant coordinate product with \(X\le N<Y\), so the
product-cover hypothesis forces \(N=X\).  Thus every \(X_i\) is
normal in \(Y_i\).
\end{proof}

We apply this lemma before taking ancestor-centralizer quotients.
An ordinary cover is a cover among the original-site product labels.
If a formal intermediate exists, the lemma forces normality of the
lower site projection in the upper one.  This normality persists on
passing to an ancestor quotient and intersecting with an actual
supported subgroup there.  The concrete nonnormality tests below can
therefore detect a formal intermediate even when that intermediate
is not itself a product label.

\begin{lemma}[The three cuts above a disjoint meet]\label{fa:cut-meet}
If \(D,D'\) are nonempty and disjoint, the only coatoms above \(C_D\wedge C_{D'}\) in the relative marker lattice are
\[
 C_D,\qquad C_{D'},\qquad C_{D\cup D'}.
\]
Each of the first two covers the meet, whereas the interval up to the third is the Young subgroup
interval \([S_D\times S_{D'},S_{D\cup D'}]\).  These three cuts and the first two cover relations persist
in the ordinary product sublattice.
\end{lemma}

\begin{proof}

In particular, in the relative label lattice, the meet of cuts on \emph{disjoint nonempty} \(D,D'\)
absorbs everything outside \(D\cup D'\), with side group \(S_D\times S_{D'}\). The only nonempty sets
available for coatom cuts above that meet are \(D,D',D\cup D'\) (unions of the unabsorbed orbits). The
first two coatoms cover the meet, since inside one such cut one would have to absorb from the other set,
hence absorb that whole set, for any strict increase. Up to the union cut instead we have the full Young
subgroup interval in \(S_{D\cup D'}\). These are formal statements before the product selection (and
within the ideal at \(M_A\)); in particular all three coatoms persist, and formal covers between products
are still covers. In the other direction when needing an intermediate such as the formal wreath label,
which need not itself be product, we will argue using product closure rather than just assuming
persistence.

\end{proof}

\begin{lemma}[Finite selection simultaneously at all ancestors]\label{fa:all-ancestors}
For the successive exclusions below, fix a bound on path-prefix size and a finite sequence of rounds, each testing boundedly many
omitted positions and finitely many structure or comparison predicates.  A sufficiently large marker
reservoir can be thinned so that, conditionally on absence of the point configuration, each round applies
simultaneously at every eligible prefix carrying the simple group of \Cref{fa:relative-labels}.  At any
such fixed prefix, a bad witness with omitted support of size \(k\) has replicas on every \(k\)-subset of
a later reservoir.  Any prescribed bounded unions of these replicas retain all properties established in
preceding rounds.  The required reservoir sizes depend only on these finite requirements, not on group
dimensions or field sizes.
\end{lemma}

\begin{proof}

Apply finite Ramsey homogenization to the bounded omission patterns. Concretely color bounded-size
ordered index subsets by all the tested label predicates (Boolean, also subdivided into the boundedly
many indicated cases) relative to their prefixes \(A\) whenever those have the required path-node normal
simple in the ordinary product interval. It is canonical if it exists, by its soluble centralizer.
Include type predicates and the bounded comparisons across subsets as used for selection. From a
sufficiently large homogeneous set, for fixed \(A\) with fresh indices after it, any bad label with small
nonempty omitted set after \(A\) thus gives replicas with \emph{every} \(k\)-set of later fresh positions
(\(k=|D|\)), including as many disjoint ones as needed. Choices of witness structures can also be further
finitely colored and thinned. The exclusions are used successively: earlier ones for larger bounded sets
as needed to make the intersections (which contain the axes omitting the support unions) common
previous-good labels, just as properties there passed upwards from anchors. Thus they apply
simultaneously relative to all indicated high-type ancestors. All bounds for combined supports in each
round need only be finite in terms of subsequent round requirements, and never involve the ambient group
dimensions or fields.

More explicitly one may use successively shorter \emph{initial segments} of a homogeneous reservoir as
rounds proceed. Assert each exclusion (conditional on absence of the point configuration) for all prefixes \(A\) of the
allowed path sizes with the stated high simple data, and omitted sets in that round's segment, \(A\)
before the further omissions in index order. If a bad full stabilizer contains a tested axis label, its
exact omitted support is a nonempty subset of the tested further omissions. Replicate this support at
fixed \(A\) using the spared tail of the preceding segment; earlier-good properties are then available
for each specified bounded union of the replicas by the result of preceding rounds on that segment. Thus
choose earlier omission bounds and segment lengths large enough \emph{backwards} (e.g. include combined
supports of the finitely many simultaneous intersections needed in a later round, and allow extra finite
Ramsey thinning there). Predicate colors on ordered subsets can include all splits/sublist choices at the
finitely bounded sizes, recording actual path-simple existence even before descending (\(B_A\) is unique
by its soluble centralizer), needed same-characteristic/type tests, and existence of a witness with
exactly those support positions in the relevant categories. Not the unbounded field sizes, group orders,
etc., only finitely many case/comparison predicates are needed. In a contradiction at a fixed prefix one
can further choose witnesses and color them by finitely many tests, without losing earlier-good
properties. All unions and links still refer to the original \(T\); in particular the residual marker
counts for the normal/missing-site tests remain arbitrarily large.

\end{proof}

\subsection{Alternating ancestors}
\label{fa:alternating}
Throughout this subsection \(B_A\) is alternating of sufficiently large degree.  The homogeneous
families, disjoint replicas, and inherited properties are supplied by \Cref{fa:all-ancestors}.  A
previous-good intersection means one to which the conclusions of all preceding pruning rounds apply.

\begin{proposition}[Subset stabilizers produce points]\label{fa:subset-points}
If a tested descendant is contained in a proper subset stabilizer, then, after the finite selections of
\Cref{fa:all-ancestors}, there are arbitrarily many distinct relative labels \(M_{\{i\}}\) that are full
stabilizers of individual natural points.  In absolute notation these are \(M_{A\cup\{i\}}\).
Consequently, in the absence of the conclusion of \Cref{r:forward-marker}, the required descendants and
their bounded intersections are transitive.
\end{proposition}

\begin{proof}

For subset stabilizers \(K_{r_D}\), among the two supported alternating groups precisely one side is
missing relative to \(K_{g_D}\) by the normal property of \Cref{fa:pure-links} (soluble quotient over
both). Use the retained side's point subset \(U_D\); \(K_{g_D}\) acts transitively there by
supplementation with the other side, even for small size. Thus retained point sets for disjoint supports
are disjoint, by mutual containment of each smaller label in the other's stabilizer and distinctness. In
a homogeneous \(k\)-set family we can take their sizes each below a small fixed fraction of the whole
(color, using arbitrarily many disjoint replicas). Their stabilizers are maximal: an increase would give
transitivity and overlapping conjugates of the large complementary supported alternating. So \(r_D=C_D\)
(relative labels; all relative coatoms are products by intersecting global cuts with \(M_A\)). The same
maximality applies to the subset union for two disjoint \(k\)-sets. The union label contains \(C_D\wedge
C_{D'}\), and therefore equals \(C_{D\cup D'}\), since the only coatoms above that meet are these three,
and the union stabilizer moves the individual subsets (correct parity outside). Consequently the union
itself is independent of the split by its small size and the stabilizer orbits.

\smallskip\noindent\emph{Decomposing the retained subsets.}

It follows that \(U_D\) splits into elementary disjoint subsets indexed by the individual \(i\in D\), the
same in every use. Indeed for any point the indicator \(p_D\) has \(p_D+p_{D'}\) fixed by the union of
disjoint index sets. Thus \(p_{F\cup i}-p_{F\cup j}\) on replacing one index is independent of the other
\(k-1\) indices \(F\), by comparing with a common disjoint auxiliary \((k-1)\)-set. For \(k\ge2\) we
obtain a sum of individual real weights plus a constant: the replacement differences obey the triangle
identity by using the same \(F\) for three positions, so are differences of weights, and replacements
connect all the \(k\)-sets. Weights differ only by zero or one up to sign; with our ample reservoir, two
or more each of a high and a low weight would produce a difference of two between indicator sums. So
either all indicators vanish (a constant one violates disjointness) or just one position is exceptional
and membership is exactly when it occurs (the reverse choice again violates disjointness). For \(k=1\)
use disjointness directly. At most one elementary subset can be empty by label distinctness. For any
nonempty one its stabilizer label is \(M_{\{i\}}\), since it contains \(C_D\wedge C_{D'}\) for pairs
meeting just in \(i\); two such meets with remaining positions disjoint already join to \(M_{\{i\}}\).
Indeed each meet fixes position \(i\) and absorbs all outside \(D\cup D'\); use four sets all meeting
pairwise just at \(i\). For \(k=1\) we already have the label \(C_{\{i\}}=M_{\{i\}}\) itself.

\smallskip\noindent\emph{Equal sizes and singletons.}

Unions of two or four of these nonempty elementary sets again have the corresponding cut labels by the
coatom-meet argument. For disjoint pairs \(P_1,P_2\), the full interval from the two cuts' meet \(J_0\)
to the union cut \(F_0\) is that from \(S_2\times S_2\) to \(S_4\), with a proper wreath intermediate. If
actual union sizes for the two differed, the ordinary interval would be a cover (project the full
stabilizers to the symmetric group on the combined set, using outside parity correction).
Apply \Cref{fa:product-cover} to the product endpoints in the original
power.  The formal wreath intermediate would force
\(K_{J_0}\cap S\trianglelefteq K_{F_0}\cap S\) before the ancestor
quotient.  Passing to that quotient and intersecting with the actual
\(B_A\) gives a contradiction: on the combined set the upper group
induces the full symmetric group by outside parity correction, whereas
the lower group induces exactly its two-part stabilizer.  For
unequal positive part sizes that stabilizer is maximal proper in the full symmetric group (an increase is
transitive, merging overlapping larger parts via conjugates of their supported symmetric group), not
normal. Thus pair-sums of sizes agree and all nonempty elementary sizes agree (compare pairs sharing one
index against a common disjoint pair). For two elementary sets the analogous formal interval is just a
cover, so their equal size must be one by block-swapping with parity correction. This gives the desired
point labels. Henceforth we can use transitivity after this exclusion if success was not already
obtained.

\end{proof}

\begin{lemma}[Exclusion of imprimitivity]\label{fa:imprimitive-exclusion}
Assume the transitivity conclusion of \Cref{fa:subset-points} for all bounded unions required in this
round.  A homogeneous family of proper full stabilizers of nontrivial equal-block partitions cannot occur
at the tested descendants.
\end{lemma}

\begin{proof}

For full stabilizers of nontrivial equal-block partitions (the imprimitivity test) let the sizes be \(a\)
per block, \(b'\) blocks. If \(a\ge3\) with \(A_a^{\,b'}\) retained, the supported primitive groups force
refinement into blocks of another system of the pattern (discrete restriction to a block is impossible by
moving a point while fixing its equivalent outside neighbor), hence equality mutually. For \(a=2\) use
instead retained even simultaneous flips of two blocks, precluding any other matching of this kind.
Otherwise \(a\ge5\) and the alternating base product is missing. Here the full stabilizer is maximal by
the supported \(A_a\)'s: an increase gives overlapping conjugate blocks, thus full alternating on merged
components, and its primitive top on blocks forces total merging. Indeed supported alternating groups of
these sizes on overlapping sets generate alternating on the union (conjugate three-cycles through a
common point; even with just one common point this brings in three-cycles using one point of the other
set and two of one's own). The full top symmetric action is available, with internal parity correction.
Take two disjoint failures, giving coatoms \(C_D,C_{D'}\) and a meet covered in each. The two block
systems have common cell size \(h\) on occupied incidences by prior transitivity. They cannot coarsen one
another by primitive top action, which is realized inside the intersection for both by supplementation
relative to the \(g\) labels. If \(h\ge2\), the cell-refinement stabilizer intersects a row label
properly (by its supported primitive block groups), strictly above the intersection because one can swap
two cells within a row while fixing points elsewhere, with internal parity correction, breaking the
columns. This contradicts covering. If \(h=1\), the intersection embeds subdirectly in \(S_{b'}\times
S_{b''}\) via the two top actions, both block counts at least five by the touching block sizes. The
kernel of projection to the first top action must still supply all even-total base sign patterns there
(supplementation with \(A_a^{\,b'}\), which acts trivially on that action's blocks and signs). But this
kernel injects as a normal subgroup of \(S_{b''}\); its abelianization has size at most two, versus its
required homomorphic sign image of size at least \(2^{b'-1}\).

\end{proof}

\begin{proposition}[Exclusion of affine and Cartesian structures]\label{fa:cartesian-exclusion}
After \Cref{fa:imprimitive-exclusion}, the tested descendants are primitive.  They cannot lie in full
normalizers of elementary abelian regular groups or in proper Cartesian-system stabilizers.
\end{proposition}

\begin{proof}

Using then common primitive previous-good labels \(E\), full normalizers of elementary abelian regular
groups are excluded by \Cref{c:alternating-affine}: the regular groups are retained and hence mutually
included in \(E\), where each is minimal normal; distinct such groups commute, whereas an abelian regular
group is self-centralizing.  For Cartesian stabilizers, \Cref{c:cartesian-retained} applies with mutual
inclusion in intersections for disjoint supports; small slot sizes \(2,3,4\) are likewise precluded by
the regular-group test. In the missing power case for two disjoint supports, the supported normal
products in \(E\) are supplied by \Cref{fa:intersection-support}, so the tuple commutator and
point-stabilizer factorization in \Cref{c:cartesian-refinement} give a common Cartesian refinement with
at most one nontrivial fine slot per pair of slot
positions. Take \(E\) to be the two-label intersection. Its primitivity gives transitivity on occupied
fine slots, say each of set size \(h\ge2\). The top action at each of the two original systems inside
\(E\) is transitive containing the alternating group by supplementation, hence primitive, so grouping
fine slots into rows and columns splits each slot position on both sides into at least two fine ones
(refinement would be improper by that primitivity). Both original slot set sizes are large by the
bounded-type consequence for the missing powers.

Let \(u\) be the intersection of the full row-system stabilizer with the fine-slot stabilizer. This is a
proper decrease by the supported natural alternating action on each row slot, and it contains \(E\). It
permits interchanging two fine slots in one row, leaving the others in place: if the parity of the
interchange needs correction then \(h\) must be odd (at least three fine slots), so use an internal
symbol transposition. Together with the full column-system stabilizer this generates the ambient group.
Here are details of the latter assertion. Local alternating gates on two overlapping sets of Cartesian
directions each of product size at least ten generate the alternating gates on their union. Split into
differences and shared part of alphabet sizes \(a,c\) and \(b\), all at least two (proper overlap). On
the union the generated group is two-transitive: ordered distinct pairs can be changed to any one fixed
pair distinct in all three positions using two-transitivity on each of the two touching position-sets
(first obtain a difference in the shared part on a position-set with some difference, then set the first
and shared values as desired, then the shared and last). For \(b\ge3\) use three-cycle gates on the two
position-sets whose lifted supports intersect in one tuple there (one shared value appearing once in each
cycle and disjoint other shared-part values). Their commutator supplies a three-cycle, sufficient by
primitivity. For \(b=2\), both \(a,c\ge5\). Flip the shared bit on two values of the first alphabet, and
use arbitrary \(A_c\) conditional on bit one. Commuting these and commuting again with conditional
\(A_c\) gives \(A_c\) conditional on bit one and membership in the chosen two values (by perfection).
Indeed the first commutation applies an element of \(A_c\) on one bit value and its inverse on the other
just for those two first-alphabet values; intersect the condition with bit one by commutators again. Two
such two-value choices intersecting once and commuting their conditional actions likewise give supported
\(A_c\) on a single fiber, again sufficient (a three-cycle there, then Jordan). All gates obtained act
independently of unused coordinates, meaning the same union permutation repeated on all outside tuples.
Now we have the column gates and those on the conjugate columns by the interchange. Indeed an altered
column and its original overlap properly, and merging them extends the latter by a slot from a second
column, next allowing merger with that entire column by the same observation. Align actual permutations
of the columns preserving fine coordinates up to a permutation (uniform fine-slot alphabets and common
fine-slot count per column); even column permutations so aligned are in the socle as well. Thus the
merged directions connect all columns by primitivity (for two columns the one interchange already
connects). The gate observation yields the whole alternating group, hence full join (the base already
supplements it). This contradicts \(u\vee r_{D_2}<M_A\).

\end{proof}

\begin{lemma}[Exclusion of primitive diagonal structures]\label{fa:diagonal-exclusion}
After the preceding alternating rounds, primitive diagonal structures are excluded as well.  Hence the
required descendants of an alternating ancestor are primitive almost simple groups.  The transitivity
conclusions continue to hold relative to all such ancestors.
\end{lemma}

\begin{proof}

By \Cref{c:primitive-diagonal}, every remaining primitive group that is not almost simple normalizes a
transitive power \(Y=B^s\), \(s\ge2\), with a full possibly twisted diagonal point stabilizer and natural
degree \(|B|^{s-1}\).  This includes two opposed regular simple factors.

First consider retention.  A retained power with transitive factor action is a minimal normal subgroup of
the common primitive intersection \(E\), with trivial centralizer because its diagonal point stabilizer
is self-normalizing.  Distinct retained powers would therefore coincide.  If the factor action is
intransitive, each factor-orbit product is transitive; in the diagonal action a transitive subproduct
uses at least \(s-1\) factors.  Thus the only possibility is \(s=2\), the two opposed regular simple
factors.  A retained regular factor is minimal normal in \(E\), and its centralizer supplies its opposite
regular group.  Distinct minimal normals commute and consequently determine the same opposed pair.  These
are exactly the retained alternatives in \Cref{c:diagonal-pruning}, and exclude distinct disjoint-support
witnesses.

It remains to consider a missing power with transitive factor action.  Take disjoint labels and their
proper previous-good intersection \(E\).  By \Cref{fa:intersection-support}, \(E\) has a missing normal
product \(R=\prod_{a=1}^s R_a\) of singly supported proper groups, with nonsoluble factors.  This
assertion does not require \(R=E\cap Y\).  Choose a minimal \(E\)-normal subgroup \(P'=C^u\) inside
\(R\).  It is nonabelian by the affine exclusion, and transitive by primitivity of \(E\).  Every factor
of \(P'\) lies in one \(R_a\): the groups \([P',R_a]\le P'\cap R_a\) are normal in \(P'\), distinct slots
commute, and these commutators generate \(P'\), since \(P'\le R\) is perfect.  The assignment is
\(E\)-equivariant.  Supplementation relative to \(g_D\) makes \(E\) transitive on the \(s\) slots, so
there is a common positive number \(h\) of factors over each slot and \(u=sh\).

The primitive-group argument of \Cref{c:primitive-diagonal}, applied to \(E\), says that the point
stabilizer \(Q=P'_x\) is a full diagonal.  If \(h>1\), the product of the projections of \(Q\) to the
\(s\) assigned batches is a proper intermediate subgroup between \(Q\) and \(P'\), normalized by \(E_x\),
contradicting primitivity.  If \(h=1\), equality of point degrees gives
\[
 |C|^{s-1}=|B|^{s-1}.
\]
But a factor \(C\) embeds in a proper subgroup \(R_a<B\), an impossibility.  Thus the diagonal case is
excluded, and the required descendants are primitive almost simple.  These conclusions apply
simultaneously at the eligible ancestors by \Cref{fa:all-ancestors}.

\end{proof}

\subsection{Classical subspaces and projector stabilizers}
\label{fa:classical-initial}
We use the natural covering groups, forms, collineations, and anti-actions of \Cref{sec:classical}.  All
dimensions required to be large can be made so by \Cref{fa:bounded-types}.

\begin{lemma}[Exclusion of parabolic stabilizers]\label{fa:parabolic-exclusion}
For a high classical ancestor, the tested descendants cannot preserve a proper linear subspace in the
absence of anti-actions, a proper signed flag in their presence, or a proper singular subspace in the
isometry cases.  The same assertion holds for the bounded intersections needed in subsequent rounds.
\end{lemma}

\begin{proof}

Soluble unipotent radicals are retained, thus each lies in all other labels on a disjoint family. For
linear subspaces (no anti-actions), signed flags (linear with anti-actions), or singular subspaces in the
isometry types, the shear arguments of \Cref{c:linear-subspaces,c:isometry-subspaces} give
a chain (using the smaller space of each signed flag, the other determined in reverse by an anti-action
of the base), up to the bounded singular-space terminal exception. Pick two in a many-member disjoint
family there (below terminal dimensions with ambiguities in orthogonal type, e.g. below Witt index minus
one in split type D). Their full stabilizers and intersection meet the ancestor socle in the
corresponding parabolics of compatible flags. By the standard parabolic overgroup theorem for the natural
\(BN\)-pair in the simple Lie-type group (overgroups of a parabolic are parabolic, with the
standard-parabolic subset order; in the classical interpretation use the indicated subspace or singular
flags, away from the split type \(D\) terminal fork here), the labels are maximal and their meet has only
the two maximal increases to coatoms: overgroups intersect the socle in parabolics obtained by dropping
flag terms, determining the overgroup since it contains the supplementing \(H\). In the signed case terms
must be kept in duality-stable pairs (possibly a single middle term). This contradicts the three coatoms
above the meet of two coatoms with nonempty disjoint supports.

\end{proof}

\begin{lemma}[Nonsingular radical lines in characteristic two]\label{fa:radical-lines}
In the orthogonal case in even characteristic, a full stabilizer of a nonsingular radical line is
maximal.  Such line stabilizers cannot occur in a homogeneous family of the tested descendants after
\Cref{fa:parabolic-exclusion}.
\end{lemma}

\begin{proof}

Indeed write the hyperplane of orthogonality to a line representative \(u\) as \(u^\perp\); the
simple-group stabilizer contains products of two nonsingular symmetries with centers therein. An increase
gives another such hyperplane \(w^\perp\); there is a common nonsingular partner for bridging the
products (high dimension). Odd and even words of symmetries using these two hyperplanes then give by
conjugacy all centers, so even products giving the simple group are contained. Indeed arbitrary even
words from the two center sets together are available by inserting a common partner for mixed products;
once all centers are obtained as symmetries by words (necessarily of odd length by Dickson parity), their
pair products thus belong and give a nontrivial normal generating set projectively by simplicity. In
detail, for nonsingular \(v\) outside both one can choose nonsingular \(z\in u^\perp\) whose symmetry
takes \(v\) into \(w^\perp\), by choosing both polar pairings \(b(z,v),b(z,w)\) nonzero there, then
ensuring
\(Q(z)=b(z,v)b(z,w)/b(v,w)\) for the quadratic form \(Q\). One can adjust this value arbitrarily with a
large nondegenerate subspace orthogonal to the indicated vectors. Neither functional vanishes identically
since the lines are distinct and \(v\notin u^\perp\). This proves the generation assertion, also
sufficient in the ambient quotient since the base supplements its socle. For two disjoint failures the
pair span \(P\) must be nondegenerate, else its singular line would be preserved and has been excluded on
the combined support. The pair intersection then agrees with the full intersection of either line
stabilizer with the plane stabilizer (by covering). Indeed the latter cuts it properly, e.g. products of
nonsingular symmetries in \(u^\perp\) move the plane (choose one center pairing nontrivially with \(w\)
and outside the plane, correcting by one fixing the plane). But this full intersection on the \(u\) side
includes the symmetry at \(u\) corrected outside \(P\), failing to preserve \(w\). This gives a
contradiction.

\end{proof}

\begin{proposition}[Maximality of an unbalanced ordered splitting]\label{fa:projector-maximal}
Let \(V=W\oplus U\), with \(0<\dim U<\dim W\), be a nondegenerate splitting in an isometry type, or an
image--kernel splitting in linear type with anti-actions.  Let \(M\) be a group of the corresponding
classical actions containing the natural special group, and let \(K\) be the full ordered splitting
stabilizer, equivalently the idempotent stabilizer for the ring action.  Assume that \(K\) supplements
the special group in \(M\), that its collineations act semilinearly irreducibly on \(U\), and that the
supported special group on \(W\) is available, as is that on \(U\) when \(\dim U\) is large.  In the
linear case assume also that \(K\) contains an anti-action fixing the projector, hence carrying its
image--kernel data into the dual annihilator data.  For sufficiently large \(\dim W\), the subgroup \(K\)
is maximal in \(M\).
\end{proposition}

\begin{proof}

Let an overgroup properly contain \(K\).  We show that the transported
copies of the supported special group on \(W\) generate the whole
ambient special group.  First we find large summands close to \(W\)
whose vectors span \(V\) in boundedly many summands.  We then collapse
those sums and recover the global elementary generators.

Throughout, supported groups use their natural matrices, equal to the
identity off the indicated summand; their projective images are the
groups in \(M\).  Transporting a splitting transports these groups
with it.  In linear type an anti-action transports the special linear
support by inverse transpose, and exchanges the corresponding primal
and dual data.

\smallskip\noindent\emph{Near splittings and bounded sums.}

Any increase contains transported copies of the supported special on \(W\) at the largest summand of any
moved projector. This copy cannot normalize the original projector: by irreducibility on its support and
dimensions it would first force its image space to agree with \(W\), then its fixed complement to be
\(U\). The high classical special groups here are generated by matrices of bounded displacement rank
(ordinary elementary/root generators, also obtained by Witt reduction). For instance take all inner
conjugates of a noncentral elementary generator (in an isometry type one can use a special linear shear
on two totally singular basis directions extended dually and identically off the associated two Witt
pairs); they suffice by projective simplicity and perfection (generation modulo the special group's
center suffices by perfection). Thus a near copy of the splitting with moved projector is attained, with
new largest space \(W'\) differing from \(W\) by bounded codimension of intersection. Use the anti-action
if necessary in linear type to get movement of the \emph{image} (also to achieve kernel movement). Use these
and their \(K\)-transports as near splittings (transports by anti-actions of these bounded-rank
displacements remain near). Every vector is a sum of boundedly many from these large spaces. Indeed
projections span \(U\) by semilinear irreducibility. If \(U\) is bounded-dimensional this suffices.
Otherwise transport a nonzero such projection by the available classical actions to get bounded-sum
generation there. In the linear or symplectic case use transitivity, and for a norm constraint use high
Witt dimension: every vector \(v\) is a sum of boundedly many vectors of a given nonzero vector's type
(norm, including isotropic nonzero). For instance prescribe \(y\) and \(v-y\) both of the given norm by
first imposing the corresponding polar pairing with \(v\ne0\) (a trace in hermitian type), then adjusting
the norm in a large orthogonal subspace, avoiding zero vectors. Here we use the usual transitivity on
nonzero vectors of given norm also under the high special by Witt extension with determinant, spinor or
Dickson corrections: a large nondegenerate subspace orthogonal to the vector to be fixed supplies
arbitrary such coset data (in unitary type norm-one determinants; in orthogonal type
reflections/symmetries at nonsingular vectors, both norm-square classes available if needed). Likewise
after prescribing the pairing to get two summands of the given norm, orthogonal adjustment can be made
perpendicular also to a first solution of the pairing, using a large nondegenerate subspace representing
every norm many times. Use \(W\) itself for the unprojected correction. In the linear setting dual
generation holds also: under transport by a fixed anti-action at the original splitting, large-space
images for near projectors correspond by a semilinear primal-to-dual identification to annihilators of
the transported kernels. The family of these initial near splittings can be taken closed under that
transport in both directions (bounded defect on both images and kernels). Thus these annihilators span
the dual, with the same bounded-sum conclusion there; in isometry type use selfadjointness and primal
spanning. In particular the common kernel of the available near projectors to the large spaces is zero.

\smallskip\noindent\emph{Collapsing bounded sums.}

One can collapse such bounded sums to one summand from some conjugate splitting, still within bounded
codimension of \(W\). Write \(W_i\) for summand spaces used and \(p_i\) for their projectors, with
bounded pairwise defect, so intersections also contain large nondegenerate spaces when relevant. If
\(s=w_i+w_j,\ w_i\in W_i,w_j\in W_j\), has \(p_i(s)\ne0\), choose \(d\in W_i\cap W_j\) so \(p_i(w_j+d)\)
is nonzero of matching type/norm as \(p_i(s)\). Use largeness of the intersection to solve the norm
condition, with nonzero choices, and the special-group orbit transitivity
with corrections established above.  The complementary components agree:
\[
 (1-p_i)(w_j+d)=(1-p_i)w_j=(1-p_i)s.
\]
The supported group on \(W_i\) can therefore send \(w_j+d\) to
\(s\), mapping \(W_j\) to another near splitting.  For sums of three or more, if all pair-collapse
projections vanish, perturb two summands oppositely by an intersection vector to make another pair
collapse. For example perturb the first and second by a common nonzero vector with opposite signs; the
first projection for the first-plus-third sum is now nonzero. Defects stay boundedly large at most
because only boundedly many steps are needed.  More precisely, write
\(\delta(E,F)=\dim E-\dim(E\cap F)\) for these equal-dimensional
spaces.  A transport \(g\) supported on \(W_i\) fixes its projector,
so
\[
 \delta(gW_j,W_i)=\delta(W_j,W_i),\qquad
 \delta(gW_j,W)\le 3d_0
\]
when both current defects from \(W\) are at most \(d_0\).
The same calculation holds on the dual supported spaces in linear
type, since \(g\) fixes that part of its splitting as well.
Thus we maintain bounded defects for both images and kernels.
Large nondegenerate parts of intersections in
these steps follow since intersection with a near space has bounded codimension in the nondegenerate
summand (and we may still take boundedly many orthogonality constraints). Finally if a nonzero pair sum
lies in both projection kernels, take a fresh near \(W_3\) with
nonzero projection of the sum.  For a nonzero summand \(w_i\),
choose in \(W_i\cap W_3\) a nonzero vector of its type and transport
it to \(w_i\) by the supported special group on \(W_i\).
This transport fixes the sum, since \(p_i(s)=0\).  Its image of
the \(W_3\)-splitting therefore covers \(w_i\) and still has
nonzero projection of \(s\).  Use that splitting to collapse the
pair.  The same construction holds dually in linear type, with
both defect bounds maintained as above.

\smallskip\noindent\emph{Global elementary generators.}

Thus all centers are available near \(W\). Supported special transvections now generate globally in
unitary and symplectic type. In orthogonal type use products of reflections/symmetries of same norm,
bridging by a common partner in the large near-space intersections; such products generate (in particular
supply global nontrivial normal generators). In linear type for \(v,f\) covered respectively in some
\(W_i\) and some dual supported \(W_j^*\), \(f(v)=0\), take a commutator of supported transvections with
tensors \(v\otimes h,\ z\otimes f\), with \(f(z)=h(v)=0,\ h(z)\ne0\), \(h\) supported in the first dual
space and \(z\in W_j\). These choices exist by the large intersections relative to the first projection
kernel (one can take \(z\) in \(W_j\cap W_i\cap\ker f\) independent of \(v\), then choose \(h\) on
\(W_i\); its scaling gives arbitrary transvection coefficient): in the linear argument keep bounded
intersection defects simultaneously on images in primal and dual. This holds for the initial near
splittings by bounded displacement (also of the contragredients); when a supported group transports
another near splitting in a collapse, it keeps its own splitting fixed, so both defects still grow only
boundedly. This gives the required full transvections. This proves maximality.\qedhere

\end{proof}

\begin{lemma}[Orthogonality of the retained projectors]\label{fa:projector-orthogonal}
After \Cref{fa:parabolic-exclusion,fa:radical-lines}, consider full stabilizers of proper nondegenerate
subspaces in isometry type, or of proper nontrivial individual projectors in linear type with
anti-actions.  Let \(e_D\) be the projector onto the side whose supported special group is retained
relative to \(g_D\).  After finite selection, witnesses on disjoint omitted supports satisfy
\(e_De_{D'}=e_{D'}e_D=0\).  For the small orthogonal exceptions the nondegenerate-line round is performed
first.
\end{lemma}

\begin{proof}

Now take isometry stabilizers of proper nondegenerate subspaces, or linear-with-anti individual proper
nontrivial projector stabilizers. Designate by \(U_D\) the side with supported special retained by the
normal property of \Cref{fa:pure-links} over \(g_D\); exactly one side is missing. Write \(e_D\) for
the projector to \(U_D\). For a pattern with large retained dimensions, disjoint supports then have
mutually orthogonal projectors as in \Cref{c:linear-subspaces,c:isometry-subspaces}, by supported
irreducibility and mutual preservation. For a pattern with bounded retained dimension (take it constant),
the full stabilizers used are maximal by \Cref{fa:projector-maximal}. We explain small
exceptions/ordering in isometry type: the full individual isometry group on \(U_D\) is available in
projection with corrections on the large complement. It acts irreducibly except possibly split orthogonal
planes over two and three. Over two there is a forbidden nonsingular radical line. Over three there are
two nonsingular axes individually invariant unless interchanged by a similarity giving semilinear
irreducibility. Thus first run the nondegenerate line exclusion, then these other dimensions with
invariant-line cases forbidden on the combined supports. There is no dependence of the line test itself
on excluding these planes.

For two disjoint bounded-dimensional occurrences the sum of the small spaces is again complemented, by
its orthogonal space or by the kernel intersection, via the prior singular/signed flag exclusions (also
nonsingular radical lines when relevant). In linear type an anti-element of \(H\) fixing both projectors
exchanges the sum with the kernel intersection, as do signed anti-actions of the pair group; a nonzero
intersection gives a signed flag. Also images determine kernels in this case. The sum projector
stabilizer cuts each individual stabilizer properly (large complementary supported group) and thus its
full intersection with either equals the pair group by covering. Consequently all linear or isometric
maps on one small summand, identity on the rest of the sum (split it using that summand's prescribed
complement) and corrected outside, preserve the other small space. The outside here is in both projector
kernels, giving a large space for arbitrary special-group coset corrections as above. Distinct small
spaces (complements determined and dimensions taken constant) cannot contain each other; this also
explains strict cutting by the sum stabilizer, since a projection of the sum to an individual large
complement is then nonzero proper, not preserved by the full supported large action. Taking differences
on the other small space under the indicated maps thus forces zero projection onto the chosen summand
whenever its individual linear/isometry group acts irreducibly and nontrivially: a nonzero projection
generates the irreducible space already by its orbit differences. Hence mutual orthogonality outside the
indicated small exceptions. For the split orthogonal plane over three with interchange, the isometries
here together with the intersection label interchanging the nonsingular axes by earlier line exclusion
suffice (the minus sign supported there on the sum splits the other small space along the projection, and
the joint action is irreducible on the summand by the interchange together with its isometries). For a
binary linear line with anti-action, use the bilinear pairing giving a fixed anti-element: each line and
kernel are mutual annihilators for this pairing, with self-value 1 at the line representative. On two
distinct such representatives the two cross-values agree, since left and right orthogonality to each line
coincide; value 1 off-diagonal would violate the complementation of the sum. This likewise gives
projector orthogonality.

\end{proof}

The orthogonal projectors admit an additive decomposition over the individual marker positions.  We
record the algebraic step separately because it works in every characteristic.

\begin{lemma}[The projector fold]\label{fa:projector-fold}
Let \(e_D\) be a family of idempotents in one associative algebra,
indexed by the \(k\)-subsets of a sufficiently large finite reservoir.
Suppose that idempotents on disjoint sets annihilate one another on both sides, and that
\(e_D+e_{D'}\), for disjoint \(D,D'\), depends only on \(D\cup D'\).  Then there are pairwise
annihilating idempotents \(h_i\) such that
\[
 e_D=\sum_{i\in D}h_i.
\]
If the \(e_D\) are selfadjoint, so are the \(h_i\).
\end{lemma}

\begin{proof}

For \(k\ge2\), \(\Delta_{ij}=e_{F\cup i}-e_{F\cup j}\) is independent of \(F\) as in the subset argument.
Disjoint-pair differences annihilate each other. Also \(e_D\) acts as identity on \(\Delta_{ij}\) on both
sides when \(i,j\in D\), and as zero when \(i,j\notin D\). The latter follows using a disjoint witness
for the difference, the former by adjoining a spare disjoint \(k\)-set to \(D\) (the union projector
supports a witness within the union). Then \(h_i=\Delta_{ij}\Delta_{il}\) is independent of distinct
auxiliary \(j,l\), including order, by disjoint annihilation.
For example, with a fourth index \(a\), the triangle identities give
\[
 \Delta_{ij}\Delta_{il}
 =\Delta_{ij}\Delta_{ia}
 =\Delta_{il}\Delta_{ia}
 =\Delta_{il}\Delta_{ij}.
\]
It meets each \(e_D\) by multiplication
either as identity-action (\(i\in D\), use a partner in \(D\)) or zero. To obtain
\(\Delta_{ij}=h_i-h_j\), choose \(D\) containing \(i,j\) and let \(D'\) replace them by two fresh
positions \(a,b\).  Then
\[
 \Delta_{ij}=\Delta_{ij}(e_D-e_{D'})
 =\Delta_{ij}(\Delta_{ia}+\Delta_{jb})=h_i-h_j.
\]
No division by the characteristic or by \(k\) occurs. Thus distinct \(h_i\)'s annihilate each other (separate by an \(e_D\)), and
\(h_i=(h_i-h_j)(h_i-h_l)=h_i^2\). They are orthogonal idempotents, selfadjoint when relevant, and the
remainder in \(e_D-\sum_{i\in D}h_i\) is constant, necessarily zero by disjointness and idempotence (each
such remainder is itself an idempotent supported by \(e_D\) on both sides). For \(k=1\) use the
projectors directly.

\end{proof}

\begin{proposition}[Exclusion of proper projectors]\label{fa:projector-exclusion}
Under the preceding classical rounds, the tested descendants cannot lie in any of the proper stabilizers
considered in \Cref{fa:projector-orthogonal}.  The remaining descendants and the required bounded
intersections therefore have irreducible anchors; in linear type their collineation parts are also
irreducible.
\end{proposition}

\begin{proof}

In the resulting homogeneous \(k\)-set patterns choose dimensions small relative to total by
orthogonality on disjoint supports. The stabilizers and disjoint-pair sum stabilizers are maximal by
\Cref{fa:projector-maximal}, and the sum label is \(C_{D\cup D'}\) as in the subset argument (it moves
each summand). For clarity the aggregate here has the requisite irreducible collineation action: split
planes over two do not arise as sums, and over three in the line test choose both norms of the same class
by coloring, avoiding the split plane. Thus the aggregate projector \(e_D+e_{D'}\) depends only on the
union; it is recoverable from its stabilizer, e.g. the supported large complementary special forces
uniqueness among projectors with small image.

Apply \Cref{fa:projector-fold} to write the projectors as sums of the \(h_i\).
In this marker configuration, at most one \(h_i\) vanishes,
and each nonzero \(h_i\) has full stabilizer label \(M_{\{i\}}\).
For \(k=1\), the full stabilizers are already the singleton
cuts \(C_{\{i\}}=M_{\{i\}}\). For \(k\ge2\), two vanishing
\(h_i\)'s would identify distinct cut projectors by replacing one index
with the other. For a nonzero \(h_i\), choose \(D,E\) meeting just
at \(i\); then \(h_i=e_De_E\). Thus the paired cut-meet argument
of \Cref{fa:subset-points} puts \(M_{\{i\}}\) in its full stabilizer.
That stabilizer is proper because \(h_i\) is a nonzero proper
projector, with nonzero disjoint witnesses remaining. Since
\(M_{\{i\}}\) is a coatom, it is the full stabilizer.

For nonzero positions, in large fixed sizes \(l,2l\) with ample further spares, all sum-projector
stabilizers on those sets \(P\) are again maximal (keep dimensions small relative to the ambient by
thinning), hence are \(C_{B(P)}\) for some nonempty \(B(P)\subseteq P\), and distinct by projector
uniqueness. Homogeneity/counting forces \(B(P)=P\) on a thinned reservoir (homogenize the decreasing-size
test; injectivity rules out strict decrease everywhere for that size). For disjoint \(l\)-sets, if the
sum dimensions differ, their meet to the combined cut is an ordinary cover, by
\Cref{fa:projector-maximal} applied within the combined aggregate (both pieces now sufficiently large).
Indeed projection to ring actions on that aggregate identifies the meet as full ordered splitting
stabilizer, with anti-action still present in the signed linear case. These applications have the
supplement hypothesis: the ancestor outer actions are supplied by the common base, and restrictions of
linear/special isometries from the inner ancestor are covered modulo the aggregate special group already
by the splitting stabilizer, using the two large pieces for arbitrary determinant etc. (orthogonal
determinant and spinor or Dickson corrections by Witt groups) and the large outside space for
compensation.  The formal wreath label is strictly between the two-cut
meet and the combined cut.  If these product endpoints formed an
ordinary cover, \Cref{fa:product-cover}, applied before the ancestor
quotient, would force normality at each original site.  This is
impossible on intersecting with the actual supported high special
group of the combined aggregate: its projective action is simple,
and the lower intersection is a proper nontrivial splitting stabilizer.
The ancestor centralizer kernel is already in the lower endpoint,
so this is also a valid nonnormality test in the quotient.  Hence dimensions agree for
disjoint \(l\)-sets, thus all individual dimensions agree (compare a pair of replacements to one common
disjoint set). Align the individual summand forms themselves by finite coloring. Now on an \(l\)-set the
sum-cut interval from the axis has at least the whole subgroup lattice of \(S_l\) just up to the
stabilizer allowing permutations of the individual projectors (kernel the axis label, and align
permutations with outside correction). The sum stabilizer exceeds this by its large full classical
action, which cannot preserve that uniform direct decomposition (same orders as in the earlier
direct-system pruning). Indeed the formal interval \([M_P,C_P]\) is exactly the subgroup lattice of
\(S_P\), with \(|P|=l\).  Its ordinary product sublattice cannot contain more labels.  Yet the
permutation stabilizer just described already supplies one ordinary label for every subgroup of \(S_l\),
and the strictly larger full sum stabilizer supplies an additional label.  This is the required
contradiction.  Thus these failures too are excluded. As
before we get prior irreducible anchors (collineation part also in linear type).

\end{proof}

\subsection{Direct systems}
\label{fa:direct-systems}

Use the direct systems and supported groups of
\Cref{c:direct-lines,c:direct-block-groups,c:direct-compatibility}, now with disjoint label
supports \(D_i\), \(K_i=K_{r_{D_i}}\) (as needed also for the following categories). Retention in the
smaller \(g_{D_i}\) labels places a group inside \emph{every} label \(K_j\) here. A previous-good anchor can
be taken to be the intersection or a large axis group below it, omitting just the joint support (all
relative to the ancestor). Thus we have transitivity and semilinear irreducibility on blocks by the
collineation part of the intersection whenever using them.

\begin{lemma}[Coincidence of exceptional matchings]\label{fa:direct-matchings}
In the retained line-system and exceptional plane-system cases of classical pruning, a family of
disjoint-support witnesses forces all the exceptional matchings to coincide.  A fixed previous-good
intersection then allows only boundedly many further systems.  Thus these cases are excluded by taking
sufficiently many replicas.
\end{lemma}

\begin{proof}

In the line-system argument of \Cref{c:direct-lines}, the diagonal groups are still mutually included by
retention, so it remains to treat its matching cases. All matchings relative to the
first system would have to agree: the signs of one of the other systems already interchange first-system
lines on a matched pair, simultaneously on a second pair if needed (choose two flips of equal norm class
for determinant/spinor correction). There are plenty of possible second pairs by transitivity and at most
two norm types (the two sign-orbit lines on a pair are of the same norm class); thus another invariant
matching must pair these same two first-system lines (otherwise choose the extra flip away from a
proposed different partner to keep that partner fixed). The choices on matched lines are unique as
before. For the exceptional plane matchings of \Cref{c:direct-block-groups}, an individual retained plane
sign likewise interchanges the
two old planes, forcing matching coincidence. The common intersection of all the labels is transitive on
the matching edges, so for that fixed intersection group bounded choices on one edge bound the number of
further systems, impossible. Here and below common intersections for any bounded family can be required
previous-good by earlier rounds at sufficiently increased omission bounds.

\end{proof}

\begin{lemma}[Support and compatibility for direct systems]\label{fa:direct-compatible}
In the remaining direct-system cases, two disjoint-support witnesses have the supported block groups
needed for the displacement-rank compatibility argument.  The block groups can be chosen as actual
block-supported matrices, with retained centers included.  Retained systems are excluded by comparable
refinements and the uniform order bound; a missing simple-site orbit has large block dimension and
bounded pair-orbit count.
\end{lemma}

\begin{proof}

In the remaining direct-system cases the supported groups of order greater than two for compatibility are
available exactly as above: use the full retained groups, or, for a missing simple-site orbit product,
\Cref{fa:intersection-support} with the previous-good intersection. (If two site orbits, can
likewise use full supports from a retained one.) Centers included with these lifts are retained. In fact
any missing simple orbit here forces large block dimension by the bounded-type consequence, with the
action on sites of bounded pair rank by supplementation relative to \(g_{D_i}\); in particular bounded
rank including the possible split four-space ambiguity cannot then occur. Block pair-orbits are bounded
by the alternating actions on at most two isometry classes (one class in linear type), and the simple
sites have fibers of size at most two over blocks. In applying the missing-product argument here the
singly supported projections modulo centers in the intersection actually yield block-supported matrices
projectively there: use representatives on that block in the supported choice of groups, since the center
kernel needed when lifting already lies in the intersection. We can take all available such block
matrices there, so the intersection's semilinear collineation stabilizer of a block normalizes the
corresponding supported group. The block module for the supported group again is semisimple without
trivial constituents by semilinear irreducibility for that stabilizer (use the module socle and the fixed
group of vectors, as before) and nontriviality of the block action. (For a missing product the singly
supported subgroups are even nonsoluble.) The displacement-rank proof of \Cref{c:direct-compatibility}
now applies: fixation of all the other system's blocks forces splitting by projections using orbit
differences there; otherwise movement both ways forces equal block dimensions, and forces any moving
block matrix to interchange a single pair of those other blocks as an involution, identity elsewhere
there and with displacement of full rank on its own support, as in \Cref{c:direct-compatibility}.  Its
retained-case proof also applies: full retained actions force comparable refinements, whose height is
absolutely bounded within the pattern by the uniform order obstruction on coarse blocks.

\end{proof}

\begin{lemma}[Merging compatible large summands]\label{fa:direct-merging}
Let \(W_1,W_2\) be compatible summands obtained by grouping blocks of one common direct splitting,
orthogonal in the isometry cases, with the corresponding complements.  If \(W_1\cap W_2\) is sufficiently
large and nondegenerate when relevant, their two supported special groups generate the supported special
group on \(W_1+W_2\).
\end{lemma}

\begin{proof}

Here the summands are direct groups of blocks of a common splitting, orthogonal in isometry cases, with
the corresponding complements. Use \(W_1\) of dimension at least \(\dim W_2\). Every vector in the sum
lies in some conjugate of \(W_1\) by the supported group of \(W_2\): for a nonzero \(W_2\)-projection
attain it by transporting a vector of the intersection (matching norm using its high rank), and the
remaining component already lies in \(W_1\). Thus we get the full symplectic or unitary transvection
generation on the sum; likewise products of orthogonal symmetries/reflections of matching norm, using
\(W_2\) and partners in its large intersections with these conjugates to bridge any two centers. These
elements give noncentral normal generation of the special. For linear type we also have dual coverage; as
in the projector argument get a transvection tensor \(v\otimes f,\ f(v)=0\), by commutators from two
supported spaces of dimension \(\dim W_1\), sufficient because twice that dimension exceeds the sum
dimension by a large amount. All assertions can use the natural covering matrices (the elementary
generators obtained give the perfect special group), with the indicated support.

\end{proof}

\begin{proposition}[Exclusion of the remaining direct systems]\label{fa:direct-exclusion}
After the preceding classical rounds, a homogeneous family of proper full direct-system stabilizers
cannot occur at the tested descendants.
\end{proposition}

\begin{proof}

Treat missing systems with two blocks first. More generally once these are excluded, two-class block
systems are impossible by coarsening, and the top action on blocks, induced inside intersections by
supplementation, is primitive (two blocks, or at least the alternating permutations with transitivity).
Thus two disjoint-support instances give incomparable compatible systems and each block splits into at
least two cells, of uniform dimension by transitivity. In the general round all those cells must in turn
be mutually isometric when appropriate (else a two-class coarsening). In the two-block round there is a
full \(2\times2\) grid of spaces, and for some choice of the pair of instances we can interchange two
same-row or same-column cells isometrically. Indeed if this fails the form must be orthogonal in odd
characteristic with odd cell dimension by similarity/transitivity. Take three failures, choosing their
common intersection to be previous-good as allowed by \Cref{fa:all-ancestors}.  Its collineation
irreducibility makes the occupied triple cells one orbit.  Since its permutation action lies in
\(S_2^3\), this orbit has size \(4\) or \(8\), and it projects onto all four cells for every pair. Size 8
violates odd pair-cell dimension by summing two equal triple-cell
dimensions. For size \(4\), the third two-block system is the unordered diagonal partition of the four cells of the first pair
(there is one third position at each occupied pair, with both third positions represented in each row and
column), so the whole intersection of that pair of stabilizers preserves it, impossible by disjoint label
supports.

Use a resolved pair, designated row and column, and let \(u\) be the intersection of the full row-system
stabilizer with the stabilizer of the simultaneous refinement. It contains the joint intersection and is
strictly smaller than the row label (uniform finer decomposition on each high-dimensional block). It
permits a swap of two fine cells in a row, leaving the others in place, or a row three-cycle of cells.
Indeed we have isometric choices as just described, and determinant, spinor or Dickson corrections inside
the cells when of large dimension. For small cell dimension here the row block dimension is large
(absolute thresholds can be made large enough to ensure at least three cells) and all cells are
equivalent, so take a three-cycle via commutators of aligned permutations. Such collineations come from
the ancestor socle.

We now apply \Cref{fa:direct-merging} to the original and altered column spaces.

A column space and its altered version lose/gain one cell by the within-row permutation and have large
intersection (at least half a block); their supported groups thus merge, extending the column \(X\) by a
cell from another column \(Y\). If that cell is large we merge with \(Y\) immediately. Otherwise use the
column label to transport within \(Y\) (alternating fine-cell permutations there while preserving \(X\))
and merge extensions along \(X\) until acquiring \(X+Y\). The resulting merged column pairs propagate to
a connected graph by block primitivity, and give full generation by repeated mergers with at least a
column in common. This contradicts the formal disjoint-support rule \(u\vee r_{D_{\rm col}}<M_A\), and
finishes direct systems.

\end{proof}

\subsection{Large tensor factors and representation bounds}
\label{fa:overgroups}

We supply maximality details for the forward field and ordered-factor tests, to avoid assumptions about
arbitrary structure stabilizers being maximal. Use ambient notation \(S\le M\le{\rm Aut}(S)\) in the
almost simple quotient at an ancestor, natural space \(V\) of dimension \(v\), \(\mathcal A={\rm
End}_F(V)\). We only use full stabilizers in groups supplementing \(S\) already by the common base. Write
\(Y\) below for the indicated supported projective special group(s), using their perfect matrix lifts
(center need not disappear in the projective field and paired constructions).

\begin{lemma}[Large-rank projective degree bounds]\label{fa:degree-gap}
For a sufficiently high-rank classical simple group of natural dimension \(m\), with the type and
characteristic conventions of \Cref{w:classical-conventions}, a nontrivial irreducible
projective representation in defining characteristic has degree at least \(m\).  Unless it is natural up
to Frobenius and duality, its degree is at least
\[
 \frac{m^2}{2}-O(m).
\]
A highest weight with \(l\) nonzero digit places has degree at least \(m^l\).  For field parameter \(P\),
every nontrivial cross-characteristic projective representation has degree at least \(P^{cm}\), for an
absolute \(c>0\); the same exponential lower bound holds for permutation degree.  These lower bounds on
nontrivial contributions apply to perfect central lifts and quotients.  Finally, a high-degree
alternating simple group acting projectively nontrivially in dimension \(v\) has natural alternating
degree \(O(v)\).
\end{lemma}

\begin{proof}

In defining characteristic, for classical simple groups of sufficiently high rank, with natural dimension
\(m\), the minimum nontrivial irreducible projective degree is at least \(m\); except for the natural
module up to Frobenius/duality, it is at least \(m^2/2-O(m)\). To see this use the restricted weights on
the universal cover and the Steinberg theory recalled in \Cref{w:classical-conventions}, using
\citet{Steinberg1963,Jantzen2003}. For a restricted digit in A the permutation orbit already
gives the latter bound except at a single endpoint only; at that endpoint with coefficient greater than
one, simple-root lowering (coefficient less than \(p\)) and the resulting orbit give it as well. In B,C,D
use signed coordinate permutations (even signs in D); at least two nonzero coordinates give the larger
bound in high rank (including spin weights), and a single first-coordinate support with coefficient
greater than one works again by lowering. Use B only in odd characteristic. Here the single first-node
coefficient one is the irreducible natural module, of degree \(m\) also in odd-dimensional B. For
\(k\ge2\) nonzero coordinates out of \(e\) in a signed word the orbit has size at least \(2^k\binom e
k\), allowing a loss of a factor two just at \(k=e\) with even-sign restriction; the stated quadratic
bound follows (including at \(k=2\)). These arguments also give the bound by \(m^l\) for \(l\) nonzero
digit places. Cross-characteristic nontrivial projective degree is at least \(P^{c m}\) for field
parameter \(P\) and absolute \(c>0\), by the argument of \Cref{w:characteristic-switch}; in particular the same exponential bound
works for permutation degree. These bounds on nontrivial contributions work also upon restriction to
perfect central lifts/quotients, pulling to the universal covers: lift projective actions there and take
a nontrivial composition factor (there is one by perfection if the projective action is nontrivial). The
center acts scalarly on irreducibles. In particular minimum permutation degree can be tested this way by
a permutation representation in another characteristic (the high perfect central lifts themselves then
have no nontrivial small permutation image). We use the classification and covering-group structure in
\Cref{w:classical-conventions}. An alternating simple of high degree acting projectively nontrivially in
degree \(v\)
has alternating degree \(O(v)\), by the following elementary argument.
Write its degree as \(n\), choose \(\ell\in\{3,5\}\) different from the
representation characteristic, and take \(r=\lfloor(n-2)/\ell\rfloor\)
disjoint \(\ell\)-cycles.  They generate \(E\cong C_\ell^r\), leaving
two spare letters.  The normalizer in \(A_n\) induces every permutation
of these generators: a permutation of the \(\ell\)-letter blocks has
its parity corrected, when necessary, on the two spare letters.
For matrix lifts of the projective action of \(E\), scalar commutators
define an alternating bicharacter \(b:E\times E\to\mu_\ell\).
Conjugation preserves \(b\).  Interchanging any two generators therefore
makes their value equal its inverse; because \(\ell\) is odd, that
value is one.  The lifts commute.  Rescaling them makes their
\(\ell\)-th powers one, and, since \(\ell\) is prime to the
characteristic, they are simultaneously diagonalizable.
The projective action of \(A_n\) is faithful by simplicity, so its
restriction embeds \(E\) in the \(\ell\)-torsion of the diagonal torus
modulo scalars.  Consequently \(r\le v-1\), and \(n\le5v+1\).

\end{proof}

\begin{lemma}[Coefficient twists on finite covering groups]\label{fa:coefficient-twists}
Let a classical finite covering group be defined by a standard pinned graph--field map of height \(f_0\).
Coefficient twisting a defining-characteristic irreducible representation by \(p^j\) is equivalent, after
restriction, to twisting its rational highest-weight lift by the source split \(p^j\)-Frobenius.  Thus
the nonzero digit positions are shifted by \(j\) modulo \(f_0\), with the diagram action at each wrap.
The same comparison may be made on composition factors.  Equivalence of fixed projective actions on a
perfect covering group implies linear equivalence of their lifts.
\end{lemma}

\begin{proof}

With a standard pinned graph/field map of height \(f_0\), coefficient \(p^j\)-twisting a
defining-characteristic irreducible is equivalent on the finite cover to twisting its rational
highest-weight lift by the split \(p^j\)-Frobenius of the source before restriction. Indeed in split
pinned coordinates, transforming coefficients in the rational representation together with the
coefficients in the argument (i.e. conjugating by coefficient Frobenius on points of target and source)
still gives an equivalent rational irreducible by highest weight. The split map permutes the standard
pinned finite fixed group; as in \Cref{w:classical-conventions}, digit positions thus move by \(j\) modulo \(f_0\), with
diagram action upon wrapping. This reasoning up to equivalence suffices (also for testing composition
factors); on a perfect covering group an equivalence of fixed projective actions lifts to linear
equivalence because scalar discrepancies are characters.

\end{proof}

\begin{proposition}[Maximality of an unequal ordered tensor factor]\label{fa:ordered-factor-maximal}
Let \(S\) be a high-rank classical simple group on \(V\), and let an ordered proper matrix tensor factor
give natural dimensions \(x<y\), both sufficiently large.  In isometry type assume that the factor is
selfadjoint.  In characteristic two assume that the bilateral factor forms are alternating, and in
orthogonal type use the canonical tensor quadratic form.  Exclude symplectic type in characteristic two.
Then the full stabilizer in \(S\) of the ordered factor is maximal.
\end{proposition}

\begin{proof}

Let \(Y_1,Y_2\) be the supported perfect projective special groups on the
two spaces, with the alternating special groups (symplectic groups) as needed on factors. They lie in
\(S\). For an overgroup within \(S\) take a minimal nontrivial normal subgroup \(N\) there. Its lift
algebra is semisimple by absolute irreducibility of the supporting action \(Y_1Y_2\), and has no multiple
center components: the larger simple action fixes all projectors by its minimum permutation degree, then
there are at most \(x\) of them (invariant spaces are modules over the larger matrix factor), so the
smaller also fixes them. The center is then scalar by perfection and full spanning by \(Y_1,Y_2\).

Any proper nonscalar matrix factor fixed individually by \(Y_1,Y_2\) gives precisely the original
factorization. Indeed lift the projective actions on factor spaces and tensor by perfection; all are
irreducible on the universal covers jointly by tensor irreducibility. The lower bounds \(x,y\) for any
respective nontrivial contributions (and \(xy\) if both in one slot) force their allocation to distinct
unique slots of a nontrivial split. Their matrix spans recover the factors. Here the separate lifts
commute as needed (any residual scalar commutators against a perfect factor give a trivial character);
joint irreducibles of the direct product over the closure are outer tensors. The supported matrices
originally span the actual factors over \(F\) by natural absolute irreducibility, and a trivial
projective action on other tensor directions means these matrices belong to the indicated one by
commuting with the others. Even for a permuted tensor system, \(Y_2\) on the larger space fixes all
slots, and acts projectively on just one (two would cost at least \(y^2\)); \(Y_1\) preserves that one
and must fix the others since their product dimension is at most \(x\). Thus a proper nonscalar lift
algebra for \(N\) already forces ordered stabilization by the overgroup. If the algebra is full and \(N\)
is a power of nonabelian simples, their lift algebras tensor to it by commuting as usual, so multiple
components also force ordered stabilization (no swap of unequal degrees). The elementary abelian
full-algebra case has the nondegenerate scalar commutator pairing, order \(v^2\) and projective linear
centralizer just itself. Its normalizer inside the simple group is far too small for \(Y_1Y_2\) (at most
\(v^{O(\log v)}\) by automorphisms of the elementary abelian group, while \(v\le y^2\)).

Thus otherwise \(N=B\) is a single simple with trivial projective linear centralizer and contains
\(Y_1,Y_2\) by Schreier. It must be high classical in the same characteristic: bounded Lie ranks are
excluded by the minimum faithful projective/linear degree of the larger simple support in any
characteristic (use also bounded-dimension adjoint realizations there); alternating type and
cross-characteristic classical type are excluded by the exponential bounds on embedding the larger
support, together with the action of \(B\) in degree \(xy\). In the alternating alternative the
alternating degree would be at most \(O(xy)\) by that action but at least exponential in \(y\) by
restriction to \(Y_2\). In the other-characteristic high classical alternative the natural dimension
there would already be at least exponential in \(y\) and its own cross-characteristic degree far exceeds
\(xy\). Its natural dimension \(w\) is at least \(x+y\): restrict its projective natural action to the
commuting supports, lift from universal covers of each (commuting linearly by perfection), and count
composition factors with nontrivial contribution from each (a joint irreducible with both nontrivial is
an external tensor over the closure, costing at least \(xy\)). Thus by the stronger degree bound, \(V\)
is projectively the natural representation of \(B\) up to twist, \(v=w\): we have absolute irreducibility
by the full matrix span, and \(w^2/2-O(w)>xy\).

In this equality case \(B\) matches the ambient root series. In ambient A, the supporting representation
is not self-dual even over the algebraic closure for the finite supporting special groups; in signed
types the form symmetry and dimension determine the series (and type A natural is not self-dual). In
characteristic two a symplectic natural \(B\) cannot preserve the orthogonal quadratic (see also the
field-transfer transitivity test in \Cref{sec:forward-completion}; over its natural finite field an
invariant quadratic, even allowing coefficients
in the closure, would have constant evaluation on nonzero vectors). Such finite-group preservation on
perfect linear lifts would be exact. Now apply \Cref{w:natural-root}: either
\(B=S\) or it comes from a proper root of the ambient Steinberg map, of field height strictly between
zero and the ambient field height \(f\), fixing the supports pointwise projectively. The latter is
impossible by restricted-weight uniqueness: on a supporting group the digit is natural at field exponent
zero (parameter \(p^f\)), not its twist through such a height even up to duality. Frobenius tests on
projective representations here and below agree linearly on universal covers by perfection. This proves
the claim.

\end{proof}

The field-transfer and exchanged-pair overgroup arguments use
\Cref{fa:degree-gap,fa:coefficient-twists,fa:ordered-factor-maximal}.  They complete the remaining
classical exclusions in \Cref{sec:forward-completion}.

\section{Completion of the forward marker theorem}\label{sec:forward-completion}

The preceding rounds have made the tested classical descendants
irreducible and excluded their direct decompositions.  We now remove
the remaining field and tensor structures.  The order matters: the
first tensor argument excludes factorizations with two large factors
but leaves a few possible retained plane factors.  The field and
exchanged-pair arguments then remove those planes, because each carries
one of the structures tested there.  Once these exclusions are complete,
every nontrivial inner subgroup normalized by a previous-good label
spans the full endomorphism algebra.  This permits the final
tensor-system exclusion and the construction of actual normal simple
descendants.

Work relative to a fixed path ancestor, with almost simple quotient
\(S\le M\le\operatorname{Aut}(S)\), and retain the ordinary subgroup
labels \(K_r\).  Its natural space is \(V\), its ground field is \(F\),
and \(\mathcal A=\operatorname{End}_F(V)\), with \(v=\dim_F V\).
The base already supplements \(S\).  The witness labels \(r_D,g_D\),
the cut labels \(C_D\), and the axis labels \(M_D\) have the relative
meaning fixed in \Cref{fa:relative-labels}.  Every group tested below
contains a previous-good axis label; its omission bound was enlarged
before the present round as in \Cref{fa:all-ancestors}.

The pure-link normal test of \Cref{fa:pure-links} applies with lower
endpoint \(g_D\), or \(g_*\) at a simultaneous intersection.  Missing
normal powers supply the independently supported subgroups of
\Cref{fa:intersection-support}; we do not require them to exhaust the
intersection with that power.  For disjoint witness supports, a proper
increase above their meet must obey the join restriction in
\Cref{fa:disjoint-join}.  These are the interval properties used in the
remaining exclusions.

All lower dimension bounds in this section are fixed absolute
thresholds, chosen before the marker reservoir.  The representation
bounds and coefficient-twist convention are those of
\Cref{fa:degree-gap,fa:coefficient-twists}.  After the structural
exclusions, we select descendants having the additional root and tensor
protections relative to every earlier member of a classical run.  The
absolute path bound will then force the desired point configuration.

\subsection{Overgroups of field and paired stabilizers}

We first establish the overgroup information needed by the field and
paired tests.  Under the following hypotheses, a prime-degree field
stabilizer or an exchanged-pair stabilizer is maximal.  For a field of
two-prime degree, every maximal overgroup must preserve one of its two
prime-degree subfields.  Thus, when two disjoint witnesses have such a
field stabilizer as their intersection, they cannot leave room for the
third cut supplied by \Cref{fa:cut-meet}.

\begin{proposition}[Field and paired overgroups in the prepared configuration]
\label{fb:field-pair-overgroups}
Use the transferred forms of \Cref{sec:classical}.
Let \(D\subseteq\mathcal A\) be a field containing \(F\), with \([D:F]=n\), where \(n\) is prime or a
product of two distinct primes, and let \(m=v/n\) be sufficiently large.
In an isometry type assume that \(D\) is selfadjoint; in orthogonal
characteristic two assume the irreducible normalization already obtained in the
preceding rounds.  Omit symplectic characteristic two with nontrivial adjoint
on \(D\).

Alternatively, in an isometry type other than symplectic characteristic two,
let \(e,e^*=1-e\) be exchanged complementary idempotents, with totally
singular images when a quadratic form is present, and let \(m=v/2\) be
sufficiently large.  Suppose that the collineation action of each overgroup
under consideration is semilinearly irreducible, as ensured here by the
previous-good axis.

Every proper overgroup in \(M\) of the full field stabilizer preserves a
nontrivial subfield algebra of \(D\) over \(F\).  Every proper overgroup of
the full paired stabilizer preserves that pair.  Consequently a prime-degree
field stabilizer and an exchanged-pair stabilizer are maximal; for a
compositum of two prime degrees its maximal overgroups can only be the
stabilizers of its prime-degree subfields.
\end{proposition}
\begin{proof}

In the field case the stabilizer includes the perfect transfer group \(Y\) over \(D\), of large natural degree \(m\). In the paired case use the paired perfect \({\rm SL}_m(F)\). Its centralizer algebra is \(Fe+F(1-e)\), by natural versus adjoint-dual representations; that for the field case is \(D\).

The full stabilizer contains a supplement to \(S\). To prove the assertion take a minimal subgroup \(N\ne1\) normal in the overgroup inside \(S\). The lift algebra is semisimple, and no center components can be permuted nontrivially by \(Y\), by its minimum permutation degree (at least exponential in an absolute positive constant times \(m f_Y\), where \(p^{f_Y}\) is its field parameter with \(m f_Y\ge v/2\); the number of center components is at most \(v\)). Thus its center lies in the displayed centralizer alternatives (perfection). Any enlargement from scalars already suffices: a larger center normalized by the overgroup then gives the indicated field (inside \(D\)) or exchanged pair. Here the high-degree centralizer facts use the transfer/paired representations as in the earlier pruning. Nor can there be a proper nonscalar central simple tensor factor, or multiple simple components of \(N\) tensoring to the full algebra: \(Y\) must individually stabilize such matrix factors in either setting (by order), acting projectively nontrivially on every factor space, as the centralizer algebras are commutative.

Indeed any such nontrivial projective representation of \(Y\) over \(F\) has degree greater than \(\sqrt v\). For the paired group this follows from \(m\); in the field case lift on its universal covering type, with some nontrivial composition factor. The factors are stable as a multiset under twist of field height \(h=\log_p |F|\). Indeed coefficient twisting the lifted representation through this height preserves the projective matrices (taken over \(F\)), hence gives the same representation by perfection on the universal cover, and in particular the same composition factors. If that constituent has orbit length \(d\), its restricted highest weight has at least \(f_Y/(d h)\) nonzero digits, since the twist circulates them modulo field height \(f_Y\) of \(Y\), where \(f_Y/h\ge n/2\). Diagram twisting does not affect the support positions. This costs dimension at least \(d m^{\lceil n/(2d)\rceil}\ge nm/2\), counting the orbit constituents with their Steinberg dimensions (use \(m^{j-1}\ge j\) for \(m\ge2,\ j\ge1\)). Thus two tensor directions are already impossible.

The full-algebra elementary abelian normalizer inside \(S\) is again too small by orders to contain \(Y\). Thus otherwise we have a full-spanning simple \(B=N\), with \(Y\le B\) as before (centralizer inside \(S\) trivial). The order/degree exclusions again force \(B\) high classical of the same characteristic, since \(Y\)'s field-height times natural degree is at least a constant times \(h v\). Let its natural dimension be \(w\ge m\), field height \(f_B\), and put \(t=f_B/\gcd(f_B,h)\). The absolutely irreducible projective representation in \(V\) costs \(v\ge w^t\) by \Cref{fa:coefficient-twists}. Orders compared with \(Y\) give
\[
             h v m\le C f_B w^2
\]
for an absolute constant \(C\). With \(m\) large this forces \(t=1\) (otherwise \(m\le C t w^{2-t}\)), then \(w^2\ge vm/C\). Again the representation must be natural up to twist by the nonnatural lower bound of \Cref{fa:degree-gap}.

We detail series and root checks here. In ambient linear type the field support uses the natural linear type over \(D\) with its embedding-conjugates and is not self-dual. In unitary type we necessarily have odd \(n\) (involution extending the ground involution); the \(n\) embedding twists of the natural unitary type occupy distinct digit positions modulo \(n f\), so again cannot match to form a self-dual module. For a pair in unitary type the natural and conjugate-dual of \({\rm SL}_m(F)\) at distinct digits likewise do not pair to a self-dual module. Signed types match by the previous tests; in the included symplectic characteristic two field case (trivial adjoint on \(D\)), \(Y\) is symplectic over \(D\) and cannot preserve a nonsingular quadratic even with extended coefficients. For example its lift is transitive on nonzero ground vectors, impossible for such a quadratic by evaluation there (also over the binary field, using polarization or a 3-space). The square-form preservation by a perfect group would be exact. Now all supporting constituent nonzero digits lie at multiples of the ambient \(f\) with \(f_Y\) a multiple of \(f\) (also in the indicated paired representations). No proper smaller positive height could fix the action projectively, even up to duality. To see the digit-position convention explicitly, the field-transfer action over the closure uses natural copies over the embeddings of \(D\) extending the ground embedding, hence twists spaced by \(h\), with \(f_Y=hn\) in the linear or trivial-adjoint case and \(hn/2\) in the involutory case. In the latter, either \(h=2f\) in unitary type, or \(h=f\) in bilateral type and the nontrivial involution fixes \(F\), so \(n\) is even there. For a pair one uses \(f_Y=h\) with natural and adjoint-dual (thus conjugate-dual in unitary type). Wrapping may diagram-twist but does not affect the positions modulo \(f\). A fixing ambient root would intertwine the representation with a coefficient twist by its height (possibly with duality) as natural projective matrices, hence on perfect lifts as above, rotating the positions to nonmultiples of \(f\). This rules out a proper root as before and proves the assertion. In particular prime field labels and exchanged-pair labels in these uses are maximal. For a two-prime compositum the only possible maximal overgroups correspond to prime subfields.\qedhere

\end{proof}

\begin{lemma}[Canonical quadratics in symplectic characteristic two]
\label{fb:canonical-quadratic-exclusion}
In the prepared symplectic configuration in characteristic two, replicated
failures whose full stabilizers preserve canonical quadratics with the ambient
polar form are impossible.  This excludes ordered tensor stabilizers with
alternating factor forms, nontrivial-adjoint field stabilizers, and exchanged
pair stabilizers as soon as the common previous-good label is irreducible.
\end{lemma}
\begin{proof}

For two disjoint failures whose stabilizers respect such quadratics (transforming as the bilinear form), irreducibility of a common previous-good label forces the same quadratic by square-difference, yet the two labels join to the top, not preserving it. This applies to tensor stabilizers (factor forms then alternating), to fields with nontrivial adjoint by hermitian trace, and exchanged pairs by their canonical hyperbolic quadratic. One can exclude all such additional quadratic stabilizers as soon as previous irreducibility on the joint support is available, also for use when a later tensor label is known to force one of them (alternation of ordered factors by previous irreducible normalization; for a slot system transitivity will give alternation).

\end{proof}

\subsection{Ordered tensor factors}
\begin{proposition}[Exclusion of two large ordered factors]
\label{fb:both-high-factors}
At the present forward-test stage, a homogeneous family of failures preserving
proper ordered matrix tensor factorizations cannot have both factor degrees
above the fixed large-degree threshold.  In isometry types the factors are
selfadjoint and have the normalized factor forms.
\end{proposition}
\begin{proof}

The factor-form arguments of \Cref{sec:classical} apply. Take a homogeneous pattern on all \(k\)-sets with designated algebra \(A_D\) on the retained side (over \(g_D\); the supported special groups are commuting normals with soluble quotient). For disjoint sets these large supported groups normalize each other; nontrivial intersection would force equality of simple groups, hence algebras and labels. Thus they commute projectively and linearly on perfect matrix lifts, so the designated algebras commute and multiply as independent factors. By coloring one can keep all their degrees small relative to total, e.g. at most \(v^{1/10}\), by using arbitrarily many disjoint sets. Now each such factor stabilizer, and the stabilizer of the product on a disjoint pair, is maximal by \Cref{fa:ordered-factor-maximal} and supplementation of the socle by the base. Thus we have labels \(C_D\) and product label \(C_{D\cup E}\) for disjoint \(D,E\), by the three-coatom argument of \Cref{sec:forward} (the aggregate supported group moves its constituent factors). The aggregate small-side algebra \(A_D A_E\) depends only on the union: among such small sides it is unique for its full stabilizer, by the supporting-group identification of individually preserved tensor factors in the observation. By repartitioning two sets at a time, products for any larger disjoint family likewise depend only on the union.

\noindent\emph{Recovering elementary factors.}

Consequently \(A_D\) tensors into elementary matrix factors \(T_i\) indexed by \(i\in D\), independent of \(D\) (scalars allowed). Here are details for \(k\ge2\). At a fixed \(B\) take external mutually disjoint \(k\)-sets \(B_i\) for \(i\in B\). Repartition using sets \(X_i\) each consisting of \(i\) and all but one entry of \(B_i\), plus the final set of leftovers. Then \(A_{X_i}\subseteq A_B A_{B_i}\); take its coefficient algebra \(L_i(B)\) in \(A_B\) relative to this tensor expression (generate from all coefficients). These coefficient algebras commute with each other since the \(A_{X_i}\) commute and the external factors are independent. They jointly generate \(A_B\), since the leftover-set factor lies entirely in the product of the \(A_{B_i}\)'s (repartitioning the auxiliary sets), and by repartition of the whole union including \(B\), \(A_B\) now lies in the product of the coefficient-generated algebra with the independent external tensor factors. Thus they are central simple tensor factors (commuting generating algebras have scalar centers and no radicals). For another choice with disjoint auxiliaries they commute across choices at different positions by the same reasoning, so coincide at matched positions by centralizers. Compare arbitrary choices through spare disjoint auxiliaries. And for a different \(B\) still containing \(i\), one can use the same \(X_i,B_i\) auxiliary choice, getting exactly the same coefficient algebra (take coefficients into the commutant of \(A_{B_i}\)). This proves the assertion; these factors are selfadjoint when needed (take coefficients also in the adjoint basis, since the algebras used are selfadjoint). They commute across arbitrary positions, hence are simultaneously tensor independent by central simplicity. At most one can be trivial by label distinctness, and each nonscalar \(T_i\) has stabilizer \(M_{\{i\}}\) as usual, since it is \(A_D\cap A_E\) for \(D\cap E=\{i\}\) (two such paired meets of cut labels, with remaining indices disjoint, join to \(M_{\{i\}}\); the stabilizer is proper by the tensor order bound). For \(k=1\) this conclusion is immediate instead.

\noindent\emph{The aggregate supplement.}

Take prescribed large fixed set sizes, on nontrivial positions with degrees again kept small enough relative to total by thinning. For all such sizes product-factor stabilizers on sets \(P\) are maximal (large unequal sides), hence \(C_{B(P)}\), \(\emptyset\ne B(P)\subseteq P\); they are distinct and thus \(B(P)=P\) after homogenizing on a sufficiently large reservoir. Use sizes including \(l,2l\) for sufficiently large even \(l\). For disjoint \(D,E\) of size \(l\), \(P=D\cup E\), write \(J=C_D\wedge C_E, F_0=C_P\). On the aggregate for \(P\) we have the full supported high special projectively, say \(R\), normal in \(K_{F_0}\) and lying inside the actual ancestor simple. There is no loss from possible outer cosets: we claim
\[
                  K_J R=K_{F_0}.
\]
Indeed transport \(R\) as a normal product over initial sites at the product label \(F_0\), and join to \(M_P\), giving the product label of \(K_{M_P}R\) in the exact interval \([M_P,F_0]={\rm Sub}(S_{2l})\). Its meet with the alternating group there is either endpoint on \({\rm Sub}(A_{2l})\), by the earlier exact normal test. If the meet is trivial, the projected normal-multiple label in \({\rm Sub}(S_{2l})\) itself has order at most two, and thus preserves some \(l\)-set there (\(l\) even). This puts \(R\) inside the actual factor stabilizer by the cut labels, impossible on its high aggregate by the tensor order bound. So it contains the alternating group formally, giving the claim by joining with \(J\).

\noindent\emph{Degree equality and the characteristic-two exception.}

In orthogonal characteristic two use instead disjoint unequal large
sizes also chosen as above, with their sum.  Then \([J,F_0]\) is
formally just a cover (Young subgroup with unequal blocks in the
combined symmetric group; any movement of the split yields a crossing
transposition by conjugation).  But \(R\) includes the supported
symplectic group on the alternating aggregate (tensor with a further
alternating space, with canonical total quadratic).  Within \(F_0\)
we have a proper intermediate between the split stabilizer \(J\) and
the upper label by imposing the canonical tensor quadratic on the
aggregate, with properness over the split by its supported high
orthogonal group and tensor orders.  Indeed \(K_J\) stabilizes both
ordered pieces and acts by semilinear similarities on them, respecting
the split's canonical aggregate quadratic (zero on pure split tensors)
with the aggregate polar multiplier.  Not all of the supported
symplectic group respects it.  This excludes the case directly.

In the remaining types use our equal-size index sets.  If the two
actual side degrees differ, \Cref{fa:ordered-factor-maximal} makes
\(K_J\cap R\) maximal and nonnormal in \(R\).  This intersection
is the full ordered stabilizer there: within \(R\), fixing one piece
also fixes its commutant piece on the aggregate.  Since
\(K_JR=K_{F_0}\), intersections with \(R\) determine all ordinary
intermediates, so \([K_J,K_{F_0}]\) would be a cover.  Yet the
formal interval has a proper wreath intermediate.
Apply \Cref{fa:product-cover} at the original sites before the
ancestor quotient.  It would make the lower site projections normal
in the upper ones, and therefore make their intersection with the
actual \(R\) normal as well, contradicting the preceding
nonnormality.  Thus all disjoint \(l\)-products of degrees agree.
Comparing replacements with spare sets makes all elementary degrees
agree; finite coloring aligns their form classes.

\noindent\emph{The final permutation count.}

Now under \(F_0\) the label allowing permutations of the chosen slots in \(P\) has kernel label \(M_P\) for that action. It induces all symmetric permutations. Indeed use the aligned tensor transpositions in the ambient socle, with corrections as follows. For common degree \(d\) the interchange has negative-eigenspace dimension \(\binom d2 v/d^2\) for sign purposes in odd characteristic, which is even when \(d\) is even by the many slots, and also when \(d\) is odd and the total dimension even. Linear/unitary determinant can always be corrected (by minus identity for odd total dimension). In orthogonal odd characteristic with even slots the tensor complement of the two has even dimension and square determinant by additional even slots. The negative-space form determinant modulo squares for the involution is then square by tensoring, giving also spinor membership. For odd slots use isometries on the large complement of all the chosen slots in \(P\) for determinant/spinor corrections: tensoring with identity on an odd-dimensional symmetric factor still gives both independent invariants (diagonalize that factor; determinant unchanged and spinor unchanged up to a determinant-dependent scalar class). The full aggregate group exceeds this permutation stabilizer by orders. Thus above \(M_P\) there are more elements than the entire formal interval to \(F_0\), contradiction.

\end{proof}

\begin{lemma}[The remaining ordered-factor possibilities]
\label{fb:nuisance-planes}
After \Cref{fb:both-high-factors}, the missing side of an ordered-factor
failure has large degree.  Its retained side can only be a linear plane over
\(\mathbf F_2\), an odd-characteristic bilateral orthogonal plane, or a
symplectic plane over \(\mathbf F_2\).
\end{lemma}
\begin{proof}

For remaining ordered-factor failures the missing side must still be of large dimension by the missing-site bound: the normal products of supported groups modulo centers have soluble remaining quotient and at most the two-slot ambiguity, so some missing simple orbit would otherwise have bounded type. (If there are two projective simple components on a side, use a missing transitive orbit thereof; supplementation over \(g_D\) gives the top action, here of bounded pair rank.) Thus it is the retained side that is now small. The component list and arguments of the earlier ordered-factor pruning then apply except for retained nuisance factors: a linear plane over two, an odd-characteristic bilinear orthogonal plane, or a symplectic plane over two. Indeed across a disjoint pattern all other retained components lie mutually in the stabilizer labels, with exactly the same normalization and spanning arguments giving commuting distinct component algebras (one or two per label); the cases giving forbidden direct systems there are already excluded. From arbitrarily many disjoint labels then take many distinct commuting slots. Label distinctness indeed requires many distinct slots since the designated slot set determines the designated algebra. The common intersection preserves each label's set of at most two slots (orbit sizes at most two even on the union), so preserves an ordered aggregate of many with still many unused, impossible by \Cref{fb:both-high-factors} on the combined support. This proves the partial exclusion needed.

\end{proof}

\subsection{Field algebras and paired projections}
\begin{lemma}[Transfer degree at full simultaneous stabilizers]
\label{fb:large-transfer}
For the prime-field witnesses in this round, the transferred natural degree
exceeds any prescribed absolute bound after the earlier missing-site tests
have been enlarged.  The same conclusion holds at the full simultaneous
field or paired stabilizers used below, relative to their lower endpoint
\(g_*\).  Prime-field labels are coatoms.  Two disjoint witnesses whose
fields commute cannot have a field compositum.
\end{lemma}
\begin{proof}

For any prime field algebra label under consideration, the large transfer degree hypothesis of the stabilizer observation holds: the transferred special inside the ancestor simple gives, modulo its soluble center, nonabelian simple orbit product(s) with soluble remaining quotient, one of them necessarily missing by the pure-link normal test, hence of high type. To spell out this use of the standard orders/structures: if the module degree were bounded, total ground dimension large would force the extension and its parameter large; permanently soluble dimensions then give only a soluble stabilizer, and otherwise the natural perfect special modulo center gives the simple(s), with at most the split orthogonal four-space ambiguity. Semilinear field automorphisms, similitudes, extra scalars or anti-actions do not affect quotient solubility. We use the same argument for the transferred groups in simultaneous intersections below \emph{when the intersection is the full simultaneous stabilizer as indicated}, over \(g_*\) there. More explicitly, any center lost on passing from such a normal transferred perfect special group to its simple site(s) is soluble normal and retained; if the simple sites are not transitive take a missing orbit (by the normal test and solubility of the rest above the supported products), with boundedly many sites and the lower pure-link endpoint supplying the orbit action by supplementation. Thus the high-type conclusion of the missing-site argument applies; we may demand the transferred natural degree exceed any needed absolute bound, by the bounded-rank test there. Thus field labels are coatoms by \Cref{fb:field-pair-overgroups,fb:canonical-quadratic-exclusion}. For two disjoint-support instances the intersection \(J\) is covered in each, with also a third coatom above it. In particular if the two fields commute with compositum a field, their intersection is the full normalizer of that field, of two-prime degree (equality would give identical labels); its large transfer and \Cref{fb:field-pair-overgroups} rule out the third coatom. For any such application in the symplectic even case with individual trivial adjoints, the compositum also has trivial adjoint.

\end{proof}

\begin{proposition}[Exclusion of fields, pairs, and nuisance planes]
\label{fb:field-pair-exclusion}
In the prepared forward-test configuration, the prime-field and exchanged-pair
failures are impossible.  The nuisance planes of \Cref{fb:nuisance-planes}
are therefore impossible as well, completing the ordered-factor exclusion.
\end{proposition}
\begin{proof}
\noindent\emph{Linear type.}

In linear type the prime-field proof with retained projective tori from \Cref{sec:classical} thus works: commutation gives a field by the direct-system exclusion; the only possible noncommuting exceptional pair gives a quaternion \(M_2\) factor over three normalized by the intersection, now forbidden by \Cref{fb:nuisance-planes}. We thereby also exclude the remaining linear ordered factors (plane over two gives a canonical quadratic field).

\noindent\emph{Fields with nontrivial adjoint in isometry types.}

In isometry types first take field patterns with nontrivial involution restriction. For two disjoint occurrences use \(R=F[D_i,D_j]\) as in the previous pruning, semisimple with either center field or two exchanged central components. We give the adaptation. If a larger center field gives an invariant central prime \(D'\), then either \(K_i\) normalizes it (forcing \(D'=D_i\) by the perfect transfer and centralizer), or the full simultaneous stabilizer of \(D_i,D'\) equals \(J\) by covering. In the latter, if \(F[D_i,D']\) is a field its large transfer inside \(J\) centralizes \(D_j\), giving commutation. If not, there must again be two exchanged components, from equal prime degrees, and the full intersection includes the paired special linear group over \(D_i\) on the components, normalized there (\(D'\) is likewise scalar over \(D_i\) on each, by the tensor of the two equal-degree fields, a product of copies of \(D_i\)). Its two spaces are dually totally isotropic and totally singular when relevant by previous irreducibility. The quotient over this paired type of group is soluble, so by the pure-link normal test and the missing-site bound its module degree on each component is again large. Indeed on a pair commuting with \(D_i\) use the transferred pairing over \(D_i\), acting by a determinant-one matrix and its inverse adjoint. In the split-algebra use this also respects \(D'\) since on each component the two prime fields act by the same field of scalars. Quadratic restrictions with zero polar on the pair spaces in characteristic two are zero by previous-good irreducibility (else the kernel sum gives an invariant proper nonzero subspace), so pairing preservation suffices there. This special lies inside the ancestor simple by perfection in the nonsoluble large-parameter/dimension range, normalized in the simultaneous stabilizer which acts by paired semilinear similarities allowing exchange (thus soluble remaining quotient; dimension one itself would just give solubility). It too centralizes \(D_j\), and its ground-field centralizer commutes with \(D_i\), as in the earlier paired situation: in the now large degree the natural and adjoint-dual constituents over the embeddings are pairwise distinguished, including across the two sides. The same argument applies directly if the center of \(R\) has two exchanged components, using their pair stabilizer which cannot contain \(K_i\) by the field centralizer; again its full intersection with \(K_i\) must be \(J\).

For center \(F\), \(R\ne\mathcal A\) since otherwise the intersection has soluble image on the fields with soluble kernel, impossible. Thus \(R\) gives an ordered tensor, necessarily with retained nuisance plane side. If \(R\) itself is that plane algebra, it cannot contain two distinct fields of the nontrivial-adjoint kind: on an odd orthogonal plane such a field has the two singular lines as eigenspaces over the closure and is uniquely recovered, and over two the quadratic field in a plane matrix algebra is unique. If its commutant is the nuisance plane instead, \(K_i\) cannot preserve the tensor (large transfer and field centralizer), so the full intersection of \(K_i\) with this ordered tensor stabilizer is \(J\). On the large factor we have the perfect transfer relative to \(D_i\), with again large degree by the missing-site argument (similarity projection on the nuisance side is soluble), supported inside \(J\). Here transfer is on that large factor space for its factor form; in the bilinear characteristic-two setting this is alternating, and the tensor with the nuisance-side identity preserves the canonical total quadratic in orthogonal type. Its centralizer in \(R\) is \(D_i\), but would include \(D_j\), impossible. In this transfer-on-a-factor use, the full simultaneous stabilizer acts on that matrix factor by semilinear similarities respecting \(D_i\), so on its factor space one again has only a soluble remaining similarity quotient over the transferred special. On the nuisance plane factor the whole similarity projection is already soluble. Likewise for the simultaneous paired transfers, after fixing field scalars and the two summands, linear action on one summand and the pairing multiplier determine the action on the other. Thus the soluble-quotient tests indeed use the corresponding supported perfect transfer images inside the simultaneous group, without needing arbitrary individual similitudes there to embed unchanged in the ancestor simple.

Consequently any remaining pair commutes. Two distinct commuting prime fields of the same degree generate a split algebra; for homogeneous adjoint type its primitive component projectors are selfadjoint (relative identifications of the fields commute with their indicated adjoints), contrary to direct-system exclusion. Distinct prime degrees give the field compositum contradiction already proved. This finishes the first prime-field case.

\noindent\emph{Exchanged pairs and trivial-adjoint fields.}

Next exclude exchanged pairs \(e,e^*=1-e\) of idempotents; total singularity when necessary follows as before. Their labels likewise are coatoms by \Cref{fb:field-pair-overgroups}. For the algebra generated by two projections from disjoint instances a larger center field would give a prime \(D'\) commuting with them and normalized by the intersection. It is not normalized by the first label (large paired group centralizer). The full intersection of the first pair stabilizer with \(N_M(D')\) thus equals the pair-label intersection \(J\). Again it includes a paired \(D'\)-linear special group of large degree by the same argument. Centralizing the other projection (perfection), it forces commutation of the two projections by distinct natural and adjoint-dual constituents. But distinct commuting exchanged pairs normalized by the intersection give all four cells (zero cells paired oppositely would instead force the same pair), whose sums on adjoint orbits give a forbidden nondegenerate direct split, just as before. A center with two exchanged components likewise would force that impossibility since its pair commutes with both. For center \(F\), their generated algebra is \(M_2(F)\) by the scalar \((e-f)^2\) argument for two idempotents \(e,f\). This must have nuisance plane type by the ordered-factor exclusion. An odd orthogonal plane gives at most one exchanged pair, and the plane over two now violates the nontrivial-adjoint field exclusion. This finishes paired exclusions, hence nuisance planes too by the canonical pair or nontrivial-adjoint quadratic field there. Remaining trivial-adjoint prime fields now have the same two-label proof as above (center \(F\) immediately impossible since all such ordered factors are excluded).

\end{proof}

\subsection{Permuted tensor systems}
We first record the generation calculation that will replace direct-space
merging when the objects being combined are sets of tensor directions.

\begin{lemma}[Merging compatible tensor directions]
\label{fb:tensor-merging}
Consider compatible sets of tensor directions \(A B\) and \(B C\) in the
prepared classical configuration, with common aggregate degree \(\dim B\)
sufficiently large.  Suppose their tensor-supported natural special groups
are available.  They generate a group containing the tensor-supported special
on \(A B C\).  In orthogonal characteristic two, use the canonical orthogonal
special on every subset with at least two slots, with product polar form and
quadratic zero on fine simple tensors.  Symplectic characteristic two has
already been excluded by \Cref{fb:canonical-quadratic-exclusion}.
\end{lemma}
\begin{proof}

On compatible sets of directions call the exclusive parts \(A,C\) and common part \(B\), also denoting their aggregate spaces by these letters. Assume \(B\) has sufficiently large degree, and the two direction sets have their full tensor-supported special groups. In orthogonal ambient characteristic two we need only the canonical orthogonal specials there and on the desired union when multislotted (still no symplectic characteristic two ancestor here). Then these groups generate so as to contain the union special; nested cases need no argument, so take both exclusive parts nonempty. Work first on \(A\otimes B_0\otimes C\) where \(B_0\) is a hyperbolic summand of \(B\) of defect at most two, or the entire \(B\) in linear type. In \(AB\) take transfers \(1+N-N^*\) with \(N\) directed between signed hyperbolic \(B\)-coordinates \(s\to r\) from distinct pairs, and arbitrary \(A\)-matrix there. These are in the natural special: the adjoint arrow uses the opposite coordinates in exchanged order, so \(X=N-N^*\) has square zero and \(1+X\) is a determinant-one isometry (in odd orthogonal type unipotent of odd order gives trivial spinor norm). For the canonical quadratic in characteristic two the displacement \(X\) has totally singular image in the two \(A\)-fibers at mutually orthogonal isotropic \(B\)-coordinates, by canonical zero on pure tensors; also \(b(x,Xx)=0\) by symmetry and adjunction for its product polar form \(b\), and displacement rank is even. In \(BC\) take similarly \(t\to s\) with arbitrary \(C\)-matrix \(L\) in the transfer arrow. Use three distinct pairs. Their commutator gives the paired transfer \(t\to r\) with product arrow tensoring the two arbitrary matrices, by elementary matrix multiplication; the higher cross terms vanish. Thus we have all such arrows on \(A\otimes C\) by linear spanning. In linear type simply use single unpaired arrows with the same calculation. In individual basis coordinates now, even short paired arrows between nonopposite (distinct-pair) coordinates within one \(B\)-pair region follow by commuting through a spare pair. These standard short transfers give the split special on \(A\otimes B_0\otimes C\) supported there: ordinary elementary/root generation (in odd symplectic type commutators give also the long roots with factor two, and in hyperbolic unitary type the commutators through an intermediate index to the opposite of the initial one give all trace-zero long parameters). In orthogonal characteristic two they give the split type D generators for the indicated quadratic, which suffices. These are assertions on natural matrices (or their projective images), with identity on the complementary summand inside \(ABC\), and acting by tensor identity in all outside directions.

To see the generator step explicitly, the arbitrary \(A\otimes C\)-arrows between distinct \(B\)-pairs are additive under transfer multiplication, and the commutator used to obtain them has just the composite arrow \(N_1N_2\) and its negative adjoint when the three pairs are distinct. Choose individual hyperbolic bases on the paired \(A\otimes C\) fibers over \(B_0\), totally singular on each coordinate side (also for the indicated quadratic). Rank-one arrows between distinct big pairs give the elementary paired updates there, then any two distinct individual pairs can use an intermediate signed coordinate from a spare big pair. Thus all distinct-pair root groups are present. If a long root is needed, commutation through an intermediate index with final index opposite to the initial one yields on that last arrow the composite minus its adjoint (higher products still zero), i.e. \(2c\) in the standard odd symplectic coordinates or \(c-\bar c\) in normalized hermitian coordinates, with arbitrary composite parameter \(c\). This supplies all long parameters even in even-characteristic unitary type. Ordinary root generation of the split special here thus applies (hyperbolic orthogonal uses just type \(D\) arrows, including at characteristic two). In linear type the unpaired calculation gives all elementary matrices analogously.

If a small anisotropic part remains, take an isometric copy of it inside \(B_0\). Using also the perfect special on \(B\) (embedded in a given supported group), conjugate to the special supported instead with \(B_0\) moved onto the orthogonal complement of that copy, by Witt transitivity with corrections (the large orthogonal complement supplies determinant/spinor data as in the subspace arguments). Here hyperbolic splitting of \(B\) uses its bilateral or hermitian pairing; in bilateral characteristic two it is alternating with no remaining part. A large enough hyperbolic \(B_0\) represents an isometric subspace copy of the small anisotropic part in the other form cases (e.g. embed together with its negative in hyperbolic pairs). The supported embedding on \(B\) used for transport lies in a given special by tensoring with identity (form preservation and high-rank perfection). These two resulting large direct spaces within \(ABC\) are compatible with large nondegenerate intersection and full sum, so \Cref{fa:direct-merging} finishes the argument. This proves the tensor merging assertion.

In applying this repeatedly in orthogonal characteristic two, use on every multislotted subset the quadratic with the product polar and zero on the fine simple tensors there. It is also zero for pure tensors across any two-part cut there (apply the canonical construction for that cut and compare on a tensor basis), so column groups or their altered versions supply the required specials. Aligned internal fine permutations carry these choices as required; when transporting whole merged sets of at least two columns the form is canonical on column tensors too, so arbitrary column similarities in those transports cause no problem.

\end{proof}

\begin{proposition}[Exclusion of permuted tensor systems]
\label{fb:permuted-tensors}
After the preceding rounds, neither a full-algebra elementary abelian
normalizer nor a nontrivial permuted tensor-system stabilizer can be a failure
label in the forward test.
\end{proposition}
\begin{proof}
\noindent\emph{Common tensor refinement.}

We now have the same full-spanning property as in \Cref{c:full-spanning} for a nontrivial inner subgroup normalized by a previous-good label. Thus the elementary abelian exclusion again works by scalar commutator pairing and inclusion of all retained elementary groups in the common intersection (normal there). Take then the tensor-system labels allowing permutations, with the very same reductions of exceptional small types as there. In the retained-product case we again have minimal normals in the common intersection recovering the slot systems. Otherwise take two disjoint occurrences, with common intersection \(E\), and missing transitive products with rich top action (thus high slot dimensions by the missing-site bound). The supported normal products in \(E\) supplied by \Cref{fa:intersection-support} contain in their images the same base-core monolith (using the lower endpoint \(g_*\)). Hence the tuple-intersection commutator argument, followed by full spanning, gives a common matrix tensor refinement by intersections assigned to slot pairs. Concretely use \(N_i=\prod_a R_{i,a}\trianglelefteq E\) of singly supported subgroups in the two missing products as supplied above (the projective tensor supports here have one nonabelian simple per slot). In the quotient by the core of \(K_{g_*}\), assign each simple factor of the monolith \(W\) for each \(i\) to some \(\bar R_{i,a}\) containing it, using \([W,\bar R_{i,a}]\) since \(W\le\bar N_i\). Then \(\overline{N_1\cap N_2}\) still contains \(W\) by commutators. Just as before take intersections of the assigned \(R_{i,a_i}\) with \(N_1\cap N_2\), normal there and their images containing the assigned simple by commutators; once again taking their mutual commutators gives its containment in the image of \(R_{1,a_1}\cap R_{2,a_2}\). Thus the tuple intersections generate a nontrivial \(E\)-normalized subgroup and jointly span fully. Distinct tuples have matrix lifts commuting linearly by their distinct slot supports in some system. So their nonscalar lift algebras are commuting central simple tensor factors as before (one per occupied slot pair), recovering each coarse algebra by containment and fullness. Use the same selfadjoint/form observations as before. The fine slots have uniform degree \(h\ge2\) and forms aligned up to scalar by previous-good transitivity; each row or column contains at least two by block primitivity in \(E\) (supplementation supplies transitive top on each system including at least all alternating permutations). Indeed a singleton width would occur throughout that system by transitivity and the other system would then group it into blocks, hence coincide by primitivity and having at least two slots in both, impossible for disjoint supports. In orthogonal characteristic two the fine forms are alternating and the total quadratic is canonical, zero on fine simple tensors; this canonical quadratic is also the one for grouping by \emph{any} bipartition of the fine slots (uniqueness from values on a tensor basis). Here \Cref{fb:canonical-quadratic-exclusion} excludes symplectic characteristic two.

\noindent\emph{A local alteration of columns.}

Use \(u\) the full row-system label intersected with the fine-system stabilizer, a proper decrease (order on any high row slot) above \(E\). Within \(u\) we can interchange two fine slots of a row or cycle three, leaving others in place as slots. Indeed use aligned permutations, commutators if at least three. For a transposition the determinant/spinor arguments given earlier work since there are at least four fine slots, using internal corrections on another slot for odd orthogonal degree (then \(h\) odd, at least three, repeated with odd multiplicity). In characteristic two orthogonal type the tensor swap has even displacement rank by the even \(h\) and at least four slots. Thus one column gets altered by a replacement, remaining in large overlap with its original version (common product degree at least the square root of the column degree).

\noindent\emph{Generation and the join contradiction.}

By \Cref{fb:tensor-merging}, adjoining \(u\) to the column label merges a column with its altered version. If the gained fine direction within another column has large degree, merge immediately with that column; otherwise transport to its other slots by aligned within-column alternating permutations (its fine-slot count then at least three) and merge extensions via their large common column. As with direct systems conjugation under the column label then gives merged pairs on a connected graph of columns by primitivity; overlapping edges (and enlarging unions) merge with at least a large column in common, giving the full special projectively, thus the ancestor simple, contradicting \Cref{fa:disjoint-join}. Thus these last systems too are excluded.

\end{proof}

\subsection{Grounded descendants and uniform thresholds}
\begin{proposition}[Grounding at bounded omissions]
\label{fb:grounding}
After the prescribed rounds, every tested descendant has an actual nested
normal simple supplement with soluble centralizer in the base.  Relative to a
high classical ancestor it is absolutely irreducible and avoids all the
excluded tensor structures.  Relative to an alternating ancestor it is
primitive and transitive unless the desired point configuration has already
been obtained.  Its simple type can be required to exceed any fixed bounded
order or bounded-rank threshold.

The same actual normal simplicity and full spanning apply to a label above a
previous-good axis and below the ancestor when it has the indicated
nonabelian part: either a contained perfect simple group as used here, or its
own pure link satisfying the normal test.  In the first case that perfect
group lies in the resulting normal simple.
\end{proposition}
\begin{proof}

First extract a normal simple inside the ancestor socle, using the
structural conclusions just proved.  In classical type this is the
minimal-normal argument of \Cref{c:classical-conclusion}: elementary
abelian normals have been excluded, while multiple nonabelian simple
factors would give a forbidden tensor system through their commuting
lift algebras.  The remaining actual normal simple spans the whole
matrix algebra.  Its centralizer in the ancestor almost simple group
therefore injects into the soluble outer quotient.  In alternating
type the primitive, non-affine, non-Cartesian and non-diagonal
conclusions give an almost simple group, as in
\Cref{c:alternating-conclusion}.

In either case the tested label is soluble over this simple, also
after pulling back the ancestor's soluble centralizer.  The simple
lifts inside the actual ancestor socle, where the quotient map is
faithful; hence successive simples really are nested subgroups.
For an axis descendant the pure-link normal test has lower endpoint
\(H\).  It puts the soluble centralizer in \(H\) and forces the
simple to be missing: if it were retained, the base-core quotient
would be soluble, contradicting its nonabelian monolith.  Thus the
axis label is supplemented by this actual simple.

For a general label above the previous-good axis, the same extraction
applies when its own pure link or a contained perfect simple group
supplies the nonabelian part.  We need no supplementation assertion
with the original base in the latter case.  Solubility of the quotient
over the extracted simple already forces every such contained
perfect group into it.  This is the conclusion used at the third
coatom in the protection arguments.

Finally, at axis descendants the missing-site product test applies
to their actual simple supplements inside the original \(S\), and
\Cref{fa:bounded-types} makes their types exceed any fixed threshold.
Every link and intersection used here has bounded omissions; the
earlier exclusion bounds were enlarged to make those intersections
previous-good.

\end{proof}

\begin{remark}[Choice of thresholds]\label{fb:uniform-thresholds}

In arranging these rounds the lower high-dimension requirements in the merging, large-factor and projection calculations are fixed thresholds, not functions of the field parameters. For instance the bounded-codimension grouping/collapse of near direct summands needs only an absolute number of transports (separate off bounded small-space dimensions first, and in the high smaller dimension use the full special orbit argument); choose the lower dimension cutoffs there accordingly. When forcing actual missing sites to be large, take the bounded-type threshold in their product tests also large enough for subsequent requirements such as large column dimensions/degrees (and hence large overlaps after moving one cell). At full simultaneous field or paired stabilizer intersections a bounded module dimension would force large field size in their transferred actions by the high total dimension; the resulting missing-site tests still have arbitrarily long pure links by the original residual marker count. Product transfers of actual simples at axis descendants likewise require only that many residual markers, not an a priori classification of all product labels at that simple. The counts of indices/replicas used together depend on fixed tests and the subsequent bounded path/omission requirements, not the sizes of those simple groups. Thus the enlargements backwards and homogeneity relative to path prefixes as at the start suffice here.

\end{remark}

\subsection{Selecting protected extensions}\label{fb:protection-setup}
We now use the actual nested simples and the path-length bound of
\Cref{w:prepared-path,w:path-bound}.  The protections to be imposed on each
continuing high-classical run are a single tensor atom and absence of a
fixing prime root, relative to every earlier member of that run.

Write \(S_i=B_D\) for any earlier member of the run, including the
current parent, and work above a large supply of fresh ordered
positions. For a trial extension at the parent with omission set \(W\supseteq D\) take increments of prescribed bounded sizes, all beyond previous positions. At trial nodes of matching (continuing) characteristic and classical type relative to \(D\), we have a single cycle under \(\sigma\) on tensor atoms as before, by the structure grounding. The orbit, label and counting arguments of \Cref{w:atom-cycle,w:coarsening-count,w:root-count} apply, only now heights up to \(K_{M_D}\) above \(K_{M_U}\) are bounded by
\[
             O(n\log\log(n+3)),\qquad n=|U\setminus D|.
\]

\begin{lemma}[Reduction to singleton protection failures]
\label{fb:protection-count}
With the bounded predicates homogenized relative to the allowed ordered
prefixes, choose trial increments \(a<b\) of matching continuation color and
with \(b\) sufficiently large compared with \((|W|+a)^4\).  If a protection
failure survives at a continuing trial extension, then one obtains arbitrarily
many coatom witnesses on distinct singleton positions, of the same protection
kind, corresponding to distinct primes.  All these witnesses are counted for
one fixed later descendant \(B_U\).
\end{lemma}
\begin{proof}

Use preliminary homogeneity for a bounded number of comparisons: existence of the actual path socles and their types/continuation relative to prefixes, the structure exclusions, and whether a bad protection label for \((D,L)\) exists for a specified ordered sublist support (and of which of the two kinds), suffices. Choose extension sizes \(a<b\) from \(W\) having matching color for continuation from \(W\), with \(b\) very large compared with \((|W|+a)^4\) (use three spaced sizes). If not continuing at \(a\), use that next size with no new run protection. Otherwise take \(W\subset L\subset U\) with these increments, both continuing. For a run ancestor \(D\), any bad prime-coarsening or prime-root label for \(B_L\) (here \(B_L\) denotes the simple at \(M_L\)) occurs among those counted for \(B_U\). Indeed \(B_U\le B_L\) by actual normal simplicity and soluble centralizer as before. Individually invariant tensor slots for \(B_L\) must group atoms for \(B_U\), with just the unique prime congruence coarsening at that prime; and fixing roots still fix \(B_U\). Their stabilizer/normalizer labels for \(B_L\) are proper below \(K_{M_D}\) and contain \(K_{M_L}\). If one such bad case occurs, we can vary the \emph{entire} list \(L\setminus D\) (even the positions from \(W\setminus D\)) within \(U\setminus D\), retaining by homogeneity the same test relative to \(D\) and ordered omitted support. These altered intermediate labels still have the indicated actual data, and their simple groups contain \(B_U\) by the same soluble-quotient argument (the altered \(M_L\) still contains \(M_U\), even when no longer under \(M_W\)). Indeed we are comparing throughout relative to the fixed \(D\), which still comes first in index order. If the nonempty omitted support of the bad label has size at least two this gives at least \((n/(|W|+a))^2\) different supports (vary \(L\setminus D\) over all subsets of that size in \(U\setminus D\); each specified support of size \(z\) can then occur with probability at most \((|L\setminus D|/n)^z\)), violating the combined atom-prime and root count \(O(1+n\log\log(n+3))\) against the stipulated spacing. If singleton it gives arbitrarily many coatom labels omitting distinct single positions \(x\), relative to \(D\), of the same kind, still from coarsenings or roots for this same \(B_U\). A singleton omitted support leaves no choice of side group on it. In the tensor case there is just one possible such full stabilizer per prime in the count for \(B_U\), and in the root case at most two (height and diagram uniqueness), so we still get arbitrarily many distinct primes.

\end{proof}

\begin{remark}[The third coatom]\label{fb:protection-third-coatom}
Take two of the singleton coatom witnesses from \Cref{fb:protection-count}.
The following common setup applies to the two possible kinds.

For two such coatom labels we also have a third distinct coatom over their intersection in the interval under \(M_D\), namely the cut on their position pair. Write \(K\) for this third coatom group in the almost simple ancestor. It lies above \(K_{M_U}\) so has full spanning and structure avoidance there. Thus it has an absolutely irreducible normal simple \(B\le S_i\) containing any perfect subgroup in \(K\cap S_i\) by Schreier and the trivial inner centralizer in \(S_i\). We will use these assertions when \(K\) contains the indicated high perfect group; they follow by \Cref{fb:grounding}. Write \(v\) for the ancestor natural dimension.

\end{remark}

\begin{lemma}[Singleton root failures]
\label{fb:singleton-roots}
The singleton coatom witnesses in \Cref{fb:protection-count} cannot be
normalizers of fixing prime roots.
\end{lemma}
\begin{proof}

In the root case take distinct primes \(r,s\). Their roots fixing \(B_U\) combine to \(\Theta\) with \(\Theta^{rs}=\Phi_i\) and the two given roots as powers as in \Cref{w:root-count}. Now \(Y=S_i(\Theta)=(G_i^\Theta)'\) lies in the intersection. Its presence forces \(B\) high classical in the same characteristic by \Cref{fa:degree-gap}, and the natural dimensions compared both ways force \(B\) to have dimension \(v\) and act naturally up to twist in the ancestor. The root series agrees by the same form checks using the natural \(Y\). In particular if the ancestor here is symplectic in characteristic two, its natural subfield \(Y\) from \(\Theta\) cannot preserve a polar-nondegenerate quadratic even over the closure, excluding orthogonal \(B\) acting naturally there. In natural subfield coordinates its symplectic matrices on the finite cover act transitively on nonzero finite vectors (invariance of a putative quadratic would be exact by perfection); a constant quadratic value on those vectors forces zero polar form (for the binary field use three independent vectors and polar bilinearity). In orthogonal ambient characteristic two the preservation test instead excludes a natural symplectic \(B\) itself the same way. Hence \(B=S_i(\Psi)\) with \(\Phi_i=\Psi^e\), \(e>1\). Applying \Cref{w:natural-root} also to \(\Psi\) fixing \(Y\) gives \(\Psi=\Theta^h,\ eh=rs\). Lifts normalizing \(B\) commute with \(\Psi\) by uniqueness with given outer part. Thus \(K\) already normalizes the simple of one of the two given roots, impossible.

\end{proof}

\subsection{Natural recognition in the tensor-coarsening configuration}\label{fb:tensor-configuration}
The root case supplied a subgroup already acting naturally on the
ancestor module.  For tensor coarsenings the subgroup in the two-cut
intersection instead acts as a tensor product of spaced natural twists.
We must prove that the normal simple group in the third coatom still
acts naturally on that module.  The following argument uses this
particular tensor configuration throughout.

Choose two distinct odd primes \(r,s\), leaving many unused primes at least as large. Put \(d=rs,\ m=v^{1/d}\); thus \(m\) is very large even compared to \(d^4\) by the remaining prime factors in the atom cycle (the \(B_U\)-atom degree is at least two, raised to the cycle length divided by \(d\)). In particular throughout the following bounds we can impose sufficiently large absolute constants and \(m\gg d^4\), by using a sufficiently large fixed number of distinct labels. In their intersection is the high derived fixed group \(Y\) on the \(d\)-slot coarsening, as in \Cref{w:coarsening-count}. On the ancestor module over the closure it acts projectively irreducibly with lift tensoring \(d\) natural modules of dimension \(m\) at twists spaced by the ancestor field height \(f\), allowing duals. Indeed the cyclically embedded fixed group projects on one slot with parameter \(p^{df}\), and the remaining slot actions are its successive coefficient twists by the ancestor map, as described there. Restricted weights have these \(d\) distinct natural (or dual) digits, hence the individual \(d\) slots are exactly its tensor atoms. If the ancestor is symplectic in characteristic two, the two coarsenings already give canonical quadratics (polar the given form, zero on their respective simple tensors). These are coefficient-ground rational by Frobenius invariance, and the respective labels respect them up to the semilinear similarity action: maps preserving a system of matrix tensor factors permute pure tensor lines (apply factor identifications to the induced semilinear matrix-algebra maps; an intertwiner after that is scalar). The quadratics are equal by square difference and absolute irreducibility under \(B_U\), using perfect lifts for exact invariance. Their joining labels cannot then generate the ancestor (the natural symplectic group here does not preserve such a quadratic). We exclude that case here.

Thus, for the remaining argument, \(d=rs\) with distinct odd primes,
\(v=m^d\), and \(m\) is larger than a sufficiently large absolute multiple
of \(d^4\).  The group \(Y\le B\le S_i\) has parameter \(p^{df}\), and
its action on \(V\) has precisely \(d\) natural or dual-natural atoms at
positions \(0,f,\ldots,(d-1)f\).  The notation \(B\) still denotes the
normal simple in the third coatom group \(K\), supplied by
\Cref{fb:grounding}.

Both \(Y\) and \(B\) act irreducibly on \(V\). We first show that
\(B\) is high classical in the same characteristic, and bound the
support of the highest weight for its action on \(V\). Those bounds
then prove irreducibility of a different module: the natural module
of \(B\), restricted to \(Y\). This permits us to lift the inclusion
to algebraic groups and compose rational highest-weight
representations. The ensuing universal cuts and weight-multiplicity
comparisons force the action on \(V\) itself to be natural.

\begin{lemma}[A single sparse restricted weight]
\label{fb:sparse-weight}
In this configuration, \(B\) is high classical in the same characteristic,
with natural dimension \(w\ge m\).  Its action on \(V\) has highest
weight \(p^x\lambda\), where \(\lambda\ne0\) is \(p\)-restricted.
Every active type-A node of \(\lambda\) defines a cut whose smaller side
has fewer than \(2d\) entries.  In a signed type, its dominant word has
fewer than \(2d\) nonzero coordinates and has no spin support.
\end{lemma}
\begin{proof}

Again \(B\) contains \(Y\) and is high classical in the same characteristic (\(v=m^d\) is small compared with cross-characteristic/permutation bounds for \(Y\)). Its natural dimension \(w\) is at least \(m\), by the exact minimum-degree bound in \Cref{fa:degree-gap}. Its own tensor atoms on the ancestor module are again one \(\sigma\)-orbit by tensor-system avoidance. They group those of \(Y\), so if nontrivially split, \(K\) already lies below one of the two given coarsening labels, impossible: a prime congruence coarsening here would have to come from the \(d\)-cycle of \(Y\)-atoms, necessarily at \(r\) or \(s\) and coinciding with the corresponding system used above. Thus its lift there has highest weight \(p^x\lambda\) with \(\lambda\ne0\) \(p\)-restricted. At every active node of \(\lambda\) in A the cut size or its complement has size less than \(2d\) (Weyl orbit with binomial lower bound versus \(m^d\)). In signed types there are fewer than \(2d\) nonzero coordinates of the dominant word, and no spin supports. Indeed with rank \(e\) and \(k\) nonzero coordinates the orbit has size at least \(2^{k-1}\binom e k\), and use \(e\ge(m-1)/2\), \(m\gg d^4\): on \(2d\le k\le e\) log-concavity puts the minimum of the bound at one of the endpoints, both too large. Likewise \(\binom{w}{2d}>m^d\) in A.

\end{proof}

\begin{lemma}[Irreducibility on the natural module]
\label{fb:natural-restriction}
Under the preceding hypotheses, the projective natural \(B\)-module remains
irreducible on the finite group \(Y\) over the algebraic closure.
\end{lemma}
\begin{proof}

Suppose the projective natural module of \(B\) is reducible on the
finite group \(Y\) over the closure.  In linear type this puts \(Y\)
in a proper parabolic.  In a form type an invariant subspace either
gives such a parabolic, gives a proper nondegenerate splitting, or
has a nonsingular radical line in the characteristic-two orthogonal
case.  We treat the parabolic case first, then the nondegenerate
splittings with two positive ranks, and finally the line cases.

A proper parabolic preserves the filtration of every nontrivial
rational highest-weight module defined by a noncentral cocharacter
(on the covering group if needed).  That filtration is nontrivial:
the projective kernel cannot contain the cocharacter torus, so
distinct cocharacter weights occur.  Parabolic containment would
therefore contradict irreducibility of \(L_B(p^x\lambda)\) on \(Y\).
For an invariant degenerate subspace in a form type, its radical
gives singular or isotropic containment; in characteristic-two
orthogonal type take the kernel of the quadratic form on the radical,
leaving only a nonsingular radical line if that kernel is zero.
It remains to exclude the nondegenerate and line cases.  We may
suppress the overall weight scale \(p^x\) in the following comparisons:
both the extremal restrictions and the multiplicities scale by the
same factor, by Steinberg's theorem.

Use the connected part of each indicated stabilizer on points, with covering groups: the projective finite source is perfect; component quotients for nondegenerate splittings here are soluble (independent isometries on complementary summands, each connected up to at most orthogonal component data; one can work on exact form-preserving lifts). Thus a connected covering subgroup above it would also have to act irreducibly. One can take the reduced identity component over the indicated connected stabilizer using the ambient central covering: it surjects onto that stabilizer (finite surjection with finitely many components and connected irreducible target), by a central isogeny, and projective invariance on its action already tests linear subspace invariance. Ordinary nondegenerate splittings use the product of the two connected classical groups on points (up to central quotients/covers, which do not change their Weyl actions), including an orthogonal two-space torus if needed. For a nondegenerate splitting with both ranks positive, maximal tori of the two parts use all ambient coordinates, except possibly one unused in two odd parts. Among restrictions of ambient orbit weights of \(\lambda\), those with zero on any missing coordinate have full maximal Euclidean norm (coordinate projection and highest-weight convexity), yet are not transitive under the two-part Weyl group, since the sparse nonempty support and spare zeros can be differently allocated. For instance one block has room (in torus rank) for the whole support, and then a nonzero coordinate can also be reassigned to the other. This contradicts irreducibility even on the corresponding connected group, whose simple has all maximal-norm weights in one orbit by the same convexity. In these coordinates one takes hyperbolic pairs on each part separately; if both are odd orthogonal (odd characteristic), combine the two leftover directions as one ambient pair over the closure, unused on restriction. The Euclidean norm in the surviving torus coordinates is Weyl-invariant for the product too, and pulling its torus back isogenously does not change this real coordinate calculation.

\noindent\emph{Line stabilizers.}

For a line in odd orthogonal ambient type B (odd characteristic), restriction of tori to D uses the same coordinates, with even-sum root lattice. If \(k\) is the last positive position in \(\lambda\), sparsity gives \(\lambda_k=a\) with \(0<a<p\). Short root lowering gives \(\lambda-\epsilon_k\) as a weight (pairing \(2a\) nonzero in the odd characteristic), already impossible by parity (the difference after scaling is still odd-sum). For a line in type D, including the nonsingular line in characteristic two, use the standard \(B_{e-1}\) stabilizer embedding, dropping the last torus coordinate \(\epsilon_e\). Here is some detail for characteristic two. In a split hyperbolic basis \(z_i,t_i\) take the nonsingular vector \(u=z_e+t_e\). On points the connected special isometry group's line stabilizer fixes \(u\) and injects on \(u^\perp/\langle u\rangle\) into the symplectic group: kernel elements must fix the whole hyperplane by the quadratic, and then only the nonsingular symmetry (odd displacement rank) could remain besides identity, outside the special group. A connected subgroup \(R\) fixing \(u\) has full image there: take the torus on the first \(e-1\) hyperbolic pairs, their orthogonal root groups, and for each \(i<e\) combine the two end-pair root groups of roots \(\epsilon_i-\epsilon_e,\epsilon_i+\epsilon_e\) with identical parameter \(a_0\) (and similarly in the opposite orientation). Their product sends \(z_e,t_e\) to \(z_e+a_0 z_i,t_e+a_0 z_i\) respectively and \(t_i\) to \(t_i+a_0u+a_0^2 z_i\), fixing the rest. This gives the long symplectic root groups on the quotient by squaring (surjective on points); all other required roots come from the orthogonal ones already used. Take the connected subgroup generated (with Zariski closure). Thus \(R\) has the whole stabilizer's points and dimension of that symplectic group, and is semisimple (a connected solvable normal would have trivial image). Relative to the indicated maximal torus the generators just given have roots \(\pm\epsilon_i,\ \pm\epsilon_i\pm\epsilon_j\) on \(e-1\) coordinates (nonzero differentials also for the combined roots), exhausting the dimension. The displayed torus is maximal by the finite kernel and rank on the symplectic quotient. So this is type \(B_{e-1}\) with the claimed coordinate restriction. We can use its simply connected cover lifting into the ambient cover (central isogeny lifting as in the path proof). In odd characteristic this is the ordinary odd-dimensional special orthogonal stabilizer on the complement up to components/centers.

For the line comparisons these character coordinates on restrictions can be used on the simply connected covers: they come from the vector weights on the indicated surviving pairs. In particular for the special stabilizer just described the characters \(\epsilon_i,\ i<e\), on its torus pull back as the short-root coordinates of type \(B_{e-1}\), and \(\epsilon_e\) becomes zero even on the cover (its vector pair is fixed by this torus). Lifting from a simply connected covering source through an ambient central covering agrees projectively and hence for these linear vector matrices by absence of scalar characters.

Again let \(k\) be the last positive position and put
\(a=\lambda_k\), so \(k\ll e\) and \(0<a<p\).  Root lowering
with pairing \(a\) gives two distinct ambient weights.  On the
line stabilizer they have the same restriction, since
\(\epsilon_e\) restricts to zero:
\[
 \begin{aligned}
 p^x(\lambda-\epsilon_k+\epsilon_e)
   &\longmapsto p^x(\lambda'-\epsilon_k),\\
 p^x(\lambda-\epsilon_k-\epsilon_e)
   &\longmapsto p^x(\lambda'-\epsilon_k).
 \end{aligned}
\]
Here \(\lambda'\) is the restricted highest weight; extend a
positive-chamber cocharacter by zero in the last coordinate to
verify that it is highest.  The two ambient weight spaces are
independent, so this restricted weight has multiplicity at least
two.  We show that an irreducible module of highest weight
\(p^x\lambda'\) has multiplicity at most one there.

After removing the scale, the relevant top Levi-face module has
highest weight \(a\omega_1\) in tail type B, since the difference
down by \(\epsilon_k\) uses just those tail nodes.  The simple
\(L(a\omega_1)\) is a composition factor of the degree-\(a\)
symmetric power of the natural B Weyl module, also in characteristic
two.  The ordinary Weyl character of that natural module has precisely
the signed coordinate pairs and one zero, independently of
characteristic.  Its symmetric power has a unique top weight, giving
the stated constituent, and multiplicity exactly one at
\((a-1)\epsilon_1\): one zero-weight term and \(a-1\) copies of the
first coordinate are necessary.  This bounds the multiplicity in
the simple constituent by one and proves the contradiction.

\end{proof}

\begin{lemma}[Rational lifting and the composed digits]
\label{fb:rational-lift}
The restriction in \Cref{fb:natural-restriction} is the restriction of a
parameter-restricted rational irreducible of the simply connected algebraic
group of \(Y\), of highest weight \(\eta\).  It gives a morphism into the
simply connected algebraic group of \(B\), preserving the required forms.
The rational composition with \(L_B(p^x\lambda)\) is irreducible.  After
removing the common scale \(p^x\), its highest weight \(\mu\) has a
single natural or dual-natural coefficient one at each nonzero digit.
\end{lemma}
\begin{proof}

This natural restriction lifts as a restricted irreducible for the parameter
of \(Y\), with rational highest weight \(\eta\). It gives a morphism
from the simply connected algebraic group of \(Y\) into that of \(B\),
as in \Cref{w:rational-composition}. The invariant bilateral form,
when relevant, is rationally invariant by self-duality and uniqueness.
For the quadratic in characteristic two, if \(\eta\) has multiple
digits, their self-dual factors have alternating forms and hence supply
the canonical quadratic, equal to the given one by uniqueness.
For a single digit, untwist to a \(p\)-restricted weight using the
standard pinned choices; split Frobenius permutes the finite points.
The cocycle for the finite-invariant quadratic then takes values in
the square module of the dual, a nontrivial \(P_Y\)-restricted simple
module because \(P_Y=p^{df}>2\). By \Cref{w:h1-injective} its rational
cohomology class is zero. Thus the cocycle is a coboundary, and that
coboundary is zero because this module has no finite-fixed vectors.
This proves rational invariance.

The composed rational module is irreducible by finite irreducibility: on the finite covering points the lift into the ambient simply connected group of \(B\) differs at most centrally from lifts of the original finite projective inclusion (same projective natural matrices by construction). Thus on \(L_B(p^x\lambda)\) one still obtains projectively the embedding action just specified with its \(d\) atoms. Its highest weight, after removing the scale \(p^x\), say \(\mu\), has only single natural (or dual-natural) coefficients one at its nonzero digits. Indeed every rational digit factor restricted after twist to the finite group is one digit there and by individual tensor invariance must be an aggregate of the \(d\) twisted natural atoms, thus exactly one by restricted-weight uniqueness (an aggregate of two or more has natural/dual digits at that many distinct positions of the finite restricted tuple). Removing the scale is integral by torus restriction.
\end{proof}

\begin{lemma}[Small cuts force endpoint support]
\label{fb:endpoint-cuts}
In this same configuration, \(\lambda\) is supported only at the two
endpoints in type A, or at the first node in a signed type.
\end{lemma}
\begin{proof}

Here is a cut argument to finish. Restrict natural torus weights of \(B\) to get the multiset \(\Omega\) of the module with highest \(\eta\). As in \Cref{w:universal-cuts}, one fixed Weyl translate of \(\lambda\) maximizes throughout the positive chamber for \(Y\), simultaneously on each active fundamental summand. The restriction of each such translated fundamental weight is itself dominant (its orbit maximum, also against \(Y\)-Weyl translates). Indeed Weyl actions on the image torus extend on it via ambient Weyl elements: after transport of the ambient maximal torus by an image normalizer element, conjugate it back in the centralizer of the image torus. Recall the simultaneous maximum exists since at an interior \(Y\)-cocharacter a maximizing orbit weight of the upper module must actually restrict to its unique highest on composition (weight convexity), hence maximizes on the whole positive chamber. Each active fundamental summand is then maximal there since for each cocharacter there is a common ambient Weyl element maximizing all summands. Since their positive integral sum is \(\mu\), each has support only at endpoints for A-type \(Y\), or the first node for signed \(Y\). Each corresponds to a strict upper cut in \(\Omega\) as in that result: all selected elements dominate all unselected and there are no identical vectors across the boundary, by highest multiplicity one (in signed type all active nodes here are early with spare zero coordinates). The sum of the selected weights is that fundamental restriction. Take the small side of size less than \(2d\); for a co-small A-cut this means the negatives of the unselected values, still with that same sum and strict dominance in the dual multiset instead. More explicitly an exchange across a strict coefficient boundary compares natural weights against each other throughout the \(Y\)-chamber. Equal-vector ties across it would produce two distinct ambient orbit weights mapping to the highest vector weight, contrary to multiplicity one. In signed type one similarly flips nonzero coefficient coordinates in checking the maximum (compensating with zero-coefficient coordinates in D), so the positive directed choices lie nonzero in the positive cone and one can make the same strict comparisons in the full natural multiset.

The middle-coordinate Weyl group in type A for \(Y\) (permuting positions other than first and last), or signed Weyl group away from the first coordinate, fixes this sum, thus preserves the selected multiset by strict domination across the cut. By its small size, every selected weight must \emph{individually} be fixed (nontrivial orbits cost at least \(m/3\)). (A permuted selection with the same sum still attains the maximal \(Y\)-chamber interior pairing, hence is the same multiset.) This includes the highest \(\eta\) or its dual, thus supported also only at the indicated endpoints or first node as appropriate. Its active simple reflections are then outside the cut, not fixed individually there. Dominating those forces each selected weight's difference down from highest to be supported only at the active node of any such reflection, which by individual fixedness allows no difference at all: the corresponding simple root (endpoint, or first in signed type) is not fixed by the subgroup either. Thus the cut takes just the highest line. We conclude \(\lambda\) has only endpoint coefficients in A, or just the first coefficient in signed types.

\end{proof}

\begin{lemma}[Multiplicity forces the natural representation]
\label{fb:natural-recognition}
In the tensor-coarsening configuration above, \(\lambda\) is a natural
weight, up to duality.  Consequently \(w=v\), and \(B\) acts naturally
on \(V\), up to Frobenius twist.
\end{lemma}
\begin{proof}

It cannot have both endpoint coefficients \(a,b>0\) (ambient for \(B\) then A), giving \(\mu=a\eta+b\eta^*\). At the lowest nonzero digit exponent \(y\) of \(\eta\), if \(Y\) is type A with just one active endpoint at that digit, \(\mu\) has both active there, impossible. In fact the two coefficients at that position are nonzero modulo \(p\) by \(0<a,b<p\) and the nonzero lone endpoint digit, with no carry into that position from below. Otherwise choose torus coordinates at the full maximum with first two values \(\eta,\eta-p^y\alpha_1\), last two values \(-\eta^*, -\eta^*+p^y\alpha_1\) (last highest-coordinate minimum first in this listing). These distinct weights occur by digit activity (also for signed type with \(\eta^*=\eta\)), and middle coefficients of \(\lambda\) are constant. Here \(\eta+\eta^*\) is not proportional to \(\alpha_1\) by the support already found in high rank, so the four entries are available distinctly at the maximum, rearranging the middle coordinates. Lowering at each of the two endpoint nodes of \(\lambda\) gives two distinct weights colliding at \(\mu-p^y\alpha_1\), up to the overall scale, impossible by simple-root multiplicity one bounds in the simple composed module.

Thus up to duality \(\lambda=z\omega_1,\ \mu=z\eta\). If \(1<z<p\), a coefficient one at any active node of the lowest \(\eta\)-digit would already contradict the required digit shape of \(\mu\). For coefficient at least two there we have three distinct nonzero nonopposite weights \(\eta,\eta-p^y\alpha,\eta-2p^y\alpha\) for that simple root (again the endpoint/first-node shape in high rank ensures nonproportionality to the root). Assign these to the first three natural coordinates of \(B\) at the maximum (in signed types choose signs using also spare zero-coefficient coordinates). The upper \(A_2\)-face for \(\lambda\) gives the two symmetric-power weights lowering twice (one step to the third coordinate, or twice using the second) colliding again along just \(\alpha\). This forces coefficient one for \(\lambda\).

All the collisions just used are comparisons in the rational composed
simple module.  With scales displayed, the first pair of distinct upper
weights restricts to the same weight
\[
             p^x\mu-p^{x+y}\alpha_1,
\]
and the symmetric-power pair in the last case restricts to
\[
             p^x\mu-2p^{x+y}\alpha.
\]
Each lies below the highest weight in a single simple-root direction, where
the divided-power PBW bound gives multiplicity at most one, as in
\Cref{w:universal-cuts}.  Hence the only remaining coefficient is one.
\end{proof}

\begin{proposition}[Singleton tensor failures]
\label{fb:singleton-tensors}
The singleton coatom witnesses in \Cref{fb:protection-count} cannot be
prime tensor-coarsening stabilizers.
\end{proposition}
\begin{proof}

By \Cref{fb:natural-recognition}, \(w=v\) and \(B\) acts naturally
up to twist. Its series matches the ancestor by the previous tests:
in type A the finite \(Y\) digit tuple is non-self-dual, and the
symplectic characteristic-two ancestor case has been removed.
The third coatom is proper and contains the supplementing base, so
\(B<S_i\). By \Cref{w:natural-root}, a proper root of smaller
height fixes \(B\), hence \(Y\).

Fixedness identifies the projective \(Y\)-action with its coefficient
twist through a height \(0<g<f\), possibly with duality, using the
matrix realization of the ambient point maps. These perfect lifts
agree linearly. But the support of the finite restricted tuple is
\[
 \{0,f,\ldots,(d-1)f\}\subseteq\mathbf Z/(df)\mathbf Z.
\]
A shift through \(0<g<f\) moves every one of these positions to a
nonmultiple of \(f\).  A diagram action on wrapping may exchange a natural
weight with its dual, but cannot change which digit positions are nonzero.
Thus no such smaller-height root fixes \(Y\), giving the required
contradiction.
\end{proof}

\subsection{Conclusion of the forward test}
\begin{proposition}[Completion of the forward marker argument]
\label{fb:forward-completion}
The forward marker assertion of \Cref{r:forward-marker} holds.  In the
prepared configuration, the assumption of no point success at all the
prescribed bounded prefixes is impossible once the initial marker set is
sufficiently large.
\end{proposition}
\begin{proof}

Fix the requested number \(m\) of point labels.  Choose the high-type
thresholds first, and let \(D_0\) be the resulting absolute bound in
\Cref{w:path-bound}.  Reserve more than \(D_0\) path nodes.  Choose
the trial omission sizes successively, with the spacing required by
\Cref{fb:protection-count}; each is a finite function of the bounds
already chosen.  This fixes a bound on every path-prefix size.
The earlier exclusion rounds can now have their further-omission
bounds enlarged backwards to cover every bounded union used by a
later round.  All such comparisons are relative to one fixed
ancestor, as in \Cref{fa:all-ancestors}.

Finally choose the initial marker reservoir large enough for those
finite Ramsey requirements, for \(m\) point labels whenever
\Cref{fa:subset-points} succeeds, and for all residual pure links and
missing-site tests.  A terminal spare reservoir, and spaced prefixes
when needed, provide the fresh indices throughout.  Only the bounded
case and comparison predicates already specified are colored; these
choices do not depend on group dimensions or field sizes.

Start with the actual simple \(S\) in the ordinary site interval.
At each subsequent bounded prefix the construction either supplies
the required \(m\) point stabilizers, or continues with an actual
nested simple supplement by \Cref{fb:grounding}.  In the second
case the socles satisfy the transitivity and full-spanning properties
relative to every ancestor, and \Cref{fa:bounded-types} excludes all
bounded types.  A continuing classical run receives both protections
by \Cref{fb:protection-count,fb:singleton-roots,fb:singleton-tensors};
an alternating or otherwise noncontinuing step needs no new run
protection.  Thus absence of point success at the bounded prefixes
would give more than \(D_0\) actual simple nodes satisfying all the
group-theoretic hypotheses of \Cref{w:path-bound}, a contradiction.
The construction must stop with point success at one of those
prefixes.  Its prefix bound and reservoir threshold depend only on
\(m\) and the fixed choices, proving the asserted uniformity.

\end{proof}

\section{Weighted templates and simultaneous size readings}\label{sec:weighted}
A weighted template starts with disjoint finite permutation sets whose
cardinalities satisfy selected numerical relations.  Its defining subgroup
interval is itself finitely representable.  Our task is to recover those
relations from an arbitrary finite representation of a nontrivial lattice
image of the template.  The recovered cardinalities may differ from the
defining ones.  What must survive is their arithmetic, read simultaneously
from the parts of one partition.  The next section will bound the number
of lattice elements independently of the defining cardinalities.

We first pass to an ordinary subgroup interval, allowing either
orientation (\Cref{t:weighted-algebra}).  Normal and projection tests then
produce a forward interval in a nonabelian simple power, with nontrivial
base projections and product subgroups at its distinguished cut coatoms
(\Cref{t:weighted-product-success}).  Applying the forward marker theorem
to a pure-marker filter places the whole remaining template in an actual
alternating power (\Cref{t:weighted-restart}).  The cut intersections
recover a single partition of its permutation domain
(\Cref{t:weighted-alignment}); equal totals and squares are read from
that partition.  Only boundedly many markers are deleted, and the order
of each simple site remains at most the original ordinary top-group order.

Literal lattice calculations throughout use the original parameter order,
even when actual subgroup containment has the reverse order.  The
distinction between subgroups in the simple power and their ordinary
one-site counterparts is fixed in
\Cref{t:weighted-power,t:weighted-site-convention}.

\subsection{The template and its nontrivial images}

\begin{definition}[Units and stock]\label{t:weighted-units}
We use disjoint transitive permutation sets called units; their index set is denoted by \(T\), and
symmetric-group notation on a set of indices (when talking about permutations of letters) uses the
entire union of their units. There are singleton markers. On the units we take independent groups
with product \(Y\), of these kinds:
\begin{itemize}
\item full symmetric groups, including on the markers;
\item primary squares of side \(x\), with group \(A_x\times A_x\) slotwise, where \(8\mid x,\ x\ge16\);
\item auxiliary units of shape \((x/4)\times2\times2\) (same side-parameter restrictions), with \emph{full symmetric groups slotwise}.
\end{itemize}

For each primary side there is its auxiliary type (total size \(x\)), and for each axis length
used there are full reference units of that size. Units are supplied in cohorts of identical type
and size. Also, for each non-marker cohort so required (or basic extra full cohort used), include
a full offset cohort of size per unit one greater than its total size per unit. These extra
offsets themselves need no new offsets. In what follows every cohort, including markers, can be
required to have at least a sufficiently large number of members. For the group-structural tests
this lower bound can be \emph{absolute}, independent of sizes and number of cohorts. We explain
the finite losses involved along with the arguments. Assume there are nonmarker units in this use.
\end{definition}

\begin{definition}[Weighted template]\label{t:weighted-template}
Adjoin a foundation set \(Z\) larger than all the letters together. Use the ordinary interval
\[
       F=[S_Z\times Y,S_{Z\sqcup(\cup T)}].
\]
(Here \(\cup T\) denotes the union of the units.) The notation for its labels is
\[
       (I,U)=S_{Z\sqcup(\cup I)}\times U,\qquad
       I\subseteq T,\quad Y_{T\setminus I}\le U\le S_{T\setminus I}.
\]
The set \(I\) is the \emph{footprint}: its units have been absorbed into the foundation orbit.
\end{definition}

\begin{lemma}[Labels, cuts, and restrictions]\label{t:weighted-labels}
The labels in \Cref{t:weighted-template} are exactly the overgroups in \(F\).
Its coatoms are the labels \(C_D\) below.  Absorption in a join follows residual
orbits.  Axis filters and axis ideals again give the weighted templates on the
units that remain.
\end{lemma}
\begin{proof}
In any overgroup the conjugates of the supported \(S_Z\) generate the full symmetric group
supported on the orbit of \(Z\), by overlap (its size before expansion already exceeds half the
ambient degree). That orbit absorbs entire units. Absorption in joins thus follows orbits; in
particular the join of the \(x_j\)'s below will be top. We shall use the following labels:
\begin{align*}
 X_D&=(D,Y_{T\setminus D}), & x_j&=X_{\{j\}}, & M_D&=X_{T\setminus D},\\
 C_D&=(T\setminus D,S_D)\quad(D\ne\emptyset),& N&=C_T,
 & J^*&=(\emptyset,\prod_{j\in T}S_{\{j\}}).
\end{align*}
The \(C_D\)'s are precisely the coatoms (here and below test by absorption, or by the unequal
pieces of each cut with the large foundation). The filter above \(X_D\) and the ideal below
\(M_D\) have the templates on the remaining units (in the filter the foundation has grown); inside
\(M_D\) the independent unchanged \(Y_D\) factors out.
\end{proof}

\begin{convention}[Images and literal lifts]\label{t:weighted-images}
We will work even with an arbitrary nontrivial lattice image \(\bar F\) of \(F\) (surjective
lattice homomorphism); often the notation for labels refers also to their images. We call the ones
before the map literal labels. Intervals between comparable literal labels map onto the
corresponding intervals (join with the lower then meet with the upper to put a lift there); all
meets and joins may be tested by first doing the literal operation and then mapping.
\end{convention}

\begin{lemma}[Footprints and faithful pure links]\label{t:weighted-footprint}
In every nontrivial lattice image \(\bar F\), the footprint is well defined.
All cut coatoms survive distinctly and are all the coatoms of the image.
Residual orbit discrepancies are detected by absorption.  The pure marker
links, and canonical permutation intervals on copies of any fixed package of
units, are faithful.
\end{lemma}
\begin{proof}
The footprint \(I\) is unambiguous in the image. In fact a collapse with changed footprint gives a
comparable one (use meets), and for \(j\) gained it then identifies \(N\vee x_j=1\) with \(N\), by
meeting first with \(x_j\). Indeed \(x_j\) is under the upper label, while its intersection with
the lower is only residual (empty footprint). Thus the identification kills every \(x_R\) for full
units \(R\) by intersection with \(N\). For a grid unit \(O\), in each direction take a label
preserving a partition of \(O\cup R\) into the parallel lines of that direction plus a reference
block \(R\) of the matching size. Use full independent bases, \emph{alternating top} on the line
blocks including \(R\), and unchanged \(Y\) elsewhere, without absorption. This contains the
bottom: on parallel lines the grid groups induce even permutations (for the auxiliary grids both
other lengths are even). It joins with \(x_R\) to absorb \(O\). Intersecting all these labels with
\(x_O\) gives bottom: preserving each parallelism forces slotwise permutations and on the primary
squares the two alternating conditions enforce both prescribed slot groups. Thus all \(x_O\) are
killed as well, impossible. In particular the cut coatoms survive, distinctly with no new ones,
and orbit discrepancies of residual labels can be detected by joining with absorptions. Here
adjoining \(x_j\) to a label for an unabsorbed \(j\) absorbs exactly its residual orbit in
addition to the old footprint. On just markers, the resulting pure lattice, in particular in the
filter absorbing everything else, is exact: for a strict change of residual group (on singleton
positions) intersect down to a cyclic group generated by a discrepancy element, holding the
footprint fixed. A proper subgroup there has a smaller orbit on some cycle (an index divisor hits
the cycle length), detected by adjoining the absorption of one singleton. The same test works
inside a permutation label acting canonically on copies of a package of units modulo the
prescribed internal bottom actions.
\end{proof}

\begin{proposition}[From algebras to ordinary intervals]\label{t:weighted-algebra}
Every nilpotent join-preserving deflation of \(\bar F\) is zero.
If some nontrivial image of \(F\) is the congruence lattice of a finite algebra
of size \(m\), there is an ordinary finite-group interval representing the
dual of a nontrivial image of \(F\), whose top group has order at most
\(60^m m!\).  The full lattice \(F\) itself has a finite algebra representation.
\end{proposition}
\begin{proof}
Let \(f\) be a nilpotent join-preserving deflation.  It kills every atom:
otherwise it fixes that atom, contrary to nilpotence.  In particular,
\(f(x_R)=0\) for every full unit.  If \(A\) is any of the parallel-line
labels used in \Cref{t:weighted-footprint}, then \(x_O\le A\vee x_R\), so
\[
 f(x_O)\le f(A\vee x_R)=f(A)\le A.
\]
Also \(f(x_O)\le x_O\).  Intersecting over the directions therefore gives
\(f(x_O)=0\).  Since all the \(x_j\)'s join to the top, \(f=0\).
The idempotent-restriction and unary-nonunit argument of
\Cref{f:prop-ordinary-interval} applies starting with a least carrier representation
among all nontrivial images of \(F\). Idempotent restriction to more than one point would still
represent such an image since it gives a surjective lattice map; nonunits thus yield the nilpotent
join deflation as there. Consequently this minimal representation gives an ordinary finite group
interval for the \emph{dual} of a nontrivial \(\bar F\), by the diagonal construction in \Cref{f:eq-diagonal-interval}. Its top order is at most \(60^m m!\) if the representing size was \(m\), taking
the simple group \(A_5\). Also \(F\) itself is representable (the transitive coset action on its
ordinary interval).
\end{proof}

\subsection{Normals: absorption, parallelism, and signs}

We next show that a normal subgroup cannot produce an intermediate label
between the base and the top.  The tests successively exclude absorption
of units, movement between units, disruption of their parallel partitions,
and finally the remaining primary-square signs.  All calculations are
made with literal labels and then passed to the image.

\begin{convention}[Normal endpoint property]\label{t:weighted-normal-convention}
For an ordinary interval \([H,G]\), property \((*)\) means that every normal
subgroup of \(G\) is either contained in \(H\) or supplemented by \(H\),
and that modulo \(\operatorname{core}_G(H)\) there is a unique minimal normal
subgroup, a power of a nonabelian simple group, supplemented by the base.
In the next three results, \([H,G]\) represents \(\bar F\) in either orientation,
and \(e\) is the label of \(HL\), with \(L\trianglelefteq G\).
All inequalities between labels use the parameter order.
\end{convention}

\begin{lemma}[Projection and intersection tests for normals]\label{t:weighted-normal-test}
A normal-multiple label cannot satisfy
\[
 e\wedge b\le a<b\le e\vee a.
\]
Its axis projections are again normal-multiple labels.  If \(L\) is minimal
normal and nontrivial, a nontrivial proper intersection \(L\cap G_b\) cannot
be normalized at a family of labels whose actual groups generate \(G\).
The equations in the proof test precisely the normalizations needed below.
\end{lemma}
\begin{proof}
Write \(e\) for (the label of) a base times a normal, always reasoning with labels in the forward
\(\bar F\) order. Dedekind's law makes it left-modular: in particular \(e\wedge b\le a<b\le e\vee
a\) is impossible in either orientation. Projections to axis filters (by joining in the label
order) again come from normals: in actual order enlarging the bottom to \(K\) projects \(HL_1\) to
\(KL_1\), and shrinking the top to \(K\) projects it to \(HL_1\cap K=H(L_1\cap K)\), for
\(L_1\trianglelefteq G\). Axis intervals used after deleting a bounded number of units are still
nontrivial images on the corresponding templates, by interval surjectivity and footprint
detection. For the normal tests below we suppose that \(e\) is a proper nonendpoint.

Here is also the intersection test used repeatedly below in the normal argument. Write \(G_b\) for
the actual ordinary overgroup at label \(b\), and \(G_e=HL\) for the actual normal multiple,
\(L\trianglelefteq G\) (similarly after any indicated base-core passage). If \(H<G_b<HL\), then
\(I_b=L\cap G_b\) is nontrivial proper in \(L\), by \(G_b=H I_b\). It is normalized at every \(s\)
for which \(\langle G_b,G_s\rangle\cap HL=G_b\), since then it is just \(L\cap\langle
G_b,G_s\rangle\). The sufficient equation in parameter order is thus \(e\wedge (b\vee s)=b\)
forward, \(e\vee(b\wedge s)=b\) reverse. If \(L\) was taken minimal normal nontrivial and the
indicated \(G_s\)'s generate \(G\) (label join forward, label meet reverse), this normalization is
impossible.
\end{proof}

\begin{lemma}[A proper normal label is internal]\label{t:weighted-normal-internal}
If \(e\) is a proper nonendpoint normal-multiple label, then
\(I(e)=\emptyset\) and \(e\le J^*\).
The conclusion remains available after any of the bounded axis restrictions
used in the proof, and when all unit actions are replaced by full symmetric
groups.
\end{lemma}
\begin{proof}
\noindent\textbf{Absorption and the offset comparisons.}
Take four disjoint equal-size unions of units in a cyclic order.
Absorb the outside and allow independent full symmetric groups on
the four blocks.  Adjoining the square of their rotation gives a
label \(a\), and adjoining the whole rotation gives \(b>a\), with
strictness detected by orbits.  Suppose the four counts of absorbed
letters from \(I(e)\) are nonconstant, with a positive count in each
parity class of positions.  An odd rotation cannot preserve these
counts, so \(e\wedge b\le a\).  But the square rotation connects
each parity pair, both of which meet the absorbed part of \(e\);
hence \(e\vee a=1\).  This is forbidden.

In any cohort, therefore, either every unit is absorbed or at most
one is: otherwise put two absorbed copies in opposite parity
positions and a nonabsorbed copy in one of the remaining positions.
Each sufficiently numerous cohort consequently has a majority
absorption bit.  Compare a typical size-\(s\) unit together with
a typical marker against a typical size-\(s+1\) offset unit.
Write their majority bits as \(\epsilon,\mu,\eta\), respectively.
The four-block test, using both block types and repeating a positive
one in both parity classes if their counts differ, forces
\[
 s\epsilon+\mu=(s+1)\eta.
\]
For bits in \(\{0,1\}\), this implies
\(\epsilon=\mu=\eta\).  Every required cohort and its offset
therefore have the marker cohort's bit.  If that bit were one,
the preceding cohort dichotomy would make every unit absorbed,
contrary to \(e<1\).  All majority bits are zero.

Finally, two absorbed units from different cohorts are impossible.
Use four blocks each containing one copy from each of those cohorts,
and put the two absorbed copies in different parity classes, choosing
the other copies nonabsorbed.  The same test applies.  Thus
\(|I(e)|\le1\).  These arguments use only sizes, offsets and
boundedly many copies, so they also hold with full unit actions
and after the indicated bounded axis deletions.

\noindent\textbf{Excluding a singleton footprint.}
If \(I=\{O\}\) and the residual group moves between other units, shift above \(x_j\) at the source
of such a move. At least two absorptions remain (\(O\) and another unit hit from \(j\)), so the
shift is top by the count just proved (the same count in a filter deleting a unit), meaning the
residual group is transitive on the complement. On four copies including \(O\) the square test
then gives total absorption on join anyway (we reach another unit besides \(O\), then its whole
transitive complement orbit), impossible. If instead the residual group preserves units
separately, use three copies including \(O\) with rotation group \(A_3<S_3\) on the blocks,
outside absorbed. This increase is strict (meet with the version allowing just one transposition
on blocks, then detect orbits); intersection with \(e\) cannot move the blocks while join absorbs
them. So \(I=\emptyset\).

\noindent\textbf{Excluding residual mixing.}
By shifting once at a mixed unit and applying the same conclusions,
mixing units now forces transitivity on all the residual letters (there is an extra absorption in
that shifted label, hence it is the filter's top). These nested temporary shifts delete only
boundedly many units before doing the counts.

For such a transitive \(e\), every marker \(x_i\) complements it. Normal complementation gives
isomorphisms in actual subgroup order (when \(K\) complements \(HL\), \(L\trianglelefteq G\),
project \(K\) to \(G/L\) to identify \([H,K]\) with \([HL,G]\)). Forward this forces \(e=N\) by
marker atomicity (it makes \(e\) a residual coatom). Replace the base by the label \(J^*\), i.e.
switch to all full unit groups, with the same transitive projected label. In the new base-core
quotient take a minimal normal inside the normal giving that label; it is still missing and
proper, so still transitive by what was proved thus far (nonmixing would now mean bottom), hence
still maps to \(N\). For a residual marker transposition \(d\) over the new base, we have
\(0<d<N,\ N\wedge(d\vee x_j)=d\) for every unit, using notation relative to the full-unit base.
Thus the intersection of the minimal normal with the subgroup labeled by \(d\) is nontrivial
proper and normalized by all labels \(x_j\) which generate the top, impossible. Indeed inside an
actual base times that normal, all overgroups of the base factor as the base times their own
intersections with the normal; and here enlarging \(d\) by any \(x_j\) leaves its intersection
unchanged.

In reverse orientation the complementation identifies \([x_i,1]\) with \([0,e]\) by meeting with
\(e\) in \(\bar F\). Thus \(X_{\{i,k\}}\wedge e>0\) on distinct markers, forcing every marker
transposition label beneath \(e\) (the literal residual action in that intersection only has the
two singleton markers available over bottom). A supported marker three-cycle group including
\(i\), which lies below \(e\) by these transpositions, then cannot be in the indicated meet image:
joining it with \(x_i\) absorbs all three and they admit transpositions below \(e\). These small
permutation labels over bottom are exact by orbit detection. So in either orientation every proper
nonendpoint normal label lies under \(J^*\).
\end{proof}

\begin{lemma}[Parallel restrictions on an internal normal]\label{t:weighted-normal-parallel}
After quotienting by the base core, suppose that a minimal nontrivial normal
subgroup gives a proper nonendpoint label \(e\).  Let \(B\) be the intersection
of \(J^*\) with all full symmetric parallel-block labels.  Then \(e\le B\).
This case is impossible in reverse orientation.  In forward orientation
\([0,B]\) is the exact lattice of subspaces of the sign space having two
coordinates for each primary square.
\end{lemma}
\begin{proof}
Pass now to the original base-core quotient and, if a proper nonendpoint occurs, take the normal
to be minimal nontrivial therein (within the purported normal), written \(L\). This still gives a
proper nonendpoint internal \(e\) (actual base times \(L\)); internal representatives can be used
since \(e\le J^*\).

For a parallel direction on \(O\), write \(A<W\) for the alternating-top parallel label used
earlier and its full symmetric-top version, now allowing full internal groups on other units. This
increase survives: adjoin a transposition \(d\) on two spare markers to both, then intersect with
a graph label coupling this swap to odd line-block action, giving different marker orbits,
detectable by absorption. Suppose \(e\not\le W\). Some element of an internal representative then
disrupts the partition; conjugating supported full bases and using primitive alternating top
generates independent full action there (a cross-transposition in addition to the within-block
symmetric groups suffices), so \(e\vee A\ge W\). Thus \(D=e\wedge W\not\le A\) by left-modularity.
Forward, \(0<D<e\) in this situation, and
\[
                  e\wedge(x_k\vee D)=D\qquad(k\ne O),
\]
since \(x_k\), like the literal internal intersection giving \(D\), preserves \(O\) and its
indicated lines; on meeting back with the internal representative one thus stays under \(W\).
Hence the nontrivial proper intersection of \(L\) at \(D\) is normalized both at \(W\) (where the
intersection is the same) and all these \(x_k\) (use \(x_k\vee D\)); together they absorb
everything, impossible by minimal normality.

In reverse orientation in fact no internal element of a representative for \(e\) can preserve the
direction's lines on \(O\) and act oddly on them. Otherwise set \(K=e\vee d\) with \(d\) just a
spare marker swap over bottom. This is strictly above \(e\) (orbit test) but residual. Take the
graph label \(P'\) coupling this swap to odd line-top on \(O\), otherwise allowing internal
actions preserving units and this parallelism. We have \(d\wedge P'=0\), and \(e\vee(K\wedge
P')=K\) using the presumed odd element times the swap. Thus the intersection of \(L\) at \(K\),
nontrivial proper by \(e<K<1\), is normalized at \(d\) (by \(K=e\vee d\)) and at \(P'\) (enlarging
actually, i.e. using \(K\wedge P'\), leaves the intersection unchanged). These generate actual top
by \(d\wedge P'=0\). This too is impossible; it rules out in particular \(D\not\le A\) obtained
above under a disruption (an odd top element in the internal intersection with \(W\) fixes the
reference block, hence acts oddly on the lines of \(O\)).

It follows in both cases that \(e\le B\), the intersection of \(J^*\) with all the \(W\)'s. This
is slotwise full on the grids and differs from bottom only by primary square sign bits. Reverse
has now been excluded (meet to place the representative under literal \(B\), then use the odd-line
test).

In the forward case \(0<e\le B\), so the image of the literal sign-space
interval is nontrivial.  The subspace lattice is simple: a nontrivial
congruence collapses an atom by atomisticity, and every two distinct atoms
lie in a plane, where such a collapse propagates to all atoms.  Hence this
interval is exact.
\end{proof}

\begin{proposition}[Normal endpoint property for weighted images]\label{t:weighted-normal}
For sufficiently large absolute cohort stock, every ordinary interval
representing a nontrivial weighted image, in either orientation, satisfies
\((*)\).  Its base acts transitively on the simple factors of its monolith
modulo the base core.  The same assertion holds for the bounded axis deletions
used here, provided the remaining stock is sufficient.
\end{proposition}
\begin{proof}

By \Cref{t:weighted-normal-internal,t:weighted-normal-parallel}, only the
forward sign-space case remains.  Take a minimal missing normal inside any
purported proper normal multiple after passing to the base-core quotient.

Forward, \([0,B]\) must now be an exact subspace lattice over \(\mathbf F_2\), since that literal
interval is simple and its image is nontrivial (\(0<e\le B\)). Here the slotwise groups over
literal bottom have quotient given by the sign space with two bits per primary square. Subspace
lattices here are simple, e.g. by atomisticity and propagation of any atom collapse across the
planes. For an active bit of the sign space of \(e\), meeting with its parallel alternating-top
test \(A\) gives its kernel \(D'\). If this is nonzero, the same formula \(e\wedge(x_k\vee
D')=D'\), \(k\ne O\), and normalization at \(A\) give the earlier forward contradiction. Thus
\(e\) has sign dimension one with vector \(v\ne0\). Swapping the two coordinates of any primary
square by \(\tau\) must fix \(v\): otherwise use as \(a\) the label adjoining \(\tau\) to bottom
and as \(b\) this label extended by the sign line \(\langle v+\tau v\rangle\); \(e\wedge b\le
a<b\le e\vee a\), strictness tested under \(B\). Hence at an active square both bits occur
together.

Choose an active square and use its \(x\) lines in one direction,
together with its size-\(x\) auxiliary unit as one additional block.
Let \(a\) have independent alternating groups inside these blocks and
the alternating group on their \(x+1\) positions.  Match coordinates
along the lines and choose any matching with the auxiliary block.
Outside this union allow full internal groups, without absorption.
The prescribed bottom lies in \(a\).  Let \(b\) instead allow all base
sign vectors of even total, with the same alternating top.  Then
\(a<b\): the base sign which is odd on every original line and even
on the extra block lies in \(b\), and its nontriviality is detected by
meeting with the exact sign interval \([0,B]\).

Write \(\alpha\) for the line-position sign and \(\beta\) for the
repeated within-line sign of the chosen square.  The preceding swap
test gives \(\alpha(v)=\beta(v)=1\).  An element of \(e\wedge b\)
has \(\alpha=0\), since \(b\) has alternating top.  As \(e\) is the
single sign line \(\langle v\rangle\), its sign component must
therefore vanish, including \(\beta\).  The auxiliary bottom action
is also even on its letters: its first-slot permutations repeat four
times, and each binary-slot swap repeats an even number of times.
Consequently \(e\wedge b\le a\).

On the other hand, a representative of the nonzero sign of \(e\)
has base-sign vector
\[
 u=(1,\ldots,1,0)
\]
on these \(x+1\) blocks, and may also permute the original lines.
Conjugate the supported matched alternating top by this representative,
then remove the top permutations using \(a\).  Modulo the alternating
bases this produces \(u-hu\) for every even block permutation \(h\).
These differences span the augmentation space: moving the unique zero
gives every pair of coordinate positions.  The conjugations remain
supported on this union regardless of the representative's actions
on other units.  Hence \(e\vee a\ge b\), contradicting the
normal-multiple test.  This proves the normal endpoint property.

The lattice is nonmodular (the exact pure marker link includes the subgroup lattice of \(S_4\),
already nonmodular: with a 4-cycle \(r\) and a dihedral reflection \(s\) inverting it, \(\langle r
s\rangle\) has trivial meet with \(\langle s,r^2\rangle\) but joins with \(\langle s\rangle\)
above it). For completeness, in the faithful base-coset action every minimal normal is
transitive by the endpoint property.  An abelian minimal normal would identify
the interval with an invariant-subgroup interval in an abelian group, which
is modular.  Two distinct minimal normals commute; supplementation by one
then makes the base induce all inner actions on the other.  Intersections
in the latter are invariant normal subproducts, again giving a modular
interval.  Thus there is a unique nonabelian monolith modulo the base core,
and it is supplemented.  The base is transitive on its simple factors,
because the full group is and the monolith acts internally on each factor.

The stock bounds thus far only spent boundedly many copies for parallel tests, rotations, and
nested shifts deleting at most a bounded number of units.
\end{proof}

\subsection{Product closures in both orientations}

\begin{convention}[Actual groups and product closure]\label{t:weighted-power}
After passing to the base-core quotient we thus have \(G=H\mathcal W\), \(\mathcal W=S^s\) a
minimal normal nonabelian power, and the labels are exact \(H\)-invariant subgroups \(K_b\) in
\([P,\mathcal W]\), where \(P=H\cap\mathcal W\), transitive operator action on the sites. Thus
\(K_b\) is written on the simple-power side (actual subgroups over the interval base correspond by
intersection with \(\mathcal W\)), not a formal permutation label. In reverse, label joins mean
actual intersections, and label bottom means actual \(\mathcal W\). We use \(c\) for coordinate
product closure, and call these its product labels when closed. Recall the extension description
in \Cref{f:lem-extension-description} of subdirect labels (site stabilizer \(D_0\le H\), action image \(\Lambda\le
{\rm Aut}(S)\), extensions to overgroups of \(D_0\) in \(H\)); above a given subdirect group in
actual order all are just restrictions on its domain interval, by reverse correspondence. Also,
intersections of comparable distinct subdirect actual groups with a fixed invariant product must
differ if nontrivial: the restrictions are independent across the tie parts, with nontrivial
projections at every site by operator transitivity, so a strict tie change distinguishes them. And
a subdirect joined (in actual order) with a nontrivial invariant product gives the whole power, by
conjugating independently supported elements and simplicity. These facts and the projection
arguments below can use the same transitive operator action on restricting a filter of actual
groups (or add the enlarged base to the operators by inner actions).
\end{convention}

\begin{proposition}[Forward proper projections]\label{t:weighted-forward-projections}
In a forward realization \([P,\mathcal W]^H\), suppose \(P>1\) and
\(c(P)<\mathcal W\).  Every cut coatom is a product label.  If
\(P<c(P)\), its closure label has empty footprint and lies below \(J^*\).
\end{proposition}
\begin{proof}
Forward, suppose first \(P<c(P)<\mathcal W\), and put \(K_p=c(P)\). Joining a label with \(p\) to
top forces subdirectness, which also suffices since \(K_p\) contains nontrivial independent
support. Thus a test \(p\wedge b\le a<b\), \(p\vee a=1\), is impossible by strictness on
intersections (both \(K_a,K_b\) then subdirect, intersections containing \(P>1\)). Exactly the
same four-block rotation test gives \(|I(p)|\le1\). In the singleton case \(I=\{O\}\), if any
other copy of \(O\) is moved in the residual action to touch a different unit, place it opposite
in parity from \(O\) in the four-copy test, with destination outside the chosen four. The join
again absorbs everything. If none are thus moved use the three-copy test instead. Hence
\(I=\emptyset\).

Shift above \(x_j\); if its closure is not full the closure label has empty extra footprint there
by the same tests (or trivially when \(x_j\) is already product). So if \(p\) mixes \(j\) with
other units then \(x_j\) must project fully by monotonicity of closure, and \(x_j\vee p=1\),
forcing transitivity of the residual group. In that case all marker \(x_i\)'s are subdirect and
their intersections with \(p\) are bottom. Meeting also with \(X_{\{i,k\}}>x_i\) forces all marker
transposition labels beneath \(p\) by strictness. The \(K_{x_i}\)'s have identical ties and
identical twists restricted to the projections of \(P\), by the common intersection \(P\) with
\(K_p\). Therefore the images of all the filters above \(x_i\) under meeting with \(p\) must agree
(ties may be split by any invariant refinement, retaining the twists, then taking the product of
the correspondingly restricted diagonals on intersection with \(K_p\)). Yet a supported marker
three-cycle group \(d'\) lies beneath \(p\) by the transpositions, and occurs in the meet image
for \(i\) outside its three positions, as \(p\wedge(x_i\vee d')\) (\(p\) is residual, \(i\) is
singleton). It cannot occur for \(i\) inside by the same absorption contradiction as before. Thus
\(p\le J^*\). Now every cut coatom is product, since subdirectness there would join to top with
\(p\). The same cut conclusion holds immediately if \(P=c(P)\ne1\).

To elaborate on the common meet-image test in that argument, put \(\mathcal T=c(P)=\prod T_o\)
with the index here ranging over sites. A subdirect \(K_{x_i}\) meeting \(\mathcal T\) in \(P\)
restricts on each tie part to the full twisted diagonal of the corresponding \(T_o\)'s in that
intersection: the projection there is exactly \(T_o\) since it comes from \(P\). These are
independent across parts. Thus whether two sites are tied, and if so the restriction of their
identifying isomorphism to \(T_o\), can be read from the two-site projection of \(P\) (graph of an
isomorphism versus a whole product; \(T_o\ne1\)). Above \(K_{x_i}\) one may split by every
operator-invariant refinement, preserving its twists within the remaining ties; this is again
invariant by equivariance of the original ties, and these are all the invariant overgroups in the
power here. Their intersections with \(\mathcal T\) depend only on that same restricted-twist data
and the refinement. This proves the claimed agreement even if the full twists were not specified
by \(P\). In marker notation with \(i\) outside the triple one can bound \(p\wedge(x_i\vee d')\)
also by \(N\wedge(x_i\vee d')=d'\) calculated literally, and it contains \(d'\) in the image.
\end{proof}

\begin{definition}[Product success]\label{t:weighted-success-definition}
We say that \emph{product success} holds when a forward realization in a
nonabelian simple power has nontrivial projections of its base intersection
at every site and every cut coatom is a product subgroup.  We allow passage
to this configuration by deleting one marker in an axis filter.
\end{definition}

We next establish in \emph{reverse} orientation that for \(P>1\) we must have \(c(P)=\mathcal W\),
and for \(P=1\) all nonbottom actual labels are subdirect. The \(P>1\) argument only needs an
invariant interval with those base and action conditions.

\begin{lemma}[Subdirect copy actions in reverse orientation]\label{t:weighted-copy-action}
In a reverse invariant realization with \(P>1\), let \(z\) canonically
permute \(k\) identical nonempty packages of units, acting as the prescribed
bottom inside each package.  For a sufficiently large absolute \(k\),
\(K_z\) is subdirect.  The same holds for its canonical alternating-copy
label once that label is sufficiently large.
\end{lemma}
\begin{proof}
Take disjoint identical nonempty packages of units and a residual label \(z=Y\rtimes S_k\)
permuting \(k\) copies canonically (fix coordinated matchings of corresponding units respecting
their types and actions). Foundation actions are implicit in residual notation; by a canonical
subgroup label here we mean the inverse image in \(z\) of a subgroup of \(S_k\). The interval
\([0,z]\) in label order is exact \({\rm Sub}(S_k)\), by the cyclic discrepancy and orbit test on
these copies. We claim \(K_z\) subdirect for sufficiently large absolute \(k\), here when \(P>1\).
Otherwise its closure corresponds to a nonidentity subgroup \(U\le S_k\). Labels indexed by
subgroups of \(U\) generated by prime order elements are product, since the corresponding atoms
(actual coatoms above the nontrivial product \(K_U\)) must be so. Hence \(U\) cannot contain a
large supported \(A_v\): that subgroup is prime-generated and would give a product base
\(K_{A_v}\ne1\) for an exact reverse \({\rm Sub}(A_v)\) inside this invariant interval, contrary
to the reverse alternating exclusion of \Cref{r:reverse-exclusion}.

If \(a<c_0\) in \({\rm Sub}(S_k)\) and \(c_0\cap U=1\), their labels are subdirect (by
monotonicity, the label of each closure is below \(U\) and below \(c_0\)), so \(\langle
a,U\rangle<\langle c_0,U\rangle\) by intersection strictness with \(K_U\). Thus any transposition
moved by conjugation by \(u\in U\), together with its conjugate, generates a group meeting \(U\)
nontrivially. We get some prime-order \(h\in U\) supported on at most four letters. Match a
nontrivial orbit of \(h\) to a fresh ordered set of the same size (2 or 3) and exchange them
simultaneously pairwise by an involution. By the very same strict-join test its group together
with its \(h\)-conjugate (dihedral, two such exchanges) meets \(U\) nontrivially. Any nonidentity
element there rotates the two sets nontrivially or exchanges them in matched fashion. By finite
Ramsey homogenize a large reservoir of fresh points so a constant such pattern on increasing
tuples lies in \(U\). Comparing two tuples differing at just the last entry \(y,y'\) gives a
supported three-cycle by taking quotients: the two permutations differ by conjugation by
\((y\,y')\) and the first (say \(w\)) fixes \(y'\) and sends \(y\) to a common anchor (possibly of
the original orbit) distinct from \(y\). Thus one quotient is \((w(y)\,y')(y\,y')\). This anchor
is fixed as a choice of position when the prefix is fixed and the last entry varies. Varying these
cycles gives a large supported alternating group, contradiction. All these Ramsey requirements are
absolute (orbit lengths two or three, bounded many permutation patterns, with sufficiently many
points for the absolute reverse alternating exclusion). This proves the symmetric-copy claim.  Its
alternating-copy subgroup also
has a subdirect actual group: in reverse order that group contains the
subdirect group just obtained.
\end{proof}

\begin{lemma}[Small labels are subdirect in the reverse proper-closure case]\label{t:weighted-small-labels}
Suppose \(P>1\) in reverse orientation and
\(K_p=c(P)<\mathcal W\).  Then \(I(p)=\emptyset\), and every label supported
on any fixed bounded number of units is subdirect.  For the subsequent tests,
a support bound of ten is enough.  The required cohort stock is absolute.
\end{lemma}
\begin{proof}
Suppose now \(K_p=c(P)<\mathcal W\). We have \(I(p)=\emptyset\): for any absorbed \(O\) use such a
canonical permutation label on many copies including it; its alternating and symmetric versions
are distinct subdirect labels, yet have the same joins with \(p\) (all copies are absorbed),
impossible.

Call a label small here when supported on a bounded number of units (literally below their
\(X\)-label); we need only fixed bounds, say up to ten. We show that every such label has a
subdirect actual group. Otherwise an atom \(e_i\) (actual coatom) beneath its nonzero closure
label in parameter order must be product, and \(e_i\le p\): the actual coatom contains the
nontrivial product closure, so cannot itself be subdirect, hence is closed by maximality. It is
residual and lifts to a literal label supported on the package in question at empty footprint.
Take there a literal cover \(B_0<C_0\) of groups giving the images bottom and \(e_i\) (groups
written on just the package). Choose many disjoint clones and write \(e_j\) for the analogous
images of \(C_0\); they too are distinct atoms and the lowers collapse. Indeed image congruences
within residual package intervals transfer both ways by canonical swaps: join with the swap label
over \(Y\), then meet with the full residual symmetric label on the other package (only \(Y\)
elsewhere). Literally this copies a supported group from one package to the other, by independent
support and conjugation. Thus both collapses and noncollapses there transfer, and each supported
interval maps onto its image interval. Distinctness of \(e_j\)'s follows also by their bottom
meets.

\noindent\textbf{Transporting product atoms.} Every \(e_j\) must likewise be product. Indeed a coatom if not product is subdirect. If subdirect here for \(j\ne i\), it contains (in actual order) both intersections of \(K_{e_i}\) with the actual groups at the canonical labels of \(\langle(i\,j)\rangle\) and \(\langle(i\,j\,l)\rangle\) on copies. This is because on the parameter side, joining \(e_i\) with either permutation label contains \(e_j\) by conjugating the literal \(C_0\). Thus its ties refine both sets of permutation-label ties, since the indicated intersections with the product \(K_{e_i}\) project nontrivially and independently across each set of parts. But their common refinement is discrete: the two canonical labels are subdirect with a common subdirect lower group \(K_z\), hence consistent twists, and generate actual top (the canonical parameter subgroups meet trivially); any sites tied in both would keep that common twist on joining. This contradicts proper subdirectness of \(K_{e_j}\). Being product and above \(P\) as an actual subgroup, each is above \(c(P)\), i.e. \(e_j\le p\).

\noindent\textbf{The faithful product interval.} Use \(k\) large such packages. Between the products of the literal covers the interval is Boolean if \(B_0\) is not normal in \(C_0\) (any full projection conjugates the supported \(B_0\) to generate \(C_0\) there); otherwise a subspace lattice over \(\mathbf F_\ell\) for the prime order of \(C_0/B_0\). This interval maps faithfully to \([0,\bigvee e_j]\) by survival or simplicity, and all entries are product (the join of the \(e_j\)'s lies under \(p\), hence all atoms under it too must be coatom products in actual order; now use atomisticity). Take \(k\) coprime to \(\ell\) in the latter case (one of two consecutive large counts suffices). Let \(E\) be the augmentation hyperplane in that case, the whole join in the Boolean case. It is nonzero and no atom beneath it is fixed by the canonical alternating copy action (label \(a\) over \(Y\)), using the faithful literal product interval. For the hyperplane an invariant line would have a fixed nonzero vector by perfection of \(A_k\), hence constant on the copies, impossible by \(\ell\nmid k\).

\noindent\textbf{The extension image.} For the subdirect \(K_a\), write \(D_a\) for its extension domain with map \(\beta\) to \({\rm Aut}(S)\). By restrictions \([D_0,D_a]\) realizes exactly \({\rm Sub}(A_k)\), so the kernel of outer action must be supplemented by \(D_0\) (the normal test in \Cref{sec:reverse}; if retained, the base-core quotient would be soluble by Schreier, impossible there). Thus \(\beta(D_a)=(\beta(D_a)\cap S)\Lambda\), identifying \(S\) with inners. Also
\[
                   (a\vee E)\wedge p=E
\]
by factoring the literal join as a set product (the lift of \(E\) is normalized by the copy
action) and using a representative of \(p\) above that lift, since \(a\wedge p=0\). Here \(a\wedge
p=0\) comes from joining a subdirect with a nontrivial product in actual order, and literally the
set product intersected with that representative of \(p\) is the lift of \(E\) times the
intersection with the lift of \(a\). Hence the product closure of \(K_a\cap K_E\) is \(K_E\) (the
equality gives generation with \(c(P)\), already under the closure by \(P\le K_a\cap K_E\)).
Writing \(T_e={\rm proj}(K_e)\) at the site, this means \(\beta(D_a)\) normalizes \(T_E\). Indeed
intersection with a product has full projection onto it precisely when its factors are matched by
the diagonal twists; by transport along the tie through the site this is precisely invariance of
the site's factor under the extension image.

Here \(T_E\) is nontrivial and proper, by \(P>1\), site transitivity
and the proper product \(K_E<\mathcal W\).  Its normalizer in \(S\)
is proper by simplicity, and is invariant under \(\Lambda\).  Transporting
it gives a product label \(0<f\le E\).  Choose an atom \(h\le f\),
and put \(B_a=\beta(D_a)\cap S\).  The normalization just proved and
the reversal of label order give
\[
 B_a\le N_S(T_E)=T_f\le T_h.
\]
Also \(\Lambda\) normalizes \(T_h\), as it does every invariant
product projection.  Thus
\(\beta(D_a)=B_a\Lambda\) normalizes \(T_h\).  The same
intersection criterion now gives
\[
 c(K_a\cap K_h)=K_h.
\]
The alternating copy action moves the literal lift of \(h\) to a
different atom \(h'\) under \(E\), so \(h'\le a\vee h\) in label
order.  Its actual group is a product and contains \(K_a\cap K_h\);
it must therefore contain the displayed product closure \(K_h\).
This contradicts the distinct actual coatoms \(K_h,K_{h'}\), proving
the small-label claim.

In this extension-image argument, \([D_0,D_a]\) really is the whole ordinary interval: all actual
overgroups of the subdirect in the invariant interval are still subdirect and correspond exactly
to restrictions of its extension map. For the normalization criterion at the tie through a site
\(o\), at \(t o\), \(t\in D_a\), the product factor of \(K_E\) is the transport by \(t\) of
\(T_E\), whereas the entry matched by \(K_a\) to \(u\) at \(o\) is the transport of
\(\beta(t)^{-1}(u)\). Thus all \(u\in T_E\) occur in the intersection exactly for stability under
\(\beta(D_a)\), and similarly for the product at \(h\); other tie parts transport. Finally the
``different atom'' test only conjugates a literal lift within the displayed Boolean/subspace
interval, where alternating copy permutations preserve the lift of \(E\) and move the given atom's
lift to a different atom there. The distinctness survives by faithfulness on this interval, and
the conjugate is contained in the literal join with the copy label, so the conclusion in the image
does not assume invariance of the entire quotient map under copy permutations.
\end{proof}

\begin{lemma}[Residual products are impossible in reverse orientation]\label{t:weighted-residual-obstruction}
Suppose all labels on the bounded supports used below are subdirect in a
reverse realization.  There is no nontrivial proper product actual group
whose parameter label \(p\) has empty footprint.  This assertion applies
both when \(P>1\) and when \(P=1\).
\end{lemma}
\begin{proof}
For small residual \(a<b\) we must have \(a\vee p<b\vee p\) by intersection strictness (actual
intersections nontrivial since joins are residual). A representative for \(p\) cannot move between
units. Otherwise for an element sending a letter of \(j\) into \(k\ne j\), choose a second letter
outside \(k\) mapping outside \(k\) (available by sizes with copies). Use a supported full
symmetric label on the one or two source units. Conjugation of the source-letter transposition in
the join gives a supported transposition connecting \(k\) with somewhere else, adjoining which to
the source label is a small strict increase by orbit detection at \(k\), impossible.

\noindent\textbf{Preserving parallelisms.} On a square both parallelisms, and on an auxiliary grid both length-two parallelisms, must be preserved. Otherwise use the parallel partition plus matching reference block with full symmetric wreath action and only \(Y\) elsewhere as \(a\), and the independent full symmetric action on that union as \(b\). Again disruption generates \(b\) in the join with \(p\); strictness is detected by meeting with the other indicated parallel label of equal length using the \emph{same} reference block, then testing whether the reference unit connects to the grid (it is the only common block of the two partitions).

\noindent\textbf{Uniform auxiliary flips.} Thus on an auxiliary unit \(h_0\times2\times2\) a representative acts by component permutations on \(h_0\), and separate bit flips possibly per component. They must actually be uniform. If for instance the first flip varies, take \(a\) acting on the (component, second bit) positions by a full symmetric group and separately by constant first flip, with only \(Y\) elsewhere. Conjugation of its supported top permutations by the varying element produces (after removing top action) first-flip differences generating all even-sum flips on those positions: differences of the nonconstant flip vector with its translates include weight two and are supported within this unit regardless of actions on other units. Adjoining those to \(a\) strictly increases it. Indeed meet with simultaneous preservation of the two parallelisms: the first flips formerly constant induce with component permutations only even actions on the (component, first bit) lines (second-bit direction, \(h_0\) even), but flipping just one component's first bit now occurs and acts oddly. As before distinguish by adjoining a spare marker swap then intersecting a parity graph label, which contains bottom. This violates strict joins again.

\noindent\textbf{The remaining primary signs.} We conclude \(p\le B\) of the earlier sign-space test, with exact \([0,B]\) since \(p>0\). On any primary slot with active sign take as \(a\) the alternating bases on those lines plus the auxiliary size-\(x\) block, with \emph{symmetric} top, outside only \(Y\). In the join with \(p\), conjugation again adds independent augmentation base signs by discrepancy with the even auxiliary action, producing the strict increase \(b\) tested under \(B\) as before. This final contradiction (or triviality already if no signs) eliminates the proposed residual product.
\end{proof}

\begin{proposition}[Reverse projections and the trivial-base cases]\label{t:weighted-reverse-projections}
In reverse orientation, if \(P>1\), then \(c(P)=\mathcal W\).
If \(P=1\), every nonidentity actual label is subdirect and the site-action
image \(\Lambda\) contains the inner automorphisms of \(S\).
In forward orientation with \(P=1\), either a marker axis filter gives
product success, or the same subdirectness and inner-action conclusions hold.
\end{proposition}
\begin{proof}

For \(P>1\) in reverse orientation, a proper closure would be residual by
\Cref{t:weighted-small-labels} and would contradict
\Cref{t:weighted-residual-obstruction}.  This proves the first assertion.

If \(P=1\), apply the result just obtained in the reversed filter with parameter range
\([0,M_{\{i\}}]\) for a spare marker \(i\) outside any tested small support. The actual base there
is nontrivial with still the same transitive operators, so it projects fully by the reverse
\(P>1\) result and the desired small labels are indeed subdirect. A proper nontrivial product
group would have residual label \(p\), since containing a singleton absorption in parameter order
would put independent nontrivial support inside that small proper subdirect group. The residual
argument rules it out, hence all nonbottom groups are subdirect as claimed. Here \(\Lambda\) must
then contain inners (otherwise its proper nontrivial inner part, or the proper-invariant-subgroup
result for soluble operators in \Cref{sec:affine}, yields a proper nontrivial invariant product).

Back forward, suppose \(P=1\). If some marker \(x_i\) does not project fully, its axis filter
gives product success by the forward argument above (we may take the enlarged actual base as
operators then). Otherwise the extension at each marker \(x_i\) has image containing inners by
atomicity and the same proper-invariant-subgroup argument, now on a diagonal part: a nontrivial
proper subgroup invariant there would transport and multiply to an invariant group strictly
between \(P=1\) and \(K_{x_i}\). Its domain interval via restriction realizes the reverse weighted
test on one less marker. If \(\Lambda\) itself does not contain inners, the kernel of that
extension map is retained by the normal test (supplementation would give image just \(\Lambda\)),
hence this gives an ordinary interval within the image in \({\rm Aut}(S)\), with base \(\Lambda\),
top containing \(S\). Missing inners there must be supplemented, so this is a reverse invariant
interval in one simple site with proper base \(\Lambda\cap S\), impossible by the reverse result
just proved: a nontrivial base cannot project fully here, and for trivial base having all
nonbottom actual groups project fully leaves no interior labels. Thus \(\Lambda\) contains inners
and again all nonbottom labels are subdirect.
\end{proof}

\subsection{Extension domains and the passage to product success}

\begin{theorem}[Product success with absolute stock]\label{t:weighted-product-success}
There are absolute stock and marker-loss bounds with the following
property.  Any ordinary realization, in either orientation, of a nontrivial
weighted image yields a forward realization of a nontrivial image of the
weighted template with at most boundedly many markers deleted, in which
product success holds.  The stock bound is independent of all numerical
unit sizes and of the number of cohorts.  The simple group at a site has
order at most the original top-group order.  A larger prescribed fixed
stock can be retained by increasing the initial absolute stock.
\end{theorem}
\begin{proof}
Take an ordinary realization of minimum top order among both orientations
of the same quotient; its order is at most the input order.  If
\(c(P)=\mathcal W\), the ties of \(P\) are fixed and every actual label
refines them.  On the tie part through a chosen site these refinements are
exactly the invariant equivalences under its part stabilizer in \(H\).
They give an ordinary realization of the opposite interval with top order
at most \(|H|<|G|\), as in the proof of \Cref{f:case-subdirect}.
This contradicts minimality. Thus absent product success we have the artificial-top extension
case in the opposite orientation, with base map containing inners.

\noindent\textbf{The extension datum and its flat labels.}
To apply \Cref{f:lem-lower-grid}, distinguish the new base from the
old operator group.  Write \(H_{\rm op}\) for that operator group,
\(A_0\le H_{\rm op}\) for its site stabilizer, and
\(\alpha:A_0\to\operatorname{Aut}(S)\) for its action on the site.
The remaining case gives \(\alpha(A_0)\ge\operatorname{Inn}(S)\).
Thus the extension datum in that lemma is
\[
 (K,H,B,\psi)=(H_{\rm op},A_0,S,\alpha).
\]
Put \(C_0=\ker\alpha\) and
\(A_{\rm in}=\alpha^{-1}(\operatorname{Inn}(S))\); these are its
base kernel and inner preimage.

Index the points of an affine four-space over \(\mathbf F_3\) by
distinct markers.  For a nonempty flat \(X\) of dimension at most two,
let \(Z_X\) be the extension label given by \(M_X\) in forward
extension order and by \(X_X\) in reverse extension order.  Write
\(D_X\) and \(E_X\) for its domain and kernel.  Restriction identifies
the entire lower ideal of this extension with the ordinary interval
\([A_0,D_X]\).  In parameter order this is the axis ideal
\([0,M_X]\), or the axis filter \([X_X,1]\) with order reversed.
It is a nontrivial weighted image deleting just the markers of \(X\),
so \(D_X>A_0\) and its normal property \((*)\) is available.

Apply that property to \(E_X\trianglelefteq D_X\).  If its kernel
were retained in \(A_0\), the whole preimage of
\(\operatorname{Inn}(S)\) would lie in \(A_0\): an element of that
preimage has the same image as some element of \(A_0\), and their
quotient lies in \(E_X\).  This normal subgroup would lie in the base
core, leaving a soluble quotient by Schreier's theorem, contrary to
the nonabelian monolith in \((*)\).  Hence the kernel supplements
the base.  Each flat extension is pure, with
\[
 D_X=A_0E_X,\qquad E_X\cap A_0=C_0.
\]
Only these flat labels need purity.

\noindent\textbf{Proper and incompatible kernel joins.}
The literal map \(D\mapsto X_D\) preserves Boolean joins and meets,
and footprints preserve its strictness in the image.  Consequently
the \(Z_X\)'s reverse flat containment, and for \(X\cap Y=\{i\}\)
their join in extension order is the point label.  Restricting the
point extension to the generated domain gives
\[
 D_i=\langle D_X,D_Y\rangle
     =A_0\langle E_X,E_Y\rangle,
 \qquad E_i=\langle E_X,E_Y\rangle.
\]
The kernel equality follows because the generated kernels lie in
\(E_i\), have the same intersection \(C_0\) with \(A_0\), and
already generate its domain together with \(A_0\).

Distinct point labels have incompatible join.  For their generated
kernel group \(E\), the equality \(E\cap A_0=C_0\) would define a
common pure extension on \(A_0E\), which is impossible.  Hence
\(E\cap A_0>C_0\).  It is normal in \(A_0\), and its nontrivial
\(\alpha\)-image contains the inner simple subgroup.  Since
\(C_0\le E\), this gives \(A_{\rm in}\le E\), exactly the
incompatible-join condition of the lower-grid lemma.

\noindent\textbf{Subdirect intersections in the point monoliths.}
Apply the projection conclusions to each ordinary point ideal
\([A_0,D_i]\), without a minimality assumption.  If they give product
success, the theorem is already proved on the corresponding bounded
marker restriction.  Otherwise a nontrivial base intersection has
full projection: this follows from the reverse projection result in
reverse orientation and from the absence of forward product success
in forward orientation.  With trivial base intersection, every
nonidentity label is subdirect by
\Cref{t:weighted-reverse-projections}, again allowing its marker-filter
success alternative.

Thus the domains of lines and planes through \(i\) intersect the
point monolith subdirectly modulo the point core.  These intersections
are strictly above the point-base intersection: each such domain is
strictly above \(A_0\), by footprints, and supplementation recovers
the domain from its intersection.  The flat ideals have \((*)\),
the proper kernel joins were just verified, and incompatible point
joins contain \(A_{\rm in}\).  All hypotheses of
\Cref{f:lem-lower-grid} therefore hold, giving the required contradiction.

\noindent\textbf{Stock and order bounds.}
Each test uses boundedly many units: four-block comparisons, packages
of at most ten units, or the fixed affine grid.  The Ramsey reservoir
for the copy-action argument concerns package indices and bounded
permutation patterns, independently of the numbers of letters inside
a package.  Cloning a bounded package therefore needs only an absolute
number of copies from any one cohort, including the two consecutive
copy counts used earlier.  Tests on different cohorts reuse these
stocks; they do not delete the tested units.  The lasting deletions
are the boundedly nested marker restrictions in the extension,
atom-filter and point-ideal steps.  Increasing the initial stock thus
leaves any prescribed fixed number of spare copies.

Only core quotients, subgroups and domain intervals have been used;
the extension domains lie in the preceding operator group.  Hence
every resulting simple site has order at most the initial minimal
top-group order, itself at most the input order.  The reverse
projection step uses \Cref{r:reverse-exclusion}.
\end{proof}

\subsection{A single alternating action for all readings}

\begin{lemma}[Restart in an actual alternating power]\label{t:weighted-restart}
From product success one obtains, after boundedly many further marker
deletions, an actual power \(A_\Omega^s\) with transitive operator action,
base intersection \(H\cap A_\Omega^s\), and product subgroups at every cut.
The degree is high, and at least forty specified surviving marker cuts are
genuine distinct alternating point stabilizers at each site.
All previously product labels below the new top remain products.
\end{lemma}
\begin{proof}
In the forward product success case, take the pure filter absorbing everything except markers (one
can leave just a sufficiently large fixed substock unabsorbed and put many spare markers among the
absorbed units). Its base \(K_{\rm sat}\) inside \(\mathcal W\) is product by intersection of
marker point cuts, nontrivial, and the filter exact pure. Apply the forward marker theorem for,
say, at least 40 distinct points. Thus at \(M_A\) for boundedly many omitted markers we have an
actual high alternating simple \(B_A\) at one site in the ordinary product group \(\Lambda
T_{M_A}\) (\(T_b={\rm proj}(K_b)\)), normal and supplemented by the pure base there. Here
\(\Lambda\cap S\le T_{\rm sat}\) as in that theorem; thus \(B_A\le T_{M_A}\), and multiplying its
transported factors gives a normal group at actual \(HK_{M_A}\), not contained in the pure base.
In the whole ideal below \(M_A\), removing these boundedly many markers, (\(*\)) shows that
already \(H\) supplements this normal group.  Explicitly, normality of
the local \(B_A\) under \(\Lambda\) makes its transport independent of
the chosen element of \(H\).  The resulting actual subgroup
\(N=\prod B_A\) lies in the old power, is normalized by
\(K_{M_A}\) and by \(H\), and is not contained in
\(HK_{\rm sat}\), hence not in \(H\).  The whole-ideal normal test
gives \(HN=HK_{M_A}\); the original transitive action of \(H\) on
the sites permutes these very factors.
Restart therefore with its power of alternating
groups \(B_A^s\), transitive operator action, and base intersection \(H\cap B_A^s\); labels that
were products below \(M_A\) remain product by intersection. In particular all relative cuts are
product (intersect original cuts with \(M_A\)). The specified genuine point stabilizers still hold
by intersection with \(B_A\). Below use notation on the remaining template (removing \(A\)) and
this new power.

If success was obtained by an earlier axis-filter shift deleting a marker, take \(H\) to be the
ordinary base of the resulting whole weighted template (thus enlarged at that prior shift). It can
remain the operator group in testing the pure filter \([K_{\rm sat},\mathcal W]^H\); that filter's
\(M_A\) is exactly the lift \(M_A\) in the whole template. In the transport just used, \(K_{M_A}\)
itself is a product by the marker-cut intersections, \(\Lambda\) normalizes the local \(B_A\), and
\(T_{M_A}\) normalizes it as well. Its power is not in the pure saturated base since \(M_A\)
strictly increases saturation and the local simple supplements it on ordinary product labels.
After supplementation by \(H\), intersection with this new power at each label is just the
intersection of the old \(K_b\) with it (the old power contained the new one). For \(M_{A\cup i}\)
from the pure theorem, local intersection with \(B_A\) can be taken on the ordinary site label as
well, since \(\Lambda\cap S\le T_{\rm sat}\); the almost simple quotient's centralizer kernel was
in that ordinary base. Thus it gives exactly the alternating point stabilizer inside the new site.
This step uses the forward marker result and its high-type reductions in
\Cref{sec:forward-completion}.
\end{proof}

\begin{convention}[Ordinary site groups and groups in the power]\label{t:weighted-site-convention}
Write \(B_A=A_\Omega\); locally let \(\Theta\) be the site action image, in \(S_\Omega\) since the
degree is high. The projection of the new base intersection lies in \(\Theta\), and all cut site
groups are nontrivial proper normalized by \(\Theta\). They contain \(\Theta\cap A_\Omega\):
adjoining it there could not reach \(A_\Omega\) by simplicity, hence could not increase by the
product coatom property. We may thus look also at the ordinary cut groups multiplying by
\(\Theta\), inside \(\Theta A_\Omega\). This respects intersections (the cut site groups all
contain \(\Theta\cap A_\Omega\)); covers between such products remain covers as ordinary site
groups since any strict intermediate, intersected with \(A_\Omega\), would transport to a strict
intermediate product.

In particular, throughout the readings \(K_b\) on this restarted template denotes an actual
intersection inside the normal power (before multiplying by any ordinary base). We use the
ordinary cut groups with \(\Theta\) mainly to test permutation actions; when using a projection or
a graph in the power keep the parity restrictions in the alternating groups. A proposed
intermediate \emph{product} obtained sitewise is a label if it contains the base intersection and
is invariant (transport to the other sites is implicit); arbitrary intermediate labels need not
themselves be products.
\end{convention}

\begin{lemma}[All cuts are full nonhalf subset stabilizers]\label{t:weighted-cuts}
In the alternating realization of \Cref{t:weighted-restart}, every cut
coatom has site group the alternating intersection of a full subset
stabilizer on \(\Omega\).  The two parts of this subset cut have unequal
sizes.  The ordinary site groups obtained by multiplication by \(\Theta\)
respect all cut intersections and the indicated covers.
\end{lemma}
\begin{proof}
Every cut coatom \(C_D\) (remaining indices) must give a full nonhalf subset stabilizer on
\(\Omega\). Indeed use at least eight genuine distinct point markers all on one side of the formal
cut, and up to six of them at a time as a prefix \(P_1\). The cut group intersected with this full
pointwise stabilizer (axis \(M_{P_1}\), by meeting the individual full point stabilizers) is
maximal proper there including in ordinary site notation with \(\Theta\). In labels it is a
relative cut coatom after deleting the prefix, and these are products; any ordinary intermediate
would transport to a product. Also the intersection with a single such point stabilizer is covered
by the cut itself, literally by the independent full groups on the two pieces in the template
(strictness as well by footprints or orbit detection, and use the cover transfer between the
indicated products as ordinary site groups). Thus if the ordinary cut group were transitive, it
would be primitive by point stabilizer maximality. Here the tested genuine points are fixed
individually by \(\Theta\) (it normalizes their respective alternating point stabilizers), so
prefix stabilizers indeed cut down inside a pointwise ambient alternating or symmetric group.
Every tested prefix stabilizer in it would have to be transitive on the remaining letters:
otherwise by its maximality in the full pointwise ambient group it would be a full subset
stabilizer there and contain a three-cycle fixing the prefix points, impossible in the original
assumed transitive primitive cut group by Jordan's Theorem (its intersection with \(A_\Omega\) is
proper). Indeed on deleting the prefix the residual cut is still proper with a nonempty unabsorbed
side by our spare markers; for a single chosen marker the upward cover back to the cut is just
restoring one singleton to a full symmetric action on its piece. Strictness if that piece is
unabsorbed follows by absorption on adjoining that marker. Taking successive nested prefixes gives
at least six-fold transitivity, impossible here by the classification of multiply transitive
groups. Thus the cut is intransitive and by maximality a full subset stabilizer (its alternating
intersection is exactly the product site group); a half subset would extend properly to the
transitive two-block stabilizer instead.
The permutation-group inputs in this argument are Jordan's three-cycle
Theorem and the classification of multiply transitive groups, in the
high-degree range, as recorded in \Cref{sec:affine}; see also
\citet[Theorem~1.1 and Section~2]{Jones2014} for the prime-cycle
statement and the exclusion of proper six-transitive groups.
\end{proof}

\begin{proposition}[Common alignment of all cuts]\label{t:weighted-alignment}
There is a single bijection from the formal units together with the
foundation position to the \(\Theta\)-orbits on \(\Omega\), under which
every cut and every intersection of cuts is the corresponding stabilizer
of unions of those orbits.  Write \(u_j>0\) for the size of the cell of unit
\(j\).  The specified genuine point markers have size one.
\end{proposition}
\begin{proof}
\noindent\textbf{Every orbit bipartition is a cut.}
The preceding lemma assigns a bipartition into unions of
\(\Theta\)-orbits to every cut.  Conversely, take any nontrivial such
bipartition.  Its individual-part stabilizer in \(A_\Omega\)
transports to an invariant product containing the base.  A coatom
above this product is a subset stabilizer by
\Cref{t:weighted-cuts}.  It must preserve the same bipartition, since
the original stabilizer together with \(\Theta\) is transitive on
each of its two parts.  Thus the product was already a coatom.
This also excludes a half-sized bipartition.  We have obtained a
bijection between all formal cuts and all bipartitions of the orbit list.

\noindent\textbf{Intersections recover the partition poset.}
At least one \(\Theta\)-orbit is nonsingleton.  Otherwise an even
degree admits a half-sized bipartition.  In odd degree, \(\Theta=1\)
and the base intersection is trivial; a full cycle then transports
to a proper transitive product, whereas every coatom above it would
be an intransitive subset stabilizer, also impossible.

For any partition of the orbit list, intersect the cuts separating
its parts.  At a site this gives the full alternating intersection
of the independent symmetric groups on their unions.  Multiplying
by \(\Theta\) leaves exactly these unions as its orbits.  The only
possible parity exception would be a part of size two with every
other part a singleton: its transposition has no outside parity
correction.  By the nonsingleton-orbit conclusion, that two-point
part is already one \(\Theta\)-orbit, so multiplication by
\(\Theta\) restores transitivity there too.  Hence distinct
partitions give distinct intersection labels, with their correct
containment order.

On the template side, intersections of cuts likewise form the
partition lattice on the units and the foundation position.  A
strict change of partition changes either the absorbed set or a
residual orbit, so footprint and orbit detection preserve strictness
in the image.  The cut-intersection poset therefore identifies these
two partition lattices.

\noindent\textbf{Recovering individual cells.}
A partition lattice on sufficiently many positions has only
position permutations as automorphisms: its merging atoms represent
pairs, their joins detect intersecting pairs, and the large incident
stars recover the individual positions.  The poset identification
thus gives one bijection of formal positions with actual orbit cells,
compatible with every cut and every intersection.  Denote the cell
of unit \(j\) by \(\Omega_j\), and put \(u_j=|\Omega_j|>0\).
Transport these cells to the other simple sites by the operator group.
This is well-defined because the site stabilizer \(\Theta\) preserves
each of its orbit cells; invariance of the cut labels gives the same
formal alignment at every site.
For a specified genuine point marker, the formal singleton cut is
an actual one-point stabilizer.  Its one-cell side must consequently
correspond to that single point, so this marker has size one.
\end{proof}

\subsection{Equality and square readings in that alignment}

The following two readings use the same alignment of \Cref{t:weighted-alignment}.

\begin{proposition}[Equal totals and equal copies]\label{t:weighted-equal-totals}
Let \(D,E\) be disjoint bounded collections of nonfoundation units with
equal formal total size, with ample spare markers outside them.  If at least
one of \(D,E\) contains more than one unit, then
\[
 \sum_{j\in D}u_j=\sum_{j\in E}u_j.
\]
In a sufficiently numerous cohort, all copies have the same actual cell
size.  In particular, all surviving markers have size one.
\end{proposition}
\begin{proof}
Put \(J=C_D\wedge C_E\) and \(F'=C_{D\cup E}\).  Equal formal
totals supply a two-block wreath label \(W\) between them.  It
survives strictly above \(J\): joining an absorption on one side
detects that \(W\) connects the two sides.  To distinguish \(W\)
from \(F'\), intersect with a cut fixing a proper nonempty
subcollection of the multi-unit side.  This prevents the wreath
exchange, while the upper label still mixes letters across the two
sides.  Thus \(J<W<F'\) in the image.

Suppose the actual totals differ.  At one site write \(J_o\) and
\(F_o\) for the projections of the product groups at \(J,F'\).
Projection onto the combined \(D,E\)-cells maps \(F_o\) onto their
full symmetric group: spare units provide outside parity correction.
Its kernel lies in \(J_o\), whose image is the stabilizer of two
unequal nonempty parts.  That stabilizer is maximal and nonnormal,
as is seen by conjugating its supported transpositions.  Therefore
\(J_o<F_o\) is a cover and \(J_o\) is not normal in \(F_o\).

The actual group \(K_W\) need not be a product.  Its coordinate
closure is nevertheless a product label between \(J\) and \(F'\),
so the site cover forces that closure to be \(K_{F'}\).  Now
\(K_W\) contains the independently supported copies of \(J_o\)
and projects onto \(F_o\) at every site.  Conjugating those copies
inside \(K_W\) therefore puts their independent normal closures
in \(F_o\) inside \(K_W\).  Nonnormality and the cover give
\(\langle J_o^{F_o}\rangle=F_o\).  Hence
\(K_{F'}\le K_W\), contradicting \(W<F'\).  The actual totals
must agree.

Apply this conclusion to disjoint pairs of copies in one cohort.
Comparing two pairs which differ at one position against a common
disjoint pair equalizes the individual cell sizes.  This works for
the marker cohort as well; its already identified genuine points
make every surviving marker size one.
\end{proof}

\begin{theorem}[Reading a primary square]\label{t:weighted-square}
For a primary square \(O\) and its full side reference \(R\), let
\(r=u_R\) in the alignment of \Cref{t:weighted-alignment}.  If \(r\) is even
and \(r\ge8\), then \(u_O=r^2\).
\end{theorem}
\begin{proof}
We first construct the two actual parallel systems, then determine their
common incidence cells and the number of blocks.

\noindent\textbf{The two parallel systems.}
Put
\[
 M=M_{\{O,R\}},\qquad F'=C_{\{O,R\}},\qquad
 J=C_{\{O\}}\wedge C_{\{R\}}.
\]
For each of the two formal parallel directions, let \(V_j\) be
the full wreath label whose blocks are the parallel lines and the
reference unit \(R\), with everything outside absorbed.  Each is
maximal below \(F'\): adjoining a permutation disrupting its blocks
and conjugating the supported base transpositions generates the full
symmetric group on their union.  The label \(V_j\) mixes \(R\)
with \(O\), whereas \(V_1\wedge V_2\le J\), since \(R\) is
the only common block of the two formal partitions.  Orbit detection
preserves these distinctions in the image.  In particular both
\(V_j\)'s remain proper and maximal below \(F'\).

First consider the product \(N_R\) of the independently supported
alternating groups on the actual cells \(\Omega_R\) at all sites.
It is normal in the ordinary group at \(C_{\{R\}}\).  Since the
formal reference unit is full, \(C_{\{R\}}=M_{\{R\}}\); the ideal
below it is a weighted image with just that reference unit deleted.
Its normal property \((*)\) applies.  But \(N_R\) lies inside the
separate-cell stabilizer at \(J<C_{\{R\}}\), so it cannot supplement
the base.  It must lie in the base, and hence is available independently
at every site in each \(K_{V_j}\).

Fix a direction \(j\).  The projection of \(K_{V_j}\), together
with \(\Theta\), is transitive on \(\Omega_O\sqcup\Omega_R\):
\(\Theta\) is transitive on each cell, and \(V_j\) mixes them.
Suppose two translates of \(\Omega_R\) overlap without coinciding.
After conjugation one of these translates is \(\Omega_R\) itself.
The component of the overlap relation containing \(\Omega_R\)
then meets \(\Omega_O\) and is \(\Theta\)-invariant.  Cell
transitivity makes it the entire union.  Alternating groups on
intersecting supports of size at least eight generate the alternating
group on their union.  Conjugating the supported factors of \(N_R\)
along the translates therefore gives, inside \(K_{V_j}\), the
independently supported alternating groups on
\(\Omega_O\sqcup\Omega_R\).  Conjugation by a site projection is
realized by an element of \(K_{V_j}\); conjugation by \(\Theta\)
uses operator invariance.  Both preserve single-site support.

Let \(N\) be the product of these larger supported groups.  It is
normal in the ordinary group at \(F'\), since the cell union is
invariant.  In the ordinary interval from the group at \(M\) to
the group at \(F'\), let \(e\) label the lower endpoint multiplied
by \(N\).  We have
\[
 M\le e\le V_j,\qquad e\not\le J.
\]
For the other direction \(k\), set \(a=J\wedge V_k\) and
\(b=V_k\).  Then \(a<b\) and \(e\wedge b\le a\).

The remaining join calculation takes place in the literal template.
Lift \(e\) between \(M\) and \(V_j\).  Because its image is not
below \(J\), this lift moves the block \(R\) to some \(j\)-line.
The literal label \(a\) contains the full supported symmetric
groups on \(R\) and on every \(k\)-line.  Conjugating its group on
\(R\) supplies the full symmetric group on that \(j\)-line.
Its crossings with the \(k\)-lines generate the full symmetric
group on \(O\), and the movement of \(R\) then generates the full
group on \(O\cup R\).  Hence \(e\vee a=F'\).  This is the
forbidden normal-multiple configuration, a contradiction.

Thus the translates of \(\Omega_R\) are disjoint or equal and
form an equivariant block partition at every site.  Its full product
stabilizer is proper, contains \(K_{V_j}\), and lies below \(K_{F'}\).
Maximality forces equality with \(K_{V_j}\).  Removing its block
\(\Omega_R\) gives a partition of \(\Omega_O\) into blocks of
size \(r\).  Doing this for both directions yields two systems,
each with \(m'=u_O/r\) blocks.

\noindent\textbf{The sign graph across all simple sites.}
Set \(a_j=V_j\wedge J\) and \(b_0=a_1\wedge a_2\).  These are full product
subgroups at the simultaneous structures.  Let \(L\) be the projection of
\(K_{b_0}\) at a site onto the actual \(O\)-cell alone.  Outside parity
correction makes this the full common stabilizer of its two systems.
We show that \(L\le A_{u_O}\).

Choose two singleton markers outside \(O,R\).  From \(b_0\) obtain \(b^-\)
by fixing the markers individually and \(b^+\) by fixing their pair setwise;
from \(J\) obtain \(j^+\) by fixing their pair setwise.  Each restriction is
an intersection with the corresponding cut labels.  Formally the group on
\(O\) at \(b_0\) is \(S_x\times S_x\) in its grid action, and every such
permutation is even because \(x\) is even.  The literal graph label \(g\)
inside \(j^+\), which couples the marker swap to the sign on \(O\), therefore
satisfies
\begin{equation}\label{t:weighted-parity-equalities}
 g\vee b^+=j^+,\qquad g\wedge b^+=b^-.
\end{equation}
It does not fix the markers individually: this survives the lattice map
by orbit detection.  This is a valid literal label before taking the
image, since the bottom action on \(O\) is even; actions on the other parts
remain independent.

Now use the actual subgroups in \(A_\Omega^s\).  Project to the \(O\)-cell
and the marker pair at every site, correcting parity on \(R\).  The three
product subgroups \(K_{b^-},K_{b^+},K_{j^+}\) have a full common kernel,
say \(Q_0\), and their quotients are respectively
\[
 L^s\times0,\qquad L^s\times C_2^s,
 \qquad S_{u_O}^s\times C_2^s.
\]
The restrictions distinguishing them concern only \(O\) and the marker
pair, so the same full kernel occurs in all three.  By the meet equality,
\(Q_0\le K_{b^-}\le K_g\).  Put \(Q=K_g/Q_0\).  The same equality gives
\[
 Q\cap(L^s\times C_2^s)=L^s\times0.
\]
In particular, the projection of \(Q\) to \(S_{u_O}^s\) is injective and its
image contains \(L^s\).  The join equality in
\Cref{t:weighted-parity-equalities}, taken with
\(L^s\times C_2^s\), forces that image to be all of \(S_{u_O}^s\).
Indeed, if the image is \(P\), then \(L^s\le P\), so the projected
join is \(\langle P,L^s\rangle=P\); it is the entire product,
not merely a subgroup surjecting to each site separately.
Consequently \(Q\) is the graph of a homomorphism
\[
 f:S_{u_O}^s\longrightarrow C_2^s,
 \qquad f(L^s)=0.
\]
No assumption is made that \(f\) acts separately at the different sites.
It is nonzero: if it vanished, \(K_g\) would lie in the product restriction
of \(K_{j^+}\) fixing the two actual points individually, the intersection
with their individual cut labels.  The literal graph connects the two
formal singleton units whereas that intersection does not, with the same
outside absorption, contradicting orbit detection.

Every homomorphism to an abelian group factors through the sign
abelianization of \(S_{u_O}^s\).  If \(L\) contained an odd permutation,
then the signs of \(L^s\) would fill \(C_2^s\), and \(f(L^s)=0\) would force
\(f=0\).  Hence \(L\) is even.  Transport identifies the actual copies of
\(L\) without changing parity.

The common incidence cells of the two systems on \(O\) now have at most one
letter each: otherwise a transposition in one such cell would belong to
\(L\).  Thus the letters of the \(O\)-cell are edges of a simple
\(r\)-regular bipartite graph, with \(m'\) vertices in each part.  The group
\(\Theta\) is transitive on edges, because it is transitive on the
\(O\)-cell.  In particular, \(m'\ge r\).

\noindent\textbf{The row chain and the column signs.}
The literal interval \([b_0,a_1]\) has exactly four elements.  On
the formal \(O\)-grid its lower endpoint contains the diagonal
base \(\Delta S_x\) and all pure row permutations \(S_x\), while
the upper endpoint is \(S_x\wr S_x\).  An intermediate is therefore
determined by its base kernel.  Intersecting this kernel with
\(A_x^x\) gives a subdirect group containing \(\Delta A_x\).
Its tie partition is invariant under all row permutations, so either
all rows are tied or all are independent.  In the first case the
intersection is \(\Delta A_x\), whose normalizer in the base is
\(\Delta S_x\), since \(A_x\) has trivial centralizer in \(S_x\).
In the second case the kernel contains \(A_x^x\).  Its quotient is
an invariant sign space containing the constant vector.  Taking
differences under transpositions shows that the only possibilities
are the constant line, augmentation, and the full sign space.
Thus the four base kernels are
\[
 \Delta S_x
 <A_x^x\Delta S_x
 <(S_x^x)_{\rm even}
 <S_x^x.
\]
Here \(x\) is even, so the constant line lies in augmentation.
The actions outside \(O\), including on \(R\), are independent
and full at both endpoints.  Passing to the lattice image leaves
at most these four elements.

In the actual realization consider the following chain on the
\(O\)-cell, where even refers to total base parity:
\[
 L<L A_r^{m'}<L(S_r^{m'})_{\rm even}<L S_r^{m'}\le S_r\wr S_{m'}.
\]
All groups in this chain lift with the full common kernel; parity can be
corrected freely outside the \(O\)-cell.  They contain the base intersection
and are \(\Theta\)-invariant, because \(\Theta\) preserves both systems and
the displayed base layers are canonical within the row system.  Thus they
are actual product labels in \([b_0,a_1]\).

The first increase is strict: a nontrivial three-cycle supported on one row
cannot preserve the column partition, since the incidences are simple and
every column has degree \(r>1\).  For the next increase, examine the kernel
of the row action of \(L\).  It can only permute columns with the same
neighborhood, called twins.  A permutation inside one twin class has the
same parity contribution on every row of that neighborhood, and contributes
zero on other rows.  Its row-sign image therefore has rank at most the
number of twin classes.  Edge transitivity implies column transitivity,
so all twin classes have the same size \(t\).  If there are no twins the
rank is zero; otherwise \(t\ge2\), and it is at most
\[
 \frac{m'}t\le\frac{m'}2<m'-1.
\]
Intersecting the displayed groups with the row base and comparing their
sign images modulo \(A_r^{m'}\) proves that the augmentation increase is
strict.  The final increase is strict as well: a row permutation has sign
\((\operatorname{sgn}\sigma)^r=1\) on letters because \(r\) is even; since
\(L\) is even, \(L(S_r^{m'})_{\rm even}\) is even, whereas
\(LS_r^{m'}\) contains an odd row-base element.

There are now four distinct labels in an interval which is an image of a
four-element chain.  They exhaust it, so
\(LS_r^{m'}=S_r\wr S_{m'}\).  In particular the row action of \(L\) is
\(S_{m'}\).  The orbit of any column neighborhood under this row action
contains every \(r\)-subset of the rows, whereas there are only \(m'\)
columns.  Hence
\[
 \binom{m'}r\le m'.
\]
As \(m'\ge r\ge8\), this forces \(r=m'\) or \(r=m'-1\).
In the latter case the missing incidences form a perfect matching:
every row and every column misses exactly one partner.  Label each
column by the row it misses.  The letters then identify with ordered
pairs \((i,j)\) of distinct indices, and every element of \(L\)
acts by the same permutation on both indices.  The full row action
\(S_{m'}\) supplies an element inducing a transposition \((a\ b)\).
On letters it exchanges \((a,b)\) with \((b,a)\), and for each
other index \(c\) it exchanges both \((a,c)\) with \((b,c)\) and
\((c,a)\) with \((c,b)\).  This is a product of
\(1+2(m'-2)=2m'-3\) transpositions, hence odd, contrary to
\(L\le A_{u_O}\).  Thus \(m'=r\), and the \(O\)-cell has
\(m'r=r^2\) letters.
\end{proof}

\Cref{fig:weighted-square} illustrates the incidence conclusion just proved.

\begin{theorem}[Weighted realization and simultaneous numerical readings]\label{t:weighted-realization}
Fix any bounded support size for the comparisons and any desired fixed
number of spare copies.  There are absolute initial cohort-stock and
marker-loss bounds, independent of the numerical unit sizes and the number
of cohorts, with the following property.  The bounds may depend on the
fixed support and spare-copy requirements.  An ordinary realization of a
nontrivial weighted image, in either orientation and with top order \(M\),
gives a forward realization after at most that many marker deletions, with
all cuts product in an actual alternating power.  A single common cell
alignment gives positive sizes \(u_j\le M\), equal on each cohort and one
on all surviving markers.  In this same alignment the equal-total reading
of \Cref{t:weighted-equal-totals} and the square reading of
\Cref{t:weighted-square} hold.

If the input is a finite algebra of size \(m\) representing a nontrivial
image, the same conclusion holds with \(M=60^m m!\), for a possibly
different nontrivial image of the original template.  The full template
itself is finitely representable.
\end{theorem}
\begin{proof}
Apply \Cref{t:weighted-algebra,t:weighted-product-success,t:weighted-restart},
then \Cref{t:weighted-alignment,t:weighted-equal-totals,t:weighted-square}.
All sizes read here are bounded by the order bound already given, since the high alternating
occurs inside an initial site. Stock bounds, including losses to \(A\) and marker requests, remain
independent of the numerical sizes read.
In particular \(u_j\le |\Omega|\le |A_\Omega|\le M\).  All comparisons are
made in this one realization; none requires a new choice of simple sites
or a new alignment.  The original stock can include the fixed losses and
the fixed spare copies required at the end.
\end{proof}

\begin{figure}[htbp]
\centering
\begin{minipage}{0.77\linewidth}
\centering
\begin{tikzpicture}[x=0.5cm,y=0.5cm,
  every node/.style={font=\small},
  letter/.style={circle,fill=black,inner sep=1.05pt}]
  \fill[black!9] (0,5) rectangle (8,6);
  \fill[black!9] (2,0) rectangle (3,8);
  \draw[step=1,black!25,thin] (0,0) grid (8,8);
  \draw (0,0) rectangle (8,8);
  \draw[thick] (0,5) rectangle (8,6);
  \draw[thick,densely dashed] (2,0) rectangle (3,8);
  \foreach \i in {0,...,7}{
    \foreach \j in {0,...,7}{
      \node[letter] at (\i+0.5,\j+0.5) {};
    }
  }
  \node[anchor=east,align=right] at (-0.6,5.5)
    {one row\\\(8\) letters};
  \draw[-{Stealth[length=3pt]}] (-0.5,5.5)--(-0.08,5.5);
  \node[align=center] at (2.5,9.15)
    {one column\\\(8\) letters};
  \draw[-{Stealth[length=3pt]}] (2.5,8.55)--(2.5,8.08);
  \node[align=center] at (4,-0.85)
    {actual \(O\)-cell\\\(u_O=8^2\)};

  \draw (10.4,0) rectangle (11.4,8);
  \foreach \j in {1,...,7}{
    \draw[black!25,thin] (10.4,\j)--(11.4,\j);
  }
  \foreach \j in {0,...,7}{
    \node[letter] at (10.9,\j+0.5) {};
  }
  \node[align=center] at (10.9,9.15)
    {reference\\cell \(R\)};
  \node[align=center] at (10.9,-0.85)
    {\(u_R=8\)};
\end{tikzpicture}
\caption{The established conclusion of \Cref{t:weighted-square},
illustrated for actual size \(r=8\) at one simple site.  In the common
cell alignment, the \(O\)-cell has \(r\) row blocks and \(r\) column
blocks, each with \(r\) letters; each intersection contains one letter
(a dot).  The separate reference cell \(R\) contains \(r\) letters.}
\label{fig:weighted-square}
\end{minipage}
\end{figure}

\FloatBarrier

\section{A uniform computable bound for weighted templates}
\label{sec:counting}

The letter sizes in a weighted template can be arbitrarily large.
We will bound its number of lattice elements using only the number of
units.  By \Cref{t:weighted-labels}, it suffices to count permutation
overgroups of the prescribed actions on each residual union of units.
Such an overgroup may have several orbits.  On each orbit we need both a
bound for the number of transitive overgroups and a bound for their
quotients by the normal closures of the prescribed actions.  The latter
bound controls how the different orbit actions can be coupled in one
subgroup.

We begin with products of natural alternating actions, allowing one
factor to act on several copies of its alphabet.  Compression in bounded
coordinate groups will let us lift a count from a block action to the
full permutation action.  We then include the private square actions
present in the weighted templates and assemble the residual orbit
actions.  All bounds are computable functions of the stated parameters;
none depends on the numerical weights of the units.

\begin{definition}[Controls]\label{n:controls}
A \emph{rank-one control} on a finite set $I$ is a group $X$ that is a
product of independent alternating groups $A_{n_i}$, with $n_i\ge7$.
Each factor acts naturally on one or more disjoint copies of its alphabet,
using the same permutation on all its copies, and fixes all other letters.
Additional fixed letters are allowed, as is the trivial group.  Its
\emph{complexity} is the number of its orbits on $I$.

A \emph{control with private squares} is the disjoint combination of a
rank-one control and square orbits, each with its own two independent
natural alternating factors.  The factors of a square act on no other
orbit.  The two alphabets of each square have the same size, at least
seven.  Its complexity again counts all orbits, including fixed letters.

A \emph{trimmed control} fixes a specified bounded number of entries in
each alphabet and retains the alternating action on the remaining entries.
If fewer than seven entries remain, the whole remaining alphabet is
declared fixed.
\end{definition}

Fixing at most $k$ entries per alphabet changes complexity by a computable
function of $k$ and the original complexity.  Indeed, each natural copy
splits into the fixed entries and at most one remaining natural orbit;
an alphabet discarded at this step has size at most $k+6$.
A private square splits into products of these alphabet orbits, so it
contributes at most $(k+6)^2$ orbits after trimming.
There is also a computable function $q_r(s)$ bounding the number of
orbits on $I^r$ for either kind of control of complexity at most $s$.
To see this, specify the orbit containing each tuple entry and then the
equality patterns of the alphabet entries in each factor.  There are at
most $2s$ factors and at most $2r$ alphabet entries to consider.
Passing from the corresponding symmetric groups to alternating groups
splits an orbit into at most $2^{2s}$ orbits.  For example, with $B_j$ the
$j$th Bell number, one may use
\begin{equation}\label{n:tuple-bound}
 q_r(s)=2^{2s}(s+1)^r B_{2r}^{\,2s}.
\end{equation}
Consequently the number of $X$-invariant $r$-ary relations is at most
$2^{q_r(s)}$.  These estimates apply to equivariant images of invariant
subsets as well: relations on such images pull back to relations on the
original tuple sets.

\begin{lemma}[Compression in a bounded base group]\label{n:compression}
Let $X$ be a rank-one control of complexity at most $s$ on $I$, and let
$D$ be a finite group of order at most $a$.  There are computable bounds,
depending only on $(s,a)$, for the following quantities:
\begin{enumerate}[label=\textup{(\roman*)}]
\item the number of $X$-invariant subgroups $J\le D^I$;
\item the number of $X$-invariant right or left cosets of each such $J$;
\item the number of tuples needed to generate $J$ under the action of
      $X$, each of which can be chosen fixed by a trimmed control of
      bounded complexity.
\end{enumerate}
Every invariant coset has a representative fixed by such a trimmed
control.  No commutativity assumption on $D$ is required.
\end{lemma}

\begin{proof}
For one factor $A_n$, combine the coordinates in its repeated natural
copies into an entry in a product $E=D^c$, where $c\le s$ is the number
of copies.  Thus the coordinates on this alphabet form a list
$(x_1,\ldots,x_n)$ in $E^n$.
Suppose first that $x$ belongs to an invariant left coset $xJ$; the case
$x\in J$ is included.  If $\sigma$ is a three-cycle of alphabet entries,
then $x^{-1}x^\sigma\in J$.  This tuple is supported on the three moved
entries, and an appropriate choice of $\sigma$ gives any prescribed
difference $x_i^{-1}x_j$ at one of them.

Let $z\in J$ be such a tuple, and let its value at a chosen entry be
$t\in E$.  Choose two fresh entries outside its support and conjugate
$z$ by a three-cycle that moves the chosen entry through those two
positions, fixing all other support entries.  If $z'$ is the resulting
tuple, then $z(z')^{-1}$ is supported on a pair and has values $t,t^{-1}$
there.  The other entries cancel coordinatewise.  Since $n\ge7$, the
fresh entries exist; since $A_n$ is two-transitive, such a pair can be
moved to any ordered pair of distinct positions.

Keep position $2$ as a reference.  For each $i\ne1,2$, use a pair tuple
in $J$ with value $x_i^{-1}x_2$ at $i$ and its inverse at position $1$.
Right multiplication changes entry $i$ to $x_2$ and accumulates the
correction at position $1$.  These are valid coordinatewise products
even when $E$ is nonabelian.  After all corrections, every entry except
possibly the first is the common reference value.  Do this independently
on every alphabet.  The resulting representative is constant beyond
the first two positions of every alphabet and is fixed by the control
fixing those positions.

Let $\mathcal C$ be the set of all tuples in $D^I$ constant beyond the
first two entries on each alphabet.  The fixed coordinates and the three
values per natural copy give, for example, $|\mathcal C|\le a^{3s}$.
Every pair multiplier used above is an $X$-translate of an element of
$J\cap\mathcal C$.  Therefore $J$ is generated by the $X$-translates of
$J\cap\mathcal C$.  This bounds its generating list and the number of
possible $J$'s, for instance by $2^{|\mathcal C|}$.  Every invariant
left coset meets $\mathcal C$, so there are at most $|\mathcal C|$
such cosets for fixed $J$.  Inversion gives the corresponding assertion
for right cosets.  The required trimming bounds were established after
\Cref{n:controls}.
\end{proof}

\begin{lemma}[Lifting across bounded fibers]\label{n:lifting}
Let $X$ be a rank-one control of complexity at most $s$ on $I$, let $D$
have bounded order, and let $U'\le S_I$ contain $X$.  Suppose $U'$ is
generated over $X$ by a bounded
list of elements $g$ such that each $X\cap gXg^{-1}$ contains a rank-one
control of bounded complexity.  Then the number of subgroups
\[
 V\le D^I\rtimes U',\qquad \operatorname{proj}(V)=U',
 \qquad (1,X)\le V,
\]
is bounded computably in $s$ and these parameters.  If
$[U':\langle X^{U'}\rangle]$ is bounded, then so is
$[V:\langle(1,X)^V\rangle]$.

In the permutation action with bounded fibers, each such $V$ is generated
over its pure copy of $X$ by boundedly many elements whose intersections
with that pure control contain rank-one controls of bounded complexity
on the expanded point set.
\end{lemma}

\begin{proof}
The base kernel $J=V\cap D^I$ is $X$-invariant, so its possibilities are
bounded by \Cref{n:compression}.  For a prescribed generator $g$, write
the fiber of the projection as $(Jf,g)$.  If
$h\in X\cap gXg^{-1}$, multiply a lift on the left by $(1,h)$ and on
the right by the matching pure lift of $g^{-1}h^{-1}g$.  The resulting
element still lies above $g$, and its base entry is $h\cdot f$.
Thus $Jf$ is invariant under the specified intersection control.
\Cref{n:compression}, applied to that control, bounds its possible
cosets and supplies a representative $f$ fixed by a trimmed subcontrol.
The kernel, these lifts, and $(1,X)$ generate $V$, proving the count.

Put $N=\langle(1,X)^V\rangle$.  In $V/N$, all $X$-translates of a
kernel generator have the same image.  Hence the image of $J$ is already
generated by the images of a bounded list of tuples, with no translates
required.  A subgroup of $D^I$ generated by $k$ tuples has order at most
\begin{equation}\label{n:base-order}
 |D|^{\,|D|^k}.
\end{equation}
Indeed, each coordinate has a type in $D^k$, and coordinates of the same
type evaluate every word in the generators identically.  Projection to
one coordinate of each type embeds the subgroup into at most $|D|^k$
copies of $D$.  The quotient of $V/N$ by the image of $J$ is
$U'/\langle X^{U'}\rangle$, which proves the index bound.

For the last assertion, let $h$ belong to the trimmed intersection
control fixing the chosen $f$.  Conjugating its pure lift by $(f,g)$
takes it back into pure $X$: the base term disappears because $h$ fixes
$f$, and $g^{-1}hg\in X$.  The compressed kernel generators likewise
commute with their trimmed pure controls.  A natural action repeated
over bounded fibers is still a rank-one control, with complexity
multiplied by a bounded factor.  These observations give all the
asserted intersections on the expanded set and remain valid after a
bounded number of such lifts.
\end{proof}

We use the following consequence of the primitive order theorem with an
explicit effective threshold.  Mar\'oti's Theorem
\citep[Theorem~1.1]{Maroti2002} says that a primitive group of degree $d$
is either one of the four Mathieu exceptions, has order at most
$d^{1+\log_2d}$, or lies between
\[
 A_\ell^e\quad\hbox{and}\quad S_\ell\wr S_e
 \quad\hbox{on degree}\quad d=\binom{\ell}{h}^{e}.
\]
Replacing subsets by their complements permits $h\le\ell/2$.

\begin{lemma}[Large primitive parts]\label{n:primitive-parts}
For each positive integer $s$ there is a computable threshold $H(s)$
with the following properties for $d>H(s)$.  A primitive group of degree
$d$ and order at least $\lfloor\sqrt{d/s}\rfloor!/2$ has one of the
alternating-product actions
\begin{equation}\label{n:three-actions}
 (e,h)=(1,1),\qquad(1,2),\qquad(2,1),
\end{equation}
with $\ell$ arbitrarily large as the threshold is increased.  Its socle
is $A_\ell^e$ and has index at most eight.  If its order is at least
$\lfloor d/s\rfloor!/2$, it contains $A_d$ in the natural action.
One may take
\begin{equation}\label{n:threshold}
 H(s)=2^{\,2^{(s+100)^{10}}}.
\end{equation}
\end{lemma}

\begin{proof}
Write $p=\log_2d$.  Outside the bounded Mathieu exceptions, the small-order
alternative has logarithm at most $p(1+p)$.  In an alternating-product
action, $e\le p$.  Unless one of \Cref{n:three-actions} holds,
$\ell\le4d^{1/3}$: if $e\ge3$, use $\binom\ell h\ge\ell$; if
$e=2,h\ge2$, use $\binom\ell h\ge\binom\ell2$; and if $e=1,h\ge3$,
use $\binom\ell h\ge\binom\ell3$.  Consequently the logarithm of
$(\ell!)^e e!$ is at most $10p^2d^{1/3}$.  For the two additional
possibilities $(1,2)$ and $(2,1)$, the same estimate with $d^{1/2}$ in
place of $d^{1/3}$ is an upper bound.

Under \Cref{n:threshold}, we have $s<p$ and
$2^{p/6}>100p^4$.  Compare the preceding upper bounds with
\[
 \log_2(v!/2)\ge v-2,
 \qquad
 v=\lfloor\sqrt{d/s}\rfloor
 \quad\hbox{or}\quad v=\lfloor d/s\rfloor,
\]
respectively.  The square-root factorial excludes the small-order
alternative and all cases outside \Cref{n:three-actions}; the linear
factorial excludes the other two actions as well.  The socles in the
three remaining cases are respectively $A_\ell$, $A_\ell$, and
$A_\ell^2$.  Their full indicated overgroups have quotient orders two,
two, and eight.  In the last case primitivity forces transitivity on
the two slots: fixing the slots would preserve the partition into
lines in either direction.  The centralizer of a natural alternating
group, and of its action on two-subsets, is trivial; a commuting
permutation is determined by the image of one point, whose stabilizer
fixes no other point in these actions.  In the product action,
commuting with each independently supported factor likewise fixes
each coordinate.  Finally, the intersection of a normal subgroup with the
alternating product is a subproduct of its simple factors, and slot
transitivity forces every nontrivial such intersection to be
the whole product.  A normal subgroup with trivial intersection
centralizes that product and is therefore trivial.  This proves the
assertions about the socle and unique minimal normal subgroup.
\end{proof}

The next overgroup argument needs bijections between equal-sized parts
that commute with one common subcontrol.  For a rank-one control fixing
each part setwise, its \emph{profile} on a part records the number of
fixed letters and the number of natural copies of each factor of that
control.  Equal profiles give equivariant bijections by matching copies
of the same factor with their alphabet coordinates, and matching the
fixed letters.  Equal cardinalities alone need not give equal profiles;
the following lemma supplies a suitable subcontrol.

\begin{lemma}[Effective equal profiles]\label{n:profiles}
Suppose a finite set is partitioned into at most $s$ equal-sized parts
$B$, each preserved by a rank-one control $X$.  Write the size of $B$ as
\[
 b_B+\sum_i k_{Bi}n_i,
\]
where $b_B$ counts its fixed letters, $k_{Bi}$ counts its natural copies
of the factor $A_{n_i}$ of $X$, and the number of factors and all
nonnegative integers $b_B,k_{Bi}$ are at most $s$.  There is a rank-one
subcontrol of $X$, of complexity bounded computably in $s$, that has
identical profiles on all parts.  If the common part size is at least
two, the subcontrol may also be chosen with at least two fixed letters
in every part.
\end{lemma}

\begin{proof}
The empty set is immediate, so suppose there is at least one part.
Subtract the expression for one part from each other expression.  This
gives a system $An=c$ with bounded integer coefficients and bounded
numbers of equations and variables.  No equation prescribes the
common total size.  Every nonnegative integer solution admits a
decomposition
\begin{equation}\label{n:profile-decomposition}
 n=u+\sum_{w\in\mathcal H}m_w w,\qquad m_w\in\mathbb Z_{\ge0},
\end{equation}
where $u$ belongs to the finite set of coordinatewise minimal
nonnegative solutions of $Au=c$, and $\mathcal H$ is the finite set
of minimal nonzero nonnegative solutions of $Aw=0$.
The finiteness is Dickson's Lemma \citep{Dickson1913}.  Given $n$, choose
such a $u\le n$; then $n-u$ is homogeneous.  Unless it is zero, subtract
a minimal nonzero homogeneous solution lying below it and repeat.
The sum of its coordinates strictly decreases, proving the decomposition.

These finite lists can be found effectively.  Enumerate solutions in
larger and larger boxes and test whether their upward orthants cover
all nonnegative solutions, using the first-order theory of addition
and order \citep{Presburger1929}.  The test has only fixed integer
coefficients: for a proposed particular list it asks whether
\[
 \forall v\in\mathbb Z_{\ge0}^r\quad
 (Av=c\ \Longrightarrow\ \bigvee_u u\le v),
\]
and the homogeneous test adds $v\ne0$.  Dickson's Lemma guarantees
termination; discarding nonminimal listed entries gives the desired
lists.  There are only boundedly many possible coefficient systems,
so the lengths and coordinates of all these lists have computable
uniform bounds.

For each original alphabet $i$, fix $u_i$ entries and split the rest
into $w_i$ sets of size $m_w$ for each $w\in\mathcal H$.
For each $w$, use one common natural alternating factor $A_{m_w}$
diagonally on all these sets, across all original alphabets and repeated
copies.  Thus the profiles on different parts refer to the same actual
factors, even when two factors have equal degrees.  If $m_w<7$, fix
those entries instead.  The number of copies on part $B$ for the new factor
is $\sum_i k_{Bi}w_i$, independent of $B$ because $Aw=0$.
The remaining fixed-letter count is also independent of $B$, by
$Au=c$.  This proves identical profiles and the complexity bound.
Suppose the common part size is at least two.  The number of fixed
letters is already the same in every part.  If it is less than two,
some new factor $A_m$, with $m\ge7$, acts on a part.  By the identical
profiles it has the same positive number of natural copies on every
part.  Fix two entries of its alphabet, simultaneously on all these
copies.  If $m\ge9$, retain $A_{m-2}$ on the remaining entries.  If
$m=7$ or $8$, declare all entries of these copies fixed, as required
by the trimming convention.  This adds at most eight fixed letters
per copy, and the number of copies is computably bounded by the
preceding construction.  The new profiles are still identical, and
every part now has at least two fixed letters.  If no factor acts,
all letters were already fixed, so the same conclusion follows from
the common part size.  The trivial control is allowed; discarding
the last active factor therefore causes no exception.  This proves
the fixed-letter assertion with computably bounded complexity.
\end{proof}

\begin{proposition}[Rank-one overgroups]\label{n:rank-one}
There are computable functions of $s$ that bound, for every rank-one
control $X$ of complexity at most $s$ on a finite nonempty set $I$,
\begin{enumerate}[label=\textup{(\roman*)}]
\item the number of transitive groups $U$ with $X\le U\le S_I$ and
      the index $[U:\langle X^U\rangle]$;
\item the length of a generating list for $U$ over $X$, chosen so
      that every generator $g$ has a rank-one control of computably
      bounded complexity inside $X\cap gXg^{-1}$.
\end{enumerate}
\end{proposition}

\begin{proof}
We prove a more precise assertion with a second parameter $h$ bounding
the number of strict steps in any chain of $U$-invariant equivalences.
Every such equivalence is a union of $X$-orbits on $I^2$, so initially
$h\le q_2(s)$ suffices.  We induct on $h$, allowing the complexity to
change by computable functions at a step.  If $|I|=1$, the assertion
is immediate.  Otherwise choose a minimal nondiscrete $U$-invariant
partition $\mathcal B$; the universal partition is allowed.  Its
parts have a common size $d\ge2$, and a part stabilizer acts
primitively on its part.  There are at most $2^{q_2(s)}$ possible
partitions.  Let $I'$ be the set of parts.

\emph{A natural copy maps injectively to $I'$.}
The map from an $X$-natural copy to $I'$ is either constant or
injective, since a natural alternating action is primitive.
If a copy for one factor maps injectively, all letters mapping to
that moving orbit of parts come from injective copies of the same
factor.  Indeed a copy on which this factor is trivial cannot have
an equivariant image in its nontrivial natural orbit.  Each
contributing copy supplies only one letter to a part, whence $d\le s$.

Retain only the natural factors of degree greater than $d+2$ to form
a smaller control $X_0$.  Discarded factors have bounded degree,
so $X_0$ has bounded complexity.  It induces a rank-one control
$X_0'$ on $I'$, whose complexity is no larger than that of $X_0$.
A part stabilizer in $X_0$ is a product of alternating
groups or natural point stabilizers.  Each nontrivial factor is
simple and has order greater than $d!$, so its action on a fiber of
size $d$ is trivial.  Label one fiber in each $X_0$-orbit of parts
and transport the labeling by $X_0$.  The trivial stabilizer action
makes this well-defined.  This labels $X_0$ as a pure control in a
bounded-fiber wreath product, and the labeling can be fixed from
$X_0$ and $\mathcal B$, independently of $U$.

The projected transitive group $U'$ has equivalence-chain bound at
most $h-1$: lift a chain on $I'$ and append the discrete partition
of $I$ at its fine end.  Apply the induction hypothesis to
$(X_0',U')$, then \Cref{n:lifting} with fiber group $S_d$.
This bounds the choices of $U$, its index modulo the normal closure
of $X_0$, and its generators with the required intersection controls
on $I$.  They also prove the assertions for $X\ge X_0$: its normal
closure is larger, generation over it needs no additional elements,
and the intersection controls remain contained in
$X\cap gXg^{-1}$.

\emph{Every natural copy maps constantly to $I'$.}
Now $X$ fixes every part, so their number $b$ is at most $s$.
If $d\le H(s)$, then $|I|\le sH(s)$ is computably bounded, and
finite subset counts bound every subgroup choice and generator
list.  The trivial subgroup is an admissible intersection control
of bounded complexity on this bounded domain.  Thus assume $d>H(s)$.
Each part contains at most $s$ natural copies or fixed letters,
so some natural axis of degree at least $d/s$ occurs in it.
\Cref{n:primitive-parts} implies that its primitive part-stabilizer
action contains $A_d$.

Let $K$ be the kernel of the action on $I'$ and put
\[
 T_0=K\cap\prod_{B\in\mathcal B}A_B.
\]
The projection of $K$ to each part is normal in its primitive
part-stabilizer group and contains the nontrivial axis action.
It therefore contains $A_B$.  The commutator subgroup of $K$ lies
in $T_0$ and still projects onto every $A_B$.  Thus $T_0$ is a
subdirect product of these nonabelian simple groups, hence a
product of full twisted diagonals over a partition of the parts.
It is normal in $U$.

For completeness, such a subdirect product is self-normalizing
inside the direct product of its simple factors.  On each diagonal
class, an element normalizing the diagonal has coordinate
conjugations that agree through all the tying isomorphisms.  Its
coordinates therefore differ from a member of that diagonal by
elements of the centers, which are trivial.  Consequently the
whole block-stabilizer normalizer satisfies
\begin{equation}\label{n:rank-one-normalizer}
 \bigl[N_{\,S_d\wr S_b}(T_0):T_0\bigr]\le2^b b!.
\end{equation}
In high degree, automorphisms of $A_d$ are induced uniquely by
permutations of its natural alphabet.  Record each twisted tie
by matching the corresponding letters across its tied parts.
The resulting relation on $I^2$ is $X$-invariant, since $U$
normalizes $T_0$.  It determines the ties and their twists, so
there are boundedly many choices of $T_0$.  For fixed $T_0$,
\Cref{n:rank-one-normalizer} bounds all choices of $U$ by subset
counts in its quotient.

Here $1<X\le T_0$: $X$ fixes each part and its natural alternating
actions are even on it.  Since $U$ is transitive on parts, it is
transitive on the diagonal simple factors of $T_0$.  A normal
subgroup of their product is a subproduct.  Thus
$\langle X^U\rangle$ contains all of $T_0$, giving the index
bound in (i).

To obtain (ii), first generate $T_0$ by supported three-cycles in
its diagonal factors.  Such a cycle is specified by at most $3b$
letters, recording its simultaneous support in all tied parts.
There are boundedly many $X$-orbits of these specifications by
\Cref{n:tuple-bound}.  Take one representative from each orbit.
Each commutes with the control trimmed to fix its support, and
these representatives generate $T_0$ over $X$.

It remains to choose suitable representatives for $U/T_0$.
Write the size of each part as
$b_B+\sum_i k_{Bi}n_i$, with fixed-letter counts and copy
multiplicities for $X$.  Their coefficients are bounded by $s$.
\Cref{n:profiles} gives a subcontrol $X_1\le X$ of bounded
complexity, with identical profiles and at least two fixed
letters on every part.  Fix a coset in $U/T_0$.  By multiplying
its representative by $T_0$, the map from one representative
part in each diagonal class to its destination may be changed
to any bijection of the prescribed parity.  The maps on its
other tied parts then correspond through the tying
identifications, since the representative transports the
diagonals.  Those identifications are $X$-equivariant, because
$X\le T_0$.  The identical $X_1$-profiles allow the chosen
bijections to be $X_1$-equivariant; a swap of two fixed letters
corrects their parity if necessary.  The resulting coset
representative centralizes $X_1$, so
$X_1\le X\cap gXg^{-1}$.  There are boundedly many cosets by
\Cref{n:rank-one-normalizer}.

This completes the induction.  All coefficient lists used for
profiles are drawn from bounded ranges, and all bounded-domain
branches use finite subset counts.  The only recursive call
strictly decreases $h$, even when the complexity grows after
small axes are discarded.  Hence the bounds proved are computable
functions of $(s,h)$, and substituting $h=q_2(s)$ proves the
Proposition.
\end{proof}

\begin{proposition}[Overgroups of controls with private squares]
\label{n:private-squares}
There are computable functions $B(s),C(s)$ such that every control
$X$ with private squares of complexity at most $s$ on a finite
nonempty set $I$ has at most $B(s)$ transitive overgroups in $S_I$,
and every such overgroup $U$ satisfies
\[
 [U:\langle X^U\rangle]\le C(s).
\]
The functions can be taken nondecreasing and at least one.
\end{proposition}

\begin{proof}
We count three successive choices: the block partition, the local
primitive structures and their diagonal ties, and the overgroup in the
resulting wreath product with bounded local groups.  The first two
choices will be encoded by
invariant relations on tuples of bounded length.

Again choose a minimal nondiscrete $U$-block partition $\mathcal B$,
with part size $d\ge2$ and part set $I'$.  The case $|I|=1$ is
immediate, and the choices of $\mathcal B$ are bounded by
\Cref{n:tuple-bound}.

The restriction of this partition to a square orbit is either
constant, a partition into lines in one direction, or discrete.
Indeed the point stabilizer in the product of its two alternating
factors is the product of two maximal nonnormal natural point
stabilizers.  If an overgroup of that stabilizer projects larger
in an axis, conjugating the independently present point
stabilizer by a mover generates that entire alternating factor
independently.  These are exactly the four resulting block
possibilities.  The discrete possibility cannot occur here:
privacy means that no other $X$-orbit can map equivariantly onto
the moving orbit of parts of this square, so a part meeting it
would have size one, contrary to $d>1$.

It follows that the image $X'$ on $I'$ is rank-one.  Each moving
part orbit is identified equivariantly with the natural alphabet
of its contributing axis.  Distinct such orbits for the same
axis are diagonal natural copies; the factor is trivial on
every other part orbit, and the active factors remain independent.
There are at most $s$ part orbits.  Moreover, an injective
rank-one copy on $I'$ can only share its parts with copies of
the same factor, one letter per copy.  Other orbits, including
all the private squares, have trivial action of that factor.
Thus an injective rank-one copy still forces $d\le s$.

If $d$ is bounded in terms of $s$, every private square has side
at most $d$, since one of its axis lines lies in a part.  Discard
its factors and declare its letters fixed.  There are boundedly
many such letters, and what remains is a rank-one control of
bounded complexity on $I$.  \Cref{n:rank-one} bounds the desired
number and index, also for the original larger $X$.

Assume henceforth that $d>H(s)$, increasing the displayed
threshold if needed only within a computable bound.  In a part
$B$, each original orbit contributes a whole orbit, an axis
line, or a singleton.  One contribution has size at least
$d/s$.  Therefore some axis of degree at least
$\sqrt{d/s}$ acts nontrivially inside $B$.  It lies in the
kernel of the action on $I'$: square axes are private, and an
injective rank-one occurrence of its factor anywhere would
force $d\le s$.

\paragraph{The local socle and its normal closure.}
By \Cref{n:primitive-parts}, the socle $S_B$ of the primitive
part-stabilizer group is $A_\ell^e$ in one of
\Cref{n:three-actions}.  It is the unique minimal normal
subgroup, has trivial centralizer, and its quotient has order
at most eight.  Let $K$ be the part kernel.  Its projection to
$B$ is a nontrivial normal subgroup of the primitive group,
so contains $S_B$.  Set
\[
 T_0=K\cap\prod_{B\in\mathcal B}S_B.
\]
This intersection projects onto every $S_B$.  To verify this
without a dependence on the number of parts, take the
$8!$-th powers of all elements of $K$; these lie in $T_0$.
In an individual simple site, every $11$-cycle is an $8!$-th
power, since $11\nmid8!$.  Such cycles generate the large
alternating group.  Since the projection of $K$ contains the
whole local socle, these powers project to a generating set
of every individual site, hence onto $S_B$.

Thus $T_0$ is a subdirect product on all the simple sites,
and is a product of full twisted diagonals.  Uniqueness of
the local socles makes it normal in $U$.  The local slots are
transitively permuted by a part stabilizer, as required by
primitivity in the two-slot case; $U$ is also transitive on
parts.  Hence $U$ is transitive on all simple sites and on
the diagonal simple factors of $T_0$.  The kernel axis in
$X$ found above has a nontrivial local projection, and the
local centralizer of $S_B$ is trivial.  Commuting that axis
with $T_0$ therefore gives a nonidentity element in
$T_0\cap\langle X^U\rangle$.  This intersection is normal
in $U$, so transitivity on the diagonal factors gives
\begin{equation}\label{n:socle-normal-closure}
 T_0\le\langle X^U\rangle.
\end{equation}

\paragraph{Encoding structures and diagonal ties.}
We next count the local structures and their diagonal ties
without assuming that $|I'|$ is bounded.  In the natural
case $(1,1)$, the site alphabet is the set of part letters.
In case $(1,2)$, choose the graph in which two distinct
two-subsets are adjacent when they intersect.  In case
$(2,1)$, choose adjacency when two distinct points share a
coordinate.  Choose this structure on one part and
transport it by $U$ to all parts; the initial part
stabilizer preserves it, so transport is well-defined.
The resulting relation on $I$ is $X$-invariant and hence
has boundedly many possibilities.

These graphs recover their natural alphabets.  In the two-subset graph,
more than three pairwise-intersecting distinct pairs have a common entry;
the large maximal cliques are exactly the stars, one for each alphabet
entry.  An ordered adjacent pair determines the star at its common entry.
In the sharing graph, a clique of size at least two lies on a unique
coordinate line, so an edge determines that line.  Disjointness groups
the lines into two directions; the lines in each direction form one
recovered site alphabet.  Thus the recovered alphabet points are stars
in the two-subset case and coordinate lines in the product case.
Their disjoint union is an equivariant image of the $X$-invariant subset
of $I^2$ consisting of ordered edges within the parts.  In the natural
case use the diagonal pairs $(p,p)$ instead.  Relations on pairs of
recovered alphabet points therefore pull back to relations on $I^4$,
whose orbit count is bounded by \Cref{n:tuple-bound}.

An isomorphism between these part structures is induced
by natural alphabet bijections, with a possible
permutation of the two slots.  Isomorphisms between high
alternating simple factors are likewise induced uniquely
by alphabet bijections.  Thus record the twisted
diagonals of $T_0$ by their matched-letter relation on
the recovered alphabets.  This relation is equivariant,
because $X$ normalizes $T_0$, and it determines the full
twisted diagonals.  Its pullback is an invariant relation
on $I^4$, so \Cref{n:tuple-bound} bounds all choices of
$T_0$ for the fixed block structure.  This explicit
four-tuple bound also covers ties between sites in
different parts.

\paragraph{Bounded-group lifts.}
Let $W$ be the stabilizer of the partition and chosen
part structures.  In fixed local coordinates its
homomorphism
\begin{equation}\label{n:normalizer-wreath}
 W\longrightarrow D^{I'}\rtimes S_{I'}
\end{equation}
records the alphabet parities and slot permutations,
where $|D|\le8$.  Its kernel is $\prod_B S_B$.
On $N_W(T_0)$ the kernel is exactly $T_0$, because a
subdirect product of nonabelian simple groups is
self-normalizing inside their direct product, as
proved in \Cref{n:rank-one}.  Hence $N_W(T_0)/T_0$
embeds in the bounded-group wreath product in
\Cref{n:normalizer-wreath}.

We can make the image of $X$ pure in this wreath product.
A part stabilizer in $X$ is a product of large
alternating groups, allowing one natural point
stabilizer $A_{n_i-1}$ when the part orbit is natural.
Since $n_i\ge7$, each nontrivial factor is simple of
order greater than eight.  Its local homomorphism to
$D$ is therefore trivial.  Fix a local quotient frame
at one representative of each $X$-orbit of parts and
transport it by $X$.  The stabilizer property makes
this unambiguous; changing frames amounts only to
changing coordinates by local elements of $D$.
All base components of $X$ are now identity.  These
frames depend only on the structures and $X$, not
on the subgroup $U$ subsequently being counted.

Apply \Cref{n:rank-one} to the rank-one part action
$X'$ on $I'$, of complexity at most $s$.  It bounds its possible transitive
overgroups and supplies the generators with
intersection controls needed in \Cref{n:lifting}.
With this complexity bound and $|D|\le8$, that lemma bounds the possible images
$U/T_0$ in \Cref{n:normalizer-wreath} and their indices
modulo the normal closure of the pure image of $X$.
For fixed structures and $T_0$, the image determines
$U$ uniquely, because $U$ contains the kernel $T_0$.
By \Cref{n:socle-normal-closure}, these quotient indices
are exactly $[U:\langle X^U\rangle]$.  Combining the
bounded choices of partitions, structures, matching
relations and lifts proves both assertions.
\end{proof}

\begin{theorem}[Uniform computable template bound]\label{n:template-bound}
There is a total computable nondecreasing function
$\Phi:\mathbb Z_{\ge0}\to\mathbb Z_{\ge1}$ such that every weighted
template
\[
 F=[S_Z\times Y,S_{Z\sqcup(\bigcup T)}]
\]
of \Cref{sec:weighted}, with at most $t$ units, excluding
the foundation $Z$ from this count, and
$|Z|>|\bigcup T|$, satisfies $|F|\le\Phi(t)$.
The bound is independent of all full-unit sizes, square sides,
auxiliary axis lengths, offset sizes, and the size of $Z$.
No stock-size hypothesis is required for this counting assertion.
\end{theorem}

\begin{proof}
On any nonempty union of units, the prescribed product $Y$
contains a control with private squares whose complexity is
bounded computably in the unit count.  On a full unit of size
at least seven, keep its natural alternating group; a smaller
full unit contributes at most six fixed letters.  On an
auxiliary unit $(x/4)\times2\times2$, if $x/4\ge7$, keep
$A_{x/4}$ on its four repeated natural copies.  Otherwise
its domain has at most $24$ letters and may be fixed entirely.
On each primary square keep its two prescribed alternating
factors; their privacy is part of the independent unit
construction.  Thus complexity at most $24t$ suffices.

Consider an arbitrary residual overgroup $U\ge Y_R$ on a
union $R$ of at most $t$ units.  Its orbits are unions of
the transitive units, so there are at most $B_t$ choices
for this orbit partition.  On every orbit $O$, its
transitive projection $U_O$ has at most $B(24t)$ choices,
by \Cref{n:private-squares}.  Let $N_O$ be the normal
closure of the prescribed unit groups on $O$ in $U_O$.
Their contained control shows that
$[U_O:N_O]\le C(24t)$.

The subgroup $N_O$, supported independently on $O$, is
actually contained in $U$: its initial unit groups are
independently supported, and conjugating them by $U$
realizes their normal closure in the projection $U_O$
without adding support on another orbit.  Consequently
$U$ contains $\prod_O N_O$.  Once its projections are
fixed, its remaining choice is a subgroup in
$\prod_O(U_O/N_O)$, a group of order at most
$C(24t)^t$.  There are at most $2^{C(24t)^t}$ such
subgroups.  The empty residual set has one choice.

By the exact literal labeling of the weighted template,
each lattice element consists of an absorbed subset
$I\subseteq T$ and one of these residual overgroups on
$T\setminus I$.  For $t\ge1$, therefore, one may take
the increasing envelope of
\begin{equation}\label{n:explicit-template-envelope}
 2^t B_t\,B(24t)^t\,2^{C(24t)^t},
\end{equation}
and put $\Phi(0)=1$.

To check effectivity of the whole construction, every
invariant-relation choice used above is bounded by a
power of two with exponent a tuple-orbit bound from
\Cref{n:tuple-bound}.  Every bounded quotient has at
most as many subgroups as subsets.  Trimming and small
degree branches use computable finite ranges; equal
profiles use maxima over finitely many bounded
coefficient systems and the terminating Presburger
procedure in \Cref{n:profiles}.  The recursion of
\Cref{n:rank-one} strictly decreases the equivalence
chain parameter.  These operations define total
computable bounds and never take a maximum by a
finite-group search over unbounded unknown degrees.
\end{proof}

\section{The Diophantine reduction}
\label{sec:diophantine}

We combine the effective counting theorem with the size readings of
\Cref{sec:weighted}.  The distinction between their quantifiers is
essential: the counting bound is computable from the number of units,
whereas the literal numerical sizes used to exhibit a template need
only exist when the Diophantine instance has a solution.

\begin{definition}[Positive circuits]\label{n:positive-circuits}
A positive circuit is a finite arithmetic circuit with positive integer
input variables, constant $1$, and addition and multiplication gates.
A \emph{value variable} records either an input or a gate value.
A circuit equation requires equality of two designated outputs.
\end{definition}

Solvability of positive circuit equations is undecidable.  Indeed,
the Davis--Putnam--Robinson--Matiyasevich Theorem gives undecidability
of polynomial solvability over the nonnegative integers
\citep{Matiyasevich1970,Matiyasevich1971}.  Replace each nonnegative variable $x$ by
$y-1$ with $y$ positive, expand the resulting integer polynomial,
and move its positive and negative coefficient terms to opposite
sides.  Add $1$ to both sides, so each side is positive even when
one collection of terms was empty.  Positive integer coefficients
are made by sums of $1$, powers by products, and the resulting
polynomials by positive circuits.  All these transformations are
effective and preserve existence of a solution.

\begin{lemma}[Encoding a circuit by weighted comparisons]
\label{n:circuit-template}
There is an absolute stock integer $S$ with the following property.
For every positive circuit equation $\mathcal C$, one can compute
an integer $t_S(\mathcal C)$ from its syntax such that every positive
solution determines a weighted template $F$ with at most
$t_S(\mathcal C)$ units.  In every size reading of this template
supplied by \Cref{t:weighted-realization}, the actual raw sizes
give a positive solution of $\mathcal C$ in the same cell alignment.
If all cells in that reading have size at most $Q$, the resulting
solution has all its value variables at most $Q$.
\end{lemma}

\begin{proof}
For every value variable $v$, including all gate values, prescribe
a full \emph{raw cohort} $R_v$ of size $v$ per unit and a full
\emph{side cohort} $E_v$ of size $16v$ per unit.  These sizes are
literal integers when a solution is given.  The relation between
the two cohorts is recorded by a comparison of one side unit
with sixteen distinct raw units.  For a constant value $v=1$,
compare one side unit with sixteen distinct singleton markers.

For each required square side, include the primary square cohort
$Q_v$ of side $16v$, its auxiliary cohort of shape
$(4v)\times2\times2$, the full axis-reference cohorts, and all
required full offset cohorts.  A unit from $E_v$ itself may serve
as the full side reference for $Q_v$.  Every extra full reference
and every raw or side cohort receives its prescribed offset
where required; an offset cohort requires no further offset.
The formal square sides satisfy $8\mid16v$ and $16v\ge16$.
The needed full references for bounded slots are included as
well.  Distinct roles may use disjoint lots even when their
numerical sizes happen to coincide.

An addition gate $c=a+b$ is recorded by the comparison
\begin{equation}\label{n:add-comparison}
 E_c=E_a+E_b,
\end{equation}
where an equality of cohort symbols here means an equality of
the indicated totals on chosen units.  For a multiplication
gate $c=ab$, introduce a fresh value $d=a+b$, with its own raw
and side cohorts and its side-to-sixteen-raw comparison.
Require \Cref{n:add-comparison} for $d=a+b$, and use the
square comparison
\begin{equation}\label{n:product-comparison}
 Q_d=Q_a+Q_b+32E_c.
\end{equation}
Whenever needed, add the square cohorts for $a,b,d$ and their
associated roles.  The output equality is recorded using
two distinct copies of each output side cohort.  Thus even
an equality between two individual values is read as an
equality between two-unit totals.

Each comparison is between disjoint collections of a fixed
bounded number of units, and at least one collection uses
more than one unit.  In particular, \Cref{n:product-comparison}
uses one unit against two square units and thirty-two side
units, so thirty-five bounds the total support of every circuit
comparison.  Repeated occurrences of the same variable within
a comparison use distinct copies.  The pair-total reading
in \Cref{t:weighted-equal-totals} equalizes all the actual
sizes in a cohort, so these choices of copies do not change
the values being read.  The same copies may be used again
in another comparison.  The comparisons for constants
use the common marker size one.

Choose once and for all a stock $S$ large enough for
\Cref{t:weighted-product-success,t:weighted-realization},
including all bounded marker losses, all marker requests
in the readings, and the bounded number of distinct copies
needed in any one comparison.  Enlarge it by a fixed amount
if necessary for the raw/side and output comparisons.
With the comparison support and spare-copy needs fixed,
neither the number of comparisons or cohorts nor any numerical
size enters this choice.  Singleton raw units may also count
as markers.  Include the global marker-loss bound in every
cohort's stock, so enough copies survive even if all deletions
fall in one raw cohort.  The same surviving copies can be
used in different comparisons.

The role list is computable from the circuit: take its
value variables, the fresh sums for multiplication gates,
the square and auxiliary roles they require, their full
axis references, and their offsets.  Put the fixed stock
in every role and include the fixed marker stock.  There
are only finitely many prescribed roles per value or
shape, and offsets introduce no further offset recursion.
Counting these listed lots gives $t_S(\mathcal C)$.
This calculation uses no solution values.  If sizes of
different roles coincide, keeping the lots separate
still gives this same bound on their total unit count.

Now start with a positive solution and construct the
literal weighted template using its values.  All the
prescribed total equalities hold.  Apply the weighted
size readings in their one common cell alignment.  Let
$r_v$ and $e_v$ be the actual raw and side sizes; the
cohort equalization makes these well-defined positive
integers.  The raw/side and constant comparisons give
\begin{equation}\label{n:read-side}
 e_v=16r_v,\qquad r_v=1\ \hbox{for every unit constant}.
\end{equation}
In particular each side reference is even and at least
sixteen, so \Cref{t:weighted-square} applies to every
required primary square and reads its actual size as
$e_v^2$.  Addition comparisons give $r_c=r_a+r_b$.
At a multiplication gate, the fresh sum gives
$r_d=r_a+r_b$ and \Cref{n:product-comparison} becomes
\begin{equation}\label{n:factor-thirty-two}
 \bigl(16(r_a+r_b)\bigr)^2
   =(16r_a)^2+(16r_b)^2+32(16r_c).
\end{equation}
Subtracting the first two squares gives
$512r_ar_b=512r_c$, hence $r_c=r_ar_b$.
This explains the factor $32$: the cross term of
the squared side sizes is $2\cdot16^2r_ar_b$, while
one side unit for $c$ has size $16r_c$.
The doubled output comparison gives the required
output equality.  All equations concern the same
actual raw sizes, because the size readings use one
alignment simultaneously.  These sizes are themselves
cell sizes, so their bound by $Q$ proves the last
assertion.
\end{proof}

\begin{theorem}[Undecidability of finite representability]
\label{n:undecidable}
There is no algorithm which, given the tables of a finite
lattice, decides whether it is isomorphic to the congruence
lattice of a finite algebra.  Equivalently, there is no
total computable bound, in terms of lattice cardinality,
for the minimum carrier size among all representable
finite lattices.
\end{theorem}

\begin{proof}
Suppose that such a decider exists.  By
\Cref{f:cor-witness-bound}, it gives a total computable
nondecreasing function $M(N)$ that bounds a representing
carrier for every representable lattice with at most
$N$ elements.  Explicitly, enumerate the finitely many
lattice tables on at most $N$ points, retain the YES
instances using the decider, and search the finite
colored-graph witnesses for each retained table until
one is found.  The graph criterion supplies a finite
algebra on that witness carrier.  Taking the maximum
over these finitely many searches produces $M(N)$.

Fix the absolute integer $S$ from \Cref{n:circuit-template}.
Its existence is sufficient: the single integer $S$
may be hardwired in the algorithm below.  This step
does not assert an algorithm for extracting a sharp,
or even a specified, sufficient stock from the
structural arguments.  Once fixed, the function
$t_S(\mathcal C)$ is computable from the circuit syntax.
For an input circuit equation $\mathcal C$, compute
\begin{equation}\label{n:computable-witness-bound}
 \begin{aligned}
  t&=t_S(\mathcal C),&
  N&=\Phi(t),&
  m&=M(N),&
  Q&=60^m m!.
 \end{aligned}
\end{equation}

We claim that if $\mathcal C$ has any positive solution,
then it has one whose value variables are all at most
$Q$.  Given an arbitrary positive solution, use
\Cref{n:circuit-template} to exhibit its weighted
template $F$.  Its literal size parameters may be much
larger than the bounds just computed, but
\Cref{n:template-bound} still gives $|F|\le N$.
The template is representable by the coset action of
its defining ordinary subgroup interval, as in
\Cref{t:weighted-algebra}.  Hence it has a finite-algebra
representation on at most $m$ points.

Among the representations of all nontrivial lattice
images of this $F$, take one with least carrier size
$m_0$.  Since $F$ itself is one such image, $m_0\le m$.
The nontrivial-image reduction of
\Cref{t:weighted-algebra} gives an ordinary interval
for the dual of a nontrivial image, with top group
order at most
\[
 60^{m_0}m_0!\le60^m m!=Q.
\]
The product-success and size-reading
\Cref{t:weighted-realization} then provides the common
actual cell alignment, after only the allowed bounded
marker deletions.  Every read cell size is positive
and at most this top-order bound; markers have size
one.  \Cref{n:circuit-template}, with its preserved
stocks, now extracts a positive solution of
$\mathcal C$ all of whose value variables are at most
$Q$.  The side sizes are multiples of sixteen by
\Cref{n:read-side}, so every use of the square reading
in this extraction satisfies its actual numerical
hypothesis.

Thus the complete quantifier implication is
\begin{equation}\label{n:quantifier-chain}
 \begin{aligned}
 &\exists\,\mathbf v>0\quad\mathcal C(\mathbf v)
 \\
 &\quad\Longrightarrow\quad
 \exists F\quad |F|\le N,\quad F\text{ representable}
 \\
 &\quad\Longrightarrow\quad
 \exists A\quad \operatorname{Con}A\cong F,\quad |A|\le m
 \\
 &\quad\Longrightarrow\quad
 \exists[H,G]\text{ as above with }|G|\le Q
 \\
 &\quad\Longrightarrow\quad
 \exists\,\mathbf r\in\{1,\ldots,Q\}^{k}
       \quad\mathcal C(\mathbf r).
 \end{aligned}
\end{equation}
where the circuit determines $k$ and all quantities in
\Cref{n:computable-witness-bound} are computed before
any solution is known.  Auxiliary fresh sums can be
included among the $k$ value variables; restricting
back to the original values preserves the conclusion.
The reverse implication from a bounded solution to
some solution is immediate.

An algorithm may therefore compute
\Cref{n:computable-witness-bound} and search all positive
assignments up to $Q$ to the finitely many value
variables, checking each gate and the output equality.
It terminates, and \Cref{n:quantifier-chain} makes its
answer correct.  It never needs to construct a weighted
template with unknown sizes: that template was used
only to prove that an existing solution has a bounded
replacement.  This would decide positive circuit
solvability, contradicting the
Davis--Putnam--Robinson--Matiyasevich Theorem.
The equivalent carrier-bound formulation follows
again from \Cref{f:cor-witness-bound}.
\end{proof}

The structural input to this reduction is precisely the
weighted success theorem, whose exclusion arguments use
the sufficiently-general subgroup theorem in the scope
established earlier.  The reduction introduces no
additional uniformity in the numerical weights beyond
the computable counting bound proved in
\Cref{n:template-bound}.

\begin{corollary}[Undecidability of finite subgroup-interval recognition]
\label{n:interval-undecidable}
There is no deterministic Turing machine which, given the order table
of a finite nonempty lattice $L$, always halts and decides whether
$L\cong[H,G]$ for some finite group $G$ and subgroup $H\le G$, where
$[H,G]$ consists of every subgroup $K$ with $H\le K\le G$, ordered by
inclusion.  Equivalently, there is no total computable bound, in terms
of lattice cardinality, on the least possible order of the top group
$G$ among lattices admitting such a realization.
\end{corollary}

\begin{proof}
Suppose first that such a decider exists.  For each $N\ge1$, enumerate
all order-relation tables on sets of at most $N$ points, retain those
which are nonempty lattice tables, and use the decider to select the
affirmative ones.  For each selected table, enumerate finite
multiplication tables in increasing carrier size.  Test the group
axioms, enumerate the subgroups $H$, and for each $H$ enumerate every
subgroup $K$ containing it.  Testing the finitely many bijections
then decides whether the resulting full inclusion interval is
isomorphic to the selected lattice.  Each selected search terminates.
The maximum of the witness group orders found and $1$, followed by
an increasing envelope if necessary, gives a total computable
nondecreasing function $B(N)$ such that every affirmative lattice
with at most $N$ elements has a realization $[H,G]$ with $|G|\le B(N)$.
Conversely, any total computable bound on the least possible $|G|$
gives a decider by the same finite tests restricted to group orders
up to that bound.  Thus it suffices to contradict the existence of $B$.

Fix the absolute stock integer $S$ of Lemma~\ref{n:circuit-template}; as
in the preceding proof, this one integer may be hardwired.  For a
positive circuit equation $\mathcal C$, compute from its syntax
\[
 t=t_S(\mathcal C),\qquad N=\Phi(t),\qquad Q=B(N).
\]
These quantities are computed before any solution is known.
If $\mathcal C$ has a positive solution, Lemma~\ref{n:circuit-template}
exhibits a weighted template $F$ with at most $t$ units, and
\Cref{n:template-bound} gives $|F|\le N$.  This literal template is
itself a full ordinary subgroup interval by Definition~\ref{t:weighted-template}.
It is nontrivial: the circuit construction includes side units of
size $16v\ge16$.  The bound $B$ therefore supplies an ordinary
realization $F\cong[H,G]$ with $|G|\le Q$.

Use the identity $F\to F$ as the nontrivial weighted image allowed
by Convention~\ref{t:weighted-images}.  Apply the ordinary-input part of
\Cref{t:weighted-realization} to this realization with its actual
top order $M=|G|$.  It gives one common size reading whose positive
cell sizes satisfy $u_j\le M\le Q$.  By Lemma~\ref{n:circuit-template},
that reading yields a positive solution of the same equation
$\mathcal C$ with every value variable at most $Q$.

Consequently, if $\mathcal C$ has any positive solution, it has one
in $\{1,\ldots,Q\}^{k}$, where $k$ is determined by the circuit and
may include its finitely many auxiliary fresh sums.  The reverse
implication is immediate.  Computing $Q$ and testing these finitely
many assignments would decide positive-circuit solvability,
contradicting the undecidability established at the start of this
section.  The argument uses the existence of $F$ and of its reading
only to prove a bounded replacement for a solution; it does not
require either to be constructed without that solution.
\end{proof}

\section{An explicit computable size ceiling}\label{sec:quantitative}

We give an ordinary computable upper bound for the cardinality of one
nonrepresentable lattice in the affine-chart family.  The bounded-rank
exclusion in \Cref{sec:bounded} used a sufficiently-general-subgroup
theorem with an unspecified uniform threshold.  To choose explicit finite
parameters here, we first prove a quantitative substitute: a finite subgroup
avoiding proper closed sets of a specified computable format lies between
the derived group and the full fixed group of a Steinberg map.  We then
bound the pruning and protected-path parameters and count the resulting
lattice.

These bounds concern the negative affine family, independently of the
weighted-template counting bound and of any hypothetical representation
decider.  They provide the fixed initial segment of lattice tables used in
\Cref{sec:minimum}.

\subsection{A uniform expanding function}

\begin{definition}[Affine ring format]\label{q:format}
For an integer $x\ge10^9$, an affine ring formula has \emph{format at most
$x$} if it has at most $x$ occurrences in total of polynomial-equality atoms
and the operations $\neg,\land,\lor,\exists,\forall$, at most $x$ variables,
and polynomial atoms of total degree at most $x$.  Occurrences are counted in
the formula tree, with a separate variable for each quantifier binder.  The
variable count includes the entire ambient affine coordinate list, even
unused coordinates.  Coefficients are arbitrary field elements, and a
polynomial is counted through its atom and degree, rather than its expanded
arithmetic expression.

Let $a(x)$ be the least positive simultaneous bound, over all algebraically
closed fields and formulas of format at most $x$, for both the format of an
equivalent quantifier-free formula in the same ambient coordinates and the
number and degrees of polynomial equations defining its Zariski closure.
\end{definition}

\begin{lemma}\label{q:acf-computable}
The integer $a(x)$ exists, is nondecreasing, and is total computable.  This
holds uniformly over the characteristic and the specialization of every
coefficient.
\end{lemma}

\begin{proof}
For fixed $x$, canonically rename the coordinate and bound variables.  There
are finitely many possible formula trees.  Replace each polynomial atom by
the polynomial containing every allowed monomial of degree at most $x$, with
an independent coefficient parameter.  This produces finitely many schemes.
Their parameter lists are finite and computable, although the specialized
coefficients need not have effective descriptions.  These parameters belong
to the tests of uniformity, not to the variable count of the specialized
input formula.

Effective quantifier elimination in the theory ACF, without fixing the
characteristic, supplies a quantifier-free formula for each scheme
\cite{CohenMahboubi2012}.  Write it as a finite disjunction of conjunctions.
For a conjunction $f_i(z)=0$, $g_j(z)\ne0$, introduce one inverse variable $u$
and the ideal
\[
 I=\bigl(f_i(z),\ u\textstyle\prod_jg_j(z)-1\bigr).
\]
The Nullstellensatz gives the closure of that conjunction as $V(I\cap k[z])$ \citep[Theorem~10.34.1]{StacksNullstellensatz}.  An elimination
Gr\"obner basis gives equations for this closure.  The closure of the original
finite union is the union of these closures, which again has a finite
equation list, for example from the product of their defining ideals.
Characteristic-independent Gr\"obner bounds are available
\cite{Dube1990}; we give also the uniformity argument needed to compute the
\emph{least} simultaneous bound.

For given bounds on variables, generators, and degrees, the assertion that a
tuple of generators has a Gr\"obner basis with certificates of size at most
$b$ is first order in its coefficients.  The certificates consist of bounded
lists of polynomials, bounded ideal combinations in both directions, choices
of leading monomials from finite lists, and Buchberger reductions or standard
representations.  The required monomial comparisons for a fixed elimination
order form finite tests; polynomial identities become coefficient equations.
We can make these assertions increasing in $b$.  Every individual tuple in
every algebraically closed field has some certificate.  If no uniform $b$
existed, the Compactness Theorem \citep[p.~26]{Marker2020}, applied to ACF with constants for the coefficients and
the negations of all these increasing assertions, would give a tuple with no
certificate at all.  Thus a uniform $b$ exists.  First-order ACF tests of
successive $b$ find one.

In such a certified elimination basis, the members involving only the
retained coordinates define the projected closure.  Independence of the
eliminated variables is itself a finite coefficient test.  Existentially
quantifying a bounded certified basis therefore gives an effectively found
formula $C_\phi(c,z)$ for the closure of the fiber defined by $\phi(c,z)$.
The coefficient specialization precedes the certificate test.  In
particular, this constructs each specialized fiber's closure, rather than
specializing the closure of the total parameter family.  Vanishing leading
coefficients, empty pieces, and empty equation lists are handled by finite
branches.

For a candidate bound $A$, enumerate the finitely many quantifier-free
templates $Q(d,z)$ of format at most $A$ in the same coordinates.  For each
input scheme, the assertion
\[
 \forall c\ \bigvee_Q\exists d\ \forall z\,
       \bigl(\phi(c,z)\longleftrightarrow Q(d,z)\bigr)
\]
tests the quantifier-free part of the bound.  Let $P_1(e,z),\ldots,P_A(e,z)$
be dense polynomials of degree at most $A$ with coefficient tuple $e$.
The additional assertion
\[
 \forall c\ \exists e\ \forall z\,
 \left(C_\phi(c,z)\longleftrightarrow
                       \bigwedge_{i=1}^{A}P_i(e,z)=0\right)
\]
tests the closure bound; zero polynomials pad shorter lists.  Conjoin these
tests over the finitely many schemes and ambient dimensions.  Validity of
the resulting sentence in all algebraically closed fields is decidable:
after effective elimination a closed sentence is a Boolean combination of
integer equalities, whose behavior is checked in characteristic zero and at
the finitely many prime divisors of the nonzero integers occurring.  Testing
$A=1,2,\ldots$ consequently finds exactly $a(x)$.  The nested input classes
make $a$ nondecreasing.  Finally, B\'ezout bounds give elementary bounds on
the number of components and their sum of degrees from these equation
bounds.  All these statements allow ambient affine matrix coordinates,
including equations specifying the group itself.
\end{proof}

\begin{definition}[The function $J$]\label{q:J}
Let $R(c,h)$ be the least positive set size forcing a monochromatic
$h$-element subset in every $c$-coloring of its unordered pairs of distinct
elements.  Let $\operatorname{HJ}(a,c,s)$ be the least positive word length
forcing a monochromatic $s$-slot word family over an $a$-letter alphabet in
every $c$-coloring.  Here some positions have fixed letters, possibly no
positions do, and the others are partitioned into $s$ disjoint nonempty
sets; each set carries one independently arbitrary common letter.  Put
\begin{align}
 J(x)=1+\max\bigl(&\{2^{2^{x!}},a(x)\}
     \cup\{R(c,h):1\le c,h\le x\}\notag\\
     &\cup\{\operatorname{HJ}(a,c,s):1\le a,c,s\le x\}\bigr).
 \label{q:J-formula}
\end{align}
The notation $J^d$ denotes iteration, not an ordinary power.
\end{definition}

The finite Ramsey and multidimensional Hales--Jewett theorems
\cite{Ramsey1930,HalesJewett1963} ensure that these least bounds exist.
For each parameter tuple they can be found, in principle, by finite coloring
tests at successively larger sizes.  Thus $J$ is nondecreasing and total
computable.  Its double-exponential factorial term ensures that an additional
iteration absorbs the elementary polynomial and exponential costs used below.
No evaluation of these finite searches is required.

\subsection{Formats for bounded root systems}

We make the genericity bound quantitative in the affine coordinates of the
adjoint action.  The conventions for positive standard point maps, including
the parameter $P$ for the special twisted types, are those of
\Cref{sec:bounded}.

\begin{lemma}[Root words and fixed points]\label{q:root-formats}
Let $10^9\le x\le y$ be integers, and let $G$ be split pinned, adjoint
simple of rank at most $x$ over an algebraically closed field of positive
characteristic.  Its adjoint representation embeds $G$ as a closed subgroup
of $\operatorname{GL}_d$, where $d=\dim G=O((x+10)^2)$.
For rational simple modules of the simply connected cover, of dimension
and highest-weight coefficients at most $y$, graphs of the induced
projective actions have format at most $J(y)$.  Stabilizers of a subalgebra
or of at most $y$ subspaces of the
corresponding matrix algebra, and normalizers of projective subgroups of
order at most $y$, have format at most $J^3(y)$.

For a positive standard map of parameter $P$, the fixed-point test has a
Boolean description with at most $J^5(y)$ atoms and degrees at most
$J^5(y)P$.  A proper closed subset of $G$ of format at most $y$ therefore
contains at most
\begin{equation}\label{q:proper-fixed-count}
                 J^{10}(y)P^{\dim G-1}
\end{equation}
fixed points.  For $P>J(y)$ the derived simple fixed group has order at least
$P^{\dim G}/(10y)^4$.  The same bounds apply to the other adjoint types in
the stated rank range.
\end{lemma}

\begin{proof}
\emph{The embedding.}
Every root is primitive in the character lattice of the adjoint torus,
which is the root lattice: it is Weyl-conjugate to a simple root.  Its
differential is therefore nonzero in every characteristic.  An element
central in the Lie algebra has no root-space components, by commuting with
toral directions.  Its remaining toral component vanishes because the
simple-root differentials form a basis.  Thus the adjoint kernel has zero
Lie algebra.  It is finite reduced, is central by connectedness, and is
trivial because $G$ is adjoint.  The separable isomorphism theorem for
algebraic groups identifies $G$ with its closed adjoint image.  The image
has determinant one.  We use a Frobenius-standard basis.

\emph{Word and stabilizer formats.}
Bruhat decomposition and the rank-one torus formulas cover both simply
connected and adjoint points by root-element words of length less than
$1000(y+10)^4$.  This includes zero parameters and words for Weyl
representatives.  Root matrices on the modules in question have polynomial
entries of bounded degree: a coefficient of $t^h$ shifts weight by $h\alpha$,
and the weights lie in the Weyl polytope.  Coroot heights in the irreducible
root lists are less than $100(y+10)$, so pairing with the coroot bounds the
rootwise weight range.  The number of word choices needed is at most
\begin{equation}\label{q:word-choices}
             \bigl(1000(y+10)^2\bigr)^{1000(y+10)^4+1}.
\end{equation}
Use entries of intermediate products as variables and the coefficients of
the root maps as parameters.  Describing a group element and its action by
the same word gives the asserted $J(y)$ format, using invertible matrices
modulo scalar matrices for a projective action.  The same bound also
dominates Weyl dimensions for the indicated highest-weight coefficients, by
the Weyl dimension formula.  Conditions that a matrix preserve a subspace,
a subalgebra, a collection of at most $y$ subspaces, or normalize a
projective group with at most $y$ listed elements add only bounded lists of
variables and equations.  Allowing matrices and translates as parameters
gives the $J^3(y)$ bound.  Images, preimages, and intersections combine these
formulas; closures use \Cref{q:acf-computable}.

\emph{Point maps and the fixed-point estimate.}
Pinned diagram maps and special isogenies in their fixed characteristics
have bounded graph format by the same words: their root parameters are
sent to $t$ or $t^p$.  Split field powers, including their inverses on
points, only twist the coefficients of sets.  Inverses of special maps use
also an inverse field power.  Images of bounded sets under these point
maps thus cost only a few further iterations of $J$.

For a positive standard map, first eliminate the bounded combination of
inner, pinned diagram, or special maps, whose graph has format at most
$J^2(y)$, and only then substitute the single split field power.  This
gives at most $J^5(y)$ atoms of degree at most $J^5(y)P$ in the fixed-point
test.  On a conjunction piece, the inequalities specify an open subset.
Since the fixed set is finite, only isolated components after the
polynomial cuts can contribute.  On a variety of dimension $d$ and sum of
component degrees $\delta$, cuts of degree at most $E$ contribute at most
$\delta E^d$ such points: choose at most $d$ generic linear combinations
of the cuts retaining the finitely many points as isolated points, and
apply affine B\'ezout for locally closed sets \citep[Definition~3, Proposition~2.9 and Theorem~2.14]{PardoSebastian2020}.  The generic-combinations step supplies the exponent $d$ locally.  For a proper subset of $G$, $d\le\dim G-1$.
The degree bound obtained from $a(y)$, the number of Boolean pieces, and
the factor $(J^5(y))^{\dim G}$ are absorbed by $J^{10}(y)$, proving
\eqref{q:proper-fixed-count}.  The ordinary order and center formulas give
the stated lower bound for the derived simple when $P>J(y)$.
\end{proof}

\begin{proposition}[Quantitative genericity]\label{q:genericity}
Let $x\ge10^9$, and let $G$ be adjoint simple of rank at most $x$ over an
algebraically closed field of characteristic $p>0$.  If a finite subgroup
$\Gamma\le G$ is not contained in any proper closed subset of affine ring
format at most $J^{1000}(x)$, then there is a positive-height standard point
map $\Psi$ such that
\begin{equation}\label{q:sandwich}
                   (G^\Psi)'\le\Gamma\le G^\Psi.
\end{equation}
Proper closed subsets in this hypothesis need not be irreducible.
\end{proposition}

\begin{proof}
Choose a fundamental simple module of the simply connected cover of $G$.
Its dimension $n$ is at most $J(x)$.  Its weights lie in one coset of the
root lattice, so the center acts by a scalar character, giving an
irreducible projectively point-faithful action of $G$.

\emph{Claim 1. The matrix lifts of every nontrivial normal subgroup of
$\Gamma$ span the full matrix algebra, and $\Gamma$ is almost simple.}
This connected action preserves no nontrivial tensor splitting individually.
Indeed the projective actions on the factors are algebraic, as can be seen
from their conjugation actions on the factor matrix algebras.  Pull back
$\operatorname{SL}\to\operatorname{PGL}$ along each map from the simply
connected source.  The reduced identity group of the pullback still
surjects onto that source: the full map is finite surjective and has only
finitely many components.  Its kernel is central of multiplicative type.
Its connected solvable radical has trivial image and hence dimension zero,
so it is semisimple.  The resulting central isogeny onto the simply
connected source is an isomorphism, providing a linear lift.  Perfection
leaves no scalar character in the tensor comparison.  Irreducibility would
make both tensor factors nontrivial irreducible modules, but their dominant
highest weights cannot add to a single fundamental weight.

An invariant unital matrix subalgebra is accordingly scalar or full.  Its
radical would otherwise give a proper invariant annihilator.  If it is
semisimple, connectedness fixes its central idempotents, so irreducibility
leaves one matrix factor.  A proper nonscalar such factor would yield the
excluded tensor splitting.

We may use $y=J^2(x)$ in the stabilizer tests of
\Cref{q:root-formats}; it covers $n^2$ and all the lists of matrix data
needed here.  Their stabilizers and the bounded unions of translates used
below have format well within $J^{1000}(x)$.  Avoidance implies
$\Gamma\ne1$.  For every nontrivial normal subgroup of $\Gamma$, the
algebra spanned by its matrix lifts has a bounded-format stabilizer
containing $\Gamma$, so is $G$-invariant.  It is nonscalar by projective
faithfulness, and hence is the full matrix algebra.

If such a normal subgroup were abelian, its distinct projective elements
would have lifts with distinct scalar conjugation characters.  Equal
characters would make the quotient of two lifts centralize the full span,
forcing equal projective elements.  Independence of the corresponding
eigenvectors bounds the order by $n^2$.  Its normalizer is then a bounded
proper closed subgroup: a finite nontrivial subgroup normal in connected
adjoint simple $G$ would be central, which is impossible.  This contradicts
avoidance.

A minimal normal subgroup is therefore a power of a nonabelian simple
group.  Lifts of its factors commute linearly across factors, by
perfection, and their algebras jointly generate the full matrix algebra.
A radical in one algebra would generate a nilpotent ideal in that full
algebra, so vanishes; the center is scalar for the same commuting-generation
reason.  Each factor algebra is consequently a full matrix algebra, and
their tensor product is simple and maps onto the full matrix span.  These
nontrivial tensor factors are permuted by $\Gamma$.  The bounded stabilizer
of their collection must contain $G$, whose connectedness fixes each
factor.  Tensor indecomposability leaves exactly one factor.  Denote the
resulting simple normal subgroup by $S$.  Fullness and perfection give it
trivial projective centralizer, and conjugation identifies
\begin{equation}\label{q:almost-simple}
                       S\le\Gamma\le\operatorname{Aut}(S).
\end{equation}

\emph{Claim 2. The simple normal subgroup $S$ is of Lie type in
characteristic $p$, of rank at most $100(n+1)$ and parameter
$P>J^{100}(x)$.}
The classification of finite simple groups, together with the usual
covering and automorphism descriptions \cite{GLS1994}, now reduces to a
large-parameter group of Lie type.  We record the numerical exclusions.
A faithful projective degree-$n$ action of $A_5^u$ forces $n\ge2u$.  Lift to the
product of the universal covers.  Each perfect factor acts nontrivially on
some composition factor, since a nonsoluble image cannot act trivially on
all composition factors.  Irreducible modules of a finite direct product
over an algebraically closed field are outer tensors, also in modular
characteristic.  A composition factor on which $a\ge1$ direct factors act
nontrivially has dimension at least $2^a\ge2a$; summing proves the claim.

In alternating degree $d$ there are $\lfloor d/5\rfloor$ disjoint copies
of $A_5$.  In high classical rank $e$ there are at least
$\lfloor(e-1)/5\rfloor$ such factors on ordinary or paired hyperbolic
coordinates; nonsplit forms still have enough coordinate pairs, and
unitary groups can use orthonormal coordinates.  These perfect permutation
groups lie in the projective special group.  Thus the relevant Lie rank
is at most $100(n+1)$, and the alternating degree is bounded as well.
For Lie parameter at most $J^{100}(x)$ in this rank range, or for these
bounded alternating groups, sporadics, and small exceptions, the ordinary
order and automorphism bounds give
$|\operatorname{Aut}(S)|<J^{110}(x)$.  A finite set of this size itself
has format below the avoidance ceiling, contradicting
\eqref{q:almost-simple}.  Hence the parameter satisfies
\begin{equation}\label{q:large-parameter}
                         P>J^{100}(x).
\end{equation}

Suppose its defining prime $p_0$ differs from $p$.  At this parameter the
standard multiplier theorem identifies the simply connected fixed group
$\mathbf H^{F_0}$ of a pinned Steinberg map with the universal central
cover of $S$.  Our projective representation lifts linearly to it and is
projectively nontrivial on noncentral unipotents.  In ordinary
field/diagram type, put $q=P$.  The highest long root $\theta$ is
diagram-stable, also in the simply laced case.  The fixed root elements
have parameters $b=cb^q$ for a fixed $c\ne0$, a one-dimensional additive
$\mathbf F_q$-space.  The fixed coroot elements $\theta^\vee(a)$,
$a\in\mathbf F_q^\times$, act by multiplication by $a^2$.

For a special twisted type write
\[
 F_0=\sigma\operatorname{Fr}_{p_0}^{m},\qquad
 \sigma^2=\operatorname{Fr}_{p_0},\qquad
 q=p_0^{2m+1}=P^2,\qquad p_0\in\{2,3\}.
\]
The highest long and short roots $\theta,\theta'$ are paired by the
special root correspondence, which preserves the positive chamber and
exchanges lengths.  The pinned root-power rules and square-Frobenius
identity hold in each of $B_2,F_4$ in characteristic two and $G_2$ in
characteristic three \citep[Sections~2--3 and Definition~5.4]{DeMedtsNaert2019}.  Their root groups commute because no positive
combination involving both is a root.  In pinned coordinates over the
split prime field,
\[
                 u(b)=x_\theta(b)F_0(x_\theta(b)),
                 \qquad b\in\mathbf F_q,
\]
is a fixed additive copy of $\mathbf F_q$.  On these rank-one parameters
the exponent of $F_0$ is $p_0^r$, where $r=m$ or $m+1$, and the root datum
also gives
\[
 F_0(\theta^\vee(a))=(\theta')^\vee(a^{p_0^r}).
\]
For $a\in\mathbf F_q^\times$, the coroot product
$t(a)=\theta^\vee(a)F_0(\theta^\vee(a))$ is fixed, since the two
coroot factors commute and $F_0^2=\operatorname{Fr}_q$.
The same identity proves fixedness of $u(b)$.  On its $\theta$-coordinate,
$t(a)$ acts by the two root--coroot pairings.  Equivariance makes the
paired coordinate the $F_0$-image of the transformed $\theta$-coordinate.
Thus $t(a)$ sends $u(b)$ to $u(cb)$, where
\[
                   c=a^{2+p_0^r\langle\theta,(\theta')^\vee\rangle}.
\]
The exponent is nonzero at our large $m$ and has absolute value at most
$2+3p_0^{m+1}$.  In ordinary type the scalar subgroup has size at least
$(P-1)/2$; in special type its size is at least
\[
 \frac{P^2-1}{2+3\sqrt{p_0}P}>\frac P{20}.
\]
Thus either case gives a noncentral additive $p_0$-group with a scalar
action subgroup of order at least $P/20$.  It acts diagonalizably in the
linear lift, with a nontrivial character.  Finite-field trace
parametrization shows that the orbit of such a character under the scalar
subgroup is free.  Consequently $n\ge P/20$, contrary to $n\le J(x)$
and \eqref{q:large-parameter}.  The defining characteristic is $p$.

\emph{Claim 3. The outer index satisfies $[\Gamma:S]\le J^3(x)$.}
This bound will control the unions of $S$-cosets used in Claim~4.
Lift to the simply connected universal cover of this large-parameter
group, now with pinned Steinberg map $F$.  The lift is unique by
perfection and irreducible by fullness.  Steinberg's Restriction Theorem,
including its twisted versions, realizes it as the restriction of a
rational simple $L(\lambda)$ of $F$-restricted highest weight
\cite{Steinberg1963,Steinberg1968,Jantzen2003}.  Write its ordinary digits
as $\lambda=\sum_h p^h\lambda_h$.  The Tensor Product Theorem shows
that at most $\log_2n$ digits are nonzero.  Each digit coefficient is less
than $n$: on a highest simple-root face a coefficient $c<p$ gives the
simple rank-one module of dimension $c+1$.

The cyclic subgroup of $\operatorname{Out}(S)$ induced by the pinned
$\operatorname{Fr}_p$ has index at most
$100(100(n+1)+10)^2$.  Every class of this cyclic subgroup realized by
$\Gamma$ preserves the lifted module.  Indeed remove its inner
discrepancy in $S$; a projective intertwining becomes a linear
intertwining on the perfect cover.  For
$F=\operatorname{Fr}_p^f\tau$, with $\tau$ a pinned diagram map, such a
shift preserves the occupied digit positions modulo $f$, by uniqueness
of $F$-restricted restriction.  Wrapping a digit past $f$ only applies a
diagram to it.  Fixing one occupied position bounds the preserving shifts
modulo $f$ by $\log_2n$.  At most three representatives per shift are
needed, since the actual order divides $f\operatorname{ord}(\tau)$.

For special types the unequal root lengths require a different slot count.
For $F=\operatorname{Fr}_p^m\sigma$ with
$\sigma^2=\operatorname{Fr}_p$ and $f=2m+1$, write
\begin{equation}\label{q:special-slot-data}
 \sigma^*\omega_i=p^{c_i}\omega_{i'},\qquad
 c_i,c_{i'}\in\{0,1\},\qquad c_i+c_{i'}=1.
\end{equation}
Here the restricted parametrization is Steinberg's special-factor
version \citep[Sections~11--12, especially Theorem~12.2]{Steinberg1963};
the conversion to the following unequal coefficient boxes and the slot
rotation are deductions from that parametrization and the displayed
root-datum pullback.  The independent universal-cover input used to
lift the projective representation has already been imposed above.
For a short node $a$ paired with a long node $a'$, the pullbacks are
$\sigma^*\omega_a=\omega_{a'}$ and
$\sigma^*\omega_{a'}=p\omega_a$.  On the finite fixed group,
$\sigma$ agrees with $\operatorname{Fr}_p^{m+1}$, by applying $\sigma$
to $F(g)=g$.  The special tensor factorization therefore converts
paired coefficients
\[
 0\le A_a<p^{m+1},\qquad 0\le B_{a'}<p^m
 \quad\longmapsto\quad C_a=A_a+p^{m+1}B_{a'}<p^f
\]
in Steinberg's short-supported parametrization.  The ordinary digit
factors occupy the disjoint ranges $0,\ldots,m$ and $m+1,\ldots,2m$.
This is a bijective change of parameters, simultaneously over all
paired nodes; it gives both existence and uniqueness for the unequal
boxes.  Put
\[
 b_i=m+c_{i'},\qquad b_{i'}=m+c_i,\qquad b_i+b_{i'}=f.
\]
The coefficient at $i$ is less than $p^{b_i}$, so the two node boxes
have $b_i$ and $b_{i'}$ slots.  Their concatenation is a circle of
length $f$.  The first exponent beyond the slots at $i$ is the
$F^*$-image of slot zero at $i'$.

Take shifts $s$ modulo $f$, discarding those for which any occupied
ordinary digit position $h$ satisfies $h+s=m\pmod f$.  At most
$\log_2n$ shifts are discarded.  A nonzero digit newly at
$e=h+s\pmod f$ is kept when $e<m$.  When $e>m$, replace it by the
module with highest weight
\begin{equation}\label{q:crossing-digit}
                         p^{e-m-1}\sigma^*\lambda_h.
\end{equation}
The pinned maps commute, and
\[
 \operatorname{Fr}_p^e=
       \operatorname{Fr}_p^{e-m-1}\sigma F.
\]
Since $F$ fixes the finite group pointwise, this replacement preserves
the finite restriction.  Isogeny pullback preserves irreducibility.
On each node pair the crossing exponent becomes
\[
               e-m-1+c_i=e-b_i
\]
at the other node, exactly rotating the concatenated slots by $s$.
Excluding $e=m$ lets the whole digit be kept or crossed simultaneously
on all its nodes.  There are no overlapping slots.  The resulting
rational tensor has irreducible finite restriction, hence is itself
irreducible, and its sum highest weight is $F$-restricted.  Uniqueness
of restricted simples on the simply connected fixed group forces the
rotation to preserve support whenever the shift is realized by
$\Gamma$.  On a nonempty paired-node cycle at most $2\log_2n$ slots
are occupied, so at most this many rotations preserve its support.
Including the discarded shifts and the bounded diagonal/graph factor
gives, in all types, the ample bound
\begin{equation}\label{q:outer-index}
                           [\Gamma:S]\le J^3(x).
\end{equation}

\emph{Claim 4. There is a positive-height standard point map $\Psi$ of $G$
such that $S=(G^\Psi)'$ and $N_G(S)\le G^\Psi$.}
If several ordinary digits occur, $S$ preserves a nontrivial tensor
splitting.  Its stabilizer in $G$ is proper by the fundamental-module
argument.  At most $[\Gamma:S]$ translates cover $\Gamma$, giving a
proper bounded-format closed set, a contradiction.  The union remains
proper because the ambient connected group variety is irreducible.

For a single digit, untwist its field power, which permutes the pinned
finite points.  The same projective $S$ is then the finite-point image of
a rational module with highest-weight coefficients below $n$.  Let $C$
be its adjoint algebraic domain.  The preimage in $C$ of the projective
image of $G$ is closed, since algebraic-group images are closed, has
format at most $J^{10}(x)$, and contains $(C^F)'$.  To obtain the format
bound use the two common-word action graphs with $y=J^2(x)$, which
covers both ranks, module dimensions, and highest-weight coefficients.
Apply \Cref{q:root-formats} with $y=J^{11}(x)$.  The inequality
$P>J^{100}(x)$ makes the derived-fixed-group order exceed the number
of fixed points allowed in a proper subset, so this preimage is all of
$C$.  The projective image of $C$ cannot be strictly smaller than that
of $G$: its pullback in $G$ would be a proper bounded-format closed
subgroup containing $S$, whose at most $[\Gamma:S]$ translates again
cover $\Gamma$.  The word-graph estimates and \eqref{q:outer-index}
keep both this test and the tensor test below $J^{1000}(x)$.

The adjoint domains are therefore identified on points through these
projective isogenies.  By the isogeny classification, transporting $F$
gives a positive-height standard point-map element $\Psi$.  Up to
inner maps, self-identifications use diagrams and positive or negative
split-field or special powers, and transport conjugates within that
point-map group.  For the dual pair $B_e,C_e$ in characteristic two,
$e\ge3$, the pinned cross isogenies commute with split powers and need
no diagram twist.  We may compose the two projective isogenies with the
adjoint quotient of their common image, which changes nothing on points,
as the maps from the adjoint domains are already point-injective.  Even
if both maps pass through a cross-type target, the cross maps and their
pointwise inverses commute with split powers.  Hence $S=(G^\Psi)'$.
If $g\in G$ normalizes $S$, then $g^{-1}\Psi(g)$ centralizes $S$,
whose inner centralizer is trivial; thus $\Psi(g)=g$.  Since $\Gamma$
normalizes $S$, this proves \eqref{q:sandwich}.
\end{proof}

\subsection{A simultaneous parameter choice}

\begin{definition}\label{q:ceiling-sequence}
Set
\begin{equation}\label{q:x-sequence}
 x_0=10^9,\qquad x_{j+1}=J^{x_j}(x_j)\quad(0\le j<5),
 \qquad M_*=x_5.
\end{equation}
These are ordinary computable integers.
\end{definition}

\begin{theorem}[Size ceiling]\label{q:size-ceiling}
The negative affine-chart construction has a choice of finite parameters
for which its lattice $L$ is nonrepresentable and
\[
                              |L|<M_*.
\]
In particular a nonrepresentable lattice exists among the tables of
cardinality at most the ordinary computable integer $M_*$.
\end{theorem}

\begin{proof}
We verify the quantitative requirements of the structural pruning,
protected-path, and bounded-root arguments.  The parameters are selected
in the order
\[
 x_0\ ;\quad D,q\ ;\quad \ell,r\ ;\quad
 \text{backward trial increments}\ ;\quad m,n\ ;\quad N,g.
\]
This order matters because only the initial graph-pattern round uses a
slot count depending on $n$.  The later-round slot counts and the chart
losses are fixed before the supply of charts is chosen, as required by
\Cref{c:parameter-order}.

\emph{High-type thresholds and structure counts.}
Use $10^8$ as the high cutoff for alternating degree or classical rank;
the remaining ranks are at most $x_0$.  The high automorphism and covering
assertions have only the usual low exceptions.  In the notation of the
direct-system pruning in \Cref{c:direct-block-groups,c:direct-compatibility},
supported classical groups of natural dimension $h$ have full special
order, also after projectivizing, at least
\[
                 s_0^{c h^2-20h},
\]
whereas the full isometry or linear order is at most
$s_0^{c h^2+20h}$, including in small dimensions.  The direct-split
upper bound
\[
                  h!s_0^{c h^2/2+20h}
\]
is strictly smaller than the supported full special order for $h>10^6$.
The linear comparison in \Cref{c:tensor-order} likewise follows from
\[
          h!|F|^{h^2/4+4}<|F|^{h^2/2-20h}
          \qquad(h>10^6).
\]
Perfection, natural irreducibility, and existence of the supported
classical groups already hold there.  In the initial subspace pruning of
\Cref{c:linear-subspaces,c:isometry-subspaces}, the retained small-space
irreducibility cutoff may be $100$.
Completing two such spaces costs dimension at most $400$.  Summing the
endomorphism dimensions at the successive span-growth steps, including
the two spans in linear type, bounds the independent-list requirement by
fewer than $1000^5$ strict steps.

In the direct-block, field, and tensor rounds, the small-special-group,
rotation, and bounded-field cases can all use $10^6$ as the large cutoff
for the indicated dimension or field size.  In particular plane rotations
then have projective order greater than four when required.  Transfer
avoids the ordinary small-special exceptions once total ground dimension
exceeds $10^6$.  This remains so on the complement of a nuisance-plane
tensor slot, by either a high transferred dimension or a large extension
degree.  The compositum-torus test in \Cref{c:isometry-fields}, in total
dimension $4s>10^6$, uses only
\[
                  \frac{|F|^s+1}{2s}>2+|F|+1.
\]
Here $s>250000$, so the inequality holds already at $|F|=2$ and
continues to hold as $|F|$ increases.
The counts here never bound arbitrary matrices over an unbounded field
merely from bounded dimension.  Instead, the order comparisons bound
refinement fibers, comparable split lengths, or the number of distinct
elementwise-commuting simple tensor slots, of which at most two are then
designated.  When matrices inside a bounded plane algebra are counted,
the field is bounded too.  Exceptional matching counts use the bounded
aligned types of a matched pair; only the plane case needs isometry
choices per edge orbit, and that case has bounded field.

For the top-action orbit count in \Cref{c:orbit-bound}, take the six
successive cutoffs
$x_0,J(x_0),\ldots,J^5(x_0)$.  There are at most four class sizes, so
one of the five gaps has no class size strictly between its endpoints.
Call sizes at most its lower endpoint small.  Small sizes are then at
most $J^4(x_0)$, while each large alternating order exceeds $1000$ times
the product of the small factorials, using $J$ at the actual lower
endpoint.  The joint action intersected with the product of the large
alternatings has index at most $2^4\cdot2^2$ times those factorials:
the factors account for parities and the two possible class swaps.
This is enough for the subdirect projection test.  Equality patterns on
at most four large-index occurrences, up to the diagonal
identifications, together with class choices and the individual small
indices, bound the orbit count by, for example,
\begin{equation}\label{q:top-orbit-bound}
                        2^4\cdot15^4\bigl(J^4(x_0)\bigr)^4.
\end{equation}
For fine-cell pairs multiply by at most $(10^6)^2$, from the aggregate
split bounds.  For the tie-chain test on simple-site pairs, fibers of
size at most $2^2$ suffice.  Point cells in the alternating partition
argument have independent supported alternating groups, giving bounded
orbits in every ordered-cell-pair fiber as well.

Invariant partition and subset counts exponentiate these orbit bounds
once.  A further isometry choice on each bounded-field plane edge orbit
costs base at most $(10^6)^4$ per orbit.  A count inside a single
bounded-field plane matrix algebra may even use the entire power-set
size.  In the shared-component ordered-factor arguments, a matrix
centralizer factor of natural degree at most $10^6$ permits at most
$\log_2(10^6)$ nontrivial elementwise-commuting new tensor slots, over
any field.  A large-space factor permits only possibly the whole factor
as a new one there.  Designating one or two components among the
original at most four and these additions has bounded cost.  Therefore
$J^{50}(x_0)$ exceeds all line-count and rich-top tie-chain bounds,
including the small-field and matching choices.  It also bounds the
retained/missing, support, group-type, large/small, and exceptional-case
subdivisions, and the categories and rounds per ancestor.  Recording
even a thousand such tests or small values up to $x_0$ is within this
bound.  No color records an arbitrary degree or an arbitrary matrix.

\emph{The protected path.}
For a nontrivial projective representation in defining characteristic
and high classical natural dimension $v$, the Weyl-orbit lower bound
$v/4$ from \Cref{w:min-degree} defeats both the smallest tensor slots
and the powers $v^{1/s}$, $s\ge2$.  These comparisons already hold at
dimension $10^6$ in \Cref{w:coarsening-count}, leaving fewer than
$10^6$ primes when that argument stops.  The positive-height fixed-order
comparison used in \Cref{w:root-count} follows from the bounded product
factors in high rank: a height step contributes at least $2^{\dim G}$.

For a change of defining characteristic in \Cref{w:characteristic-switch},
the paired linear subgroup has dimension $h\ge v_j/3$.  Its column-shear
characters give $v_i\ge P_j^{h-1}-1$, and hence, for $v_j\ge10^6$,
\begin{equation}\label{q:characteristic-change}
 \log_2v_i\ge\frac{v_j}{5}\log_2P_j.
\end{equation}
A second change would give $v_l\log_2P_l\le5\log_2v_j$.
Using the uniform order bound $|S_l|\le P_l^{3v_l^2}$, we obtain
\[
 \log_2|S_l|\le
 \frac{75(\log_2v_j)^2}{\log_2P_l}
 \le75(\log_2v_j)^2<\frac{2v_j}{5}\le2\log_2v_i,
\]
contradicting the required representation-degree comparison.  The strict
inequality holds at $v_j=10^6$ and thereafter because
$v_j/(\log_2v_j)^2$ increases in this range.  For
\Cref{w:alternating-switch}, the block count and Bochert bound give index
at least $2^{e/7}$ for a proper transitive subgroup in large degree $e$.
The numerical comparison is
\[
                       e!<2^{(\binom e2-1)/7},
\]
since $e!\le e^e$ and $14\log_2e<e-1-2/e$ for $e\ge10^6$.
The classification of $2$-transitive groups has no nonnatural
alternating exception at these degrees.

The word tests in \Cref{w:a-obstruction,w:signed-obstruction} require
only a high constant rank: there must be enough
prefix positions and neighbors, and in the spin case the exponentially
many sign choices must exceed the rank.  Fewer than $100$ pair colors
at a fixed series and a homogeneous set of size ten suffice.  Thus a
high protected path has at most two high alternating nodes and at most
one characteristic switch on its classical subsequence, giving at most
four consecutive same-characteristic classical runs.  Within each run
pigeonhole among the four series, and then apply pair Ramsey to the
shapes of the weight digits in that series.  The total bounds are far
below $J^2(x_0)$.  We may therefore take
\begin{equation}\label{q:Dq}
                  D=J^{1000}(x_0)<x_1,
                  \qquad q=7^D<x_1.
\end{equation}
The extra stages allow a displaced point start and a stage with a
bounded-type parent.  The field order $q$ exceeds every absolute
line-count requirement, and $q-1$ has the distinct prime divisors $2$
and $3$, as required by the affine-chart construction.

\emph{Bounded-type descents and residual length.}
Use $x_1$ as the rank bound in \Cref{q:genericity}.  If its avoidance
hypothesis fails, there is a proper closed $X$ of format at most
$J^{1000}(x_1)$.  By \Cref{w:point-maps}, every point map used in
\Cref{w:witness} has one of the forms
\[
 \operatorname{Int}(g)\tau\operatorname{Fr}_p^h,
 \qquad
 \operatorname{Int}(g)\sigma^\epsilon\operatorname{Fr}_p^h,
 \qquad h\in\mathbf Z,\quad\epsilon\in\{0,1\},
\]
where $\tau$ is a pinned diagram map.  Indeed
$\sigma^2=\operatorname{Fr}_p$ and
$\sigma^{-1}=\sigma\operatorname{Fr}_p^{-1}$.  In the fixed split
models, positive and negative field powers change only equation
coefficients.  The remaining inner, diagram, and special maps have
bounded-format graphs, independently of $h$; special maps occur only
in characteristics two and three.  Thus images and inverse images
are given by bounded formulas after the coefficient twist.  These
point maps are Zariski homeomorphisms, so their images of $X$ are
closed.  Applying the fiberwise equation bound of
\Cref{q:acf-computable} gives equations of degree less than
$J^{1100}(x_1)$ for all the translates and transforms defining $Z$.

For $u$ ambient coordinates and equation degree at most $d_0$, their
entire equation span has dimension at most $\binom{u+d_0}{u}$.
Selecting a basis therefore replaces the possibly infinite
intersection defining $Z$ by that many equations.  This elementary
cost keeps its format below $J^{1200}(x_1)$.  Its conjugation
normalizer is defined by the bounded-format condition
\[
 \forall z\in G\quad
       \bigl(z\in Z\longleftrightarrow gzg^{-1}\in Z\bigr).
\]
Taking equations for this closed subgroup and its monomial-space
dimension keeps the latter below $J^{1400}(x_1)$.  The independent-list
strict descents securing the sandwich are bounded by this dimension;
one or two more list members suffice.

All large-parameter requirements may be imposed by
\begin{equation}\label{q:bounded-cutoff}
                             P>J^{2000}(x_1).
\end{equation}
This follows from the fixed-point versus degree comparison, including
the zero-height fixed loci.  To compare a smaller positive height whose
parameter is at least $\sqrt P$, the usual product formulas, including
the actual central index, give relative errors at most
\[
                          \frac{100(x_1+10)^4}{\sqrt P}.
\]
The minimal leading full-order growth factor per height has integer
square and exceeds the relevant center bound
$z\le10(x_1+10)^2$.  The displayed error is therefore small enough
for the strict order ratio.  For parameters below $\sqrt P$ the gross
order comparison suffices.

All fixed-group orders below the cutoff in
\eqref{q:bounded-cutoff}, as well as bounded alternating, sporadic, and
exceptional-small socle orders, are below $J^{2010}(x_1)$.  Reserving
residual length $J^{2020}(x_1)$ enforces the cutoff by the available
index over the base, both below current path nodes and below tested
bounded-type flats.  In the final bounded-type label estimate $C2^t$,
the sandwich quotient permits $C=2^z$.  Hence the increment
\begin{equation}\label{q:t}
                              t=J(x_1)
\end{equation}
suffices.

\emph{Flag selection and backward pattern increments.}
Only the two continuation colors in \Cref{w:selection} are needed.  Set
\begin{equation}\label{q:ell-r}
                 \ell=J^{2030}(x_1),\qquad r=3\ell.
\end{equation}
Then $Dr<x_2$ and $t<r$.  The tensor-prime and root counts give the
failure union bound
\begin{equation}\label{q:flag-error}
                     2q^{-\ell}D(3Dr+10^6),
\end{equation}
which is smaller than the two-color matching fraction $1/3$.
Indeed $q^\ell>6D(9D\ell+10^6)$ at the displayed choices.

The total number of successive high-pruning rounds in
\Cref{c:pruning-schedule} at one stage is
less than $x_2/100$: these comprise the simultaneous initial subspace
round and the ordered categories and cases relative to at most $D$
ancestors.  Every subsequent round uses at most $x_2$ slots.  Compute
the required increments backwards, beginning with $s=r$.  Suppose the
desired next increment satisfies $s\le y$, where $y\ge x_2$.
The current path rank is at most $Dr$, and the alphabet of coordinate
functionals has size at most $q^s$.  A failure color records an
ancestor, a bounded category and case, and a pullback subset of at most
$q^{Dr+s}$ points.  These use fewer than $J(y)$ colors.  For example,
the subset contribution alone is bounded by
\[
       2^{q^{Dr+s}}\le2^{y^{2y}}<2^{2^{y!}}<J(y),
\]
with ample room for the bounded ancestor and case factors.  The
alphabet and number of slots also fit below $J(y)$.  The definition
of $J$ thus bounds the Hales--Jewett vertical length by $J^2(y)$,
and $J^3(y)$ accommodates the preceding increment including the
horizontal directions.  Fewer than $x_2/100$ backward rounds require
fewer than $3x_2/100$ iterations, strictly below the $x_2$ iterations
defining $x_3$.  In particular the increment sought in the initial
round is below $x_3$.

\emph{Truncation, initial slots, and chart supply.}
Choose
\begin{equation}\label{q:mnNg}
 \begin{gathered}
       m=2x_3,\qquad n=4x_3,\qquad y=J(x_3),\\
       N=J^4(y),\qquad g=J^4(y).
 \end{gathered}
\end{equation}
Path ranks are below $x_2$ and trial increments below $x_3$, so all
path and trial ranks lie strictly below $m$.  The additional at most
$x_2$ independent-list and pair-span allowance still lies below $m$;
residual heights are greater than $2x_3$.  These exceed the reserved
bounded-type residual length.  Initial subspace pruning may ask for
more labels than the interval height, but its argument never requires
the whole list to have a proper span: only pair spans and the bounded
independent lists must fit the truncation.

After $m,n$ are fixed, the initial graph patterns may use $n+x_2$
slots.  This fits below $y=J(x_3)$; the color and alphabet bounds fit
below $J(y)$ as before.  These bounds also cover the direct trial at
a bounded-type parent, without any earlier anchor.  Thus the value
of $N$ in \eqref{q:mnNg} suffices.  The chart loss in any round is
less than $J(y)$, since it is bounded by the ancestor/category/case
colors and two matching failed charts suffice for contradiction.
Keep one successful flat in each surviving chart for the next round.
There are fewer than $x_2/100$ rounds per node, so the choice of $g$
leaves successful charts after all of them.  The supply is private
to each path node, rather than a common pool for the whole tree.
Both $N$ and $g$ are less than $x_4$.

These choices implement the pattern procedure: in a failed chart use
graphs of horizontal-to-vertical linear maps, or affine maps at the
empty start, and color them by the failure data and pullback flat.
The Hales--Jewett family fixes the pullback.  Varying its slots by
scalar multiples of a horizontal functional nonzero at a point of
that flat gives distinct witness flats indexed by an affine parameter
space.  Pair spans add one dimension and contain the whole parameter
line; bounded independent lists add one dimension per new point.
The chart intersection rules therefore turn these patterns into the
line and descent configurations used in the exclusions.  Witnesses
from two distinct failed charts have label intersection exactly the
base group.  Successive rounds retain earlier exclusions because
the later chosen labels contain the earlier successful anchor.
Consequently at each high node a pruned increment-$r$ flat gives the
next almost simple socle with its ancestor protections, and the flag
selection preserves the required root and tensor data.  A bounded-type
parent instead gives the direct bounded-complexity contradiction.
Thus a putative representation either has that contradiction or
forces a protected high path beyond the chosen $D$.  The resulting
finite lattice is nonrepresentable.

\emph{Counting every lattice element.}
Take one leaf for each required leaf attachment.  A chart has at most
\begin{equation}\label{q:flat-count}
                                F_0=2^{q^N}
\end{equation}
flat subsets.  There are $g$ initial charts, branching at most $gF_0$
at each of $D$ generations, and at most $F_0$ leaf attachments per
trunk chart.  Thus the number of all charts is at most
\begin{equation}\label{q:chart-count}
                  C=(D+1)g(1+gF_0)^D(1+F_0).
\end{equation}
Both orientations together have at most $2CF_0$ interior labels, and
there are two common endpoints.  Put $X=x_4$ and $A=J(X)$.  Since
$q,N<X$, we have
\[
                  q^N\le X^X<2^{X!},\qquad
                  F_0<2^{2^{X!}}<A.
\]
Also $D,g<X<A$, so $1+gF_0<A^2$ and
$(D+1)g(1+F_0)<A^4$.  Hence $C<A^{2X+4}$ and
\begin{equation}\label{q:strict-ceiling}
             |L|\le2CF_0+2<A^{2X+7}<J(A)=J^2(X)<x_5.
\end{equation}
The middle inequality follows from $J(A)>2^{2^{A!}}$ and
$A>X\ge10^9$; the last follows from the $x_4$ iterations defining
$x_5$.  This counts the endpoints as well as every interior label and
proves the claimed ceiling.
\end{proof}

\section{A finite description of a minimum lattice}\label{sec:minimum}

The computable ceiling $M_*$ from \Cref{q:size-ceiling} can be combined
with one ordinary Busy Beaver constant to specify a nonrepresentable
lattice of minimum cardinality.  The constant is defined independently
of lattices and representations.  Once it is given, the carrier and
every order-table bit are specified by bounded finite arithmetic.
All arithmetic and all assertions of termination in this section refer
to the ordinary standard natural numbers.

\subsection{The register language and its constant}

\begin{definition}[Register programs and size]\label{q:register-language}
Registers $v_i$, $i\ge0$, hold natural numbers and initially all contain
zero.  Terms use these registers, the constants $0,1$, addition,
truncated subtraction $u\mathbin{\dot-}v$, multiplication, exponentiation
$u^v$, and integer floor quotient.  We set $0^0=1$ and give quotient by
zero the value zero.  Bounded formulas use the binary predicates $=,<$,
the connectives $\neg,\land,\lor$, and quantifiers
$\exists v_i<t$, $\forall v_i<t$, with $v_i$ absent from the bound term
$t$.  A variable bound within a formula is local to that formula and
does not write to a register.

Programs are formed by assignments of terms to registers, binary
sequencing, if-then-else with a bounded condition and two program
branches, and while with a bounded condition and a program body.  They
have the usual deterministic semantics.  A terminating program outputs
its final value of $v_0$.

The size of a program is its syntax-tree cost: each operation, constant,
predicate, connective, quantifier, assignment, sequencing,
if-then-else, or while node costs one.  Each occurrence of $v_i$ costs
$1+i$, including assignment targets and quantifier binders.  Parentheses
cost nothing, and no extra cost is charged for the use of a quantifier's
term as its bound.
\end{definition}

\begin{definition}[The Busy Beaver cutoff]\label{q:BB}
With $M_*=x_5$ as in \eqref{q:x-sequence}, set
\begin{equation}\label{q:BB-bound}
 \begin{split}
 k&=2^{(M_*+10)^4},\\
 B&=1+\max\Bigl(\{0\}\cup
       \{\operatorname{output}(p):p\text{ terminates and has size }
                                      \le k\}\Bigr).
 \end{split}
\end{equation}
The programs in this definition use exactly
\Cref{q:register-language}; in particular $B$ is not a symbol in their
language.
\end{definition}

This is a maximum-output variant of Rad\'o's finite-program extremal
construction \citep{Rado1962}, with the register syntax and cost fixed
here.  It is not the classical two-symbol state-count function.
There are finitely many programs of bounded size, since the cost of a
register occurrence bounds its index.  Thus $B$ is a definite finite
integer.  It is an ordinary, unevaluated Busy Beaver constant, defined
by termination and output of programs without an oracle.

\subsection{Bounded arithmetic for lattice tables and witnesses}

\begin{definition}[Digits and order tables]\label{q:table-predicate}
For $b\ge1$, put
\begin{equation}\label{q:digit}
 \operatorname{dig}_b(z,i)=
   \left\lfloor\frac{z}{b^i}\right\rfloor
       \mathbin{\dot-}
       b\left\lfloor\frac{z}{b^{i+1}}\right\rfloor.
\end{equation}
Base $1$ is used only where the intended digit is zero, in which case
this formula indeed gives zero.  For $n\ge1$ and $0\le t<2^{n^2}$,
let $T_{n,t}$ be the relation on $\{0,\ldots,n-1\}$ defined by
\begin{equation}\label{q:order-bit}
          a\preceq b\quad\Longleftrightarrow\quad
                         \operatorname{dig}_2(t,an+b)=1.
\end{equation}
Write $E(a,b)$ for the right side, with $n,t$ understood.  The bounded
predicate $P(n,t)$ says that $T_{n,t}$ is a lattice: it is the conjunction
of
\begin{align*}
 &\forall a<n\ E(a,a),\\
 &\forall a,b<n\ \bigl((E(a,b)\land E(b,a))\Longrightarrow a=b\bigr),\\
 &\forall a,b,c<n\ \bigl((E(a,b)\land E(b,c))\Longrightarrow E(a,c)\bigr),
\end{align*}
and the following two assertions for every $a,b<n$:
\begin{align*}
 \exists u<n\ \bigl(&E(a,u)\land E(b,u)\land
   \forall v<n\,((E(a,v)\land E(b,v))\Longrightarrow E(u,v))\bigr),\\
 \exists l<n\ \bigl(&E(l,a)\land E(l,b)\land
   \forall v<n\,((E(v,a)\land E(v,b))\Longrightarrow E(v,l))\bigr).
\end{align*}
Repeated bounded variables denote successive quantifiers with the
displayed bound.  Implication, equivalence, and non-strict comparison
throughout this section abbreviate their expressions using the
primitive connectives and $=,<$.
\end{definition}

The next definition gives a bounded formula for the graph criterion,
including the admissible placements and the paths.  All bound variables
are renamed when abbreviations are expanded, to avoid capture.

\begin{definition}[The witness predicate]\label{q:witness-predicate}
For $j\ge1$, the formula $Q(n,t,j)$ asserts that there is
\begin{equation}\label{q:color-encoding}
 z<n^{j^2},\qquad d(p,q)=\operatorname{dig}_n(z,pj+q)\quad(p,q<j),
\end{equation}
with the following properties.

First, the colors are symmetric, bottom occurs exactly on loops, and
the triangle inequality holds.  Explicitly, set
$\operatorname{Bot}(a)\equiv\forall b<n\ E(a,b)$ and require
\begin{align}
 &\forall p,q<j\quad d(p,q)=d(q,p),\notag\\
 &\forall p,q<j\quad
        \bigl(\operatorname{Bot}(d(p,q))\longleftrightarrow p=q\bigr),
        \label{q:graph-bottom}\\
 &\forall p,q,u<j\ \forall b<n\quad
 \bigl(E(d(p,u),b)\land E(d(u,q),b)\bigr)
                       \Longrightarrow E(d(p,q),b).\notag
\end{align}
Second, require order reflection:
\begin{equation}\label{q:graph-reflect}
 \forall a,b<n\quad \neg E(a,b)\Longrightarrow
  \exists p,q<j\ \bigl(E(d(p,q),a)\land\neg E(d(p,q),b)\bigr).
\end{equation}

For the third requirement, let $0<s<2^{j^2}$ be a nonempty ordered-pair
mask and write
$\operatorname{Sel}_s(p,q)\equiv\operatorname{dig}_2(s,pj+q)=1$.
A number $f<j^j$ specifies a vertex map by
$f[p]=\operatorname{dig}_j(f,p)$.  Set
\begin{equation}\label{q:placement-arithmetic}
 \operatorname{Adm}(f)\equiv
   \forall p,q<j\ E\bigl(d(f[p],f[q]),d(p,q)\bigr).
\end{equation}
The derived undirected adjacency is the bounded predicate
\begin{equation}\label{q:adjacency-arithmetic}
 \begin{split}
 \operatorname{Adj}_s(r,r')\equiv\exists f<j^j\ \bigl(
   &\operatorname{Adm}(f)\land
       \exists a,b<j\ \bigl(\operatorname{Sel}_s(a,b)\land\\[-2pt]
   &((f[a]=r\land f[b]=r')\lor
                     (f[a]=r'\land f[b]=r))\bigr)\bigr).
 \end{split}
\end{equation}
For $h<j^{j+1}$, put $h[i]=\operatorname{dig}_j(h,i)$, and define
\begin{equation}\label{q:path-arithmetic}
 \begin{split}
 \operatorname{Path}_s(p,q)\equiv\exists h<j^{j+1}\ \bigl(
      &h[0]=p\land h[j]=q\land\\[-2pt]
      &\forall i<j\ (h[i]=h[i+1]\lor
                            \operatorname{Adj}_s(h[i],h[i+1]))\bigr).
 \end{split}
\end{equation}
Write
\begin{align}
 \operatorname{Up}_s(b)&\equiv
   \forall a,c<j\,
       \bigl(\operatorname{Sel}_s(a,c)\Longrightarrow E(d(a,c),b)\bigr),
           \label{q:mask-upper}\\
 \operatorname{BelowJoin}_s(p,q)&\equiv
   \forall b<n\,
       \bigl(\operatorname{Up}_s(b)\Longrightarrow E(d(p,q),b)\bigr).
           \notag
\end{align}
The third requirement is
\begin{equation}\label{q:graph-connect-arithmetic}
 \forall s<2^{j^2}\quad s\ne0\Longrightarrow
   \forall p,q<j\,
     \bigl(\operatorname{BelowJoin}_s(p,q)
                          \Longrightarrow\operatorname{Path}_s(p,q)\bigr).
\end{equation}
The conjunction of all three initial color requirements---symmetry,
the bottom condition \eqref{q:graph-bottom}, and the triangle condition---
with \eqref{q:graph-reflect} and \eqref{q:graph-connect-arithmetic},
under the existential quantifier in \eqref{q:color-encoding},
is $Q(n,t,j)$.
\end{definition}

Every expression just given expands to the register language's bounded
arithmetic.  All element and vertex variables have bounds $n$ and $j$;
the remaining bounds $n^{j^2}$, $2^{j^2}$, $j^j$, and $j^{j+1}$ are
terms.  In particular the definition uses neither an unbounded path
quantifier nor a quantifier over a function as a separate sort.

\begin{lemma}[Meaning of the arithmetic witness]\label{q:arithmetic-criterion}
For every table satisfying $P(n,t)$,
\begin{equation}\label{q:arithmetic-R}
 \mathcal R\bigl((\{0,\ldots,n-1\};\preceq)\bigr)
       \quad\Longleftrightarrow\quad
                     (\exists j\ge1)\ Q(n,t,j).
\end{equation}
Moreover $Q(n,t,j)$ holds exactly when the colored-graph criterion has
a witness with $j$ vertices.
\end{lemma}

\begin{proof}
For a table satisfying $P(n,t)$, the digit encoding lists all color
tables on $j$ vertices.  The predicate $\operatorname{Bot}$ identifies
the least lattice element regardless of its numerical label.  Being
below every common upper bound is equivalent to being below the join,
so the triangle clause and $\operatorname{BelowJoin}_s$ express the
corresponding lattice inequalities in \Cref{f:criterion}.

The placement codes list all selfmaps of the vertex set.  The nonzero
ordered-pair masks select exactly the same edge sets as undirected
masks: colors are symmetric, and $\operatorname{Adj}_s$ permits either
ordering of the endpoints.  Every connected pair has a path of at most
$j-1$ edges; equality steps pad it to the $j$ steps encoded by $h$.
Conversely, every encoded path certifies connectivity.  The base-one
digit convention covers $n=1$ and $j=1$, and the masks include loops.
Thus $Q(n,t,j)$ holds exactly when there is a graph witness with $j$
vertices.  The representability equivalence now follows from
\Cref{f:criterion}.
\end{proof}

\subsection{The finite table selection}

\begin{definition}[Carrier and order]\label{q:minimum-table}
Let
\[
 \mathcal I=\{(n,t):1\le n\le M_*,\ 0\le t<2^{n^2}\}
\]
have increasing pair-lexicographic order, with $n$ compared first.
Define
\begin{equation}\label{q:selection-z}
 z_{n,t}=\begin{cases}
  1,& P(n,t)\ \text{and}\ (\forall j=1,\ldots,B)\ \neg Q(n,t,j),\\
  0,&\text{otherwise},
 \end{cases}
\end{equation}
and, with the empty product equal to one,
\begin{equation}\label{q:selection-w}
        w_{n,t}=z_{n,t}
                \prod_{\substack{(m,u)\in\mathcal I\\(m,u)<(n,t)}}
                              (1-z_{m,u}).
\end{equation}
Set
\begin{equation}\label{q:a0t0}
       a_0=\sum_{(n,t)\in\mathcal I}n w_{n,t},\qquad
       t_0=\sum_{(n,t)\in\mathcal I}t w_{n,t}.
\end{equation}
The structure $L_0$ has carrier $\{i:0\le i<a_0\}$ and order
\begin{equation}\label{q:minimum-order}
               i\le_{L_0}l\quad\Longleftrightarrow\quad
                    \operatorname{dig}_2(t_0,ia_0+l)=1.
\end{equation}
\end{definition}

All products, sums, and tests in this definition are bounded finite
constructions once $B$ is fixed.  We next prove that the cutoff is
exhaustive on this entire initial segment, so that a failed bounded
test selects a negative instance.

\begin{lemma}[Size of a graph-witness search]\label{q:search-size}
For $1\le n\le M_*$ and $0\le t<2^{n^2}$, there is a program of
\Cref{q:register-language}, of size less than $k$, which terminates
exactly when $Q(n,t,j)$ holds for some positive $j$.  On termination
it outputs the least such $j$.  The program contains neither $B$ nor
any oracle instruction.
\end{lemma}

\begin{proof}
Fix this particular pair $n,t$.  Consider the following program schema,
where the first two right-hand sides are arithmetic terms for the
specified binary numerals and $Q$ is expanded as in
\Cref{q:witness-predicate}:
\[
 \begin{array}{l}
 v_1:=\text{the numeral for }n;\\
 v_2:=\text{the numeral for }t;\\
 v_3:=1;\\
 \textbf{while }\neg Q(v_1,v_2,v_3)\textbf{ do }v_3:=v_3+1;\\
 v_0:=v_3.
 \end{array}
\]
Binary sequencing makes this a program of
\Cref{q:register-language}.  Each test of $Q$ is bounded and therefore
terminates.  The loop tests the positive integers in order, proving the
claimed termination and output properties without assuming $P(n,t)$.
Writing the numerals by binary Horner
expressions uses only $0,1$, addition, and multiplication.  Since
$t<2^{n^2}$, the two numeral terms cost at most $100(n+1)^3$ nodes,
with ample allowance for their fixed register targets.

The bounded formula templates and the rest of the program schema are
fixed independently of $n,t$.  Their expansion can be counted directly.
Use $2=(1+1)$, replace $A\Rightarrow B$ by $\neg A\lor B$, replace
$A\Leftrightarrow B$ by
$(\neg A\lor B)\land(\neg B\lor A)$, and use binary conjunctions
and one binder for each successive quantifier.  Substitute every
abbreviation separately, without shared subexpressions or loop
unrolling.  Counting a variable occurrence initially as one gives:
\begin{center}
\begin{tabular}{lrrr}
\toprule
Clause & Nodes & Variable occurrences & Binders\\
\midrule
Symmetry & 53 & 26 & 2\\
Bottom condition & 173 & 68 & 4\\
Triangle condition & 240 & 92 & 4\\
Order reflection & 198 & 72 & 4\\
Mask and connectivity & 717 & 281 & 13\\
\midrule
Full $Q$ & 1394 & 542 & 28\\
\bottomrule
\end{tabular}
\end{center}
The last line adds four conjunction nodes and the outer existential:
its binder and the seven-node bound $n^{j^{1+1}}$ contribute nine
nodes with the quantifier itself.  Renaming the expanded binders
consecutively as $v_{100},\ldots,v_{127}$ avoids all capture; the
bound of each quantifier is renamed in the preceding environment.
Only $v_1,v_2,v_3$ remain free.  Every node then costs at most $128$,
so $Q$ costs at most $128\cdot1394=178432$.
Outside the two numeral assignments, the rest of the program costs
$29$: initialization costs $6$, the while and negation nodes cost $2$,
increment costs $11$, final output costs $6$, and the four sequencing
nodes cost $4$.  In particular the following much larger allowance
bounds the entire program:
\begin{equation}\label{q:program-size}
 100(n+1)^3+(2\cdot10^6)^{204}
                      <2^{(M_*+10)^4}=k.
\end{equation}
The second summand exceeds $178461$.  For the strict inequality,
$M_*+1\le2^{M_*}$ bounds the first summand by $2^{3M_*+7}$,
and $2\cdot10^6<2^{21}$ bounds the second by $2^{4284}$.
Since $M_*>10^9$, their sum is less than $2^{3M_*+8}$, which
is less than $2^{(M_*+10)^4}$.

\end{proof}

\begin{lemma}[Uniform Busy Beaver witness cutoff]\label{q:witness-cutoff}
If $1\le n\le M_*$, $0\le t<2^{n^2}$, and $P(n,t)$ holds, then
\begin{equation}\label{q:finite-cutoff}
 \mathcal R(T_{n,t})\quad\Longleftrightarrow\quad
                   (\exists j\in\{1,\ldots,B\})\ Q(n,t,j).
\end{equation}
In fact every representable such table has a witness of size strictly
less than $B$.
\end{lemma}

\begin{proof}
For a representable table, \Cref{q:arithmetic-criterion} makes the
program in \Cref{q:search-size} terminate with a witness size $j$.
Its size is less than $k$, so the definition of $B$ gives $j<B$.
Conversely, any $j\le B$ satisfying $Q(n,t,j)$ supplies a
representation by \Cref{q:arithmetic-criterion}.
\end{proof}

\begin{theorem}[Finite description of a minimum example]\label{q:minimum}
The bounded formulas \eqref{q:selection-z}--\eqref{q:minimum-order},
with the ordinary unevaluated Busy Beaver constant
\eqref{q:BB-bound}, specify the full finite order table of a
nonrepresentable lattice $L_0$ of minimum cardinality.  Its cardinality
is the integer $a_0$ in \eqref{q:a0t0}, and
\[
                              7<a_0\le M_*.
\]
This is a finite order description and a proof of minimum cardinality;
it does not supply a numerically evaluated value of that cardinality.
\end{theorem}

\begin{proof}
By \Cref{q:witness-cutoff}, for every index in $\mathcal I$ the bit
$z_{n,t}$ is one exactly when the table is a nonrepresentable lattice.
By \Cref{q:size-ceiling}, some such bit is one.  Let $(n_*,t_*)$ be
the first index for which it is one.  The product in
\eqref{q:selection-w} is one at that index; at each later index it
contains the zero factor $1-z_{n_*,t_*}$, and at each earlier index
the factor $z_{n,t}$ is zero.  Hence exactly one $w_{n,t}$ equals
one.  Formula \eqref{q:a0t0} gives $a_0=n_*$ and $t_0=t_*$, so
\eqref{q:minimum-order} is precisely this lattice's full order table.
It is nonrepresentable.

Every nonempty lattice of cardinality smaller than $a_0$ has an
order table among earlier indices in $\mathcal I$.  Its selection
bit is zero and its lattice predicate is true, so the cutoff lemma
makes it representable.  Thus $L_0$ has minimum cardinality among all
finite nonrepresentable lattices.  The represented range through
seven elements in \Cref{f:small-lattices} gives $a_0>7$, and its
index gives $a_0\le M_*$.  Nothing in this argument evaluates $B$:
all its uses concern the outputs of the terminating programs in its
definition.  The result is therefore the claimed finite description
with an exhaustive bound, rather than a numerical enumeration.
\end{proof}

The constants and entries above are unambiguous in the standard natural
numbers.  Their definition supplies no ordinary computable procedure
for evaluating the Busy Beaver constant, no evaluated decimal value
of $a_0$, and no conventional small diagram of the minimum lattice.
Those would be stronger forms of identification than the finite
description proved here.

\begin{proof}[Proof of \Cref{m:main}]
Part~\textup{(i)} is \Cref{f:criterion}; Part~\textup{(ii)} is
\Cref{n:undecidable}.  The finite description and its minimum-cardinality
property are proved in \Cref{q:minimum}.  The strict lower bound of seven
follows from \Cref{f:small-lattices}.  Part~\textup{(iv)} is
Corollary~\ref{n:interval-undecidable}.
\end{proof}
\begingroup
\raggedright
\renewcommand{\bibsection}{\section*{\refname}\addcontentsline{toc}{section}{\refname}}
\bibliographystyle{plainnat}
\bibliography{refs}
\endgroup
\end{document}
